% Twenty Firmware Projects on the NUCLEO-H7A3ZI-Q, from the Bench Inventory
% Build:  python build.py            (PDF + HTML)
%         python build.py --check sections/c07.tex   (compile one chapter alone)
%         python lint.py                             (house rules)
%         python build.py --drift ../EmbeddedLinux_Top20/build.py
\documentclass[11pt,a4paper]{article}
\usepackage[utf8]{inputenc}
\usepackage[T1]{fontenc}
\usepackage{lmodern}
\usepackage{microtype}
\usepackage[margin=21mm,top=24mm,bottom=24mm]{geometry}
\usepackage{xcolor}
\usepackage{graphicx}
\usepackage{booktabs}
\usepackage{tabularx}
\usepackage{array}
\usepackage{enumitem}
\usepackage{float}
\usepackage{caption}
\usepackage{listings}
\usepackage{tcolorbox}
\usepackage{tikz}
\usepackage[european]{circuitikz}
% Shared TikZ and circuitikz setup, used by main.tex and by the standalone
% figure wrapper that build.py generates for the HTML/SVG output.
\usetikzlibrary{arrows.meta,positioning,calc,shapes.geometric,shapes.misc,shapes.symbols,fit,backgrounds,chains,decorations.pathreplacing,decorations.markings,patterns,shadows,matrix}

% ---------------------------------------------------------------- colours
\definecolor{hwcol}{RGB}{255,236,179}   % hardware, boards, sensors
\definecolor{kercol}{RGB}{197,222,255}  % kernel space
\definecolor{usrcol}{RGB}{205,240,205}  % user space
\definecolor{extcol}{RGB}{228,228,228}  % external systems, cloud, PC
\definecolor{fwcol}{RGB}{245,210,235}   % MCU firmware
\definecolor{notecol}{RGB}{255,250,205} % notes
\definecolor{linecol}{RGB}{60,60,60}

% ---------------------------------------------------------------- block diagrams
\tikzset{
  every picture/.style={font=\sffamily\small, line width=0.6pt, >=Latex},
  blk/.style={draw=linecol, rounded corners=2pt, align=center, minimum height=8mm,
              inner sep=4pt, text width=28mm, fill=white},
  hw/.style={blk, fill=hwcol},
  kern/.style={blk, fill=kercol},
  user/.style={blk, fill=usrcol},
  ext/.style={blk, fill=extcol},
  fw/.style={blk, fill=fwcol},
  wide/.style={text width=44mm},
  narrow/.style={text width=18mm},
  layer/.style={draw=linecol, dashed, rounded corners=3pt, inner sep=6pt, fill=none},
  layerlabel/.style={anchor=north west, font=\sffamily\scriptsize\bfseries, inner sep=2pt},
  flow/.style={->, thick, draw=linecol},
  bus/.style={<->, thick, draw=linecol},
  lbl/.style={font=\sffamily\scriptsize, fill=white, inner sep=1.5pt, midway},
  irq/.style={->, thick, draw=red!70!black, dashed},
  % ---------------------------------------------------------------- UML
  umlclass/.style={draw=linecol, fill=white, align=left, inner sep=0pt, font=\sffamily\scriptsize},
  lifeline/.style={draw=linecol, dashed},
  actor/.style={draw=linecol, fill=extcol, rounded corners=2pt, align=center, minimum height=7mm, inner sep=4pt, font=\sffamily\small},
  msg/.style={->, draw=linecol, thick, font=\sffamily\scriptsize},
  reply/.style={->, draw=linecol, dashed, font=\sffamily\scriptsize},
  activation/.style={draw=linecol, fill=kercol, inner sep=0pt, minimum width=2.4mm},
  state/.style={draw=linecol, rounded corners=3mm, fill=usrcol, align=center, minimum height=8mm, inner sep=5pt},
  initial/.style={circle, fill=linecol, inner sep=0pt, minimum size=2.4mm},
  final/.style={circle, draw=linecol, fill=white, inner sep=0pt, minimum size=3.4mm,
                path picture={\fill[linecol] (path picture bounding box.center) circle (1mm);}},
  trans/.style={->, draw=linecol, thick, font=\sffamily\scriptsize},
  umlnote/.style={draw=linecol, fill=notecol, align=left, font=\sffamily\scriptsize, inner sep=3pt},
  % ---------------------------------------------------------------- bench art
  board/.style={draw=linecol, fill=green!25!black!60, rounded corners=1.5mm, minimum height=20mm, minimum width=32mm,
                text=white, font=\sffamily\bfseries\small, align=center},
  hat/.style={board, fill=blue!40!black!60},
  breadboard/.style={draw=linecol, fill=white, minimum height=12mm, minimum width=40mm, font=\sffamily\scriptsize},
  pinrow/.style={draw=linecol, fill=black, minimum height=1.6mm, inner sep=0pt},
  cable/.style={line width=1.2pt},
  usbcable/.style={line width=1.4pt, draw=black!70},
  ledsym/.style={circle, draw=linecol, minimum size=3mm, inner sep=0pt},
}

% ================================================================ MCU figures
% Appended for the firmware volume. Same colour vocabulary as above:
% hwcol = silicon, kercol = RTOS and kernel, usrcol = application task,
% fwcol = firmware module, extcol = host and debugger.
% Solid fills only: figures are built latex -> dvi -> dvisvgm, a path on which
% PDF transparency does not survive.
\usetikzlibrary{shapes.multipart,decorations.pathmorphing}

\tikzset{
  % ------------------------------------------------------------ memory map
  mmap/.style={draw=linecol, anchor=south west, align=center, inner sep=2pt,
               minimum width=34mm, text width=34mm, font=\sffamily\scriptsize},
  mmflash/.style={mmap, fill=hwcol},
  mmram/.style={mmap, fill=kercol},
  mmperiph/.style={mmap, fill=fwcol},
  mmext/.style={mmap, fill=extcol},
  mmres/.style={mmap, fill=black!6},
  mmaddr/.style={font=\scriptsize\ttfamily, anchor=east, inner sep=2pt},
  mmgap/.style={draw=linecol, dashed, line width=0.4pt},
  mpu/.style={draw=red!70!black, line width=0.8pt, dashed, rounded corners=1pt, fill=none},
  % ------------------------------------------------------------ register bit field
  regcell/.style={draw=linecol, fill=white, anchor=south west, inner sep=0pt,
                  minimum height=6mm, align=center, font=\sffamily\scriptsize},
  regfield/.style={regcell, fill=kercol},
  regrsvd/.style={regcell, fill=black!6},
  regro/.style={regcell, fill=extcol},
  regset/.style={regcell, fill=hwcol},
  regbit/.style={font=\sffamily\tiny, anchor=south, inner sep=1pt},
  regname/.style={font=\sffamily\bfseries\small, anchor=east, inner sep=3pt},
  regreset/.style={font=\tiny\ttfamily, anchor=north, inner sep=1pt},
  % ------------------------------------------------------------ clock tree
  osc/.style={blk, fill=hwcol, text width=20mm, minimum height=7mm},
  pll/.style={blk, fill=fwcol, text width=22mm, minimum height=9mm},
  clkmux/.style={draw=linecol, fill=white, trapezium, trapezium angle=70,
                 trapezium stretches body, shape border rotate=270,
                 minimum width=9mm, minimum height=8mm, inner sep=1pt, font=\sffamily\tiny},
  clkdiv/.style={draw=linecol, fill=kercol, minimum width=12mm, minimum height=6mm,
                 inner sep=1pt, align=center, font=\sffamily\tiny},
  clkgate/.style={draw=linecol, fill=white, circle, inner sep=0pt, minimum size=3.2mm,
                  font=\sffamily\tiny},
  clkdom/.style={blk, fill=usrcol, text width=24mm, minimum height=6mm},
  clkline/.style={-Latex, draw=linecol, line width=0.7pt},
  clkfreq/.style={font=\sffamily\tiny, fill=white, inner sep=1pt},
  % ------------------------------------------------------------ NVIC and vectors
  vecslot/.style={draw=linecol, fill=white, anchor=south west, align=left,
                  minimum width=44mm, minimum height=5mm, inner xsep=3pt,
                  font=\sffamily\scriptsize},
  vecsys/.style={vecslot, fill=kercol},
  vecirq/.style={vecslot, fill=hwcol},
  vecfree/.style={vecslot, fill=black!6},
  vecnum/.style={font=\tiny\ttfamily, anchor=east, inner sep=2pt},
  prio/.style={draw=linecol, fill=fwcol, minimum width=26mm, minimum height=6mm,
               inner sep=2pt, align=center, font=\sffamily\scriptsize},
  prioaxis/.style={-Latex, draw=linecol, line width=0.7pt},
  preempt/.style={-Latex, draw=red!70!black, thick},
  tailchain/.style={-Latex, draw=linecol, dashed, thick},
  % ------------------------------------------------------------ DMA and cache
  master/.style={blk, fill=fwcol, text width=22mm, minimum height=7mm},
  busmatrix/.style={draw=linecol, fill=extcol, minimum width=72mm, minimum height=7mm,
                    align=center, font=\sffamily\small},
  memblk/.style={blk, fill=kercol, text width=24mm, minimum height=10mm},
  cache/.style={blk, fill=usrcol, text width=20mm, minimum height=8mm},
  cacheword/.style={draw=linecol, fill=white, anchor=south west, inner sep=0pt,
                    minimum width=8mm, minimum height=4.5mm, font=\sffamily\tiny},
  cwclean/.style={cacheword, fill=usrcol},
  cwdirty/.style={cacheword, fill=orange!35},
  cwstale/.style={cacheword, fill=red!25},
  dmaflow/.style={-Latex, thick, draw=blue!55!black},
  cpuflow/.style={-Latex, thick, draw=linecol},
  incoherent/.style={draw=red!70!black, line width=0.8pt, decorate,
                     decoration={zigzag, amplitude=0.6mm, segment length=2mm}},
  % ------------------------------------------------------------ timing diagram
  signal/.style={draw=linecol, line width=0.8pt, line cap=round, line join=round},
  signalb/.style={signal, draw=blue!55!black},
  signame/.style={font=\sffamily\scriptsize, anchor=east, inner sep=3pt},
  tick/.style={draw=linecol!45, dotted, line width=0.4pt},
  tspan/.style={Latex-Latex, draw=linecol, line width=0.5pt},
  tlabel/.style={font=\sffamily\tiny, fill=white, inner sep=1pt},
  sample/.style={circle, fill=red!70!black, inner sep=0pt, minimum size=1.4mm},
  hiz/.style={draw=linecol, dashed, line width=0.6pt},
  % ------------------------------------------------------------ RTOS schedule
  taskrun/.style={draw=linecol, fill=usrcol, anchor=south west, inner sep=1pt,
                  minimum height=5mm, align=center, font=\sffamily\tiny},
  taskready/.style={taskrun, fill=hwcol},
  taskblock/.style={taskrun, fill=black!8},
  taskisr/.style={taskrun, fill=fwcol},
  tasklane/.style={font=\sffamily\scriptsize, anchor=east, inner sep=3pt},
  taxis/.style={-Latex, draw=linecol, line width=0.6pt},
  release/.style={-Latex, draw=linecol, line width=0.6pt},
  deadline/.style={draw=red!70!black, dashed, line width=0.6pt},
  rtstate/.style={state, fill=usrcol, text width=18mm, minimum height=7mm,
                  font=\sffamily\scriptsize},
}

% ---------------------------------------------------------------- helper macros
% Lengths are plain numbers in millimetres.

% \memregion{<y>}{<height>}{<style>}{<name and size>}{<start address>}
\newcommand{\memregion}[5]{%
  \node[#3, minimum height=#2mm] at (0mm,#1mm) {#4};%
  \node[mmaddr] at (-0.5mm,#1mm) {#5};}
% \memgap{<y>}  a "not to scale" break across the column
\newcommand{\memgap}[1]{\draw[mmgap] (-2mm,#1mm) -- (36mm,#1mm);}

% Register bit fields. Bit \regwidth-1 is leftmost, as in every reference manual.
\newcommand{\regwidth}{32}
\newcommand{\regunit}{4}
% \bitfield{<hi>}{<lo>}{<style>}{<label>}
\newcommand{\bitfield}[4]{%
  \node[#3, minimum width={(#1-#2+1)*\regunit mm},
        text width={(#1-#2+1)*\regunit mm-2pt}]
    at ({(\regwidth-1-#1)*\regunit mm},0mm) {#4};}
\newcommand{\bitscale}{%
  \foreach \b in {0,...,\the\numexpr\regwidth-1\relax}
    \node[regbit] at ({(\regwidth-0.5-\b)*\regunit mm},6mm) {\b};}

% \vecentry{<row from bottom>}{<style>}{<exception number>}{<handler name>}
\newcommand{\vecentry}[4]{%
  \node[#2] at (0mm,{#1*5mm}) {#4};%
  \node[vecnum] at (-1mm,{#1*5mm+2.5mm}) {#3};}

% \cacheline{<x>}{<y>}{<style>}{<label>}
\newcommand{\cacheline}[4]{\node[#3] at (#1mm,#2mm) {#4};}

% \clocktrain{<y>}{<x start>}{<cycles>}{<period>}  50 percent duty, 4 mm high
\newcommand{\clocktrain}[4]{%
  \draw[signal] (#2mm,#1mm) \foreach \i in {1,...,#3}
      { -- ++(0,4mm) -- ++({0.5*#4mm},0) -- ++(0,-4mm) -- ++({0.5*#4mm},0) };}
% \busval{<y centre>}{<x>}{<width>}{<text>}  the usual bus-value lozenge
\newcommand{\busval}[4]{%
  \draw[signal] (#2mm,#1mm) -- ++(1mm,2mm) -- ++({#3mm-2mm},0) -- ++(1mm,-2mm)
                            -- ++(-1mm,-2mm) -- ++({-#3mm+2mm},0) -- cycle;%
  \node[font=\sffamily\tiny] at ({#2mm+0.5*#3mm},#1mm) {#4};}

% \slice{<lane y>}{<t start>}{<duration>}{<style>}{<label>}
\newcommand{\slice}[5]{\node[#4, minimum width=#3mm] at (#2mm,#1mm) {#5};}
% \lane{<lane y>}{<name>}
\newcommand{\lane}[2]{\node[tasklane] at (0mm,{#1mm+2.5mm}) {#2};}
% \deadlinemark{<t>}{<height>}
\newcommand{\deadlinemark}[2]{\draw[deadline] (#1mm,-1mm) -- (#1mm,#2mm);}

\usepackage{fancyhdr}
\usepackage[hidelinks,pdftitle={Twenty Firmware Projects on the NUCLEO-H7A3ZI-Q, from the Bench Inventory},pdfauthor={Joseph Ambrose Pagaran}]{hyperref}

% ---------------------------------------------------------------- unicode helpers
\DeclareUnicodeCharacter{2192}{\ensuremath{\rightarrow}}
\DeclareUnicodeCharacter{2190}{\ensuremath{\leftarrow}}
\DeclareUnicodeCharacter{2264}{\ensuremath{\leq}}
\DeclareUnicodeCharacter{2265}{\ensuremath{\geq}}
\DeclareUnicodeCharacter{2248}{\ensuremath{\approx}}
\DeclareUnicodeCharacter{03A9}{\ensuremath{\Omega}}
\DeclareUnicodeCharacter{2126}{\ensuremath{\Omega}}
\DeclareUnicodeCharacter{2026}{\ldots}
\DeclareUnicodeCharacter{00D7}{\ensuremath{\times}}
\DeclareUnicodeCharacter{00B1}{\ensuremath{\pm}}
\DeclareUnicodeCharacter{00B5}{\textmu{}}
\DeclareUnicodeCharacter{2013}{-}
\DeclareUnicodeCharacter{2014}{-}

% ---------------------------------------------------------------- layout
% Sources lists carry long addresses. Let them break at more places than the
% default, which is only the hyphen, or they run into the margin.
\setlength{\Urlmuskip}{0mu plus 1mu}
\def\UrlBreaks{\do\/\do\-\do\.\do\_\do\?\do\&\do\=\do\#\do\:\do\~}
\setlength{\emergencystretch}{3em}
\hyphenpenalty=800
\tolerance=1500
\setlength{\parskip}{4pt}
\setlength{\parindent}{0pt}
\captionsetup{font=small,labelfont=bf,skip=4pt}
\setlist{itemsep=2pt,topsep=3pt,parsep=0pt}
\pagestyle{fancy}
\fancyhf{}
\fancyhead[L]{\sffamily\small Twenty Firmware Projects on the NUCLEO-H7A3ZI-Q}
\fancyhead[R]{\sffamily\small\nouppercase{\benchmark}}
% \benchmark is a short running label: "Project NN" inside a project, the
% section name elsewhere. Long project titles would otherwise run into the
% left-hand header.
\newcommand{\benchmark}{\leftmark}
\fancyfoot[C]{\sffamily\small\thepage}
\renewcommand{\headrulewidth}{0.3pt}

% ---------------------------------------------------------------- listings
\lstset{basicstyle=\ttfamily\footnotesize,columns=fixed,basewidth=0.48em,keepspaces=true,
  breaklines=true,frame=single,rulecolor=\color{gray!60},backgroundcolor=\color{gray!6},
  xleftmargin=3pt,framexleftmargin=3pt,showstringspaces=false,tabsize=4,
  commentstyle=\color{green!40!black},keywordstyle=\color{blue!60!black}\bfseries,
  stringstyle=\color{red!50!black},extendedchars=false,aboveskip=6pt,belowskip=6pt}
\lstnewenvironment{asciiart}{\lstset{language={},frame=none,backgroundcolor=\color{white},
  breaklines=false,basicstyle=\ttfamily\scriptsize,keywordstyle=,commentstyle=,stringstyle=}}{}
\lstnewenvironment{ccode}{\lstset{language=C}}{}
\lstnewenvironment{shellcode}{\lstset{language=bash}}{}
\lstnewenvironment{pycode}{\lstset{language=Python}}{}
\lstnewenvironment{makecode}{\lstset{language=make}}{}
\lstnewenvironment{plaincode}{\lstset{language={}}}{}
\lstnewenvironment{dtscode}{\lstset{language={},morecomment=[l]{//},morecomment=[s]{/*}{*/}}}{}
\lstnewenvironment{yamlcode}{\lstset{language={},morecomment=[l]{\#}}}{}
\lstnewenvironment{cppcode}{\lstset{language=C++}}{}
\lstnewenvironment{rustcode}{\lstset{language={},morecomment=[l]{//},morecomment=[s]{/*}{*/},
  morestring=[b]",morekeywords={fn,let,mut,const,static,struct,enum,impl,trait,pub,use,mod,
  match,if,else,loop,while,for,in,return,unsafe,async,await,dyn,where,type,as,ref,move,Self}}}{}
\lstnewenvironment{asmcode}{\lstset{language=[x86masm]Assembler,morecomment=[l]{//},
  morecomment=[l]{;},morekeywords={ldr,str,ldrex,strex,dmb,dsb,isb,wfi,wfe,bx,blx,svc,
  cpsid,cpsie,msr,mrs,push,pop,bl,cbz,cbnz,it,ite}}}{}
\lstnewenvironment{ldcode}{\lstset{language={},morecomment=[s]{/*}{*/},
  morekeywords={MEMORY,SECTIONS,ENTRY,PROVIDE,KEEP,ALIGN,ORIGIN,LENGTH,AT,PHDRS}}}{}

% ---------------------------------------------------------------- authoring macros
% \project{NN}{Title}{Primary board}{Theme}
% \project{NN}{Title}{Target board}{Theme}
% The macro keeps its name because build.py's converter keys off it; it prints
% "Chapter". The label prefix sec:c must match DOC["label_pfx"] in build.py.
\newcommand{\project}[4]{\clearpage\section[#2]{#2}\label{sec:c#1}%
  \renewcommand{\benchmark}{Chapter #1}%
  {\color{black!60}\sffamily\small Target board: #3 \quad|\quad Theme: #4}\par\medskip}
% key facts box: \begin{keyfacts} \item[Boards] ... \end{keyfacts}
\definecolor{volaccent}{HTML}{7A3B8F}
\newenvironment{keyfacts}{\begin{tcolorbox}[colback=volaccent!5,colframe=volaccent!80!black,
  title=Key facts,fonttitle=\sffamily\bfseries,left=2pt,right=2pt,top=2pt,bottom=2pt]
  \begin{description}[style=sameline,leftmargin=34mm,labelwidth=32mm,font=\sffamily\bfseries\small,itemsep=1pt]}
  {\end{description}\end{tcolorbox}}
% note box: \begin{note}{Title} ... \end{note}
\newtcolorbox{note}[1]{colback=yellow!8,colframe=orange!75!black,title=#1,fonttitle=\sffamily\bfseries,
  left=3pt,right=3pt,top=2pt,bottom=2pt}
% numbered steps
\newenvironment{steps}{\begin{enumerate}[label=\textbf{Step \arabic*.},leftmargin=*,itemsep=5pt]}{\end{enumerate}}
% \diagram{figure-name}{Caption}: inputs figures/<name>.tex, scaled down only when too wide
\newsavebox{\diagbox}
\newcommand{\diagram}[2]{\begin{figure}[H]\centering
  \sbox{\diagbox}{\input{figures/#1.tex}}%
  \ifdim\wd\diagbox>\linewidth\resizebox{\linewidth}{!}{\usebox{\diagbox}}\else\usebox{\diagbox}\fi
  \caption{#2}\label{fig:#1}\end{figure}}

% \pubdate{April 2010}: a publication date that the publisher gives to month
% precision only. The house rule is that every date this volume states carries
% weekday, day, month and year, and that rule cannot apply to a day a publisher
% never published. Writing one would be inventing a fact, which is the one thing
% this volume refuses to do. crosscheck.py recognises this macro and exempts it,
% so every month-and-year that is NOT wrapped is still reported as a defect.
\newcommand{\pubdate}[1]{#1}

\title{\sffamily\bfseries Twenty Firmware Projects on the NUCLEO-H7A3ZI-Q\\[2pt]\large from the Bench Inventory}
\author{Joseph Ambrose Pagaran}
\date{Sunday 20 September 2026}

\begin{document}
\maketitle
\thispagestyle{empty}
\section*{About this document}
\addcontentsline{toc}{section}{About this document}

This is a build plan for twenty firmware projects on one microcontroller board that is already on the bench, with three sensor shields, one power instrument, one data acquisition hat and a drawer of cables. Nothing here needs to be bought. Each chapter is self-contained and has the same shape: why it exists, the prior art and what is taken from it, the parts it uses, a system architecture figure, the peripheral configuration, the wiring, a memory and timing budget, a software design in UML, a data-flow sketch, a repository layout, numbered steps with real commands and code, measurable acceptance criteria, the variants it touches, pitfalls, best practices, stretch goals, a sourced roadmap, the evidence to publish, and its sources.

The twenty were chosen so that no two teach the same thing, and so that together they cover what a firmware engineer is expected to have done: a toolchain and a linker script owned end to end, buffering and interrupts, transfers by engine, framing and checksums, timing proved against an external instrument, low-power modes with a charge number, an application state machine, a codec with a host twin, energy accounted per phase, a command engine that never blocks, a test rig, three sensor shields with their bus topologies, the placement question for a classifier, the argument between a vendor layer and the registers, transforms with a numerical acceptance test, the cache and protection-unit trap this family is known for, and a real-time kernel with a measured latency distribution.

\begin{note}{What is deliberately left out}
The STWIN.box is on this bench and appears in no chapter, because it is the committed subject of a separate condition-monitoring build and splitting that subject across two documents would serve neither. The Linux side of the same cable belongs to the sibling volume of twenty embedded Linux projects, whose Project 12 owns the service running on the Raspberry Pi; this volume owns the firmware side and does not repeat it. The kit-lab workbook puts this same board on the bench four times but always flashes a vendor sample in order to teach the Linux host, so its text is not repeated either, although its address map and its one-shield-at-a-time rule are authoritative here and are reused.
\end{note}

Where the chapters came from is worth one sentence, because it explains the numbering. Chapters 1 to 7 are the seven exercises of a workbook written for an interview on Wednesday 2 September 2026. Chapters 8 to 18 come from a scored project list, and appendix A prints every entry of that list which did not become a chapter, keeping its original number and score and saying which chapter absorbed it or why it stayed out. Chapters 19 and 20 are new: the cache and protection-unit chapter because this family's worst trap deserves one, and the real-time kernel chapter because a job-fit note records the absence of a shipped kernel application as the one standing gap in the portfolio.

\section*{The trap that runs through this book}
\addcontentsline{toc}{section}{The trap that runs through this book}

\begin{note}{This part is not the STM32H7 the internet is about}
The STM32H7A3 is documented by reference manual RM0455. The popular member of the family, the STM32H743, is documented by RM0433. The clock tree, the power and regulator configuration and the memory map all differ. Most tutorials, blog posts, video series and example projects that say "STM32H7" were written against the other one. Code copied from them does not run slowly; it produces a board that does not boot.
\end{note}

Every well-engineered project encodes this distinction deliberately, and the ones that do not show exactly what it costs. The Rust hardware crate routes this part through its twin's peripheral module and selects the right manual, so the pin mapping is correct. MicroPython ships a linker script written for this part while borrowing the twin's alternate-function table. By contrast, the one board another language offers in this family hardcodes a different member's memory size in its only linker script for the family, and its own comments admit that two memory domains are unmapped and therefore skipped silently by cache maintenance.

The rule that follows is short. No linker script, clock configuration or memory map is inherited from a sibling part without being checked against RM0455 and the datasheet for this part. Any chapter that cites material written for the other member says so in the sentence that cites it. Where a fact has not been checked yet it is written as a question for the bench and not as an assertion, and the reader can tell the two apart at a glance.

\section*{The board, and what it does not have}
\addcontentsline{toc}{section}{The board, and what it does not have}

The NUCLEO-H7A3ZI-Q carries an STM32H7A3ZIT6Q: a Cortex-M7 at 280 MHz with 2 MB of flash, a switched-mode supply, and an on-board ST-LINK V3E that is a debug probe, a virtual COM port and a drag-and-drop disk at the same time. It has Zio and ST morpho connectors and an Arduino-compatible header that the three shields plug onto.

What it does not have matters more, because most of the published material assumes otherwise. There is no Ethernet: the upstream board description lists no network controller and no media access controller anywhere, so every tutorial that pairs this family with a network stack over Ethernet is for a different part. Any network path in this book goes through a coprocessor or a cellular modem. There is no on-board card socket, no external flash and no bus transceiver for the automotive bus. There is no cryptographic accelerator.

Four facts about the board are settled from two independent machine-readable sources that name this exact board, and are written as fact throughout: the 32.768 kHz crystal is fitted, so the real-time clock is usable and chapters 7, 10 and 20 rely on it; the high-speed clock is supplied by the debugger in bypass mode at 8 MHz, and 8 divided by 2, multiplied by 140, divided by 2 is the 280 MHz every chapter assumes; the user LEDs are LD1 green on PB0, LD2 yellow on PE1 and LD3 red on PB14; and the user button is on PC13. The virtual COM port is USART3, but the pins PD8 and PD9 are a convention of this board family rather than a checked fact, and are confirmed in the board manual before they are printed as one.

\section*{How to read and how to work}
\addcontentsline{toc}{section}{How to read and how to work}

\begin{itemize}
\item \textbf{Order.} Chapter 1 gates everything, because nothing else is meaningful without a build you trust. Chapters 2 to 7 are the workbook and prepare the node. Chapters 8 to 12 build the node, which chapter 20 then rebuilds on a kernel. Chapter 19 gates every chapter that moves a buffer with a transfer engine, which is 13, 15 and 18. Figure~\ref{fig:front_map} shows those dependencies and nothing else, because the rest of the order is a preference.
\item \textbf{Effort.} Each chapter's key-facts box gives a difficulty from 1 to 5 and an estimate in evenings of about four hours. The total is roughly eighty evenings. Pick by interest and by what the next application asks for.
\item \textbf{Evidence first.} Every chapter ends with the artefacts to publish: a repository whose README opens with a figure, a console log, a measurement table generated from data rather than typed, and a short note on what the measurement found. A chapter is finished when the evidence exists, not when the code runs once.
\item \textbf{Measurements carry their instrument.} Every number in this book names the instrument that produced it, or says "not measured". There are no invented microamps and no borrowed cycle counts.
\item \textbf{Licences before code.} Every dependency names its licence before a reader writes anything around it, because finding out afterwards is expensive.
\end{itemize}

\diagram{front_map}{The twenty chapters against the peripherals and subsystems they exercise, with the four real dependencies drawn and nothing else. The red arrows are the one gate that is easy to miss: any chapter that moves a buffer with a transfer engine depends on chapter 19.}

\section*{The bench inventory}
\addcontentsline{toc}{section}{The bench inventory}

\begin{table}[H]\centering\small
\begin{tabular}{p{42mm}p{74mm}p{34mm}}
\toprule
Item & What it is & Used in chapters \\
\midrule
NUCLEO-H7A3ZI-Q & STM32H7A3ZIT6Q, Cortex-M7 at 280 MHz, 2 MB flash, switched-mode supply, ST-LINK V3E & every chapter \\
X-NUCLEO-IKS4A1 & MEMS shield with three bus topologies selected by jumper & 13, 16, 20 \\
X-NUCLEO-IKS5A1 & Industrial MEMS shield: ISM6HG256X, ISM330IS with its in-sensor processor, IIS2DULPX, IIS2MDC, ILPS22QS & 14 \\
X-NUCLEO-53L8A1 & Eight by eight multizone time-of-flight shield & 15 \\
DFRobot SEN0032 & ADXL345 three-axis accelerometer breakout, I2C or SPI & 19 \\
MCC 118 DAQ HAT on a Raspberry Pi & 8 channels, 12 bit, 100 kS/s aggregate, the calibrated analogue reference & 6, 12 \\
nRF-PPK2 & Power Profiler Kit II, ammeter or source meter to 1 A, with eight digital inputs & 7, 10, 12 \\
Raspberry Pi 3, 3B+, 4 & Linux hosts for the edge side and for the test rig & 9, 12, 16 \\
SIM7020E, SIM7070G, SIM7600E-H & Cellular modems on their own supply, never fed from the board & 11 \\
Renkforce USB/TTL cable & 3.3 V serial console cable, red 5 V lead unconnected & several \\
Breadboard, jumpers, LEDs, resistors & Wiring and indicators. Every header on the bench is female & 19 and others \\
Windows 11 host with the Linux subsystem & Build host, and the Linux boards themselves & every chapter \\
\bottomrule
\end{tabular}
\caption{The inventory, and where each item is used. Small parts such as resistors, microSD cards and USB supplies are assumed to be in the drawer. The STWIN.box is on the bench and appears in no chapter, which is deliberate: it is the committed subject of a separate condition-monitoring build.}
\end{table}

\section*{The twenty chapters at a glance}
\addcontentsline{toc}{section}{The twenty chapters at a glance}

\begin{table}[H]\centering\small
\begin{tabular}{p{6mm}p{74mm}p{34mm}p{12mm}p{16mm}}
\toprule
No. & Chapter & Target & Level & Evenings \\
\midrule
1 & The toolchain, first light, and printf over the ST-LINK & Nucleo alone & 3 & 3 \\
2 & A single producer, single consumer ring buffer & Nucleo alone & 3 & 3 \\
3 & Receiving on interrupt without losing bytes & Nucleo alone & 3 & 4 \\
4 & Circular DMA and the idle line & Nucleo alone & 4 & 4 \\
5 & Framing and the hardware CRC unit & Nucleo alone & 3 & 4 \\
6 & Sampling on a timer at exactly 1 kHz & Nucleo + MCC 118 & 4 & 4 \\
7 & Stop mode, RTC wake, and a battery number & Nucleo + PPK2 & 4 & 4 \\
8 & The node's state machine, with a transmit that is a stub & Nucleo alone & 3 & 4 \\
9 & The payload codec and its Python twin & Nucleo + host & 3 & 4 \\
10 & Where the energy goes: wake, sense, compute, send & Nucleo + PPK2 & 4 & 4 \\
11 & An AT engine that never blocks & Nucleo + SIM7020E & 4 & 5 \\
12 & Energy as a regression test, and the rig that runs it & Nucleo + Pi + PPK2 & 5 & 5 \\
13 & The IKS4A1: FIFO, watermark, interrupt & Nucleo + IKS4A1 & 3 & 4 \\
14 & The IKS5A1: low and high g at once, wide pressure & Nucleo + IKS5A1 & 4 & 4 \\
15 & Eight by eight time of flight, decided on the MCU & Nucleo + 53L8A1 & 4 & 4 \\
16 & Where the classifier runs: sensor, MCU or host & Nucleo + IKS4A1 + Pi & 5 & 5 \\
17 & The same driver twice: vendor layer against registers & Nucleo + a shield & 4 & 4 \\
18 & Transforms and filters with a numerical acceptance test & Nucleo alone & 4 & 4 \\
19 & Caches, the MPU, and why DMA reads stale bytes & Nucleo + ADXL345 & 5 & 4 \\
20 & Three jobs at once: the node as an RTOS application & Nucleo + IKS4A1 & 5 & 6 \\
\bottomrule
\end{tabular}
\caption{Overview. Level is 1 for routine and 5 for research-grade; evenings are about four hours each and are estimates rather than promises. Where this table and a chapter's key-facts box differ, the key-facts box is authoritative.}
\end{table}

\section*{Variants, and how to read along an axis}
\addcontentsline{toc}{section}{Variants, and how to read along an axis}

The same exercise solved several ways is the most useful thing a book like this can offer, and twenty chapters multiplied by seven axes is not a book. So each variant is built in full exactly once, in the chapter where it teaches most, and is referenced from everywhere else it applies. Appendix F prints the whole matrix, and every chapter's variants table names the chapter where each row is built in full.

The seven axes are: the execution model, from a polled loop through interrupts and transfers to a kernel task; synchronisation, from a volatile flag through semaphores, mutexes, queues, notifications, event groups and stream buffers; time and safety, covering the tick, the general and low-power timers, the real-time clock, tickless idle and the two watchdogs; peripheral substitution, meaning the same job done on different silicon; the operating system, bare metal against two kernels; the language, with C as the baseline and C++, Rust and MicroPython on the target; and intelligence and reach, from local only to a classifier in the sensor, on the microcontroller or on a host.

To read along an axis rather than front to back, start at the chapter that builds the variant in full and follow its variants table outward. The language axis, for example, starts at chapter 1 for the five-minute interpreter path, goes to chapter 9 for C++ and the host twin, and ends at chapter 17 for Rust. The synchronisation axis is almost entirely chapter 20.

\section*{The instruments, and what the bench does not have}
\addcontentsline{toc}{section}{The instruments, and what the bench does not have}

There is no oscilloscope, no logic analyser and no Joulescope on this bench, confirmed on Sunday 20 September 2026. There is also no soldering iron, so nothing that needs a solder bridge moved is proposed anywhere in this book. Every timing and energy claim therefore comes from one of three instruments: the nRF-PPK2, whose current trace and eight digital inputs are time-aligned and which is this book's logic recorder for slow signals; the MCC 118 on a Raspberry Pi at 100 kS/s aggregate, which is the calibrated analogue reference; and the board's own timers and cycle counter. A chapter that would be easier with an oscilloscope says so in one sentence and names the substitute.

\section*{Conventions and assumptions}
\addcontentsline{toc}{section}{Conventions and assumptions}

\begin{itemize}
\item \textbf{Host.} A Windows 11 machine with the Linux subsystem is the build host, and the Raspberry Pi boards are the second host for the test rig.
\item \textbf{Console.} 115200 baud, 8N1, no flow control, over the probe's virtual COM port. The maximum reliable rate on that port is confirmed in chapter 3 rather than assumed.
\item \textbf{Logic levels.} 3.3 V only. Nothing on this bench puts 5 V into a processor pin. The cellular modems need peaks of about 2 A from their own supply and are never fed from the board.
\item \textbf{Pin numbers.} Port and pin as the board manual gives them, with the Arduino header name alongside where a shield is involved. Anything taken from the Nucleo-144 convention rather than the manual is marked as convention until it is checked.
\item \textbf{Code.} C is the baseline, Python is the host language for design, analysis and test harnesses, and shell is the glue. Every snippet is a starting point rather than a tested build; the acceptance criteria are what count.
\item \textbf{Generated code.} Nothing generated by a configuration tool is presented as reviewed source. Where a chapter uses a generator it names the peripherals to configure and the constraints they satisfy, then shows the code written on top.
\item \textbf{Naming.} Repositories are lower case with hyphens and begin with the board name, one repository per chapter, each with a licence file and a README that opens with the architecture figure.
\end{itemize}

\section*{How this document is built}
\addcontentsline{toc}{section}{How this document is built}

The source is LaTeX with one file per chapter under \texttt{sections/} and one TikZ or circuitikz file per figure under \texttt{figures/}. \texttt{build.py} runs pdflatex for the PDF, renders every figure through latex and dvisvgm to SVG, and converts the same LaTeX subset to a single self-contained HTML file with the figures embedded. \texttt{python build.py -\/-check sections/c19.tex} compiles one chapter alone for quick iteration, and \texttt{python lint.py} enforces the house rules a compiler cannot see. The authoring contract is in \texttt{AUTHORING.md} and the prior-art pool with its licence categories is in \texttt{SOURCE.md}.

\clearpage
\tableofcontents
\clearpage

\part{The workbook}
\project{01}{The toolchain, first light, and printf over the ST-LINK}{NUCLEO-H7A3ZI-Q}{Cross toolchain, startup, linker script, flashing}

\begin{keyfacts}
\item[Board] NUCLEO-H7A3ZI-Q, and nothing else
\item[Peripherals] RCC and PWR for the clock tree, GPIO for three LEDs, USART3 for the console
\item[Toolchain] arm-none-eabi-gcc with CMake, plus an open flashing tool. The vendor IDE is the alternative, not the baseline
\item[Operating system] Bare metal
\item[Difficulty] 3 of 5
\item[Effort] 3 evenings of about four hours
\item[Deliverable] A repository whose every byte you can account for: hand-written vector table, startup code and linker script, a build that runs from a clean clone in two commands, a blinking LED and a line of text on the console
\end{keyfacts}

\subsection*{Why this project}

Every later chapter needs a build you trust and a way to see what the board is
doing. This chapter builds both, and does it without a code generator, because
the one thing worth knowing before anything else is what is actually in the
image. A generated project is several thousand lines pinned to one version of
one vendor's tool. It compiles, and it teaches nothing about the first
instruction the processor executes.

The deliverable is small on purpose: one blinking light and one line of text.
What makes it worth three evenings is that you will have written the vector
table, the startup code that copies initialised data into memory and zeroes the
rest, the linker script that decides where each of those lives, and the clock
configuration that takes the part from its reset state to full speed. Chapter 2
onwards assumes all of that exists and is understood.

There is a second reason, specific to this board, and it is the reason this
volume exists at all. Most published material for this family was written for a
different member of it. Starting from a hand-written, checked configuration is
the only way to be sure you are not inheriting somebody else's part.

\begin{note}{This part is not the STM32H7 the internet is about}
The STM32H7A3 is documented by reference manual RM0455. The popular member of
the family, the STM32H743, is documented by RM0433, and its clock tree, power
configuration and memory map all differ. Code copied from an H743 tutorial does
not run slowly here; it produces a board that does not boot. Every linker
script, clock configuration and memory map in this book is checked against
RM0455, and every chapter that cites material written for the other part says so
in the sentence that cites it.
\end{note}

\subsection*{Prior art and what to reuse}

\begin{table}[H]\centering\small
\begin{tabular}{p{34mm}p{46mm}p{38mm}p{20mm}}
\toprule
Source & What it gives & What it does not & Licence \\
\midrule
Lyubka's bare-metal programming guide & The clearest open walkthrough of vector table, startup code and linker script, building up to a working clock tree & Its board templates stop at smaller cores. The clock tree, the caches and the memory regions of this part are yours & MIT \\
Oblaukhov's CMake package layer & A mature toolchain file and target handling that already knows this silicon family & It pulls in the vendor pack, so it is not the hand-written path this chapter takes & MIT \\
Two register-level projects for this exact board & One brings the part to 280 MHz with correct flash wait states; the other builds with a modern toolchain and no assembly file & Neither states a licence, so both are read and neither is copied & None stated \\
The MicroPython board definition & A working, checked clock configuration for this exact part, in a file you can read in a minute & It is a whole interpreter, not a starting point for C & MIT \\
The vendor's firmware package & Reference startup and linker scripts for this part, and a basic example tree & Its middleware sits under a licence tied to this vendor's silicon & Mixed \\
\bottomrule
\end{tabular}
\caption{Prior art for chapter 1. The two unlicensed repositories are read as worked examples and re-derived, because a repository with no licence file is not open source.}
\end{table}

What is left to write is the whole of it, and that is the point. The sources
above settle questions rather than supply code: what the startup sequence must
do, which multiplier and dividers reach 280 MHz, and which pins the three LEDs
are on. The code is yours, which is what makes the repository portfolio evidence
rather than a fork.

\subsection*{Parts from the inventory}

\begin{table}[H]\centering\small
\begin{tabular}{p{40mm}p{66mm}p{40mm}}
\toprule
Part & Role & Interface \\
\midrule
NUCLEO-H7A3ZI-Q & The whole chapter. Its on-board probe is also the programmer, the serial port and a drag-and-drop disk & Micro USB to the host \\
A USB data cable & Power, programming and console on one lead & Micro USB \\
Host PC & Build, flash and terminal & Windows with the Linux subsystem, or a Raspberry Pi \\
\bottomrule
\end{tabular}
\caption{Inventory items used in chapter 1. Nothing is wired and nothing is bought.}
\end{table}

\subsection*{System architecture}

\diagram{c01_arch}{Everything in this chapter happens on one board and one cable. The probe on the left half of the board programmes the microcontroller on the right half and carries its console back to the host.}

The probe is a separate microcontroller on the same board. It presents three
things to the host at once: a debug interface for programming and halting, a
serial port bridged to USART3 on the target, and a small disk you can copy a
binary onto. This chapter uses the first two and mentions the third, because
knowing all three exist saves an hour the first time one of them misbehaves.

\subsection*{Peripheral configuration}

\begin{table}[H]\centering\small
\begin{tabular}{p{22mm}p{26mm}p{24mm}p{30mm}p{30mm}}
\toprule
Peripheral & Mode & Clock source & Pins and function & Interrupt and transfers \\
\midrule
RCC and PWR & Bypass input, then phase-locked loop & 8 MHz from the probe & None & None \\
GPIOB, GPIOE & Push-pull output & Peripheral bus & PB0, PE1, PB14 & None \\
USART3 & Asynchronous, 115200 8N1 & Peripheral bus & Confirm in the board manual & Polled in this chapter \\
GPIOC & Input & Peripheral bus & PC13, user button & None in this chapter \\
\bottomrule
\end{tabular}
\caption{Peripheral configuration. The LED and button pins are confirmed from two independent sources; the console pins are the usual convention for this board family and are checked in the board manual before you trust them.}
\end{table}

\subsection*{Wiring}

\diagram{c01_wiring}{Nothing is wired in this chapter. The three LEDs and the button are on the board, and their pins are confirmed.}

\subsection*{Memory and timing budget}

\diagram{c01_mem}{The clock tree this chapter configures. The multiplier and dividers are taken from a working configuration for this exact part, not from a tutorial written for its sibling.}

\begin{table}[H]\centering\small
\begin{tabular}{p{44mm}p{28mm}p{28mm}p{30mm}}
\toprule
Quantity & Budget & Measured & Margin \\
\midrule
Flash used & 12 kB, revised from 8 kB & 11\,012 B & +1\,276 B, 10.4 per cent \\
Static memory used & 4 kB & 484 B & +3\,612 B, 88 per cent \\
Reset to first printed line & 100 ms & 2.489 ms & +97.5 ms, 97.5 per cent \\
\bottomrule
\end{tabular}
\caption{The budget table, which every chapter carries. Measured on the board at commit 8ecd613 on Saturday 3 October 2026, with \texttt{arm-none-eabi-size -A} for the first two rows and the part's own cycle counter for the third.}
\end{table}

\textbf{The flash budget was revised upward and the original is kept in the table,
because a budget quietly replaced by whatever was built is not a budget.} 8 kB was
chosen before anything had ever been compiled. Measured, \texttt{printf} and its
dependencies from newlib-nano are 3\,920 B on their own, which is 48 per cent of
that original figure, in a chapter whose title is ``the toolchain, first light, and
printf over the ST-LINK''. The budget was set without measuring the thing the
chapter is named after, which is exactly the mistake the article cited in this
chapter's prior art table exists to warn about.

Where the 11\,012 B goes, attributed from the link map rather than estimated:

\begin{table}[H]\centering\small
\begin{tabular}{p{74mm}p{24mm}p{24mm}}
\toprule
Source & Bytes & Share \\
\midrule
newlib-nano \texttt{libc}, almost all of it \texttt{printf} & 3\,920 & 36 per cent \\
this repository's own code & 4\,557 & 41 per cent \\
string literals and the remainder & 2\,535 & 23 per cent \\
\bottomrule
\end{tabular}
\caption{Every byte of the image attributed to where it came from, read out of the
link map. The compiler is not asked to estimate and neither is the author.}
\end{table}

The alternative to revising was to cut the explanatory text this project prints
about its own provenance. That saves at most the 2\,580 B of \texttt{.rodata} and
would land at 8\,432 B, still over the original 8 kB, so it would not even achieve
the budget while making the volume worse at what it is for. On a part with 2 MB of
flash the image is 0.53 per cent of what is available, so the budget's purpose here
is to notice unintended growth rather than to fit, and 12 kB with 10 per cent
headroom serves that.

\textbf{Static memory is 484 B, which is \texttt{.data} plus \texttt{.bss} and is
what static conventionally means.} It is reported separately from the 9\,220 B that
the linker script reserves for the stack and heap, because that is a reservation
rather than a set of variables: 8 kB of stack and 1 kB of heap, named in
\texttt{c/ld/stm32h7a3zi.ld}. All of it together is 9\,704 B, 7.40 per cent of the
128 kB of DTCM. Collapsing the two would have shown the 4 kB budget exceeded by a
factor of two for no real reason.

\textbf{The last row was expected to wait for chapter 6 and did not.} The caption
used to say so. Measuring the time from reset to the first printed character needs a
clock running before anything else happens, and the Cortex-M7 has one:
\texttt{DWT\_CYCCNT} is architectural, defined by ARM for ARMv7-M, so RM0455 has
nothing to say about it and it was available in this chapter all along. It is
started at the top of \texttt{Reset\_Handler}, before \texttt{.data} is copied,
which is safe only because those are peripheral registers rather than RAM.

So the figure excludes only the hardware's own reset sequence and the vector fetch.
Of the 2.489 ms, about 2.4 ms is the delay calibration deliberately spinning 8\,000
iterations to measure itself, so this row is mostly the cost of a measurement rather
than of startup. An earlier version spent 11.989 ms there, which was 99 per cent
calibration, and the probe was reduced once that was visible.

\subsection*{Firmware design (UML)}

\diagram{c01_uml}{From power-on to the first printed line. Everything before main is code written in this chapter.}

The processor does very little for you at reset. It reads two words from the
start of flash, the initial stack pointer and the address of the first
instruction, and jumps there. Everything a C programmer takes for granted,
namely that initialised globals hold their values and uninitialised ones are
zero, happens because the startup code copies one region and clears another
before calling main. The linker script is what tells the startup code where
those regions are, which is why the two are written together and neither makes
sense alone.

\subsection*{Data flow (ASCII)}

\begin{asciiart}
  host                                   board
  +------------------------------+       +--------------------------------+
  | cmake --build build          |       | reset: MSP and PC from flash   |
  |   |                          |       |   |                            |
  |   v                          |  SWD  |   v                            |
  | firmware.elf ----------------+-----> | startup: copy .data, zero .bss |
  |                              |       |   |                            |
  | terminal at 115200           |       |   v                            |
  |   ^                          | USART | clocks: 8 MHz -> 280 MHz       |
  |   |                          |  3    |   |                            |
  |   +--------------------------+-------+ main: blink LD1, print a line  |
  +------------------------------+       +--------------------------------+
\end{asciiart}

\subsection*{Repository layout}

\begin{plaincode}
nucleo-h7a3-firstlight/
  CMakeLists.txt
  cmake/arm-none-eabi.cmake     # the toolchain file, not the project file
  ld/stm32h7a3zi.ld             # written here, checked against RM0455
  src/startup.c                 # vector table and reset handler, in C
  src/system.c                  # clock tree, flash latency, voltage scaling
  src/uart.c                    # USART3 at 115200, polled
  src/retarget.c                # printf to USART3
  src/main.c
  tools/reconcile_size.py       # sums the map file against the size output
  README.md                     # starts with the architecture figure
\end{plaincode}

\subsection*{Steps}

\begin{steps}
\item \textbf{Install the toolchain and prove it runs.} You need a cross
compiler, CMake and a flashing tool. Nothing here needs the vendor IDE.
\begin{shellcode}
arm-none-eabi-gcc --version
cmake --version
probe-rs --version          # or: openocd --version
\end{shellcode}
The flashing tool knows this part by name. Confirm it before going further,
because a tool that does not know your die will fail in a way that looks like a
hardware fault.
\begin{shellcode}
probe-rs chip list | grep -i H7A3
\end{shellcode}

\item \textbf{Write the linker script.} Two regions to begin with. The sizes
come from the datasheet, not from a sibling part.
\begin{ldcode}
MEMORY
{
  FLASH (rx) : ORIGIN = 0x08000000, LENGTH = 2048K
  RAM  (rwx) : ORIGIN = 0x24000000, LENGTH = 1024K   /* AXI SRAM, confirm */
}
ENTRY(Reset_Handler)
SECTIONS
{
  .isr_vector : { KEEP(*(.isr_vector)) } > FLASH
  .text   : { *(.text*) *(.rodata*) } > FLASH
  _sidata = LOADADDR(.data);
  .data   : { _sdata = .; *(.data*) . = ALIGN(4); _edata = .; } > RAM AT> FLASH
  .bss    : { _sbss = .; *(.bss*) *(COMMON) . = ALIGN(4); _ebss = .; } > RAM
  . = ALIGN(8);
  _estack = ORIGIN(RAM) + LENGTH(RAM);
}
\end{ldcode}
The three symbol pairs are the contract with the startup code. Nothing else in
the file matters yet.

\item \textbf{Write the startup code.} It can be C. Assembly is traditional, not
required.
\begin{ccode}
#include <stdint.h>

extern uint32_t _sidata, _sdata, _edata, _sbss, _ebss, _estack;
extern int main(void);
void SystemInit(void);

void Reset_Handler(void)
{
    uint32_t *src = &_sidata, *dst = &_sdata;
    while (dst < &_edata) *dst++ = *src++;      /* initialised globals */
    for (dst = &_sbss; dst < &_ebss; ) *dst++ = 0;   /* the rest */

    SystemInit();
    main();
    for (;;) { }                                 /* main must not return */
}

static void Default_Handler(void) { for (;;) { } }

__attribute__((section(".isr_vector"), used))
void (* const vector_table[])(void) = {
    (void (*)(void)) &_estack,
    Reset_Handler,
    Default_Handler,        /* NMI */
    Default_Handler,        /* HardFault, replaced properly in chapter 3 */
};
\end{ccode}

\item \textbf{Bring the clock up to 280 MHz.} The input is 8 MHz from the probe,
in bypass mode because it is a clock signal and not a crystal. Divide by two,
multiply by 140, divide by two. Before raising the frequency, raise the core
voltage and add flash wait states, in that order. Doing it in the wrong order is
the classic way to produce a board that runs at reset speed or not at all.

\item \textbf{Blink one LED.} LD1 is on PB0. Enable the port clock first; a pin
configured before its clock is running does nothing and reports no error.

\item \textbf{Bring up the console and retarget printf.} USART3 at 115200. Write
the character output hook so the standard library finds it.
\begin{ccode}
int _write(int fd, const char *buf, int len)
{
    (void) fd;
    for (int i = 0; i < len; i++) uart_putc(buf[i]);
    return len;
}
\end{ccode}
Link against the small standard library and ask for floating point support only
if you need it, because it is not free.
\begin{shellcode}
arm-none-eabi-gcc ... --specs=nano.specs --specs=nosys.specs
\end{shellcode}

\item \textbf{Account for every byte.} Compare the linker map against the size
output and make a script do it, so the check survives being forgotten.
\begin{shellcode}
arm-none-eabi-size -A build/firmware.elf
python tools/reconcile_size.py build/firmware.map build/firmware.elf
\end{shellcode}

\item \textbf{Prove the clean-clone build.} Delete the build directory, clone
into a fresh folder, and build and flash with two commands on a machine with no
vendor IDE installed. If that does not work, the repository is not evidence of
anything.
\begin{shellcode}
cmake -B build -DCMAKE_TOOLCHAIN_FILE=cmake/arm-none-eabi.cmake
cmake --build build -j && probe-rs run --chip STM32H7A3ZITx build/firmware.elf
\end{shellcode}
\end{steps}

\subsection*{Build, flash and debug}

\diagram{c01_timing}{What the linker placed and where, and the boot timeline whose intervals are still unknown. Do not guess them; chapter 6 builds the instrument that measures them.}

Three ways to get the image onto the board, in the order you should try them:
the flashing tool above, the open debug server, or copying the binary onto the
disk the probe presents. The last needs no tools at all and is the fastest way
to prove a board is alive.

\begin{shellcode}
openocd -f interface/stlink.cfg -f target/stm32h7x.cfg \
  -c "program build/firmware.elf verify reset exit"
\end{shellcode}

\begin{note}{When the board enumerates but says nothing}
In order of likelihood: a charge-only cable, which powers the board and never
appears as a device; the wrong serial device, because the probe presents more
than one and the wrong one is silent rather than an error; a host service that
probes new serial devices and corrupts the first frames; and a peripheral whose
clock was never enabled. If the first bytes after a reset are rubbish, the probe
buffered output from before the reset.
\end{note}

\subsection*{Verification and acceptance criteria}

\begin{itemize}
\item A clean clone builds and flashes in two commands on a machine with no
vendor IDE, proven on a second machine and not only on yours.
\item The size output and the map file reconcile to the byte, and a script
prints the reconciliation rather than you doing it by eye.
\item A breakpoint on the first instruction of the reset handler is reached
before any library code, and two sentinel variables show that initialised data
was copied and uninitialised data was zeroed.
\item The clock configuration is asserted at startup: a frequency above what the
selected voltage scale allows stops the board with a named error rather than
producing an intermittent fault later.
\item The console prints a line within a second of reset, at 115200, and
survives twenty scripted replug cycles.
\end{itemize}

\subsection*{Variants}

\begin{table}[H]\centering\small
\begin{tabular}{p{24mm}p{26mm}p{44mm}p{22mm}p{20mm}}
\toprule
Axis & Variant & What changes & Cost & Built in full in \\
\midrule
Language & MicroPython & Copy the published firmware for this exact board onto the probe's disk and you have a prompt in five minutes, with no build at all & Determinism and direct control of transfers & Here, as the five-minute path \\
Language & C++ & The same startup code, plus what you disable: exceptions cost roughly 10 to 30 kB of tables, which on 2 MB of flash is a choice rather than a necessity & A decision to justify & Chapter 9 \\
Language & Rust & The hardware crate names this part and routes it through its twin with the right manual selected; the async framework names this exact package and pin count & The vendor library is replaced entirely & Chapter 17 \\
Operating system & Bare metal & The baseline here & None & Here \\
Execution model & Polled loop & The baseline here. Everything is in main & Nothing runs while you wait & Chapter 3 \\
Peripheral & Probe disk instead of a debugger & Drag a binary onto a disk. No tools, no drivers & No debugging & Here, in the step above \\
\bottomrule
\end{tabular}
\caption{Variants for chapter 1. The interpreter path is worth trying first even if you never use it again, because it proves the board and the cable before you have written a line.}
\end{table}

\subsection*{Pitfalls}

\begin{itemize}
\item Inheriting a linker script from the better-known member of this family.
Its memory sizes and layout are different, and the failure looks like a hardware
fault.
\item Raising the clock before the voltage scale and the flash wait states.
\item Configuring a pin before enabling its port clock. No error is reported.
\item Believing a summarised read of a large generated file. One check of the
flashing tool's target list reported that this part was absent; searching the
raw file found the exact die. Search the file, not a summary of it.
\item Letting main return. The loop at the end of the reset handler is not
decoration.
\item Assuming the console pins from a tutorial. The peripheral is confirmed;
the pins are convention until you read the board manual.
\end{itemize}

\subsection*{Best practices applied}

\begin{itemize}
\item Nothing generated is presented as reviewed source. Every line in the image
was written or read by the author.
\item The build is reproducible from a clean clone, which is the only definition
of a build that means anything to someone else.
\item Facts about the board carry their evidence. Pins confirmed from two
independent sources are written plainly; pins taken from convention are marked
as such.
\item A measurement stays blank until it has been taken.
\end{itemize}

\subsection*{Stretch goals}

\begin{itemize}
\item Add a second build configuration that places the hot loop in
tightly coupled memory and compare the size output, which is the thread chapter
19 picks up.
\item Replace the polled console with the interrupt-driven one from chapter 3
and measure the difference in time spent waiting.
\item Build the same project with the modern LLVM-based toolchain. Note the
rename: the older embedded toolchain for this architecture ended at version
19.1.5 and is superseded.
\end{itemize}

\subsection*{Roadmap and next steps}

Chapter 2 gives the console somewhere to put bytes that arrive faster than they
are consumed, which is the first thing that breaks once the link is faster than
a human reading it.

For going deeper on what happens before main, the standard progression is the
series that starts at bare-metal C and works through the standard library and a
bootloader. For the build system, the vendor's own tool has generated this kind
of project since version 6.11 of its configurator, and the distinction it draws
between tool-owned and author-owned files is worth reading even if you never use
it. Appendix H lists the courses; the one that walks a single problem from a
polled loop to an event-driven architecture across fifty-six lessons is the best
free companion to this whole volume, though its code licence means it is
recommended and never quoted.

\subsection*{Portfolio evidence}

\begin{itemize}
\item A public repository whose README opens with the architecture figure and
states, in one sentence, that no line of the image was generated.
\item The reconciliation output showing the map file and the size output
agreeing to the byte.
\item A terminal recording of a clean clone building, flashing and printing, on
a machine that has never had the vendor IDE installed.
\end{itemize}

\subsection*{Sources}

Normative references:
\begin{itemize}
\item Reference manual RM0455, for this part. Not RM0433, which is its sibling.
\item The board user manual for MB1363, for the connectors, the jumpers and the
console pins.
\item The STM32H7A3xI datasheet, for memory sizes and the clock limits at each
voltage scale.
\end{itemize}

Reusable implementations:
\begin{itemize}
\item Lyubka's bare-metal programming guide, MIT.\newline
\url{https://github.com/cpq/bare-metal-programming-guide}
\item Oblaukhov's CMake package layer for this silicon family, MIT.\newline
\url{https://github.com/ObKo/stm32-cmake}
\item The flashing and debugging tool whose target list includes this die.\newline
\url{https://github.com/probe-rs/probe-rs}
\item Memfault's article on printf for embedded systems, which is the definitive
treatment of the retarget hook and its costs.\newline
\url{https://interrupt.memfault.com/blog/printf-on-embedded}
\item The MicroPython board definition for this exact board, MIT, which is the
shortest readable statement of a working clock configuration.\newline
\url{https://github.com/micropython/micropython}
\end{itemize}
   % Toolchain, first light, printf over the ST-LINK
\project{02}{A single producer, single consumer ring buffer}{NUCLEO-H7A3ZI-Q}{Lock-free buffering, memory ordering on the M7}

\begin{keyfacts}
\item[Board] NUCLEO-H7A3ZI-Q, and nothing else
\item[Peripherals] USART3 receive interrupt as the producer, the cycle counter for the cost measurement, one spare pin as a marker
\item[Toolchain] arm-none-eabi-gcc with CMake for the target, and the host compiler for the same two source files
\item[Operating system] Bare metal
\item[Difficulty] 3 of 5
\item[Effort] 2 evenings of about four hours
\item[Deliverable] One C file and its header that compile unchanged on the target and on the host, a test that pushes tens of millions of bytes through the structure off-target, a barrier you can justify line by line, and a counter that reports how many bytes were dropped
\end{keyfacts}

\subsection*{Why this project}

Chapter 1 left you with a console that prints. The moment anything arrives
faster than the code that reads it, you need somewhere to put bytes that is not
a single variable. That place is a ring buffer, and it is the single most reused
structure in this book: chapter 3 fills it from an interrupt, chapter 4 fills it
from a transfer engine, chapter 5 parses frames out of it, and chapter 11 runs a
command queue on top of it. Getting it right once is worth two evenings.

The structure itself is twenty lines. What takes the two evenings is the
argument for why those twenty lines are correct when one side of them runs in an
interrupt handler and the other runs in the main loop. That argument has three
parts, and only the first is widely understood. The first part is that a
power-of-two capacity turns a modulo into a mask. The second is that free running
indices, never wrapped and only masked at the point of use, give you the
occupancy for free and remove a whole class of off-by-one faults. The third part
is the ordering of the data write against the index publication, which is where
the word \texttt{volatile} is usually written and where it does not do the job
people think it does.

This chapter also builds the habit that makes every later chapter testable. The
buffer compiles for the host with no change, so the logic is exercised by a test
that runs in a second on a laptop rather than by flashing a board and watching a
terminal. What the host cannot test is named plainly rather than glossed over:
the host is strongly ordered, so it will never reproduce a reordering fault. The
ordering is argued from the architecture and from the published correctness
proof, not from a passing test.

\begin{note}{What volatile does and what it does not do}
Declaring the head and tail indices \texttt{volatile} stops the compiler holding
them in a register across the loop and stops it reordering them against each
other. It says nothing about the ordinary store to the data array that has to be
visible before the index that publishes it. A compiler is free to sink that
store past a volatile one, and the C standard gives you no ordering between a
volatile access and a non-volatile access. The fix is an explicit release on the
publishing store and an acquire on the consuming load, which is what the rest of
this chapter is about.
\end{note}

\subsection*{Prior art and what to reuse}

\begin{table}[H]\centering\small
\begin{tabular}{p{34mm}p{46mm}p{38mm}p{20mm}}
\toprule
Source & What it gives & What it does not & Licence \\
\midrule
Majerle's lightweight ring buffer & A mature, maintained implementation with zero-copy regions designed for handing a contiguous span to a transfer engine, which is exactly what chapter 4 needs & It assumes a coherent view of memory. With the data cache on and the buffer in main memory, the maintenance step is yours and is chapter 19 & MIT \\
Quantum Leaps' lock-free ring buffer & The smallest readable version, with its atomicity assumptions stated in the header rather than implied, and a test suite you can run & Single byte elements and no zero-copy path. It is a teaching reference more than a dependency & MIT \\
Le, Guatto, Cohen and Pop, SBAC-PAD 2013 & The correctness argument under the C11 memory model, with the precise placement of the acquire and the release, and the batching variants that reduce how often the two sides look at each other's index & It is a paper about queues on cache-coherent multiprocessors, not about an interrupt handler on one core. The mapping to this case is yours & Paper \\
Arm application note DAI0321A & What the barrier instructions actually order on this profile, and when a data synchronisation barrier is needed rather than a data memory barrier & It predates this core's revision and says nothing about the cache maintenance operations chapter 19 needs & Vendor note \\
\bottomrule
\end{tabular}
\caption{Prior art for chapter 2. Both libraries are permissively licensed and either could be a dependency. This chapter writes the structure anyway, because the ordering argument is the deliverable and you cannot make that argument about code you have not written.}
\end{table}

What is left to write is small and specific: the structure, the two access
functions, the placement of the two barriers, the overflow counter, and the host
test harness. If you would rather depend on a library in production, depend on
the first one in the table; it is maintained and it has the zero-copy path this
book needs by chapter 4. Write it once yourself first, because the alternative
is a project where the most load-bearing twenty lines are the only ones you
cannot explain.

\subsection*{Parts from the inventory}

\begin{table}[H]\centering\small
\begin{tabular}{p{40mm}p{66mm}p{40mm}}
\toprule
Part & Role & Interface \\
\midrule
NUCLEO-H7A3ZI-Q & Target build, the interrupt that fills the buffer, and the cycle counter that costs it & Micro USB to the host \\
A USB data cable & Power, programming and console on one lead & Micro USB \\
Host PC & The second build of the same source, and the test that runs in a second & Windows with the Linux subsystem, or a Raspberry Pi \\
\bottomrule
\end{tabular}
\caption{Inventory items used in chapter 2. Nothing is wired and nothing is bought. The host build is half the chapter.}
\end{table}

\subsection*{System architecture}

\diagram{c02_arch}{One producer, one consumer, one buffer. The two marked points are where the ordering is established: the producer releases the data before it publishes the head, and the consumer acquires the head before it reads the data.}

Only two pieces of code touch the structure, and they touch different halves of
it. The producer writes the data array and advances the head. The consumer reads
the data array and advances the tail. Neither writes the other's index. That
restriction is the whole reason the structure needs no lock, and it is also the
restriction that a second producer quietly removes: the moment a second
interrupt source calls the put function, everything in this chapter stops being
true. Chapter 20 is where that case is handled properly, with a queue owned by a
kernel.

\subsection*{Peripheral configuration}

\begin{table}[H]\centering\small
\begin{tabular}{p{22mm}p{26mm}p{24mm}p{30mm}p{30mm}}
\toprule
Peripheral & Mode & Clock source & Pins and function & Interrupt and transfers \\
\midrule
USART3 & Asynchronous, 115200 8N1 & Peripheral bus & Confirm in the board manual & Receive interrupt as the producer; the handler is chapter 3's subject \\
DWT cycle counter & Free running & Core clock & None & None. Confirm it runs with no debugger attached \\
GPIO, one Zio pin & Push-pull output & Peripheral bus & Chosen from the pins a shield leaves free & Marker for the power instrument \\
Memory & Buffer placement & Not applicable & None & DTCM first, main memory once a transfer engine is the producer \\
\bottomrule
\end{tabular}
\caption{Peripheral configuration. The buffer itself is not a peripheral, but where it lives is a configuration decision and it is the one that chapter 19 revisits. The cycle counter behaviour without a debugger is on the confirm list and is checked before the measurement in the last step is trusted.}
\end{table}

\subsection*{Wiring}

\diagram{c02_wiring}{Nothing is wired in this chapter. The only physical connection is the one cable from chapter 1, and the only optional addition is a marker pin for the power instrument, which chapter 10 uses and this chapter merely leaves room for.}

\subsection*{Memory and timing budget}

\diagram{c02_mem}{The buffer as a strip of slots. The upper strip is the ordinary case and the lower strip is the wrapped case, which is the same arithmetic and not a special case in the code. The indices are free running counts of bytes ever written and ever read; only the array subscript is masked.}

\begin{table}[H]\centering\small
\begin{tabular}{p{50mm}p{24mm}p{30mm}p{26mm}}
\toprule
Quantity & Budget & Measured & Margin \\
\midrule
Buffer capacity & 256 bytes & 256 bytes & exact by construction \\
Structure overhead & 16 bytes & 12 bytes & +4 bytes \\
Added cycles for one DMB in put and one in get & 10 & 7 & +3 \\
Added flash for the same & no budget set & 28 bytes & \\
Bytes dropped & 0 in one hour & 0 in 3.2 million & duration short, see below \\
Mismatched bytes delivered & 0 & 0 in 12.2 million & across all four modes \\
\bottomrule
\end{tabular}
\caption{Measured on the board on Saturday 3 October 2026 at 64 MHz, with the cycle counter for the cycles and \texttt{arm-none-eabi-size} for the flash. Two rows the budget asked for are missing on purpose and the reason is below.}
\end{table}

\textbf{Two rows the budget asked for are not filled, and dividing a measurement
in half to fill them would be inventing a number.} The budget asks for cycles per
put and cycles per get separately. What was measured is a put and a get as a pair,
inside a loop, including the loop control and two function calls that do not inline
across translation units. Splitting that into two halves would assume they cost the
same, which is the thing a measurement is supposed to settle. Isolating them needs
a further measurement that subtracts an empty loop of the same shape, and it has
not been made.

The absolute figure is therefore a harness figure: 141 cycles for a pair at mode 0.
All four modes carry the identical harness, so the harness itself cancels in a
comparison. What does NOT cancel is where each build happens to place the code, and
that is the subject of the hedge below. This paragraph said the differences were
clean before that was measured, and they are not.

\begin{table}[H]\centering\small
\begin{tabular}{p{8mm}p{52mm}p{20mm}p{20mm}p{22mm}}
\toprule
Mode & Ordering & Cycles & Flash & Mismatches \\
\midrule
0 & none, not shippable & 141 & 8\,552 & 0 in 3.0 M \\
1 & compiler only, a signal fence & 141 & 8\,560 & 0 in 3.2 M \\
2 & compiler and processor, one DMB & 148 & 8\,580 & 0 in 3.2 M \\
3 & acquire load and release store & 158 & 8\,604 & 0 in 2.8 M \\
\bottomrule
\end{tabular}
\caption{The four ordering choices measured on the part. Cycles are for a put and get pair including the harness; flash is the \texttt{text} section of the whole image. The producer runs in thread mode and the consumer in the SysTick handler.}
\end{table}

\textbf{The compiler barrier costs nothing in time and is not a no-op.} Modes 0 and
1 agree to the hundredth of a cycle, which is expected because a signal fence emits
no instruction. But the two images are not identical: mode 1 is 8 bytes larger, so
the compiler did generate different code. Something was being reordered, or could
have been, and the fence prevented it at no cost in time. That is the one
unambiguous recommendation this chapter can make from its own numbers: take mode 1,
because it is free.

\textbf{The DMB appears to cost 7 cycles per pair and definitely costs 28 bytes.
The cycle figure is not safe and the flash figure is.} Acquire and release appear
to cost 17 cycles and definitely cost 52 bytes.

The cycles need that hedge because of something measured a few hours later, on
chapter 9's encoder. Two builds of it differed only by a renamed variable and three
extra \texttt{printf} lines, and the encoder came out 7.8 per cent slower in the
second. Nothing about the encoder changed. The new string literals grew
\texttt{.rodata}, the code after them shifted, and with no instruction cache
enabled every instruction is fetched from flash, so what the loop costs depends on
where the linker put it.

An earlier version of this paragraph said the loop had landed differently against
32 byte flash fetch lines. The 32 bytes was invented. It is the right figure for a
Cortex-M7 cache line, which ARM fixes by the architecture and which chapters 4 and
19 cite correctly, and it was carried across to the flash interface's fetch width,
which is a different quantity, lives in RM0455, and has not been read. The
observation stands without it: code moved and the cost changed. The width of the
granule that makes cost depend on address is not established here.

The four builds in the table above are four separate images with four different
sizes, 8\,552 to 8\,604 bytes, so their code sits at four different places. A 5 per
cent difference between them cannot be attributed to a barrier when an unrelated
edit produced 7.8 per cent on another project the same day. The honest statement is
that the DMB costs something between nothing and about 15 cycles, and this
experiment cannot narrow it further.

One detail pushes the other way and is worth keeping. Modes 0 and 1 differ by 8
bytes of code and measured identically, to the hundredth of a cycle. So placement
does not always matter: it matters when a shift moves a loop across a fetch line
boundary and not otherwise. That is why the effect is erratic rather than
proportional, and why it cannot be corrected for by assuming a percentage.

What would separate the barrier from its placement is enabling the instruction
cache, and on Friday 2 October 2026 that turned out to be available immediately
rather than in chapter 19. ARM's \texttt{cachel1\_armv7.h} in the Cube pack holds
the enable sequence, the device header declares both caches present on this
Cortex-M7 r1p2, and every name in the sequence is ARMv7-M architectural, so RM0455
governs none of it. It is in \texttt{c/board/icache.c}.

That also makes a better experiment available than the one in the table above. The
cache is off at reset, so a single image can measure a function cold, enable the
cache, and measure the same function again at the same address in the same build.
Nothing has moved by a byte between the two figures, so the difference is the cache
alone, which no comparison of separate builds could establish. Chapter 9 now does
exactly that. The four modes here should be redone the same way, each image
reporting both states, and until they are, the flash column is the measured result
of this chapter and the cycle column is an indication.

\textbf{Nothing failed, in any mode, including the one with no barrier at all.}
Twelve point two million bytes passed through the ring across the four builds with
zero mismatched bytes. The reading that follows from this is narrower than it
looks, and the firmware prints the caveat rather than a tick.

It does not show the barriers are unnecessary. It is consistent with the
single-core argument, that a core sees its own stores in program order and that
taking an exception on the same core is a context-synchronising event, so the
producer's data write cannot be observed by its own handler after the index write
that followed it. It is equally consistent with the compiler simply not having
reordered anything that mattered on these three million bytes, and mode 1 being 8
bytes different says the compiler was doing something.

An absence over three million bytes is not proof of correctness. The case where a
barrier is expected to earn its 7 cycles is the one where the other observer is a
bus master rather than an interrupt, which is chapter 4 with a transfer engine and
chapter 19 with the cache. Neither is built, so this chapter reports what it
measured and leaves that question where it belongs.

\textbf{The drop row is short of its budget and says so.} The budget asks for a one
hour soak. Each mode ran about fourteen seconds, which is 3.2 million bytes and zero
drops, with the consumer comfortably ahead of the producer throughout. An hour is a
different claim and has not been made.

\subsection*{Firmware design (UML)}

\diagram{c02_uml}{The put and the get as a sequence. The three notes sit beside the lifelines rather than on them, and each marks a point where the order of two adjacent operations is the whole argument.}

The design has exactly one invariant and it is worth writing down in the header
above the structure: the consumer may read a slot only after the producer has
published a head that includes it, and the producer may write a slot only after
the consumer has published a tail that releases it. Everything else follows.
The occupancy is the difference of the two free running counts, which is correct
across the wrap of the counters themselves because unsigned arithmetic wraps in
a defined way. The buffer is full when that difference equals the capacity, and
empty when it is zero, which means no slot is sacrificed to distinguish the two
states and the capacity is the number you wrote and not one less.

\subsection*{Data flow (ASCII)}

\begin{asciiart}
  producer side                      shared state                consumer side
  (USART3 interrupt)                                             (main loop)

  read RDR ---------> byte
                        |
                        v
                   buf[head & MASK] = byte        buf[256]
                        |                         head: free running count in
                   release barrier                tail: free running count out
                        |                         drops: bytes refused
                        v
                   head = head + 1  ------------> head ---+
                                                          |
                                                    acquire barrier
                                                          |
                                                          v
                        +---------------------  used = head - tail
                        |                                 |
                   tail = tail + 1  <------------ tail <--+
                        ^                                 |
                        |                                 v
                        +--------------------- byte = buf[tail & MASK]
\end{asciiart}

\subsection*{Repository layout}

\begin{plaincode}
nucleo-h7a3-ring/
  CMakeLists.txt                # target build, arm-none-eabi
  cmake/arm-none-eabi.cmake     # reused unchanged from chapter 1
  src/ring.c                    # the whole structure, no target headers
  src/ring.h                    # the interface and the invariant, in a comment
  src/barrier.h                 # the three ordering choices, one macro each
  src/main.c                    # target: fill from USART3, drain in the loop
  test/Makefile                 # host build: cc -std=c11 -O2 -Wall -Wextra
  test/test_ring.c              # single-threaded property test
  test/test_ring_threads.c      # two-thread soak, with its limits documented
  tools/soak_report.py          # reads the console counters into a table
  README.md                     # opens with the architecture figure
\end{plaincode}

\subsection*{Steps}

\begin{steps}
\item \textbf{Write the interface before the implementation.} Two functions,
one structure, no dynamic allocation and no target headers. The header is where
the invariant is recorded, because a comment in the header is read and a comment
in the source is not.
\begin{ccode}
#ifndef RING_H
#define RING_H
#include <stddef.h>
#include <stdint.h>

/* Single producer, single consumer. Exactly one context may call ring_put
 * and exactly one other may call ring_get. Two producers are not supported
 * and the structure gives no diagnostic if you try.
 *
 * head and tail are free running counts of bytes ever written and ever read.
 * They are never wrapped; only the array subscript is masked. Unsigned
 * overflow is defined, so head - tail stays correct across the wrap.
 */
#define RING_SIZE 256u                    /* power of two, checked below */
#define RING_MASK (RING_SIZE - 1u)

typedef struct {
    uint8_t  buf[RING_SIZE];
    uint32_t head;                        /* written by the producer only */
    uint32_t tail;                        /* written by the consumer only */
    uint32_t drops;                       /* producer only: refused bytes  */
} ring_t;

void     ring_init(ring_t *r);
int      ring_put(ring_t *r, uint8_t b);  /* 1 accepted, 0 refused */
int      ring_get(ring_t *r, uint8_t *b); /* 1 delivered, 0 empty  */
uint32_t ring_used(const ring_t *r);
#endif
\end{ccode}
The capacity check belongs in the source, not in a comment, so that a later
edit to an awkward number stops the build rather than producing a structure
whose mask silently no longer matches its size.
\begin{ccode}
_Static_assert((RING_SIZE & RING_MASK) == 0u, "RING_SIZE must be a power of two");
_Static_assert(RING_SIZE <= 0x80000000u, "capacity must stay under half the counter");
\end{ccode}

\item \textbf{Understand why the mask is not an optimisation.} A modulo by a
runtime value is a division, and on this core a division is not single cycle. A
modulo by a power of two known at compile time is a mask, which is one
instruction. That much is ordinary. The reason the power of two is structural
rather than an optimisation is the free running index: if the capacity were not
a power of two, the counters would have to be wrapped explicitly at the
capacity, and then the difference between them would no longer be the
occupancy, and you would be back to sacrificing a slot to tell full from empty.
The mask and the free running counters are one decision, not two.

\item \textbf{Put the barrier where the argument says it goes.} Write the data,
then release, then publish the head. Read the head, then acquire, then read the
data. Three implementations of that pair of barriers are worth having in one
header so you can choose with your eyes open.
\begin{ccode}
/* barrier.h: three choices, in rising order of strength and cost. */

/* 1. Compiler only. Correct when the counterpart is code on this same core,
 *    because exception entry and return are context synchronising: the
 *    handler cannot observe a partially retired program order. It is not
 *    correct when the counterpart is a transfer engine or another master. */
#define RING_RELEASE()  __atomic_signal_fence(__ATOMIC_RELEASE)
#define RING_ACQUIRE()  __atomic_signal_fence(__ATOMIC_ACQUIRE)

/* 2. Compiler and processor. One DMB. Correct in every case in this book,
 *    including a transfer engine as the counterpart, and this is the
 *    default. Cost is a handful of cycles, measured in the last step. */
/* #define RING_RELEASE()  __atomic_thread_fence(__ATOMIC_RELEASE) */
/* #define RING_ACQUIRE()  __atomic_thread_fence(__ATOMIC_ACQUIRE) */

/* 3. The index accesses themselves carry the ordering. Clearest to read,
 *    and what the C11 correctness argument in the paper is written against.
 *    Requires head and tail to be plain uint32_t, not volatile. */
/* store: __atomic_store_n(&r->head, h + 1u, __ATOMIC_RELEASE);            */
/* load:  uint32_t t = __atomic_load_n(&r->tail, __ATOMIC_ACQUIRE);        */
\end{ccode}
Note what is not in that list: \texttt{volatile} on its own. It appears in most
published versions of this structure and it is not sufficient, for the reason in
the box above. Note also what the third choice requires. If you mark the indices
\texttt{volatile} and then also use the atomic builtins on them, you have
written something the compiler is allowed to treat as undefined; pick one.

\item \textbf{Write the two functions.} They are short because the decisions
were made before the typing started.
\begin{ccode}
#include "ring.h"
#include "barrier.h"

void ring_init(ring_t *r)
{
    r->head = r->tail = r->drops = 0u;
}

/* Producer context only. Returns 0 and counts the refusal when full. */
int ring_put(ring_t *r, uint8_t b)
{
    uint32_t h = r->head;                    /* our own index: plain read  */
    uint32_t t = r->tail;                    /* theirs: see note below     */
    RING_ACQUIRE();                          /* their tail, then our write */
    if ((uint32_t)(h - t) >= RING_SIZE) {
        r->drops++;                          /* policy, counted not hidden */
        return 0;
    }
    r->buf[h & RING_MASK] = b;               /* the data                   */
    RING_RELEASE();                          /* data visible before index  */
    r->head = h + 1u;                        /* publish                    */
    return 1;
}

/* Consumer context only. Returns 0 when empty. */
int ring_get(ring_t *r, uint8_t *b)
{
    uint32_t t = r->tail;
    uint32_t h = r->head;
    RING_ACQUIRE();                          /* their head, then our read  */
    if (h == t)
        return 0;
    *b = r->buf[t & RING_MASK];
    RING_RELEASE();                          /* read done before release   */
    r->tail = t + 1u;
    return 1;
}

uint32_t ring_used(const ring_t *r)
{
    return (uint32_t)(r->head - r->tail);    /* correct across wrap        */
}
\end{ccode}
The stale read in each function is safe in one direction only, and that is worth
saying out loud. The producer may read a tail that is older than the truth; the
consequence is that it believes the buffer is fuller than it is and refuses a
byte it could have taken. The consumer may read a head that is older than the
truth; the consequence is that it believes the buffer is emptier than it is and
returns early. Both errors are conservative. Neither can produce a read of a
slot that has not been written, which is the only outcome that would be a fault.

\item \textbf{Decide the overflow policy and then count it.} Two policies are
defensible. Refuse the newest byte, which keeps the oldest data and is right for
a command stream where the front of a frame matters. Or overwrite the oldest,
which keeps the most recent data and is right for a sample stream where the
latest reading is the useful one. This chapter refuses the newest, because
chapter 5 parses frames and a frame with its head removed is worse than a frame
that never arrived.

What makes it a design decision rather than an accident is the counter. A buffer
that silently refuses bytes looks exactly like a link that works, until the day
it does not. Print the counter with every status line and make a non-zero value
visible.
\begin{ccode}
printf("rx ok=%lu drops=%lu used=%lu peak=%lu\n",
       (unsigned long) stats.accepted, (unsigned long) rx.drops,
       (unsigned long) ring_used(&rx), (unsigned long) stats.peak_used);
\end{ccode}
The peak occupancy is the number that tells you how close to the edge the
system runs. Record it in the consumer, which is the only context allowed to
call the occupancy function without racing itself, and reset it when the line is
printed.

\item \textbf{Build the same source for the host.} No target headers were
included, so this is a two line makefile and not a porting exercise. This is the
habit that makes chapters 5, 9 and 12 possible.
\begin{shellcode}
cd test && make && ./test_ring
\end{shellcode}
\begin{makecode}
CFLAGS = -std=c11 -O2 -Wall -Wextra -Werror -I../src
test_ring: test_ring.c ../src/ring.c
	$(CC) $(CFLAGS) -o $@ $^
\end{makecode}

\item \textbf{Write the test that would find an off-by-one.} A test that puts
three bytes and gets three bytes proves nothing. The test that matters drives
random burst sizes through the structure for tens of millions of bytes and
checks the sequence on the far side, so that every boundary is crossed many
times. Two hundred thousand steps with burst sizes uniform over 0 to 263, and
two loops per step, accepts about 21.6 million bytes, not the million a first
estimate of this suggests.

The generator matters more than it looks, and the obvious choice is wrong here.
A linear congruential generator with a power-of-two modulus has low bits of very
short period: its lowest three bits repeat with period eight. Taking the burst
size as that value modulo \texttt{RING\_SIZE + 8}, which is 264, which is
\(8 \times 33\), therefore gives burst sizes with a period-eight structure, and
the test crosses the buffer boundaries far less randomly than it appears to.
Use a generator with no low-bit weakness and record the seed.
\begin{ccode}
/* xorshift32. Not a linear congruential generator: the low bits of one with a
 * power-of-two modulus have period 8, and the burst size below is taken modulo
 * 264, which is 8 times 33, so the sizes would carry that period. */
static uint32_t xs32(uint32_t *s)
{
    uint32_t x = *s;
    x ^= x << 13; x ^= x >> 17; x ^= x << 5;
    *s = x;
    return x;
}

int main(void)
{
    ring_t r; ring_init(&r);
    uint32_t seed = 0x2026091u;                  /* recorded on purpose */
    uint64_t offered = 0, in = 0, out = 0, drops = 0;
    for (uint32_t step = 0; step < 200000u; step++) {
        uint32_t n = xs32(&seed) % (RING_SIZE + 8u);      /* over capacity */
        for (uint32_t i = 0; i < n; i++) {
            offered++;
            if (ring_put(&r, (uint8_t)(in & 0xffu))) in++; else drops++;
        }
        n = xs32(&seed) % (RING_SIZE + 8u);
        for (uint32_t i = 0; i < n; i++) {
            uint8_t b;
            if (!ring_get(&r, &b)) break;
            if (b != (uint8_t)(out & 0xffu)) { puts("ORDER FAULT"); return 1; }
            out++;
        }
        if (ring_used(&r) != in - out) { puts("OCCUPANCY FAULT"); return 1; }
    }
    /* The reconciliation, asserted and not merely printed. A buffer that
     * silently refuses bytes looks exactly like a link that works. */
    if (offered != in + drops)        { puts("POLICY FAULT"); return 1; }
    if (ring_drops(&r) != drops)      { puts("POLICY FAULT"); return 1; }
    printf("offered=%llu in=%llu out=%llu drops=%llu used=%u\n",
           offered, in, out, drops, ring_used(&r));
    return 0;
}
\end{ccode}
Three properties are checked and each one catches a different class of fault:
the byte sequence catches a masking or wrap fault, the occupancy catches an
index fault, and the drop count reconciling against the difference between bytes
offered and bytes accepted catches a policy fault.

\item \textbf{Be honest about what the host test cannot show.} Run the two
thread version as well, because it costs ten lines and it will find a genuine
logic fault if one exists. Then write in the file, at the top, that it cannot
find an ordering fault: the host is strongly ordered and will happily run a
version with both barriers removed for a week without complaint. The ordering
argument comes from the architecture reference and the published proof, and the
only machine that can falsify it is the target.

\item \textbf{Connect it to the target and measure the cost.} The producer is
the receive interrupt, which chapter 3 builds properly; here it can be the
simplest possible handler that reads the data register and calls the put
function. Measure with the cycle counter, around a block of puts rather than
around one, because a single call is shorter than the measurement overhead.
\begin{ccode}
DWT->CYCCNT = 0;                                  /* see the confirm list  */
for (int i = 0; i < 1000; i++) ring_put(&rx, (uint8_t) i);
uint32_t cyc = DWT->CYCCNT;
printf("put: %lu cycles per byte at 280 MHz\n", (unsigned long)(cyc / 1000u));
\end{ccode}
Run it three times: with the compiler-only barrier, with the processor barrier,
and with no barrier at all. The third is not a configuration you would ship; it
is there so the cost of correctness is a number in your table rather than a
belief.

\item \textbf{Record where the buffer lives.} Put it in the tightly coupled data
memory for now and write one line in the README saying so, because the moment a
transfer engine becomes the producer that choice reverses: the main engines
cannot reach that memory at all. That is chapter 19, and the note you write now
is what stops the bug in chapter 4.
\end{steps}

\subsection*{Build, flash and debug}

\diagram{c02_timing}{The two indices as an event timeline: values on the left, events across the bottom. The occupancy is the gap between the traces, the refusal is where the gap reaches the capacity, and the wrap is invisible because the counters never wrap in the region that matters.}

The build is chapter 1's, with two source files added and one new target for the
host. Keep the two builds in one repository and run the host one on every
commit, because it costs a second.

\begin{shellcode}
cmake -B build -DCMAKE_TOOLCHAIN_FILE=cmake/arm-none-eabi.cmake
cmake --build build -j && probe-rs run --chip STM32H7A3ZITx build/firmware.elf
make -C test && ./test/test_ring && ./test/test_ring_threads
\end{shellcode}

To look at the structure on a halted target, print it as a whole rather than
field by field, and print the occupancy as an expression so that you are reading
the same arithmetic the code uses.

\begin{shellcode}
(gdb) print rx
(gdb) print rx.head - rx.tail
(gdb) print/x rx.buf[rx.tail & 0xff]
\end{shellcode}

\begin{note}{When a board stops responding after this chapter}
A put function called from two different interrupt priorities is the likely
cause, and it presents as data that is fine for hours and then wrong once. The
recovery for a board that no longer enumerates is chapter 1's: hold the reset,
connect under reset with the flashing tool, and erase. Nothing in this chapter
touches option bytes, and nothing in this book does until the question on the
confirm list about unrecoverable option-byte writes has been settled.
\end{note}

\subsection*{Verification and acceptance criteria}

\begin{itemize}
\item The host test runs two hundred thousand random bursts, which accepts about
21.6 million bytes, and reports zero order faults and zero occupancy faults. It
runs in under two seconds and is part of the build. Write the loop in C rather
than driving the structure from a scripting language through a foreign function
interface: the same two hundred thousand steps cost about two minutes that way,
and a test that takes two minutes is a test somebody switches off.
\item Bytes offered equals bytes accepted plus bytes dropped, \emph{asserted} by
the test rather than printed by it, and against the structure's own counter as
well as the caller's tally, which are two different numbers that must agree. A
count that is only printed is not a check.
\item The structure compiles for the host and for the target from the same two
files with no conditional compilation in either.
\item The capacity assertion stops the build when the capacity is changed to a
number that is not a power of two. Verified on every test run by a script that
compiles the structure with 100, 255 and 1 and requires each to fail on the
assertion, then with 2, 256 and 4096 and requires each to succeed. Both halves
matter: a check that only ever expects failure would pass against a compiler
that refused everything. Verifying it by hand once and then trusting it is a
check nobody is running.
\item The cost of each barrier choice is a measured number in the budget table,
taken with the cycle counter, with the no-barrier case included for comparison
and marked as not shippable.
\item An overnight soak at the highest rate chapter 3 establishes reports a drop
count of zero and a peak occupancy well under the capacity.
\end{itemize}

\subsection*{Variants}

\begin{table}[H]\centering\small
\begin{tabular}{p{24mm}p{26mm}p{44mm}p{22mm}p{20mm}}
\toprule
Axis & Variant & What changes & Cost & Built in full in \\
\midrule
Execution model & Interrupt into a ring buffer & The producer is an interrupt handler and the consumer is the main loop. This is the baseline here and the structure is the deliverable & The handler must stay short & Here and chapter 3 \\
Execution model & Interrupt with a flag & One byte and a flag instead of a buffer. Simpler, and it loses data the moment two bytes arrive between visits to the loop & Data loss at any real rate & Chapter 3 \\
Execution model & DMA & The producer becomes a transfer engine, the head is derived from the transfer counter rather than written by code, and the compiler-only barrier stops being sufficient & Cache maintenance appears & Chapter 4 \\
Synchronisation & A volatile flag & The minimal case, useful only to show what it cannot do & Correct for one byte and nothing else & Chapter 3 \\
Synchronisation & Queue owned by a kernel & The restriction to one producer is lifted, and the kernel takes the cost of the critical section & Several hundred cycles and a dependency & Chapter 20 \\
Synchronisation & Stream buffer & The kernel's own single producer single consumer structure, which is this structure with a kernel's blocking semantics around it & Blocking needs a scheduler & Chapter 20 \\
Overflow policy & Overwrite the oldest & The producer advances the tail as well when full, which breaks the rule that each side owns one index and needs rethinking rather than a one-line edit & Correctness argument changes & Here, as the discussion in the step \\
Language & C++ & A template over the element type and the capacity, with the capacity check as a static assertion on the type & Nothing at run time & Chapter 9 \\
Intelligence and reach & Local only & The buffer feeds a parser on the same board & None & Here \\
\bottomrule
\end{tabular}
\caption{Variants for chapter 2. The overwrite-the-oldest policy is listed as a discussion rather than as code because it is not a variant of this structure; it is a different structure with a different correctness argument, and pretending otherwise is how the two get mixed in one file.}
\end{table}

\subsection*{Pitfalls}

\begin{itemize}
\item Believing \texttt{volatile} is a barrier. It orders volatile accesses
against each other and nothing else, and the access that matters here is the
ordinary store to the data array.
\item Wrapping the indices at the capacity instead of masking at the point of
use. It works, and then the occupancy is no longer a subtraction, and then one
slot is sacrificed, and then somebody changes the capacity to the number they
first wrote in the specification.
\item Drawing the burst size from the low bits of a linear congruential
generator. With a power-of-two modulus its lowest three bits repeat with period
eight, and the burst size here is taken modulo 264, which is \(8 \times 33\), so
the sizes inherit that period and the test crosses the buffer boundaries far
less randomly than it appears to. The symptom is a test that looks thorough and
is not.
\item Printing the drop count instead of asserting on it. Reconciling bytes
offered against bytes accepted plus bytes dropped is only a check if something
fails when it does not hold.
\item A capacity that is not a power of two. The mask silently stops matching
the size, and the failure appears as data corruption under load rather than as
an error.
\item A second producer. Two interrupt sources calling the put function is the
most common way this structure fails in the field, and it presents as rare
corruption that a test never reproduces.
\item Reading the occupancy from the producer context for a control decision.
The value is conservative in one direction only, and code that treats it as
exact will eventually act on a stale number.
\item Leaving the buffer in tightly coupled memory when a transfer engine
becomes the producer. The main engines cannot reach that memory, and the
symptom is a transfer that never completes rather than an error.
\item Testing only on the host and concluding the ordering is proven. The host
is strongly ordered. It cannot tell you anything about this.
\end{itemize}

\subsection*{Best practices applied}

\begin{itemize}
\item The invariant is written in the header, above the structure, in the place
a reader reaches before writing code against it.
\item The overflow policy is counted rather than hidden, and the counter is
printed whether or not it is zero.
\item The capacity constraint is a static assertion, so an edit that breaks it
stops the build instead of producing a silent fault.
\item The same source compiles for the host and for the target with no
conditional compilation, which keeps the test honest.
\item The limit of the test is documented in the test file itself, so the next
reader does not infer a guarantee the test cannot give.
\item The cost of correctness is measured rather than assumed, including the
case that is not shippable.
\end{itemize}

\subsection*{Stretch goals}

\begin{itemize}
\item Add the zero-copy pair of functions that hand out a contiguous span and
then commit it, which is what chapter 4 needs to give a transfer engine a
destination without an intermediate copy. The maintained library in the prior
art table is the reference for the interface.
\item Implement the batching variant from the correctness paper, where each side
caches the other's index and refreshes it only when its own arithmetic says the
buffer looks full or empty. Measure whether it is worth anything at this scale,
and publish the answer either way.
\item Add a second structure holding fixed-size records rather than bytes, and
measure whether the parser in chapter 5 gets simpler or merely different.
\item Run the target soak with the barriers removed and see whether a fault ever
appears on this core. A negative result after a week is a publishable finding
about this part, not a proof that the barriers are unnecessary.
\end{itemize}

\subsection*{Roadmap and next steps}

Chapter 3 fills this buffer from a real interrupt handler, which is where the
vendor library's two documented failure modes appear and where the rate at which
bytes begin to be lost is measured rather than asserted. Chapter 4 replaces the
handler with a transfer engine and the head with a hardware counter. Chapter 5
parses frames out of the buffer and is the first place the structure is load
bearing for something other than a console.

For the ordering material itself, the architecture reference for this profile
and the barrier application note are the primary sources, and the correctness
paper in the prior art table is the only place the acquire and release placement
is derived rather than asserted. The community roadmap referenced in appendix H
places this structure in the firmware track rather than the Linux one, and its
central advice, that projects teach more than reading, is the reason this
chapter ends with a soak and not with a summary.

\subsection*{Portfolio evidence}

\begin{itemize}
\item A repository whose README opens with the architecture figure and states
the invariant in two sentences before any code appears.
\item The host test output showing tens of millions of bytes through the structure with
zero faults, run from the build system rather than by hand.
\item The budget table with the measured column filled in for all three barrier
choices, naming the cycle counter as the instrument.
\item An overnight soak log with the drop counter and the peak occupancy, and a
sentence saying what rate produced it.
\end{itemize}

\subsection*{Sources}

Normative references:
\begin{itemize}
\item The architecture reference manual for this profile, for what the memory
barrier instructions order and for the rule that exception entry and return are
context synchronising.
\item Arm application note DAI0321A, on the memory barrier instructions and when
each is required.
\item Reference manual RM0455, for this part, when the buffer's placement and
the memories reachable by each transfer engine become the question in chapter
19. Not RM0433, which is its sibling.
\item The C11 standard, for the memory model the atomic builtins implement and
for the definition of unsigned wraparound the free running indices rely on.
\end{itemize}

Reusable implementations:
\begin{itemize}
\item Majerle's lightweight ring buffer, MIT, with the zero-copy regions that
chapter 4 needs.\newline
\url{https://github.com/MaJerle/lwrb}
\item Quantum Leaps' lock-free ring buffer, MIT, which states its atomicity
assumptions in the header.\newline
\url{https://github.com/QuantumLeaps/lock-free-ring-buffer}
\item Le, Guatto, Cohen and Pop, Correct and Efficient Bounded FIFO Queues,
SBAC-PAD 2013, the correctness argument under the C11 memory model.\newline
\url{https://www.irif.fr/~guatto/publications/sbac13.pdf}
\end{itemize}
   % A single producer, single consumer ring buffer
\project{03}{Receiving on interrupt without losing bytes}{NUCLEO-H7A3ZI-Q}{Interrupt-driven receive, overrun, and two documented library failure modes}

\begin{keyfacts}
\item[Board] NUCLEO-H7A3ZI-Q, and nothing else
\item[Peripherals] USART3 receive with its interrupt, the NVIC, the cycle counter, one spare pin as a marker
\item[Toolchain] arm-none-eabi-gcc with CMake on the target, Python with pyserial on the host
\item[Operating system] Bare metal
\item[Difficulty] 3 of 5
\item[Effort] 3 evenings of about four hours
\item[Deliverable] Three receive paths built in full and compared on one board: a polled loop, an interrupt with a flag, and an interrupt into chapter 2's ring buffer. Both documented library failure modes reproduced on purpose and then repaired. A host harness that raises the rate until loss begins, and a measured number for where that is
\end{keyfacts}

\subsection*{Why this project}

This is the only chapter in the book with no clean reference implementation to
start from. What exists instead is better: a well documented set of failure
modes in the silicon vendor's own community, reported repeatedly over years,
with the register-level explanation in the threads. Two of them matter here.
The first is that the blocking-free receive call disables its own receive
interrupt when the requested count completes, so every byte that arrives between
the completion callback and the next call to re-arm it is lost, and the loss is
silent. The second is that the overrun path returns the peripheral to its ready
state without clearing the overrun flag, so the interrupt re-asserts
immediately and the handler runs continuously, which presents as a board that
has stopped rather than as an error.

Reproducing a documented bug on purpose and then repairing it teaches more than
starting from working code, because the repair is only convincing once you have
seen the symptom. This chapter therefore has two halves. The first half builds
three receive paths and the failure modes, all on one board with one switch
between them. The second half measures: a host harness raises the line rate in
steps and reports where bytes first go missing, and the counters on the target
separate loss inside the target from loss in the bridge between the host and the
board. Without that separation the measurement means nothing, because the most
common finding, that a link fails above some rate, is usually a statement about
the bridge.

The chapter also owns four entries in the book's variant matrix. The polled loop
and the interrupt with a flag are built in full here and referenced from
everywhere else. The interrupt into a ring buffer is shared with chapter 2,
which built the structure; this chapter supplies the producer. And the host side
in Python is built in full here, so later chapters that need a host script have
a pattern to follow rather than a new invention each time.

\begin{note}{Two failure modes are reproduced here and they are documented not invented}
Both come from the vendor's own community and both are register-level
consequences of how the library's receive path is written, not of the silicon.
Neither is reproduced from library source in this book. The call pattern that
triggers each one is shown, the symptom is described, and the repair is written
at register level so that it can be read and checked. Never present generated
or vendor-supplied code as reviewed source; name the behaviour, then show the
code you wrote.
\end{note}

\subsection*{Prior art and what to reuse}

\begin{table}[H]\centering\small
\begin{tabular}{p{34mm}p{46mm}p{38mm}p{20mm}}
\toprule
Source & What it gives & What it does not & Licence \\
\midrule
controllerstech's interrupt ring buffer for this vendor & The most widely read worked example of exactly this problem, and a clear statement of the shape of the solution & It has \textbf{no licence file}, so default copyright applies and it cannot be vendored or adapted. It is also written for an older family whose flag clearing procedure differs from this one & None stated \\
The vendor's community threads on the two failure modes & The reproduction conditions, the register-level cause, and the confirmation that these are known rather than local & No fix you can depend on. Thread code is unreviewed and often written for a different family & Forum posts \\
The vendor's basic example tree for this board & A working receive example at both the library and the register level, permissively licensed, which is the precedent for this book's variants format & Its receive is a fixed count, which is the shape that loses bytes between calls. It is a starting point and not an answer & BSD-3-Clause \\
pyserial & The host side: a portable serial port with settable rate, timeouts and buffer control & No measurement method. The rate ramp, the counter reconciliation and the separation of target loss from bridge loss are the author's & BSD-3-Clause \\
Chapter 2 of this book & The ring buffer, its ordering argument and its drop counter & Nothing about the peripheral, the flags or the NVIC & This volume \\
\bottomrule
\end{tabular}
\caption{Prior art for chapter 3. The most useful repository in this row has no licence file, which means all rights reserved. It is read as a worked example and re-derived, and that is stated here rather than discovered later.}
\end{table}

What is left to write is the handler, the flag handling, the switch between the
three paths, and the whole of the measurement. The handler is fifteen lines. The
measurement is the chapter.

\subsection*{Parts from the inventory}

\begin{table}[H]\centering\small
\begin{tabular}{p{40mm}p{66mm}p{40mm}}
\toprule
Part & Role & Interface \\
\midrule
NUCLEO-H7A3ZI-Q & The receiver under test, and its own counters & Micro USB to the host \\
A USB data cable & Power, programming and the console under test & Micro USB \\
Host PC & Runs the load harness and reads the counters & Python with pyserial \\
Renkforce USB/TTL cable, optional & A second serial path when the question is whether the on-board bridge is the limit. \textbf{Disconnect the red 5 V lead} & 3.3 V logic only \\
\bottomrule
\end{tabular}
\caption{Inventory items used in chapter 3. The USB to serial cable is optional and is the control experiment: if a rate fails on the on-board bridge and passes on the cable, the bridge was the limit.}
\end{table}

\subsection*{System architecture}

\diagram{c03_arch}{The three receive paths share everything except the block in the middle. One build switch selects which one is compiled, so the comparison is between three programs that differ in one place.}

The switch matters more than it looks. A comparison between three
implementations that also differ in their initialisation, their console output
and their counters is not a comparison. Everything outside the marked block is
identical across the three builds, including the status line format, so the
harness on the host does not know or care which path is running.

\subsection*{Peripheral configuration}

\begin{table}[H]\centering\small
\begin{tabular}{p{22mm}p{26mm}p{24mm}p{30mm}p{30mm}}
\toprule
Peripheral & Mode & Clock source & Pins and function & Interrupt and transfers \\
\midrule
USART3 & Asynchronous receive, 8N1, rate set by the harness & Peripheral bus & Confirm in the board manual before printing them & Receive interrupt enabled once and never disabled \\
USART3 receive FIFO & Present on this generation; depth and threshold encoding checked in RM0455 & Peripheral bus & None & Drained in a loop inside the handler \\
NVIC & Priority grouping set once at startup & Not applicable & None & USART3 at a middle priority, chosen in the last step \\
DWT cycle counter & Free running & Core clock & None & Measures the handler, entry to exit \\
GPIO, one Zio pin & Push-pull output & Peripheral bus & Chosen from the pins a shield leaves free & Raised on handler entry, lowered on exit \\
\bottomrule
\end{tabular}
\caption{Peripheral configuration. The FIFO row is written as a check rather than as a fact: this family added a receive FIFO to the USART and the threshold encoding is not inherited from an older part. Confirm it in RM0455 for this device, not in a document for the H743.}
\end{table}

\subsection*{Wiring}

\diagram{c03_wiring}{Nothing is wired for the main experiment. The optional second path, drawn faintly, is the control: a separate USB to serial cable on a second USART, used only to answer the question of whether the on-board bridge is the limit. Its red 5 V lead is disconnected.}

\begin{note}{Three point three volt logic only}
If the optional cable is used, its signal lines are 3.3 V and its red lead
carries 5 V. Disconnect the red lead before the cable is anywhere near the
header. Nothing in this chapter needs 5 V, the board supplies its own power over
the debug connection, and a 5 V line into a general purpose pin damages the
part in a way no software repairs.
\end{note}

\subsection*{Memory and timing budget}

\diagram{c03_mem}{Where a byte can be while it waits, and how long each place buys you. The shift register holds one byte, the peripheral FIFO holds a few, and the ring buffer holds the rest. Each stage is a time budget, not just a capacity.}

\begin{table}[H]\centering\small
\begin{tabular}{p{44mm}p{28mm}p{28mm}p{30mm}}
\toprule
Quantity & Budget & Measured & Margin \\
\midrule
Byte time at 115200 8N1 & 86.8 \textmu{}s & computed, not measured & not applicable \\
Byte time at 921600 8N1 & 10.9 \textmu{}s & computed, not measured & not applicable \\
Core cycles per byte at 921600 & 3038 & computed at 280 MHz & not applicable \\
Handler duration, entry to exit & 200 cycles & not measured & not measured \\
Worst case interrupt latency & 1000 cycles & not measured & not measured \\
Highest rate with zero loss & 921600 & not measured & not measured \\
Overruns in a ten minute run & 0 & not measured & not measured \\
\bottomrule
\end{tabular}
\caption{The budget table. The first three rows are arithmetic and are marked as such; the rest wait for the cycle counter and the harness. The headroom is the story: at the highest rate on the list the handler has about three thousand cycles per byte, so the handler's own length is not the limit and the latency under a longer interrupt is.}
\end{table}

\subsection*{Firmware design (UML)}

\diagram{c03_uml}{The receive path as a state machine, with the two documented failure modes drawn as the transitions that should not exist. The repaired path never leaves the enabled state and never returns to ready with a flag still set.}

The state machine is worth drawing because the repair is a statement about
states rather than about lines of code. The library's receive call is a
transaction: it moves to a receiving state, counts down a requested number of
bytes, and on reaching zero moves back to ready and disables the interrupt that
was feeding it. Everything about that is reasonable for a protocol where you
know the length in advance. For a console, a command line or any stream whose
length is not known, it is the wrong shape, and the bytes that arrive during the
ready state are not delayed, they are gone.

The repaired path has one state. The receive interrupt is enabled at
initialisation and is never disabled. There is no count, no completion and no
re-arming, so there is no window. The handler drains whatever the peripheral
holds into the ring buffer and returns. The overrun condition is not an error
that ends a transaction; it is a counter that goes up.

\subsection*{Data flow (ASCII)}

\begin{asciiart}
  host                          bridge                target
  +--------------------+       +---------+      +-----------------------------+
  | loadgen.py         |       | ST-LINK |      | USART3 shift register       |
  |  9600 -> 921600    | ----> | V3E     | ---> |   |                         |
  |  known pattern     |  USB  | virtual |      |   v                         |
  |  counts bytes sent |       | COM     |      | receive FIFO, drained in a  |
  +--------------------+       +---------+      | loop inside the handler     |
           ^                                    |   |                         |
           |                                    |   v                         |
           |   status line once a second        | ring_put -> ring buffer     |
           +------------------------------------+   |                         |
                                                |   v                         |
   sent == accepted + drops + overruns ?        | main loop consumes, counts  |
   no  -> the bytes never reached the target    +-----------------------------+
   yes -> the target is where they were lost
\end{asciiart}

\subsection*{Repository layout}

\begin{plaincode}
nucleo-h7a3-rx/
  CMakeLists.txt                 # -DRX_PATH=polled|flag|ring selects one
  cmake/arm-none-eabi.cmake      # reused unchanged from chapter 1
  src/ring.c, src/ring.h         # reused unchanged from chapter 2
  src/usart3_ll.c                # register level: init, flags, ICR, NVIC
  src/rx_polled.c                # variant 1, built in full
  src/rx_flag.c                  # variant 2, built in full
  src/rx_ring.c                  # variant 3, the one that keeps the bytes
  src/faults.c                   # the two documented failure modes, on purpose
  src/stats.c                    # one status line format for all three
  src/main.c
  host/loadgen.py                # the rate ramp, pyserial
  host/report.py                 # turns a run into the results table
  docs/results.md                # the measured table, with the date of each run
  README.md
\end{plaincode}

\subsection*{Steps}

\begin{steps}
\item \textbf{Build the polled loop first and make it lose bytes.} This is the
first of the two variants this chapter owns, and it is worth building properly
rather than as a straw figure, because it is genuinely the right answer when the
loop has nothing else to do.
\begin{ccode}
/* rx_polled.c: correct as long as the loop returns here in under one byte
 * time. At 115200 that is 86.8 us. At 921600 it is 10.9 us. */
void rx_poll_step(void)
{
    if (USART3->ISR & USART_RXNE_FLAG) {      /* see the note on the name */
        uint8_t b = (uint8_t) USART3->RDR;    /* the read clears the flag */
        stats.accepted++;
        consume(b);
    }
    if (USART3->ISR & USART_ISR_ORE) {
        USART3->ICR = USART_ICR_ORECF;        /* clear it or it re-asserts */
        stats.overruns++;
    }
}
\end{ccode}
Now add the line that is present in every real program and absent from every
tutorial: some other work in the loop. A single blocking print at 115200 takes
longer than a byte time at any rate above about 9600, and the polled loop starts
losing bytes at exactly that point. Measure it rather than asserting it.
\begin{ccode}
while (1) {
    rx_poll_step();
    housekeeping();        /* 2 ms: a print, a sensor read, anything real */
}
\end{ccode}

\item \textbf{Confirm the flag name before you rely on it.} This family added a
receive FIFO to the USART, and the device header renamed the receive-not-empty
flag when it did. Both the older name and the newer combined name appear in
published code for parts in this family, and the wrong one is a compile error
rather than a silent fault, which is the one piece of luck in this chapter.
Define it once and check it against the device header and RM0455.
\begin{ccode}
/* One place for the name, because it differs across this family's headers. */
#if defined(USART_ISR_RXNE_RXFNE)
#define USART_RXNE_FLAG USART_ISR_RXNE_RXFNE   /* parts with the receive FIFO */
#else
#define USART_RXNE_FLAG USART_ISR_RXNE
#endif
\end{ccode}

\item \textbf{Build the interrupt with a flag, the second variant this chapter
owns.} It is the smallest step away from polling and it is where most projects
stop. It is also the version that loses data at any rate, silently, and the
point of building it is to see that.
\begin{ccode}
/* rx_flag.c. One byte of storage and one flag. */
static volatile uint8_t rx_byte;
static volatile uint8_t rx_ready;

void USART3_IRQHandler(void)                  /* the flag variant */
{
    if (USART3->ISR & USART_RXNE_FLAG) {
        rx_byte  = (uint8_t) USART3->RDR;
        rx_ready = 1;                         /* the previous byte is gone */
    }
}
\end{ccode}
There is no counter that can be added here without changing the design, and
that is the lesson. The handler cannot tell whether the byte it is overwriting
had been read, because a flag carries one bit of state and the question needs
two. A version that counts the loss is a version that has a queue of length two,
which is a ring buffer with the capacity written badly.

\item \textbf{Build the interrupt into the ring buffer.} Chapter 2 built the
structure and argued its ordering. This is its producer, and it is the shortest
of the three.
\begin{ccode}
/* rx_ring.c. The receive interrupt is enabled once, at init, and never
 * disabled. There is no transaction, no count and no re-arming, so there
 * is no window in which bytes arrive with the interrupt switched off. */
void USART3_IRQHandler(void)
{
    uint32_t isr = USART3->ISR;

    if (isr & USART_ISR_ORE) {                /* check before draining */
        USART3->ICR = USART_ICR_ORECF;        /* write one, clear it     */
        stats.overruns++;
    }
    while (USART3->ISR & USART_RXNE_FLAG) {   /* drain the whole FIFO    */
        if (!ring_put(&rx, (uint8_t) USART3->RDR))
            break;                            /* full: counted in ring   */
    }
}
\end{ccode}
Three details carry the repair. The interrupt is enabled once and never
disabled. The flag is drained in a loop rather than once, because a FIFO may
hold several bytes by the time the handler runs. And the overrun is cleared
through the clear register, which is how this generation of the peripheral
works: on older parts the flag was cleared as a side effect of reading the
status register and then the data register, and code copied from those parts
leaves the flag set here.

\item \textbf{Reproduce the first documented failure mode on purpose.} The
receive call that takes a count disables the receive interrupt when the count
reaches zero and calls a completion callback. The usual pattern is to re-arm it
inside that callback. Between the disable and the re-arm there is a window, and
at any rate where a byte can arrive inside that window, it is lost with no
indication.
\begin{ccode}
/* faults.c, mode 1. The call pattern is shown; the library's internals are
 * described in prose above and are not reproduced here. */
void on_receive_complete(void)        /* runs from the library's handler */
{
    stats.accepted++;
    consume(one_byte);
    start_receive(&one_byte, 1);      /* re-arm: the window is above this */
}
\end{ccode}
Run the harness at a low rate and then raise it. The loss is not gradual in the
way a buffer overflow is; it appears as a steady fraction of bytes missing that
grows with the rate, because the window is a fixed length of time and the byte
time is shrinking. Record the fraction at three rates, because that shape is the
signature that distinguishes this failure mode from a full buffer.

\item \textbf{Reproduce the second documented failure mode.} Send a burst faster
than the handler is allowed to run, so the overrun flag sets, and then take the
error path that returns the peripheral to a ready state without writing the
clear register. The flag stays set, the interrupt condition stays asserted, and
the handler is entered again immediately on return. The board stops making
progress in the main loop and the console goes quiet.
\begin{ccode}
/* faults.c, mode 2. Removing one line is the whole reproduction. */
void USART3_IRQHandler(void)
{
    uint32_t isr = USART3->ISR;
    if (isr & USART_ISR_ORE) {
        /* USART3->ICR = USART_ICR_ORECF;   <- the line that is missing */
        state = STATE_READY;                 /* back to ready, flag still set */
        return;
    }
    while (USART3->ISR & USART_RXNE_FLAG)
        ring_put(&rx, (uint8_t) USART3->RDR);
}
\end{ccode}
Confirm it is the condition you think it is rather than a fault elsewhere: halt
the target and read the status register. The overrun bit is set, and the
interrupt is pending in the NVIC, and the main loop's counter has stopped
advancing. That is three independent observations of the same cause, which is
the standard to hold yourself to before writing a repair.
\begin{shellcode}
(gdb) print/x USART3->ISR
(gdb) print/x NVIC->ISPR[1]
(gdb) print loop_counter
\end{shellcode}

\item \textbf{Write the host harness.} This is the Python variant, built in full
here and referenced from every later chapter that needs a host script. It walks
the rate upwards, sends a known pattern for a fixed time at each rate, and reads
the target's own counters afterwards.
\begin{pycode}
#!/usr/bin/env python3
"""loadgen.py: raise the line rate until bytes go missing, and say where."""
import argparse, sys, time
import serial

RATES = [9600, 19200, 38400, 57600, 115200, 230400, 460800, 921600]
PATTERN = bytes(range(256))          # one lap of every byte value

def status(port, timeout=2.0):
    """Ask the target for its counters. One line: ok=N drops=N ovr=N."""
    port.reset_input_buffer()
    port.write(b"?\n")
    deadline = time.time() + timeout
    while time.time() < deadline:
        line = port.readline().decode("ascii", "replace").strip()
        if line.startswith("ok="):
            return dict(kv.split("=") for kv in line.split())
    return None

def run_one(dev, baud, seconds):
    with serial.Serial(dev, baud, timeout=0.5, write_timeout=5.0) as p:
        time.sleep(0.2)                       # let the bridge settle
        before = status(p)
        sent, end = 0, time.time() + seconds
        while time.time() < end:
            p.write(PATTERN)
            sent += len(PATTERN)
        p.flush()
        time.sleep(0.3)                       # drain the bridge
        after = status(p)
    got  = int(after["ok"])   - int(before["ok"])
    drop = int(after["drops"])- int(before["drops"])
    ovr  = int(after["ovr"])  - int(before["ovr"])
    return sent, got, drop, ovr

def main():
    ap = argparse.ArgumentParser()
    ap.add_argument("device")
    ap.add_argument("--seconds", type=float, default=5.0)
    a = ap.parse_args()
    print(f"{'baud':>8} {'sent':>9} {'accepted':>9} {'drops':>7} "
          f"{'overruns':>9} {'lost':>7} verdict")
    for baud in RATES:
        sent, got, drop, ovr = run_one(a.device, baud, a.seconds)
        lost = sent - got - drop
        if lost == 0 and ovr == 0:
            verdict = "clean"
        elif ovr > 0:
            verdict = "target overran"
        else:
            verdict = "never arrived: suspect the bridge"
        print(f"{baud:>8} {sent:>9} {got:>9} {drop:>7} {ovr:>9} "
              f"{lost:>7} {verdict}")

if __name__ == "__main__":
    sys.exit(main())
\end{pycode}
The three-way verdict is the part worth keeping. A run where bytes are missing
and the target reports no overrun and no drop did not lose them in the target at
all, and reporting that as a target limit would be wrong.

\item \textbf{Separate the target from the bridge.} If any rate produces the
third verdict, repeat that rate on the optional USB to serial cable, on a second
USART, with the red lead disconnected. If the same rate is clean there, the
on-board bridge was the limit. Write both numbers in the results file with the
full date of each run, because the bridge's behaviour is a property of the host,
its driver and its operating system as much as of the board, and a number
without a date and a host is not a result.

\item \textbf{Set the interrupt priority deliberately and measure the
handler.} Set the priority grouping once at startup and give the receive
interrupt a middle priority, so that a faster interrupt added later can preempt
it and a slower one cannot delay it. Then measure the handler with the cycle
counter and the marker pin together: the counter gives the duration, and the pin
gives the power instrument something to align against in chapter 10.
\begin{ccode}
NVIC_SetPriorityGrouping(3);                 /* fix the split once, at boot */
NVIC_SetPriority(USART3_IRQn, 6);            /* middle: room either side    */
NVIC_EnableIRQ(USART3_IRQn);

void USART3_IRQHandler(void)                 /* instrumented build only     */
{
    MARKER_GPIO->BSRR = MARKER_PIN;          /* pin high on entry           */
    uint32_t t0 = DWT->CYCCNT;
    rx_drain();
    stats.handler_cycles = DWT->CYCCNT - t0;
    MARKER_GPIO->BSRR = MARKER_PIN << 16;    /* pin low on exit             */
}
\end{ccode}
Compare the measured duration against the byte time in the budget table. If the
handler is two hundred cycles and a byte at the highest rate is three thousand,
the handler is not the constraint, and the thing to look at next is the longest
interrupt in the system that can delay it.

\item \textbf{Publish the results table.} Three paths, eight rates, one table,
with the date of the run and the host named. This is the deliverable, and it is
the thing that does not exist anywhere else for this board.
\begin{shellcode}
python host/loadgen.py /dev/ttyACM0 --seconds 5 | tee docs/results_ring.md
\end{shellcode}
\end{steps}

\subsection*{Build, flash and debug}

\diagram{c03_timing}{The vector table and the priority axis. Only the entries this chapter uses are drawn. The preemption arrow shows a higher priority interrupt interrupting the receive handler, and the tail chain shows the case that costs less than a full exit and entry.}

One switch selects the path, so a full comparison is three builds and three
runs.

\begin{shellcode}
for p in polled flag ring; do
  cmake -B build-fw/$p -G Ninja -DCMAKE_TOOLCHAIN_FILE=cmake/arm-none-eabi.cmake -DRX_PATH=$p
  cmake --build build-fw/$p -j
done
cp build-fw/ring/firmware.bin "$PROBE_DISK"/   # onto the probe disk
\end{shellcode}

\begin{note}{When the board goes quiet and the debugger still connects}
That is the signature of the second failure mode: the processor is alive and is
spending all of its time entering one handler. Halt and read the peripheral's
status register and the NVIC pending register. If a flag is set that nothing
clears, you have found it. If instead the debugger cannot connect at all, that
is a different fault: hold the reset, connect under reset, and erase. Nothing in
this chapter writes option bytes, and nothing does until the question on the
confirm list about unrecoverable option-byte writes is settled.
\end{note}

\subsection*{Verification and acceptance criteria}

\begin{itemize}
\item The three paths produce the same status line format, so the host harness
does not know which is running. Verified by running the harness against all
three without editing it.
\item At the highest rate that is clean, bytes sent equals bytes accepted plus
drops plus overruns, exactly, over a ten minute run.
\item The polled loop's loss threshold is measured with the housekeeping work
present and again with it removed, and the two numbers differ, which is the
demonstration that the loop's other work is the variable.
\item Both documented failure modes are reproduced deliberately, each behind a
build switch, each with the symptom recorded: a growing fraction of missing
bytes for the first, and a stalled main loop with a set flag for the second.
\item The repaired path runs for ten minutes at the highest clean rate with zero
overruns and zero drops.
\item The handler duration is a measured number from the cycle counter, and it
is compared in the results file against the byte time at each rate.
\item Every row of the results file carries the full date of the run and the
name of the host machine.
\end{itemize}

\subsection*{Variants}

\begin{table}[H]\centering\small
\begin{tabular}{p{24mm}p{26mm}p{44mm}p{22mm}p{20mm}}
\toprule
Axis & Variant & What changes & Cost & Built in full in \\
\midrule
Execution model & Polled loop & The loop reads the flag itself. Correct while the loop returns within one byte time, and the measurement here says where that stops being true & Nothing else can run & Here \\
Execution model & Interrupt with a flag & One byte and one flag. The smallest step from polling, and it loses data at any real rate with no way to count the loss & Silent data loss & Here \\
Execution model & Interrupt into a ring buffer & The handler is the producer of chapter 2's structure. The interrupt is enabled once and never disabled & A structure to reason about & Here and chapter 2 \\
Execution model & DMA, circular, idle line & The handler disappears and a transfer engine writes the buffer. Variable length frames come from the idle line rather than from a count & Cache maintenance appears & Chapter 4 \\
Execution model & RTOS task & The handler gives a notification and a task drains the buffer, which moves the work out of interrupt context & A scheduler & Chapter 20 \\
Synchronisation & A volatile flag & Built here as the second variant, and shown to be sufficient for exactly one byte & Correct only at one byte & Here \\
Language & Host in Python & The harness, pyserial, the rate ramp and the three-way verdict. Every later host script follows this pattern & None & Here \\
Language & Host in Node.js or Java & The same harness on a different runtime, which is worth doing once to prove the target does not care & A second runtime & Chapter 12 \\
Peripheral substitution & A different USART & The same three paths on another instance, which changes the pins and nothing else & A second set of pins & Chapter 17 \\
\bottomrule
\end{tabular}
\caption{Variants for chapter 3. Four entries in the book's matrix are built in full here, which is why this chapter is longer in code than its neighbours.}
\end{table}

\subsection*{Pitfalls}

\begin{itemize}
\item Re-arming a counted receive call inside its own completion callback. The
window between the disable and the re-arm is where the bytes are lost, and
nothing reports it.
\item Returning from an error path without clearing the flag that caused it. On
this generation the overrun flag is cleared through the clear register, not by
reading the status and data registers as on older parts.
\item Reading the flag once per handler entry instead of draining in a loop.
With a receive FIFO the peripheral may hold several bytes, and a handler that
takes one leaves the rest to become an overrun.
\item Copying the flag and clearing procedure from a tutorial written for an
older family. The registers have the same names and different semantics, which
is worse than having different names.
\item Concluding that a rate is the target's limit when the target's own
counters say the bytes never arrived. Measure both ends before naming a cause.
\item Printing from inside the handler while debugging the handler. A blocking
print is longer than several byte times and creates the overrun you are looking
for.
\item Leaving the instrumented build's marker pin code in the shipping build.
It is two instructions, but it is two instructions you did not intend to ship.
\end{itemize}

\subsection*{Best practices applied}

\begin{itemize}
\item The receive interrupt is enabled once and never disabled, so there is no
window in which bytes arrive with the peripheral quiet.
\item Every loss has a counter: drops in the buffer, overruns in the peripheral,
and the difference between bytes sent and bytes accounted for on the host.
\item The three paths differ in exactly one file, so the comparison is a
comparison.
\item A failure mode is reproduced before it is repaired, and the reproduction
stays in the repository behind a build switch so the next reader can see it.
\item A claim about the register names is written as a check against the
reference manual for this part, not as a fact inherited from a sibling.
\item Every published number carries its instrument, its full date and its host.
\end{itemize}

\subsection*{Stretch goals}

\begin{itemize}
\item Add the receive timeout, which fires after a chosen number of idle bit
times and gives a natural frame boundary without the line-idle interrupt that
chapter 4 uses. Compare the two on the same traffic.
\item Drive the harness from the Raspberry Pi rather than the Windows host and
see whether the bridge result moves. If it does, the earlier number was about
the host and the results file should say so.
\item Add a deliberate long interrupt at a higher priority and find the duration
at which the receive path begins to overrun. That number is the real headroom
and nobody publishes it for this part.
\item Run the same three paths at the register level and through the vendor
library and compare code size and handler duration, which is the measurement
chapter 17 generalises.
\end{itemize}

\subsection*{Roadmap and next steps}

Chapter 4 removes the handler entirely: a transfer engine writes into the buffer
and the write position comes from the transfer counter rather than from code,
and the line-idle condition gives frame boundaries for variable length data.
Chapter 5 puts a frame format and a checksum on top of this byte stream.
Chapter 12 turns the harness written here into a rig that runs on every commit.

For the material behind the two failure modes, the vendor's community threads
are the primary source and are cited individually in the repository rather than
in this book, because thread addresses move. The community roadmap referenced in
appendix H places interrupt handling and the vendor library's shape squarely in
the firmware track, and the free course that walks one problem from a polled
loop to an event-driven architecture covers this same progression across its
early lessons; its code licence means it is recommended and never quoted.

\subsection*{Portfolio evidence}

\begin{itemize}
\item A repository with three receive paths behind one build switch and both
documented failure modes reproducible on demand.
\item The results table: three paths, eight rates, with the drop, overrun and
never-arrived columns filled in, each row carrying its full date and host.
\item A short written finding on where loss begins on this board and whether the
limit is the target or the bridge, which is a number that is asserted in several
places and measured in none.
\item The handler duration from the cycle counter, next to the byte time at each
rate, so the headroom is visible rather than claimed.
\end{itemize}

\subsection*{Sources}

Normative references:
\begin{itemize}
\item Reference manual RM0455, for this part: the USART register set, the
receive FIFO, the flag clearing procedure through the clear register, and the
receive timeout. Not RM0433, which is its sibling and differs.
\item The board user manual for MB1363, for the console pins and for what the
on-board bridge does and does not guarantee about rate.
\item The architecture reference manual for this profile, for interrupt
priority, preemption and tail chaining.
\item The device header for this part, for the flag name that changed when the
receive FIFO was added.
\end{itemize}

Reusable implementations:
\begin{itemize}
\item pyserial, BSD-3-Clause, the host side of the harness.\newline
\url{https://github.com/pyserial/pyserial}
\item The silicon vendor's basic example tree for this board, BSD-3-Clause,
which carries the same problem solved polled, on interrupt and by transfer
engine in sibling directories.
\item controllerstech's interrupt receive ring buffer, which has \textbf{no
licence file} and is therefore read as a worked example and never vendored.\newline
\url{https://github.com/controllerstech/stm32-uart-ring-buffer}
\item Chapter 2 of this volume, for the ring buffer this chapter fills.
\end{itemize}
   % Receiving on interrupt without losing bytes
\project{04}{Circular DMA and the idle line}{NUCLEO-H7A3ZI-Q}{Transfer-driven receive, variable length frames, the transfer counter}

\begin{keyfacts}
\item[Board] NUCLEO-H7A3ZI-Q, and nothing else
\item[Peripherals] USART3 receive, DMA1 stream 0 through the request multiplexer, the line-idle interrupt, the cycle counter
\item[Toolchain] arm-none-eabi-gcc with CMake on the target, the chapter 3 harness on the host
\item[Operating system] Bare metal
\item[Difficulty] 4 of 5
\item[Effort] 3 evenings of about four hours
\item[Deliverable] A receive path in which no interrupt runs per byte, frames of any length arrive whole, the write position is read out of a hardware counter, the wraparound is handled as two spans rather than as a special case, and the cache maintenance that the canonical reference leaves out is present and argued
\end{keyfacts}

\subsection*{Why this project}

Chapter 3 ended with a handler that runs once per byte. At 115200 that is about
eleven thousand interrupt entries a second, each one costing the entry, the
drain and the exit. It works, and the measurement in that chapter says by how
much. This chapter removes the per-byte cost entirely: a transfer engine writes
every byte into memory with no processor involvement at all, and the processor
finds out where the engine has got to by reading a hardware counter.

Two ideas carry the chapter, and neither is obvious the first time. The first is
that in circular mode the transfer counter is the write position: it counts down
from the buffer size, so the number of bytes written since the engine last
restarted is the buffer size minus the counter. Nothing writes a head index;
the head index is a register you read. The second idea is that the
peripheral's line-idle condition is a frame boundary for free. A protocol with a
length prefix tells you when a frame ends, and chapter 5 builds one; a line that
goes quiet tells you the same thing without any protocol at all, which is what
makes this the right path for a command line, a modem's replies, or anything
whose length you do not know in advance.

The third thing this chapter carries is a gap in its own prior art, and it is
worth naming at the top rather than in a footnote. The canonical reference
implementation for this problem is excellent, permissively licensed, and has no
cache handling of any kind. With the data cache on and the buffer in main
memory, the processor can read a line it cached before the engine wrote it. The
maintenance step that fixes that is this chapter's work, the argument for it is
chapter 19's, and a chapter that copied the reference and said nothing would
have shipped a fault.

\begin{note}{This chapter reverses a decision from chapter 2}
Chapter 2 put the ring buffer in tightly coupled memory, which is the right
place while the producer is code on the core. It is the wrong place now. The
main transfer engines are expected not to reach the tightly coupled memories at
all on this part, and the symptom is a transfer that never starts or never
completes rather than an error. Which engine reaches which memory is item 16 on
the confirm list, is settled against RM0455 and not against a document for the
H743, and is the subject of chapter 19. The buffer moves to the main memory in
step 1 and a comment in the linker script says why.
\end{note}

\subsection*{Prior art and what to reuse}

\begin{table}[H]\centering\small
\begin{tabular}{p{34mm}p{46mm}p{38mm}p{20mm}}
\toprule
Source & What it gives & What it does not & Licence \\
\midrule
Majerle's transfer-driven receive and transmit examples & The canonical treatment of this exact problem, with working examples across eight silicon families, the position arithmetic, the wraparound split and the three-event structure all laid out clearly & \textbf{No cache handling anywhere}, and its H7 example is for a different member of this family, not for this board. Under the rule in the front matter, nothing from it is inherited without being checked against RM0455 & MIT \\
Majerle's lightweight ring buffer & The zero-copy pair of functions that hand out a contiguous span, which is the shape this chapter's consumer wants & The same coherence assumption as above & MIT \\
The silicon vendor's basic example tree for this board & A transfer-driven receive at both abstraction levels, permissively licensed, and the request multiplexer configuration for this part & Fixed length transfers, no idle line, no frame concept & BSD-3-Clause \\
Chapters 2 and 3 of this book & The consumer side, the drop accounting and the measurement harness & Nothing about transfer engines & This volume \\
Chapter 19 of this book & The cache argument in full, with the region attributes and the three ways to make a buffer safe & Nothing until you read it; this chapter uses the narrowest of the three fixes and says so & This volume \\
\bottomrule
\end{tabular}
\caption{Prior art for chapter 4. The first row is close to the whole chapter in a box and is worth reading before writing a line. What it does not give is the part that matters most on this silicon, which is why this chapter exists as more than a summary.}
\end{table}

What is left to write is the cache maintenance, the placement argument, the lap
guard described in step 7, and the measurement that compares this path against
chapter 3's on the same board with the same harness. That comparison does not
exist anywhere for this part.

\subsection*{Parts from the inventory}

\begin{table}[H]\centering\small
\begin{tabular}{p{40mm}p{66mm}p{40mm}}
\toprule
Part & Role & Interface \\
\midrule
NUCLEO-H7A3ZI-Q & The receiver, the transfer engine and the counters & Micro USB to the host \\
A USB data cable & Power, programming and the link under test & Micro USB \\
Host PC & Sends frames of deliberately varied length with deliberate gaps & Python with pyserial, from chapter 3 \\
\bottomrule
\end{tabular}
\caption{Inventory items used in chapter 4. Nothing is wired and nothing is bought. The host script is chapter 3's with one new mode.}
\end{table}

\subsection*{System architecture}

\diagram{c04_arch}{Nothing on the byte path touches the processor. The engine writes, the counter records where it got to, and three different events cause the same function to run and work out what is new.}

The structure worth noticing is that the three events are not three code paths.
Half transfer, transfer complete and line idle all call one function, which
reads the counter, compares it with what it saw last time, and hands the new
span to the consumer. A design that handles the three separately acquires three
subtly different copies of the same arithmetic, and the one that is used least
is the one that is wrong.

\subsection*{Peripheral configuration}

\begin{table}[H]\centering\small
\begin{tabular}{p{22mm}p{26mm}p{24mm}p{30mm}p{30mm}}
\toprule
Peripheral & Mode & Clock source & Pins and function & Interrupt and transfers \\
\midrule
USART3 & Asynchronous receive, 8N1, receive request to the engine enabled & Peripheral bus & Confirm in the board manual & Line-idle interrupt only. No per-byte interrupt at all \\
Request multiplexer & USART3 receive request routed to the chosen stream & Peripheral bus & None & Request number read from RM0455 for this part \\
DMA1 stream 0 & Circular, peripheral to memory, byte width, memory address incrementing & Bus clock & None & Half transfer and transfer complete \\
Buffer memory & Main memory, not tightly coupled & Not applicable & None & Reachable by this engine: confirm in RM0455 \\
Data cache & On, which is the interesting case & Core clock & None & Invalidated by address over the span before each read \\
DWT cycle counter & Free running & Core clock & None & Measures the check function and the invalidate \\
\bottomrule
\end{tabular}
\caption{Peripheral configuration. Two rows are written as checks rather than as facts: the request number and the reachability of the chosen memory. Both are read from RM0455 for this device. A request number copied from the H743 lands on a different peripheral.}
\end{table}

\subsection*{Wiring}

\diagram{c04_wiring}{Nothing is wired. The only change from chapter 3's bench is on the host, which now sends frames of varied length with deliberate gaps between them, because a link that is never quiet never produces an idle condition and the chapter's main event would never fire.}

\subsection*{Memory and timing budget}

\diagram{c04_mem}{The buffer, the two positions, and the two spans a wraparound produces. Above the strip, the cache line grid that the invalidate operates on, which is why the buffer is aligned and its size is a whole number of lines.}

\begin{table}[H]\centering\small
\begin{tabular}{p{44mm}p{28mm}p{28mm}p{30mm}}
\toprule
Quantity & Budget & Measured & Margin \\
\midrule
Buffer size & 512 bytes & not applicable & not applicable \\
Cache line length & 32 bytes & from the architecture manual & not applicable \\
Longest frame that fits & 512 bytes & not applicable & not applicable \\
Idle detection delay & one character time & 86.8 \textmu{}s at 115200, computed & not applicable \\
Interrupts per second at 115200 & about 25 & not measured & not measured \\
Interrupts per second, chapter 3 & about 11520 & not measured & not measured \\
Cycles in the check function & 120 & not measured & not measured \\
Cycles in the invalidate & 60 & not measured & not measured \\
Last byte to consumer & 100 \textmu{}s & not measured & not measured \\
\bottomrule
\end{tabular}
\caption{The budget table. The two interrupt-rate rows are the headline of the chapter and the reason to build it, and both are estimates until the cycle counter and a counting build have produced them. The idle delay is arithmetic from the line rate and is marked as computed.}
\end{table}

\subsection*{Firmware design (UML)}

\diagram{c04_uml}{Three events, one function. The sequence shows the line-idle case, which is the one that delivers a short frame promptly; the other two are the same messages with a different trigger.}

The consumer contract is worth stating before the code. The check function
hands the consumer a pointer and a length into the receive buffer, not a copy.
That is deliberate: copying every byte a second time would give back a large
part of what the transfer engine just saved. It also means the consumer must
finish with the span before the engine laps the buffer and writes over it, which
is the constraint that replaces chapter 2's full-buffer test and is the subject
of step 7.

\subsection*{Data flow (ASCII)}

\begin{asciiart}
  USART3        DMA1 stream 0            rx_buf[512] in main memory
  +--------+    +---------------+        +---------------------------------+
  | RDR    |--->| circular      |------->| .....DATA DATA DATA.............|
  +--------+    | NDTR counts   |        +---------------------------------+
      |         | down from 512 |          ^              ^
      |         +---------------+          |              |
      |                 |               old_pos     pos = 512 - NDTR
      | IDLE            | HT, TC           |              |
      v                 v                  +---- new ----+
  +-------------------------------+                |
  | rx_check(): pos = 512 - NDTR  |                v
  |   if pos == old_pos: return   |   invalidate the cache over the span
  |   if pos >  old_pos: one span |                |
  |   else:              two spans|                v
  +-------------------------------+        consumer reads in place
\end{asciiart}

\subsection*{Repository layout}

\begin{plaincode}
nucleo-h7a3-dma-rx/
  CMakeLists.txt
  cmake/arm-none-eabi.cmake       # reused unchanged from chapter 1
  ld/stm32h7a3zi.ld               # gains one section for the buffer
  src/dma_rx.c                    # stream setup, the check function, the spans
  src/dma_rx.h
  src/cache.c                     # invalidate by address, aligned and argued
  src/usart3_ll.c                 # from chapter 3, minus the per-byte handler
  src/frames.c                    # the consumer: counts frames and lengths
  src/main.c
  host/framegen.py                # chapter 3's harness with a frame mode
  docs/compare_c03_c04.md         # the two paths, same harness, same board
  README.md
\end{plaincode}

\subsection*{Steps}

\begin{steps}
\item \textbf{Move the buffer and say why in the linker script.} A comment in
the file is the only place this decision survives a year. The alignment and the
size are both driven by the cache line length, which is read from the
architecture manual for this core rather than assumed.
\begin{ldcode}
/* The transfer engine writes here, so it cannot be tightly coupled memory:
   the main engines are not expected to reach DTCM on this part. Confirm the
   reachable regions in RM0455, not in a document for the H743. The buffer is
   also cache-line aligned and a whole number of lines long, so that an
   invalidate over it can never touch a line shared with anything else. */
SECTIONS
{
  .dma_rx (NOLOAD) : ALIGN(32) { KEEP(*(.dma_rx)) . = ALIGN(32); } > AXISRAM
}
\end{ldcode}
\begin{ccode}
#define RX_BUF_SIZE 512u                    /* a whole number of 32 B lines */
__attribute__((section(".dma_rx"), aligned(32)))
static volatile uint8_t rx_buf[RX_BUF_SIZE];
\end{ccode}

\item \textbf{Route the request and start the stream in circular mode.} Three
settings make this chapter work and each of them is the one people leave out.
Circular mode, so the engine restarts at the end instead of stopping. Memory
address incrementing and peripheral address fixed, because one end moves and the
other does not. And the receive request enabled in the peripheral, because the
engine waits to be asked.
\begin{ccode}
/* Request number read from RM0455 for this device. A number copied from a
 * sibling part selects a different peripheral and the transfer never moves. */
DMAMUX1_Channel0->CCR = (USART3_RX_DMA_REQUEST << DMAMUX_CxCR_DMAREQ_ID_Pos);

DMA1_Stream0->CR  = 0;                       /* stop before configuring    */
while (DMA1_Stream0->CR & DMA_SxCR_EN) { }   /* the disable is not instant */
DMA1_Stream0->PAR = (uint32_t) &USART3->RDR;
DMA1_Stream0->M0AR= (uint32_t) rx_buf;
DMA1_Stream0->NDTR= RX_BUF_SIZE;
DMA1_Stream0->CR  = DMA_SxCR_CIRC            /* wrap, do not stop          */
                  | DMA_SxCR_MINC            /* memory address advances    */
                  | DMA_SxCR_HTIE            /* half transfer interrupt    */
                  | DMA_SxCR_TCIE;           /* transfer complete too      */
DMA1_Stream0->CR |= DMA_SxCR_EN;
USART3->CR3 |= USART_CR3_DMAR;               /* now the peripheral asks    */
\end{ccode}
The wait loop on the enable bit is not decoration. Disabling a stream is a
request, not an instruction, and configuring registers while the bit is still
set produces a stream that runs with a mixture of the old and new settings.

\item \textbf{Enable the line-idle interrupt and nothing else on the
peripheral.} There is no receive interrupt any more. That is the point.
\begin{ccode}
USART3->ICR = USART_ICR_IDLECF;              /* clear a stale flag first   */
USART3->CR1 |= USART_CR1_IDLEIE;
NVIC_SetPriority(USART3_IRQn, 6);
NVIC_EnableIRQ(USART3_IRQn);
NVIC_SetPriority(DMA1_Stream0_IRQn, 6);      /* same level: they tail chain */
NVIC_EnableIRQ(DMA1_Stream0_IRQn);
\end{ccode}
Giving the two sources the same priority is deliberate. They never need to
preempt each other, they call the same function, and at equal priority two
pending interrupts tail chain rather than paying a full exit and entry, which
chapter 3's figure shows.

\item \textbf{Read the write position out of the counter.} This is the idea the
chapter turns on and it is one line.
\begin{ccode}
static inline uint16_t dma_write_pos(void)
{
    /* The counter counts down from the size, so the number of bytes written
     * since the last wrap is the size minus the counter. There is no head
     * index anywhere in this program: the head index is a register. */
    return (uint16_t)(RX_BUF_SIZE - (uint16_t) DMA1_Stream0->NDTR);
}
\end{ccode}
Read it once into a local. Reading it twice in one function gives you two
different answers, because the engine does not stop while you think.

\item \textbf{Handle the wraparound as two spans and not as a special case.}
The whole of the difficult part is nine lines, and it is difficult only until
it is written down.
\begin{ccode}
static uint16_t old_pos;

void rx_check(void)                     /* called from all three events */
{
    uint16_t pos = dma_write_pos();     /* read once, exactly once      */

    if (pos == old_pos)
        return;                         /* nothing new: the common case */

    if (pos > old_pos) {                /* one contiguous span          */
        rx_deliver(&rx_buf[old_pos], (uint16_t)(pos - old_pos));
    } else {                            /* wrapped: the tail, then the head */
        rx_deliver(&rx_buf[old_pos], (uint16_t)(RX_BUF_SIZE - old_pos));
        if (pos > 0u)
            rx_deliver(&rx_buf[0], pos);
    }
    old_pos = pos;
}
\end{ccode}
Note what is absent. There is no modulo, no mask and no test for the buffer
being full, because the engine does not ask permission. There is also no case in
which the consumer is handed a span that crosses the end of the array, which is
the property that lets it work on a plain pointer and length.

\item \textbf{Invalidate before you read, and know why the alignment matters.}
This is the step the canonical reference does not have, and on this part with
the data cache on it is the difference between a receive path that works and one
that delivers whatever the processor happened to cache earlier.
\begin{ccode}
static void rx_deliver(volatile uint8_t *p, uint16_t len)
{
    /* The engine wrote this span through the bus, not through the cache.
     * Anything the core already holds for these addresses is out of date,
     * so discard it before reading. The operation works on whole cache
     * lines; the buffer is aligned and a whole number of lines long, so a
     * line is never shared with a variable whose newer value would be
     * discarded with it. Chapter 19 argues this properly and gives two
     * other ways to arrive at the same guarantee. */
    SCB_InvalidateDCache_by_Addr((uint32_t *)(uintptr_t) p, (int32_t) len);
    consume((const uint8_t *) p, len);
}
\end{ccode}
The sharing hazard is worth one more sentence because it is the part that
produces a fault somewhere else entirely. An invalidate discards whole lines. If
the last line of the buffer also holds an ordinary variable, and the core had
written that variable but the write was still in the cache, the invalidate
discards it and the variable reverts. The alignment and the size are what make
that impossible, and they are declared in two places so that neither can drift
alone.

\item \textbf{Face the counter you cannot have.} Chapter 3 had an overrun flag:
the peripheral told you when a byte was lost. Here there is no such flag. If the
consumer falls behind, the engine simply writes over data that was never read,
and the counter arithmetic cannot tell a lap from no movement at all. That is a
real loss of visibility and the honest response is a guard rather than a claim.
\begin{ccode}
/* Two guards, neither of which is a detector, because a detector is not
 * available. The first bounds how long the consumer may be away; the second
 * reports how close the system came. */
uint16_t advance = (uint16_t)((pos - old_pos) & (RX_BUF_SIZE - 1u));
if (advance > stats.peak_advance)
    stats.peak_advance = advance;              /* publish this every second */
if (advance > (RX_BUF_SIZE / 2u))
    stats.near_lap++;                          /* half a buffer in one gap  */
\end{ccode}
The half transfer and transfer complete interrupts are the other half of the
guard. Between them they guarantee the function runs at least twice per lap of
the buffer whatever the traffic looks like, so a lap can only happen if the
consumer is blocked for longer than half a buffer's worth of line time. At
115200 with a 512 byte buffer that is about twenty two milliseconds, which is a
number to put in the README rather than a hope.

\item \textbf{Add the three handlers, which are all the same handler.} Clear the
flag that caused the entry, then call the one function.
\begin{ccode}
void USART3_IRQHandler(void)
{
    if (USART3->ISR & USART_ISR_IDLE) {
        USART3->ICR = USART_ICR_IDLECF;   /* clear, or it re-asserts */
        rx_check();
    }
}

void DMA1_Stream0_IRQHandler(void)
{
    uint32_t lisr = DMA1->LISR;           /* stream 0 flags live in LISR */
    if (lisr & DMA_LISR_HTIF0) { DMA1->LIFCR = DMA_LIFCR_CHTIF0; rx_check(); }
    if (lisr & DMA_LISR_TCIF0) { DMA1->LIFCR = DMA_LIFCR_CTCIF0; rx_check(); }
    if (lisr & DMA_LISR_TEIF0) { DMA1->LIFCR = DMA_LIFCR_CTEIF0;
                                 stats.transfer_errors++; }
}
\end{ccode}
The transfer error flag is in that handler for a reason. It is the one thing
that fires when the address you gave the engine is not reachable from it, which
is exactly the fault the note at the top of this chapter is about, and without
this branch it presents as silence.

\item \textbf{Send frames whose length varies, with gaps.} Chapter 3's harness
gains one mode. The gaps are what produce the idle condition, and a test that
sends a continuous stream never exercises the chapter's main event.
\begin{pycode}
def frames(dev, baud, lengths, gap_s=0.005, rounds=200):
    """Send frames of varied length with a quiet gap after each one.
    The gap must exceed one character time or no idle condition occurs."""
    with serial.Serial(dev, baud, timeout=0.5) as p:
        for _ in range(rounds):
            for n in lengths:
                p.write(bytes((i & 0xff) for i in range(n)))
                p.flush()
                time.sleep(gap_s)
        return status(p)          # ok, frames, lengths seen, near_lap
\end{pycode}
Use lengths that cross the interesting boundaries deliberately: one byte, the
cache line length, one byte more than a cache line, half the buffer, and one
byte less and one byte more than the whole buffer. The last two are the ones
that find a wraparound fault.

\item \textbf{Measure both paths with the same harness.} This is the
deliverable. Run chapter 3's build and this one against the same host script at
the same rates, and record interrupts per second, processor time spent in
interrupt context, and the delay from the last byte of a frame to the consumer
seeing it.
\begin{shellcode}
python host/framegen.py /dev/ttyACM0 --baud 115200 --mode frames \
  | tee docs/compare_c03_c04.md
\end{shellcode}
Expect the interrupt count to fall by roughly three orders of magnitude and the
latency for a short frame to rise by about one character time, because the idle
condition cannot be known until the line has been quiet for that long. Both
directions are the point: this path is not better, it is a different trade, and
the table is what lets a reader choose.
\end{steps}

\subsection*{Build, flash and debug}

\diagram{c04_timing}{Bytes arriving, the transfer counter stepping down as each one lands, and the line-idle flag setting once the line has been quiet for a character time. The counter is the only thing in this picture the processor reads.}

\begin{shellcode}
cmake -B build-fw -G Ninja -DCMAKE_TOOLCHAIN_FILE=cmake/arm-none-eabi.cmake
cmake --build build-fw -j && cp build-fw/firmware.bin "$PROBE_DISK"/   # onto the probe disk
\end{shellcode}

Three things are worth looking at on a halted target, in this order. The stream
enable bit, because a stream that was configured while it was still running does
not report that. The transfer counter, which should be somewhere between one and
the buffer size and never zero on a circular stream. And the buffer itself,
which tells you at a glance whether the engine is writing at all.

\begin{shellcode}
(gdb) print/x DMA1_Stream0->CR
(gdb) print DMA1_Stream0->NDTR
(gdb) x/64xb rx_buf
\end{shellcode}

\begin{note}{When the buffer stays empty and nothing reports an error}
In order of likelihood: the receive request was never enabled in the peripheral,
so the engine is armed and nobody is asking; the request number in the
multiplexer belongs to a different peripheral, which is what happens when it is
copied from a sibling part; the buffer is in a memory this engine cannot reach,
which the transfer error branch in step 8 turns from silence into a counter; or
the stream was configured while its enable bit was still set. Read the enable
bit and the counter before changing anything, because three of these four leave
a distinct signature in those two registers.
\end{note}

\subsection*{Verification and acceptance criteria}

\begin{itemize}
\item Every frame the host sends arrives whole and in order, for a length set
that includes one byte, one cache line, one more than a cache line, half the
buffer, one less than the buffer and one more than the buffer.
\item A frame that straddles the end of the buffer is delivered as two spans and
reassembled correctly by the consumer, proven by sending a length set whose
running total is chosen to land the boundary inside a frame.
\item With the data cache on, the delivered bytes match what the host sent. The
same test with the invalidate removed is run once, deliberately, and its failure
is recorded, because a cache fix that has never been seen to be necessary tends
to be removed later.
\item Interrupts per second at 115200 is measured in a counting build and is
compared in one table against chapter 3's figure for the same rate.
\item The delay from the last byte of a frame to the consumer is measured and is
within a character time of the computed idle delay.
\item Over a ten minute run at the highest rate chapter 3 found clean, the
near-lap counter stays at zero and the peak advance stays below half the buffer.
\item The transfer error counter is zero, and has been seen to be non-zero once
on purpose by pointing the engine at a memory it cannot reach.
\end{itemize}

\subsection*{Variants}

\begin{table}[H]\centering\small
\begin{tabular}{p{24mm}p{26mm}p{44mm}p{22mm}p{20mm}}
\toprule
Axis & Variant & What changes & Cost & Built in full in \\
\midrule
Execution model & DMA, circular, idle line & The baseline here. No per-byte interrupt, the write position comes from the counter, and frame boundaries come from the quiet line & Cache maintenance, and no overrun flag & Here \\
Execution model & DMA double buffered & Two buffers with the engine swapping between them, so the consumer always owns a whole one and the span arithmetic disappears & Twice the memory & Chapter 6 \\
Execution model & Interrupt into a ring buffer & Chapter 3's path. Lower latency on a short frame, far more interrupts, and an overrun flag you do not get here & About eleven thousand entries a second at 115200 & Chapters 2 and 3 \\
Execution model & Polled loop & Read the counter from the main loop with no interrupt at all, which works when the loop is fast and predictable & Latency is the loop period & Chapter 3 \\
Time and safety & Receiver timeout instead of the idle line & The peripheral counts a programmable number of idle bit times rather than one character, which lets the frame gap be tuned to the protocol & One more register to configure & Here, as the stretch goal \\
Synchronisation & A volatile flag & The check function sets a flag and the main loop does the delivery, moving work out of interrupt context & Latency of one loop & Chapter 3 \\
Operating system & RTOS task & The handler notifies a task which does the delivery and the cache maintenance & A scheduler & Chapter 20 \\
Language & Host in Python & The frame generator is chapter 3's harness with one mode added & None & Chapter 3 \\
\bottomrule
\end{tabular}
\caption{Variants for chapter 4. Only the first row is built in full here. The double buffered case belongs with the timer sampling of chapter 6, where the fixed block size makes it the natural shape rather than an extra buffer for its own sake.}
\end{table}

\subsection*{Pitfalls}

\begin{itemize}
\item Leaving the buffer in tightly coupled memory. The engine cannot reach it,
and without the transfer error branch the symptom is silence.
\item Copying the request number from a sibling part. It selects a different
peripheral and the engine waits for a request that never comes.
\item Configuring a stream while its enable bit is still set. Disabling is a
request that completes when the engine is ready, not when you wrote the zero.
\item Reading the transfer counter twice in one function. The two reads give
different answers and the arithmetic between them is wrong in a way that appears
once an hour.
\item Omitting the cache maintenance because the reference implementation does.
It is correct there because those examples run with the cache off or on parts
without one. It is not correct here.
\item Invalidating a span that shares a cache line with something else. The
invalidate discards whole lines, and the fault appears in the other variable.
\item Assuming an overrun flag exists. It does not on this path, and a program
that waits for one waits forever while data is quietly replaced.
\item Testing with a continuous stream. The line never goes quiet, the idle
condition never fires, and the chapter's main event is never exercised.
\end{itemize}

\subsection*{Best practices applied}

\begin{itemize}
\item The placement decision is recorded in the linker script, in a comment,
next to the thing it constrains.
\item The alignment and the size are both tied to the cache line length and are
declared in two places that cannot drift apart without a build failure.
\item Three events call one function, so there is one copy of the arithmetic.
\item The absence of an overrun flag is named and answered with a bounded guard
and a published worst case, rather than left as an assumption.
\item The cache fix is proven necessary once, on purpose, so that its removal
later is a visible regression rather than a tidy-up.
\item The comparison with the previous chapter uses the same harness, the same
board and the same rates, so it is a comparison.
\end{itemize}

\subsection*{Stretch goals}

\begin{itemize}
\item Replace the idle line with the peripheral's receiver timeout, which counts
a programmable number of idle bit times, and measure how short the inter-frame
gap can be before frames start merging. That number is a property of the
protocol you are about to design in chapter 5.
\item Add the zero-copy commit interface from the maintained ring buffer library
so the consumer can take a span and release it later, and measure whether the
extra bookkeeping costs more than the copy it removes.
\item Run the same path with the data cache off and measure the throughput and
the latency. The result feeds chapter 18's question about whether memory
placement matters more than anything else on this core.
\item Move the buffer into each memory the part offers in turn and record which
ones this engine reaches. That table does not exist for this device and it is
chapter 19's central artefact.
\end{itemize}

\subsection*{Roadmap and next steps}

Chapter 5 puts a frame format on this byte stream: a length prefix, byte
stuffing, and a checksum computed by the hardware unit. It takes the span this
chapter delivers and turns it into messages, and it is where the idle line stops
being the only thing that says a frame has ended.

Chapter 6 uses the same engine in its double buffered form for timer-triggered
sampling, where the block size is fixed and the wraparound arithmetic of this
chapter is unnecessary. Chapter 19 is where the cache argument used in step 6 is
made properly, with the region attributes and the two alternatives to
invalidating by address.

The canonical reference in the prior art table is the best single piece of
writing on this topic and is worth reading in full before and after this
chapter. The community roadmap referenced in appendix H puts transfer engines
and cache coherency in the firmware track and weights them heavily, and the
written series on this core's cache, listed in chapter 19's sources, is the
clearest teaching treatment of why step 6 exists.

\subsection*{Portfolio evidence}

\begin{itemize}
\item A repository in which the linker script, the alignment and the invalidate
are visibly one decision, with the reason written where the decision is.
\item The comparison table: interrupts per second, processor time in interrupt
context, and short-frame latency, for chapter 3's path and this one, on the same
board with the same harness and the full date of the run.
\item A recorded failure: the same build with the invalidate removed, showing
the delivered bytes disagreeing with what was sent, which is the evidence that
the fix is load bearing.
\item The near-lap and peak-advance counters over a ten minute run, with the
computed worst case next to them.
\end{itemize}

\subsection*{Sources}

Normative references:
\begin{itemize}
\item Reference manual RM0455, for this part: the transfer engine's registers
and flag groups, the request multiplexer numbering, the memories each engine can
reach, the line-idle flag and its clear bit, and the receiver timeout. Not
RM0433, which is its sibling and differs in the memory map.
\item The architecture reference manual and the core's technical reference, for
the cache line length and the invalidate-by-address operation.
\item The STM32H7A3xI datasheet, for the memory sizes and base addresses the
linker script places the buffer in.
\end{itemize}

Reusable implementations:
\begin{itemize}
\item Majerle's transfer-driven receive and transmit examples, MIT, the
canonical reference for this problem, with examples across eight families and
none for this exact board.\newline
\url{https://github.com/MaJerle/stm32-usart-uart-dma-rx-tx}
\item Majerle's lightweight ring buffer, MIT, for the zero-copy span
interface.\newline
\url{https://github.com/MaJerle/lwrb}
\item pyserial, BSD-3-Clause, for the frame generator on the host.\newline
\url{https://github.com/pyserial/pyserial}
\item Chapters 2, 3 and 19 of this volume, for the consumer, the harness and the
cache argument respectively.
\end{itemize}
   % Circular DMA and the idle line
\project{05}{Framing and the hardware CRC unit}{NUCLEO-H7A3ZI-Q}{Byte stuffing against a length prefix, CRC-16, and the peripheral's reversal settings}

\begin{keyfacts}
\item[Board] NUCLEO-H7A3ZI-Q, and nothing else
\item[Peripherals] The hardware checksum unit, USART3 with the receive path of chapter 4, the cycle counter
\item[Toolchain] arm-none-eabi-gcc with CMake on the target, Python on the host, and the same C files built for both
\item[Operating system] Bare metal
\item[Difficulty] 3 of 5
\item[Effort] 3 evenings of about four hours
\item[Deliverable] A frame format with a stated reason for every field, an encoder and decoder that build unchanged on the host, a 16-bit checksum computed by the peripheral that agrees with a software reference on every input tested, and a Python twin that corrupts frames on purpose and counts what gets through
\end{keyfacts}

\subsection*{Why this project}

The previous three chapters delivered bytes. Bytes are not messages. A message
has a beginning, an end, and some way of telling you that what arrived is what
was sent, and none of those three things exists in a serial byte stream. This
chapter supplies all three and then does the part that almost every project
skips: it proves that the checksum the silicon computes is the same number a
published reference computes, on every input that was tried, rather than on the
one input somebody tested first.

There are two honest ways to find the end of a frame and they trade against each
other in a way worth understanding before choosing. A length prefix is cheap and
obvious, and it fails badly: corrupt the length and the receiver skips an
arbitrary number of bytes, so a single bad byte costs several frames and there
is no reliable way back into step. Byte stuffing removes one value from the
payload so that value can be an unambiguous delimiter, which costs a pass over
the data and about one byte in every two hundred and fifty four, and buys
resynchronisation that is exactly as hard as looking for the next delimiter.
This chapter builds the second and analyses the first, because the analysis is
what makes the choice a decision.

The checksum half of the chapter is where the specific engineering is. This
family carries a programmable checksum peripheral, and the vendor's own
application note exists because getting it to agree with a standard reference is
a known trap with a trail of community questions behind it. Three settings have
to be right, and they are not the three you would guess: the polynomial and the
initial value, which are obvious, and the input and output bit reversal, which
are not. There is a fourth trap that no setting fixes, and it is the width of
the write you use to feed the data register.

\begin{note}{This chapter depends on the confirm list}
Item 18 asks whether this part has a checksum peripheral and how programmable it
is. The register layout drawn in this chapter's figure is read from RM0455 for
this device and not inherited from a sibling. If a check shows the unit here is
fixed rather than programmable, the chapter does not collapse: the software
reference in step 3 becomes the implementation, the measurement in step 9 becomes
a comparison of two software routines, and the finding is written down. What
must not happen is a chapter that assumes the feature and a reader who discovers
otherwise.
\end{note}

\subsection*{Prior art and what to reuse}

\begin{table}[H]\centering\small
\begin{tabular}{p{34mm}p{46mm}p{38mm}p{20mm}}
\toprule
Source & What it gives & What it does not & Licence \\
\midrule
Majerle's lightweight packet library & Close to this whole chapter in a box: framing, an address field, a command field, flags, a selectable checksum, and it is built on the ring buffer of chapter 2, which is a pleasing fit & It offers an 8-bit and a 32-bit checksum and \textbf{not a 16-bit one}, so the 16-bit variant and its routing to the peripheral is the work this chapter does & MIT \\
McQueen's byte stuffing library & The best implementation of consistent overhead byte stuffing there is, with the encode and decode both in a form you can vendor, and clear handling of the worst case expansion & No frame concept, no checksum and no state machine for receiving. It solves one problem completely and stops & MIT \\
The vendor's application note AN4187 & The reason this chapter exists: what each configuration bit of the checksum peripheral does, and how to set them so the result matches a standard reference. It is the document the community threads all end at & It is written for the family rather than for this device, so its register layout is checked against RM0455 before use & Vendor note \\
Chapters 2 and 4 of this book & The buffer the frames arrive in, and the span delivery that hands the decoder a pointer and a length & Nothing about frames & This volume \\
\bottomrule
\end{tabular}
\caption{Prior art for chapter 5. The first row would be the right dependency in a product. It is not used here because the 16-bit checksum it lacks is precisely the part that has to be built against the peripheral, and a chapter that vendored the library would have skipped its own subject.}
\end{table}

What is left to write is the 16-bit checksum in both forms, the configuration
that makes them agree, the receive state machine, and the Python twin that
corrupts frames deliberately. The byte stuffing is vendored and wrapped rather
than rewritten, because it is permissively licensed, complete, and there is
nothing to learn from typing it again.

\subsection*{Parts from the inventory}

\begin{table}[H]\centering\small
\begin{tabular}{p{40mm}p{66mm}p{40mm}}
\toprule
Part & Role & Interface \\
\midrule
NUCLEO-H7A3ZI-Q & The checksum peripheral, the decoder and the counters & Micro USB to the host \\
A USB data cable & Power, programming and the framed link & Micro USB \\
Host PC & The Python twin: encodes, decodes, and corrupts on purpose & Python, no extra packages beyond pyserial \\
\bottomrule
\end{tabular}
\caption{Inventory items used in chapter 5. Nothing is wired and nothing is bought. The interesting hardware in this chapter is a peripheral, not a part.}
\end{table}

\subsection*{System architecture}

\diagram{c05_arch}{The framing layer sits between the span delivery of chapter 4 and the application. The checksum peripheral is the only piece of silicon added, and it is used by both the encoder and the decoder.}

Two things about the shape are deliberate. The encoder and the decoder share the
checksum routine rather than each having one, because two implementations of a
checksum in one program is two chances to be wrong and only one of them gets
tested. And the peripheral is behind a two-function interface, so the software
reference can be substituted at build time, which is what makes the agreement
test in step 6 possible at all.

\subsection*{Peripheral configuration}

\begin{table}[H]\centering\small
\begin{tabular}{p{22mm}p{26mm}p{24mm}p{30mm}p{30mm}}
\toprule
Peripheral & Mode & Clock source & Pins and function & Interrupt and transfers \\
\midrule
Checksum unit & 16-bit polynomial, programmable initial value, reversal as the reference requires & Peripheral bus, clock enabled before any access & None & None. It is a register, not an event source \\
USART3 & As chapter 4 left it & Peripheral bus & Confirm in the board manual & Line idle only \\
DMA1 stream 0 & As chapter 4 left it & Bus clock & None & Half transfer and transfer complete \\
DWT cycle counter & Free running & Core clock & None & Measures the four checksum implementations \\
\bottomrule
\end{tabular}
\caption{Peripheral configuration. The checksum unit has no interrupt and no transfer capability: you write bytes to it and read a result. Its clock enable is the one thing that can make every configuration look correct and every result be zero.}
\end{table}

\subsection*{Wiring}

\diagram{c05_wiring}{Nothing is wired. The only structure this chapter adds to the link is the frame, so the frame is what the figure draws: the delimiter, the stuffing overhead byte, the payload and the checksum, with the sizes each field costs.}

\subsection*{Memory and timing budget}

\diagram{c05_mem}{The checksum peripheral's control register, and the two value registers beside it. Four settings decide whether the result matches a published reference, and two of them are the ones people do not think to check.}

\begin{table}[H]\centering\small
\begin{tabular}{p{44mm}p{28mm}p{28mm}p{30mm}}
\toprule
Quantity & Budget & Measured & Margin \\
\midrule
Frame overhead, stuffing & 1 byte per 254 & from the format, not measured & not applicable \\
Frame overhead, fixed & 4 bytes & from the format, not measured & not applicable \\
Checksum table in flash & 512 bytes & not measured & not measured \\
Cycles, bit at a time, 64 bytes & 3100 & not measured & not measured \\
Cycles, table driven, 64 bytes & 400 & not measured & not measured \\
Cycles, peripheral, byte writes & 280 & not measured & not measured \\
Cycles, peripheral, word writes & 90 & not measured & not measured \\
Decoder state, per link & 32 bytes & not measured & not measured \\
\bottomrule
\end{tabular}
\caption{The budget table. Every cycle figure is an estimate from the instruction counts until the cycle counter has run, and the two peripheral rows are the interesting comparison: the width of the write matters more than the fact that it is hardware at all.}
\end{table}

\subsection*{Firmware design (UML)}

\diagram{c05_uml}{The receive framer as a state machine. There are only two states and one of them does all the work, which is the property that makes resynchronisation after a corrupt frame trivial rather than an exercise.}

The state machine is small because the frame format was chosen to make it small.
A delimiter that cannot occur inside a frame means the decoder never has to
guess: any delimiter ends the current frame and starts the next one, and a
receiver that has just powered up in the middle of a transmission discards
whatever it has and is in step from the next delimiter onwards. Compare that
with the length-prefixed alternative, where recovering from a corrupt length
means guessing where the next frame might start and checking whether the
checksum happens to pass, which is a guess with a one in sixty five thousand
chance of being wrong each time you try it.

\subsection*{Data flow (ASCII)}

\begin{asciiart}
  transmit                                      receive
  payload  [ 41 00 42 7E ]                      span from chapter 4
     |                                             |
     v  append CRC-16 over the payload             v  bytes until a delimiter
  [ 41 00 42 7E ][ B1 29 ]                      [ 03 41 03 42 7E B1 29 ]
     |                                             |
     v  byte stuffing removes every 00             v  unstuff
  [ 03 41 03 42 7E B1 29 ]                      [ 41 00 42 7E ][ B1 29 ]
     |                                             |
     v  add the delimiter                          v  CRC-16 over the payload
  00 [ 03 41 03 42 7E B1 29 ] 00                compare, accept or count
     |                                             |
     +----------------> USART3 ------------------->+
\end{asciiart}

\subsection*{Repository layout}

\begin{plaincode}
nucleo-h7a3-framing/
  CMakeLists.txt                  # -DCRC_IMPL=hw|table|bitwise
  cmake/arm-none-eabi.cmake       # reused unchanged from chapter 1
  src/crc16.c, src/crc16.h        # the software reference, host and target
  src/crc16_hw.c                  # the peripheral, same two-function interface
  src/cobs.c, src/cobs.h          # vendored, MIT, copyright header kept
  src/frame.c, src/frame.h        # encode, and the two-state decoder
  src/main.c
  test/test_crc_agree.c           # host: reference against known vectors
  test/test_frame_roundtrip.c     # host: encode then decode, random payloads
  host/twin.py                    # the Python twin, and the bit flipper
  docs/crc_settings.md            # the four settings and why each one is that
  README.md
\end{plaincode}

\subsection*{Steps}

\begin{steps}
\item \textbf{Choose the frame shape and record the reason.} Two candidates,
and the table belongs in the repository as well as here, because the next
person to touch this will otherwise redo the analysis or, worse, change the
format without it.

\begin{table}[H]\centering\small
\begin{tabular}{p{30mm}p{48mm}p{48mm}}
\toprule
Property & Length prefix & Byte stuffing with a delimiter \\
\midrule
Overhead & Fixed, 2 bytes & About 1 byte in 254, plus 1 \\
Finding the end & Trust the length & Look for the delimiter \\
After one corrupt byte & The length may be wrong, so an arbitrary number of following frames are lost & The next delimiter re-synchronises, always \\
Joining a stream late & Cannot be done reliably & Discard to the next delimiter \\
Cost & None & One pass over the payload each way \\
Knowing the size early & Yes, which suits a fixed buffer & No, which suits a stream \\
\bottomrule
\end{tabular}
\caption{The framing decision for chapter 5. The deciding row is the third one: on a link that shares a cable with a debugger and is unplugged regularly, recovering in step matters more than two bytes.}
\end{table}

\item \textbf{Vendor the byte stuffing and wrap it.} It is permissively
licensed and complete. Keep the copyright header, record the version and the
date in the README, and write a thin wrapper so that the rest of the program
calls your names and not a third party's.
\begin{ccode}
/* cobs.c is vendored unchanged. Upstream is MIT licensed; the copyright
 * header at the top of that file stays there. This wrapper exists so that
 * frame.c never names a third party symbol directly, which makes the
 * dependency one file to replace rather than a search across the tree. */
size_t frame_stuff(const uint8_t *in, size_t n, uint8_t *out, size_t cap);
size_t frame_unstuff(const uint8_t *in, size_t n, uint8_t *out, size_t cap);
\end{ccode}

\item \textbf{Write the software reference and pin it to a published check
value.} The reference is the arbiter for everything that follows. Its check
value is the checksum of the nine characters \texttt{123456789}, which every
published catalogue of these polynomials lists, and it is the single most useful
line in the test file.
\begin{ccode}
/* CRC-16 with polynomial 0x1021, initial value 0xFFFF, no input or output
 * reversal, no final exclusive-or. The published check value over the nine
 * bytes "123456789" is 0x29B1. If your implementation does not produce
 * that, nothing further in this chapter is worth doing. */
uint16_t crc16_ref(const uint8_t *d, size_t n)
{
    uint16_t c = 0xFFFFu;
    while (n--) {
        c ^= (uint16_t)((uint16_t) *d++ << 8);
        for (int i = 0; i < 8; i++)
            c = (c & 0x8000u) ? (uint16_t)((c << 1) ^ 0x1021u)
                              : (uint16_t)(c << 1);
    }
    return c;
}
\end{ccode}

\item \textbf{Configure the peripheral one setting at a time.} Four settings
and each one has a reason. Enable the peripheral clock first: without it every
register reads back as it was written in your imagination and the result is
always zero.
\begin{ccode}
void crc16_hw_init(void)
{
    RCC->AHB4ENR |= RCC_AHB4ENR_CRCEN;    /* the clock, before anything   */
    (void) RCC->AHB4ENR;                  /* read back: the write posts   */

    CRC->POL  = 0x1021u;                  /* the polynomial               */
    CRC->INIT = 0x0000FFFFu;              /* the initial value            */
    CRC->CR   = CRC_CR_POLYSIZE_0         /* 01: a 16-bit polynomial      */
              | CRC_CR_RESET;             /* load INIT, start clean       */
    /* REV_IN stays 00 and REV_OUT stays 0: this reference reverses
     * neither. A reference that does reverse needs 01 or 11 here, and
     * that single choice is what the application note is about. */
}
\end{ccode}
The read back after enabling the clock is not superstition. The write to the
enable register is posted on the bus, and an access to the peripheral issued in
the next instruction can arrive before the clock does. The read forces the write
to complete and costs a handful of cycles once at startup.

\item \textbf{Feed it one byte at a time and know why the write width
matters.} This is the trap that no configuration bit fixes and that produces a
result which is wrong in a way that looks almost right.
\begin{ccode}
uint16_t crc16_hw(const uint8_t *d, size_t n)
{
    CRC->CR |= CRC_CR_RESET;              /* reload INIT for this message */

    /* An 8-bit write feeds exactly one byte. Writing the same value as a
     * 32-bit word feeds four bytes, three of them zero, and gives a
     * different and entirely plausible looking answer. The cast is the
     * whole point of this line and deserves the comment. */
    volatile uint8_t *dr8 = (volatile uint8_t *) &CRC->DR;
    while (n--)
        *dr8 = *d++;

    return (uint16_t) CRC->DR;            /* low half for a 16-bit poly   */
}
\end{ccode}
Once agreement is established, add the faster path: feed four bytes at a time
with 32-bit writes for the aligned middle of the buffer and bytes for the ends.
Do it second, not first, because the byte order within a word is one more thing
that has to agree, and adding two unknowns at once makes neither findable.

\item \textbf{Prove agreement rather than assuming it.} One test, on the host,
comparing the reference against the peripheral is impossible, because the
peripheral is on the target. So the test runs in two halves: the host proves the
reference against the published vectors, and the target proves the peripheral
against the reference over a large input space and reports a single number.
\begin{ccode}
/* On the target. Both implementations are compiled in for this build. */
uint32_t mismatches = 0;
for (uint32_t len = 0; len <= 300u; len++) {
    for (uint32_t trial = 0; trial < 64u; trial++) {
        fill_pseudorandom(buf, len, len * 1000u + trial);
        if (crc16_ref(buf, len) != crc16_hw(buf, len))
            mismatches++;
    }
}
printf("crc agree: %lu mismatches over %lu inputs\n",
       (unsigned long) mismatches, 301UL * 64UL);
\end{ccode}
Lengths from zero upwards matter. The zero-length case is where an
implementation that resets in the wrong place shows itself, and lengths around
multiples of four are where the word-feeding path in the previous step will fail
if the byte order is wrong.

\item \textbf{Write the decoder as two states.} It takes the span that chapter
4 delivers, which may be any length and may arrive in pieces, and turns it into
whole frames.
\begin{ccode}
void frame_feed(const uint8_t *p, size_t n)     /* from rx_deliver */
{
    while (n--) {
        uint8_t b = *p++;
        if (b == FRAME_DELIM) {
            if (rx.len > 0u)
                frame_complete();               /* may accept or count   */
            rx.len = 0u;                        /* either way, start over */
            continue;
        }
        if (rx.len < sizeof rx.buf) {
            rx.buf[rx.len++] = b;
        } else {
            rx.len = 0u;                        /* too long: drop it     */
            stats.oversize++;                   /* and say so            */
            rx.dropping = 1u;
        }
    }
}
\end{ccode}
The oversize counter is there for the same reason chapter 2's drop counter is.
A decoder that silently discards a frame longer than its buffer looks exactly
like a link that works until somebody sends a longer message.

\item \textbf{Complete the frame and be strict about what you accept.} Three
checks in order, cheapest first, and each with its own counter so that a link
that is failing tells you how.
\begin{ccode}
static void frame_complete(void)
{
    uint8_t out[FRAME_MAX];
    size_t  m = frame_unstuff(rx.buf, rx.len, out, sizeof out);

    if (m == 0u)               { stats.bad_stuffing++; return; }
    if (m < 2u)                { stats.too_short++;    return; }

    size_t   plen = m - 2u;
    uint16_t want = (uint16_t) out[plen] | (uint16_t)(out[plen + 1u] << 8);
    uint16_t got  = crc16(out, plen);       /* hw or reference, one name */

    if (want != got)           { stats.bad_crc++;      return; }

    stats.accepted++;
    app_on_frame(out, plen);
}
\end{ccode}
Note the single name \texttt{crc16}. Which implementation it resolves to is a
build switch, which is what lets the same decoder run on the host in the test
and on the target with the peripheral.

\item \textbf{Write the Python twin and make it corrupt frames on purpose.}
The twin is the second opinion. A bug that exists in both the C and the Python
is possible but unlikely; a bug in one of them is found the first time they
disagree.
\begin{pycode}
#!/usr/bin/env python3
"""twin.py: encode, decode and deliberately corrupt frames."""

DELIM = 0x00

def crc16(data: bytes, init: int = 0xFFFF, poly: int = 0x1021) -> int:
    c = init
    for b in data:
        c ^= b << 8
        for _ in range(8):
            c = ((c << 1) ^ poly) & 0xFFFF if c & 0x8000 else (c << 1) & 0xFFFF
    return c

assert crc16(b"123456789") == 0x29B1, "the reference itself is wrong"

def encode(payload: bytes) -> bytes:
    body = payload + crc16(payload).to_bytes(2, "little")
    return bytes([DELIM]) + cobs_encode(body) + bytes([DELIM])

def flip(frame: bytes, bit: int) -> bytes:
    """Invert one bit of a frame, counting bits from the first byte."""
    b = bytearray(frame)
    b[bit // 8] ^= 1 << (bit % 8)
    return bytes(b)

def sweep(payload: bytes):
    """Flip every single bit of a frame in turn and report what the
    receiver did with each one. Nothing here should ever be accepted."""
    frame = encode(payload)
    for bit in range(len(frame) * 8):
        send(flip(frame, bit))
    return read_counters()          # accepted must still be zero
\end{pycode}
The single-bit sweep is the acceptance test. Every frame in it is corrupt, and
the accepted counter on the target must not move. Then run a two-bit sweep,
which is a longer job and a more interesting result, and finally a random
corruption run that flips a random number of bits in random positions. The last
one is where a 16-bit checksum starts to let frames through, and the rate at
which it does is a measurement to report rather than a property to claim.

\item \textbf{Measure the four implementations.} Bit at a time, table driven,
peripheral with byte writes, peripheral with word writes, over a 64 byte
payload, with the cycle counter.
\begin{ccode}
for (int impl = 0; impl < 4; impl++) {
    crc16_select(impl);
    uint32_t t0 = DWT->CYCCNT;
    for (int i = 0; i < 1000; i++) (void) crc16(payload, 64);
    printf("%s: %lu cycles per frame\n", crc16_name(impl),
           (unsigned long) ((DWT->CYCCNT - t0) / 1000u));
}
\end{ccode}
The result worth writing up is the comparison between the two peripheral rows
rather than the comparison between hardware and software. A peripheral fed one
byte at a time spends most of its advantage on bus accesses, and the interesting
question is how much of that a wider write recovers. Report the table's 512
bytes of flash alongside, because that is the other side of its speed.
\end{steps}

\subsection*{Build, flash and debug}

\diagram{c05_timing}{Where a frame's time goes, for a 64 byte payload. Every bar is an estimate from the instruction counts until the cycle counter has run, and the figure says so rather than presenting estimates as results.}

\begin{shellcode}
cmake -B build-fw -G Ninja -DCMAKE_TOOLCHAIN_FILE=cmake/arm-none-eabi.cmake -DCRC_IMPL=hw
cmake --build build-fw -j && cp build-fw/firmware.bin "$PROBE_DISK"/   # onto the probe disk
make -C test && ./test/test_crc_agree && ./test/test_frame_roundtrip
python host/twin.py /dev/ttyACM0 --sweep single
\end{shellcode}

\begin{note}{When the checksum is always zero or always the same value}
In order of likelihood: the peripheral clock was never enabled, so every write
lands nowhere and every read returns the same thing; the data register was
written as a 32-bit word rather than a byte, which gives a stable wrong answer
that changes with the data and therefore looks plausible; the reset bit was not
set between messages, so the second frame continues the first one's computation;
or the initial value was written after the reset rather than before it. Read the
control register and the polynomial register back on a halted target before
changing any code.
\end{note}

\begin{shellcode}
(gdb) print/x CRC->CR
(gdb) print/x CRC->POL
(gdb) print/x CRC->INIT
\end{shellcode}

\subsection*{Verification and acceptance criteria}

\begin{itemize}
\item The software reference produces 0x29B1 over the nine bytes
\texttt{123456789}, checked by an assertion in both the C test and the Python
twin, so that neither can drift alone.
\item The peripheral and the reference agree on every input in the sweep:
lengths zero to three hundred, sixty four pseudorandom inputs at each length,
zero mismatches reported as one number.
\item The word-feeding path agrees with the byte-feeding path for every length,
including lengths that are not a multiple of four.
\item A round trip on the host, encode then decode, returns the original
payload for ten thousand random payloads including empty ones, payloads that are
all delimiter bytes, and payloads at the maximum length.
\item The single-bit sweep over a whole frame produces zero accepted frames and
a bad-checksum counter equal to the number of bits flipped, minus any flips that
the stuffing rejects first, which are counted separately.
\item A random corruption run reports the rate at which corrupt frames are
accepted, as a measured number with the number of trials beside it.
\item The four implementations are timed with the cycle counter and the table's
flash cost is reported next to its speed.
\end{itemize}

\subsection*{Variants}

\begin{table}[H]\centering\small
\begin{tabular}{p{24mm}p{26mm}p{44mm}p{22mm}p{20mm}}
\toprule
Axis & Variant & What changes & Cost & Built in full in \\
\midrule
Framing & Byte stuffing with a delimiter & The baseline here. Resynchronisation is finding the next delimiter & One pass each way & Here \\
Framing & Length prefix & Two bytes and no pass over the data, and no reliable recovery from a corrupt length & Frames lost after one bad byte & Here, as the analysis in step 1 \\
Checksum & Peripheral, 16-bit & The baseline here, with the four settings that make it agree with the reference & Four settings that must be right & Here \\
Checksum & Software, table driven & 512 bytes of flash buys most of the speed with none of the configuration & Flash & Here \\
Checksum & Software, bit at a time & No flash, no peripheral, and the slowest by a wide margin. The arbiter for the other three & Cycles & Here \\
Checksum & 8-bit or 32-bit & What the packet library in the prior art table already offers, which is why this chapter built the one it does not & None & Named, not built \\
Peripheral substitution & Checksum in a transfer & Feeding the unit from a transfer engine rather than by hand, which only pays on long payloads & An engine and a stream & Chapter 17 \\
Language & C++ & The frame type as a template over the payload size, with the maximum stuffed size computed at compile time & Nothing at run time & Chapter 9 \\
Language & Host in Python & The twin, the bit flipper and the sweeps & None & Chapter 3 \\
Intelligence and reach & Edge host over a framed link & This frame format is what chapters 9 and 11 send over & None here & Chapters 9 and 11 \\
\bottomrule
\end{tabular}
\caption{Variants for chapter 5. Three checksum implementations are built in full here rather than one, because the chapter's subject is that they agree, and agreement needs more than one.}
\end{table}

\subsection*{Pitfalls}

\begin{itemize}
\item Writing the data register as a word when you mean a byte. The answer is
stable, depends on the data, and is wrong.
\item Forgetting the peripheral clock. Every register reads back as zero and
every result is the same value, which reads as a configuration problem and is
not.
\item Setting the initial value after the reset bit rather than before it. The
reset loads whatever is in the initial value register at that moment.
\item Assuming the reversal settings are zero because the reference did not
mention them. Most published 16-bit checksums reverse both input and output, and
the one used here is chosen partly because it does not, which keeps the first
agreement test simple.
\item Reusing the unit between two interrupt priorities without protecting it.
It has one accumulator, and a handler that computes a checksum in the middle of
a longer one returns a corrupted result to both.
\item Trusting a length field. One corrupt byte in the length costs an
unbounded number of following frames.
\item Reporting that a checksum detects all errors. It detects a large fraction,
the twin measures which, and the number belongs in the README.
\item Two checksum implementations in one program, one for sending and one for
receiving. Only one of them ends up tested.
\end{itemize}

\subsection*{Best practices applied}

\begin{itemize}
\item The choice of framing is a table with a deciding row, kept in the
repository, so that the next change to the format is a decision and not a
preference.
\item Every implementation of the checksum is behind one name, chosen at build
time, which is what makes them comparable and testable against each other.
\item Agreement is proven over a range of inputs and reported as a single
number, not demonstrated on one example.
\item The vendored library keeps its copyright header and its version is
recorded, and it is reached through a wrapper so the dependency is one file.
\item Every rejection has its own counter, so a failing link says how it is
failing rather than only that it is.
\item The corruption rate that gets through is measured and published rather
than asserted from a property of the polynomial.
\end{itemize}

\subsection*{Stretch goals}

\begin{itemize}
\item Add the reversed configuration and reproduce a second published checksum
with the same peripheral, changing only the two reversal fields. That exercise
is the application note's whole content and it takes twenty minutes once the
first one agrees.
\item Feed the unit from a transfer engine over a long payload and find the
length at which it beats the word-writing loop, which is a crossing point nobody
has published for this part.
\item Extend the twin into a property test that generates payloads and
corruptions from a specification, and run it as part of the build in the way
chapter 12 runs its rig.
\item Measure the undetected error rate against payload length for this
polynomial on this format, and publish the curve. It is a small piece of real
work and the answer is specific to the frame shape rather than to the
polynomial alone.
\end{itemize}

\subsection*{Roadmap and next steps}

Chapter 6 leaves the serial link and starts sampling on a timer, where the
frames built here become the way a block of samples leaves the board. Chapter 9
puts a payload codec inside this frame and gives it a Python twin of its own,
built on the pattern established here. Chapter 11 sends these frames to a modem,
where the reason for choosing resynchronisation over two bytes of overhead
becomes concrete.

The vendor's application note in the prior art table is the document to read
alongside this chapter, and the packet library named there is the right
dependency for a product that needs addressing and a command field rather than a
bare payload. The community roadmap referenced in appendix H places protocol
work in the firmware track and pairs it with the testing material that chapter
12 draws on; its argument that projects teach more than reading is the reason
this chapter ends with a corruption sweep rather than a summary of checksum
theory.

\subsection*{Portfolio evidence}

\begin{itemize}
\item A repository containing a frame format, a settings document that explains
each of the four peripheral settings, and a test that fails if the software
reference drifts from its published check value.
\item The agreement result: zero mismatches over the full input sweep, printed
by the target itself, with the full date of the run.
\item The corruption sweep output: single-bit, double-bit and random, with the
accepted count for each and the number of trials beside it.
\item The four-way timing table with the table implementation's flash cost next
to its speed, naming the cycle counter as the instrument.
\end{itemize}

\subsection*{Sources}

Normative references:
\begin{itemize}
\item Reference manual RM0455, for this part: the checksum peripheral's
registers, the meaning of each control field, the polynomial size encoding and
the reset behaviour. Not RM0433, which is its sibling.
\item The vendor's application note AN4187, on configuring the checksum
peripheral so that its result matches a standard reference, which is the
document behind every community thread on this subject.
\item The published catalogue of checksum polynomials and their check values,
for the constant 0x29B1 that pins the software reference.
\end{itemize}

Reusable implementations:
\begin{itemize}
\item McQueen's byte stuffing library, MIT, vendored with its copyright header
intact.\newline
\url{https://github.com/cmcqueen/cobs-c}
\item Majerle's lightweight packet library, MIT, which is the right dependency
for a product and which offers 8-bit and 32-bit checksums but not the 16-bit one
this chapter builds.\newline
\url{https://github.com/MaJerle/lwpkt}
\item Majerle's lightweight ring buffer, MIT, which that packet library is built
on and which chapter 2 of this volume re-derives.\newline
\url{https://github.com/MaJerle/lwrb}
\item pyserial, BSD-3-Clause, under the Python twin.\newline
\url{https://github.com/pyserial/pyserial}
\end{itemize}
   % Framing and the hardware CRC unit
\project{06}{Sampling on a timer at exactly 1 kHz}{NUCLEO-H7A3ZI-Q with an MCC 118 on a Raspberry Pi}{Timer-triggered acquisition, proving the rate externally}

\begin{keyfacts}
\item[Board] NUCLEO-H7A3ZI-Q, plus a Raspberry Pi 4 carrying an MCC 118 DAQ HAT as the external witness
\item[Peripherals] TIM6 as the trigger, ADC1 with hardware trigger, DMA in double-buffer mode, TIM2 and TIM5 as instruments, one free GPIO as the marker
\item[Toolchain] arm-none-eabi-gcc with CMake as in chapter 1; Python 3 with the HAT library on the Pi
\item[Operating system] Bare metal on the board, Raspberry Pi OS on the witness
\item[Difficulty] 4 of 5
\item[Effort] 4 evenings of about four hours, one of them entirely on the Pi
\item[Deliverable] Three sampling builds that share one application, an external measurement that states the sample rate with an uncertainty, and the two instruments the rest of this book measures with
\end{keyfacts}

\subsection*{Why this project}

Every chapter after this one makes a claim about time. Chapter 7 claims a wake
takes so many milliseconds, chapter 10 splits a duty cycle into phases, chapter
18 counts cycles per transform and chapter 20 reports a latency distribution.
None of those claims mean anything until there is an instrument, and this bench
has no oscilloscope and no logic analyser. So this chapter does two jobs at
once. It builds a sampling front end that acquires at a fixed rate while the
processor is busy elsewhere, and it builds the instruments: a cycle counter
wrapper with a fallback, and a gated frequency counter that can be pointed at
any periodic signal on the board.

The second job is the reason the chapter is worth four evenings rather than
one. A sample rate that is configured is not a sample rate that is proven. The
usual practice is to write the prescaler and the reload value, do the
arithmetic, and call the result 1 kHz. That arithmetic is worth doing and this
chapter does it, but it establishes only that the divider chain is exact. It
says nothing about the oscillator the chain divides, nothing about whether the
interrupt actually fires when the timer says it does, and nothing about jitter.
The only honest way to settle those is to put a second, independent clock next
to the first and compare them. The MCC 118 on a Raspberry Pi is that second
clock: it samples a pin the board toggles from its own sampling interrupt, at a
rate high enough that each edge lands between two known instants.

The third reason is structural. A sampling loop that does its work inside the
interrupt looks correct at 1 kHz and fails silently at 10 kHz, because the
overrun has nowhere to be reported. Decoupling acquisition from processing is
the pattern the rest of Part I and all of Part II depend on, and this chapter
builds it three ways and measures the difference.

\begin{note}{This part is not the STM32H7 the internet is about}
The timer inventory, the transfer request numbers that connect the converter to
a transfer stream, and the rule about which transfer engine reaches which memory
all differ between RM0455, which documents this part, and RM0433, which
documents its better-known sibling. The prior art named below was written on a
Cortex-M7 Nucleo of the other kind. Its structure transfers; its request
numbers do not. Read them out of RM0455's request table for this part, not out
of a working project for another one.
\end{note}

\subsection*{Prior art and what to reuse}

\begin{table}[H]\centering\small
\begin{tabular}{p{32mm}p{46mm}p{40mm}p{20mm}}
\toprule
Source & What it gives & What it does not & Licence \\
\midrule
Ian Ross's five-step progression on a Cortex-M7 Nucleo & Exactly this chapter's shape: the same acquisition solved in rising order of sophistication, with a written series explaining each step's reason & It is small, unmaintained, and written for a different member of this family. No double buffering, no external verification, no instrument & MIT \\
The vendor's basic examples for this board & Working register settings for a timer trigger into the converter, in a tree that also solves the same job polled and on interrupt & Its example projects sit under a licence tied to this vendor's silicon; only the basic examples are permissive. No verification of the resulting rate & BSD-3-Clause for the basic examples \\
Majerle's transfer-driven serial work & The half-transfer and transfer-complete pattern this chapter reuses for the converter, proven across eight families & It is a serial receive problem, not acquisition, and it does no cache maintenance & MIT \\
GorgonMeducer's performance counter & A cycle-accounting library that deliberately runs on the system tick rather than the core's trace counter, because the tick is universally present & It is a whole library where this chapter needs a twenty-line wrapper with a self-test & Apache-2.0 \\
The HAT library for the acquisition board & A Python interface to the external witness with a scan API and published examples & Its licence is not recorded in this volume's prior-art pool. Read the repository's licence file before vendoring any of it. This chapter only calls it & Read before use \\
\bottomrule
\end{tabular}
\caption{Prior art for chapter 6. The five-step progression is the closest thing to this chapter that exists, and its written series is worth reading before you start; the part it is written for is not this one.}
\end{table}

What is left to write is the verification and the instruments. No source in the
pool measures the rate it configures. The external witness method below, the
edge-interpolation arithmetic that turns a sampled square wave into a rate with
an uncertainty, the cycle counter wrapper with a runtime self-test and a timer
fallback, and the gated counter and its reciprocal upgrade are all the author's
work, and they are the reason this chapter comes before the seven that cite it.

\subsection*{Parts from the inventory}

\begin{table}[H]\centering\small
\begin{tabular}{p{40mm}p{66mm}p{40mm}}
\toprule
Part & Role & Interface \\
\midrule
NUCLEO-H7A3ZI-Q & The subject. Samples on a timer and toggles a marker pin from the sampling interrupt & Micro USB to the host \\
Raspberry Pi 4 & The witness. Runs the scan and the analysis & Ethernet or wireless to the host \\
MCC 118 DAQ HAT & The calibrated reference, 8 channels, 12-bit, 100 kS/s aggregate & 40-pin header on the Pi, screw terminals to the board \\
Two jumper wires & Marker pin and a common ground reference & Screw terminal at one end, Zio header at the other \\
A 220 \textohm{} resistor & In series with the marker pin, so a wiring error costs nothing & Breadboard \\
\bottomrule
\end{tabular}
\caption{Inventory items used in chapter 6. Nothing is bought and nothing is soldered. Whether male-to-male jumper wires exist in the drawer is item 36 on the confirm list; the screw terminals on the HAT accept a bare wire end either way.}
\end{table}

\subsection*{System architecture}

\diagram{c06_arch}{Two independent clocks, one signal between them. The board's divider chain produces the sample instants; the witness has its own oscillator and never sees the board's. What the comparison proves is agreement, which is a stronger and more modest claim than accuracy.}

The board's side is a chain with no software in it: a timer reaches its reload
value, emits a trigger event, the converter starts on that event, and the
transfer engine moves the result to memory. The processor is told only when a
half of the buffer is complete, and between those notifications it is free. The
witness is deliberately dumb: it samples one channel continuously at a high
rate, writes the samples to a file with a start timestamp, and the analysis
happens afterwards on the same Pi. Nothing on it is synchronised to the board,
which is the point.

\subsection*{Peripheral configuration}

\begin{table}[H]\centering\small
\begin{tabular}{p{20mm}p{28mm}p{24mm}p{30mm}p{30mm}}
\toprule
Peripheral & Mode & Clock source & Pins and function & Interrupt and transfers \\
\midrule
TIM6 & Up counter, update event on TRGO & APB1 timer clock & None & No interrupt in the timer build \\
ADC1 & Single channel, hardware trigger on the timer's TRGO & ADC kernel clock, confirm the source in RM0455 & One analogue input, channel confirmed in the datasheet & Transfer request to a DMA stream \\
DMA1 stream & Double buffer, two memory pointers & Bus clock & None & Half and full transfer interrupts \\
GPIO marker & Push-pull output, toggled in the sampling interrupt & Peripheral bus & A free Zio pin, item 13 on the confirm list & None \\
TIM2 & Free running 32-bit, or input capture for the reciprocal counter & APB1 timer clock & Capture pin from the alternate-function table & Capture interrupt, reciprocal build only \\
TIM5 and TIM7 & External clock mode 2 counting edges, and an up counter as the gate & The signal under test, and the APB1 timer clock & External trigger pin from the alternate-function table & Gate interrupt only \\
SysTick & 1 kHz reload, the simplest build & Core clock & None & SysTick exception \\
\bottomrule
\end{tabular}
\caption{Peripheral configuration. Two 32-bit timers exist on this part and both are spent on instruments rather than on the application, which is a deliberate allocation and is worth stating in the repository README.}
\end{table}

\subsection*{Wiring}

\diagram{c06_wiring}{One signal wire and one ground wire. The series resistor is not protection for the board, which drives the pin, but insurance against the wire finding the wrong terminal.}

Two safety points before the wires go in. The board's logic is 3.3 V and
nothing from the Pi or the HAT may be driven back into a board pin; the marker
is an output on the board and an input on the HAT, and it stays that way. The
grounds must be common or the witness measures an undefined potential, so the
ground wire is not optional and goes in first. The HAT's input range and
whether it offers an external trigger are item 35 on the confirm list: read the
HAT's user guide before assuming a 3.3 V swing is comfortably inside the range,
even though every range this HAT is likely to offer contains it.

\subsection*{Memory and timing budget}

\diagram{c06_mem}{Where 1 kHz actually comes from. The divider chain is exact, and that is a different claim from accurate. Everything to the left of the timer is inherited from chapter 1's clock configuration and its accuracy is the probe's.}

The arithmetic is worth writing out because it is the part that is exact. With
the peripheral bus prescaler set so the timer kernel clock is 280 MHz, a
prescaler value of 279 divides by 280 and gives a 1 MHz counter. A reload value
of 999 counts 1000 of those and gives an update event every millisecond. Both
divisions are integer, so there is no residue and no slow drift from rounding.
The same is true of the tick: a reload of 279999 on a 280 MHz core clock is
exactly one millisecond and fits in the tick's 24 bits with room to spare.

What is not exact is the 8 MHz the whole chain starts from, which comes from the
probe and carries the probe's tolerance. That number is not in this volume and
is not guessed at here. It is what the external measurement below reports.

\begin{table}[H]\centering\small
\begin{tabular}{p{44mm}p{28mm}p{28mm}p{30mm}}
\toprule
Quantity & Budget & Measured & Margin \\
\midrule
Sample period & 1.000000 ms by construction & not measured & not measured \\
Interrupt to marker edge & under 2 \textmu{}s & not measured & not measured \\
Edge to edge jitter & under 1 \textmu{}s peak & not measured & not measured \\
Cycles in the sampling interrupt & under 2000 of 280000 & not measured & not measured \\
Buffer, 128 samples per half & 512 B in a region the stream can reach & 512 B by construction & none needed \\
Processing budget per half & 128 ms & not measured & not measured \\
\bottomrule
\end{tabular}
\caption{The budget for chapter 6. Every measured cell is filled by the end of the chapter, by the witness for the first three rows and by this chapter's own cycle counter for the fourth. Nothing is filled by arithmetic.}
\end{table}

\subsection*{Firmware design (UML)}

\diagram{c06_uml}{One half of the buffer filling, with nothing to do in the interrupt but set a flag. The processing runs on the half that is no longer being written, which is what makes the block size a budget rather than a deadline.}

The design has one invariant and it is worth stating before the code. At any
instant the transfer engine is writing one half of the buffer and the
application may read the other. The interrupt's only job is to record which
half just became readable, and the only failure mode is that the application
has not finished with a half by the time the engine comes back to it. That
failure is detectable: the interrupt can see that the previous half was never
consumed and increment an overrun count. A design where the overrun is not
countable is the design this chapter replaces.

\subsection*{Data flow (ASCII)}

\begin{asciiart}
  board                                                      witness
  +-----------------------------------------------+          +----------------------------+
  | TIM6 update -> TRGO                            |          | Raspberry Pi 4             |
  |   |                                            |          |   MCC 118 DAQ HAT          |
  |   v                                            |          |   1 channel, 100 kS/s      |
  | ADC1 conversion (no software in the path)      |          |                            |
  |   |                                            |          |   scan -> samples.npy      |
  |   v                                            |          |        |                   |
  | DMA1 stream, double buffer                     |          |        v                   |
  |   +--> half A [128] --+                        |          |   threshold crossings      |
  |   +--> half B [128] --+                        |          |   linear interpolation     |
  |            |                                   |          |        |                   |
  |            v  half/full transfer interrupt     |          |        v                   |
  |   flag = which half is readable                |          |   straight-line fit        |
  |            |                                   |          |   slope = sample period    |
  |            v                                   |          |   residual = jitter        |
  |   main loop: feature over the readable half    |          +----------------------------+
  |            |                                   |                       ^
  |            v                                   |   3.3 V marker        |
  |   GPIO marker toggled in the interrupt --------+-----------------------+
  |                                               (220 ohm in series, common ground)
  +-----------------------------------------------+
\end{asciiart}

\subsection*{Repository layout}

\begin{plaincode}
nucleo-h7a3-sampling/
  CMakeLists.txt                  # three build targets from one application
  cmake/arm-none-eabi.cmake       # unchanged from chapter 1
  src/main.c                      # the application, identical in all three
  src/acq_systick.c               # build 1: sample in the tick exception
  src/acq_timer.c                 # build 2: timer TRGO starts the converter
  src/acq_dma_double.c            # build 3: double-buffered transfer
  src/acq.h                       # the one interface all three satisfy
  src/marker.c                    # the GPIO the witness watches
  instr/cyccnt.c                  # cycle counter wrapper, DWT or timer
  instr/cyccnt.h
  instr/freqcount.c               # gated counter, and the reciprocal upgrade
  instr/freqcount.h
  witness/scan.py                 # runs on the Pi, writes samples and metadata
  witness/rate.py                 # edges, fit, rate with an uncertainty
  witness/test_rate.py            # the analysis tested on synthetic input
  docs/measurement.md             # the method, written before the first run
  README.md
\end{plaincode}

\subsection*{Steps}

\begin{steps}
\item \textbf{Write the interface all three builds satisfy.} Decide this first,
because it is what makes the comparison honest. Three implementations behind one
header, selected by the build, so the application source is byte identical
across them.
\begin{ccode}
/* acq.h: one interface, three implementations. */
#include <stdint.h>
#include <stddef.h>

typedef struct {
    const uint16_t *samples;   /* the half that is safe to read */
    size_t          count;     /* samples in that half */
    uint32_t        seq;       /* increments once per half, never wraps in a run */
} acq_block_t;

void     acq_start(void);
int      acq_take(acq_block_t *out);   /* 1 if a block is ready, else 0 */
uint32_t acq_overruns(void);           /* halves the application never read */
\end{ccode}

\item \textbf{Build the cycle counter wrapper before anything else.} It is the
cheapest instrument on the board and the rest of the book uses it. Whether the
core's trace counter runs with no debugger attached is item 25 on the confirm
list, so the wrapper does not assume: it enables the counter, checks that it
moved, and reports failure rather than returning zeros.
\begin{ccode}
/* instr/cyccnt.c */
#include <stdint.h>
#include "stm32h7xx.h"          /* CMSIS device header for this part */

/* The debug block's lock access register. Some Cortex-M7 implementations
   require this key before the trace registers accept writes; on others the
   register reads as zero and the write is harmless. Do not decide which
   applies by reading a forum post: the self-test below decides it. */
#define DWT_LAR (*(volatile uint32_t *) 0xE0001FB0u)
#define DWT_LAR_KEY 0xC5ACCE55u

static int cyccnt_ok;

int cyccnt_init(void)
{
    CoreDebug->DEMCR |= CoreDebug_DEMCR_TRCENA_Msk;
    DWT_LAR = DWT_LAR_KEY;
    DWT->CYCCNT = 0u;
    DWT->CTRL  |= DWT_CTRL_CYCCNTENA_Msk;

    uint32_t a = DWT->CYCCNT;
    for (volatile int i = 0; i < 64; i++) { }
    cyccnt_ok = (DWT->CYCCNT != a);
    return cyccnt_ok;                 /* 0 means: use the timer backend */
}

uint32_t cyccnt_now(void) { return DWT->CYCCNT; }
\end{ccode}
If the self-test returns zero, the fallback is a free-running 32-bit timer with
its prescaler set to one and its reload value left at the maximum, read through
the same header. It is one of only two 32-bit timers on this part, which is why
chapter 20 is careful about who owns it. The two backends differ in one way
that matters and the header says so: the trace counter counts core clocks and
the timer counts timer clocks, and on this part those are the same number only
because chapter 1 configured the buses that way. Store the tick rate next to
the reading rather than assuming it.

\item \textbf{Build the first acquisition: the system tick.} This is the
simplest thing that can work and it is a real variant, not a strawman. The tick
exception fires at 1 kHz, starts a conversion, waits for it, stores it, and
toggles the marker.
\begin{ccode}
/* acq_systick.c */
void acq_start(void)
{
    SysTick->LOAD = 280000u - 1u;     /* 280 MHz core clock, exactly 1 ms */
    SysTick->VAL  = 0u;
    SysTick->CTRL = SysTick_CTRL_CLKSOURCE_Msk | SysTick_CTRL_TICKINT_Msk
                  | SysTick_CTRL_ENABLE_Msk;
}

void SysTick_Handler(void)
{
    marker_toggle();                  /* first, so the witness sees the edge */
    ADC1->CR |= ADC_CR_ADSTART;
    while ((ADC1->ISR & ADC_ISR_EOC) == 0u) { }   /* the cost of this build */
    ring_put((uint16_t) ADC1->DR);
}

/* marker.c: one store, no read-modify-write, no branch. The witness can only
   see what this function shows it, so its own cost is part of the measurement.
   Measure it with the wrapper from step 2 and record the number: every later
   chapter that uses a marker pin inherits it. */
void marker_toggle(void)
{
    static uint32_t level;
    level ^= 1u;
    MARKER_PORT->BSRR = level ? (1u << MARKER_PIN)
                              : (1u << (MARKER_PIN + 16));
}
\end{ccode}
The waiting loop is the whole lesson. Measure it with the wrapper from step 2
and write the number in the budget table. At 1 kHz it is invisible; the point of
measuring it now is that chapter 13 samples faster.

\item \textbf{Build the second acquisition: a timer starts the converter.} Now
no software sits between the timer and the conversion. The timer's update event
is routed to its trigger output and the converter is configured to start on it.
\begin{ccode}
/* acq_timer.c */
void acq_start(void)
{
    RCC->APB1LENR |= RCC_APB1LENR_TIM6EN;
    TIM6->PSC = 280u - 1u;            /* timer clock 280 MHz -> 1 MHz */
    TIM6->ARR = 1000u - 1u;           /* 1 MHz -> 1000 Hz, exactly */
    TIM6->CR2 = (2u << TIM_CR2_MMS_Pos);          /* TRGO on update */
    TIM6->EGR = TIM_EGR_UG;

    /* ADC1: external trigger, rising edge, source number from RM0455's
       trigger table for THIS part. Do not copy it from an H743 project. */
    ADC1->CFGR = (ADC1->CFGR & ~ADC_CFGR_EXTSEL_Msk)
               | (ADC_TRIGGER_TIM6_TRGO << ADC_CFGR_EXTSEL_Pos)
               | (1u << ADC_CFGR_EXTEN_Pos);      /* rising edge */
    ADC1->CR  |= ADC_CR_ADSTART;
    TIM6->CR1  = TIM_CR1_CEN;
}
\end{ccode}
The marker now moves into the end-of-conversion interrupt, which is one step
further from the timer and so one step worse as a witness of the timer. Toggle
it in both places under a build switch and let the witness tell you how much
that step costs. That comparison is a measurement this book has not seen
published for this family.

\item \textbf{Build the third acquisition: double-buffered transfer.} The
stream controller on this family carries two memory pointers and a bit saying
which one it is currently filling. Set both pointers, enable double-buffer
mode, and take the half and full transfer interrupts.
\begin{ccode}
/* acq_dma_double.c */
#define HALF 128u
static uint16_t buf_a[HALF], buf_b[HALF];
static volatile uint32_t ready_seq, taken_seq;
static const uint16_t *ready_ptr;

void acq_start(void)
{
    DMA1_Stream0->PAR  = (uint32_t) &ADC1->DR;
    DMA1_Stream0->M0AR = (uint32_t) buf_a;
    DMA1_Stream0->M1AR = (uint32_t) buf_b;
    DMA1_Stream0->NDTR = HALF;
    DMA1_Stream0->CR   = DMA_SxCR_DBM            /* two pointers, ping-pong */
                       | DMA_SxCR_MINC
                       | (1u << DMA_SxCR_MSIZE_Pos)   /* 16-bit memory */
                       | (1u << DMA_SxCR_PSIZE_Pos)   /* 16-bit peripheral */
                       | DMA_SxCR_TCIE
                       | DMA_SxCR_EN;
    /* The request number that connects ADC1 to this stream comes from the
       DMAMUX request table in RM0455. It is not the same number as RM0433. */
}

void DMA1_Stream0_IRQHandler(void)
{
    DMA1->LIFCR = DMA_LIFCR_CTCIF0;
    /* CT says which pointer the engine is filling NOW, so the other is ready */
    ready_ptr = (DMA1_Stream0->CR & DMA_SxCR_CT) ? buf_a : buf_b;
    ready_seq++;
    marker_toggle();
}
\end{ccode}
Two traps live in these twenty lines and both belong to this family. The first
is the cache: with the data cache enabled and the buffers in the main memory
region, the processor can read a stale line that the engine has already
overwritten. The fix is to invalidate the block before reading it, and the
complete treatment is chapter 19. The second is reach: the tightly coupled
memories are fast and the main transfer engines cannot see them at all, so a
buffer moved there for speed produces a transfer that appears to run and writes
nothing. That is item 16 on the confirm list and the most common fault in this
family.

\item \textbf{Run the witness.} On the Pi, scan one channel continuously and
write the samples with the metadata that makes them interpretable later. The
API names below follow the HAT library's published examples; check them against
the version you install rather than against this page.
\begin{pycode}
# witness/scan.py -- runs on the Raspberry Pi
import json, time
import numpy as np
from daqhats import mcc118, OptionFlags, hat_list, HatIDs

RATE = 100000.0          # requested; the library reports what it granted
SECONDS = 60
CHANNEL = 0

board = mcc118(hat_list(filter_by_id=HatIDs.MCC_118)[0].address)
actual = board.a_in_scan_actual_rate(1, RATE)
board.a_in_scan_start(1 << CHANNEL, int(actual * SECONDS), actual,
                      OptionFlags.CONTINUOUS)
t0 = time.time()
chunks = []
while sum(len(c) for c in chunks) < actual * SECONDS:
    r = board.a_in_scan_read_numpy(-1, 1.0)
    if r.hardware_overrun or r.buffer_overrun:
        raise SystemExit("witness overran; lower the rate or the duration")
    chunks.append(r.data)
board.a_in_scan_stop()
board.a_in_scan_cleanup()

v = np.concatenate(chunks)
np.save("samples.npy", v)
open("meta.json", "w").write(json.dumps(
    {"requested_rate": RATE, "actual_rate": actual,
     "started_unix": t0, "channel": CHANNEL, "samples": int(v.size)}))
\end{pycode}
Record the rate the library granted rather than the rate you asked for. A HAT
that quietly rounds the scan rate to something it can divide exactly turns a
clean measurement into a systematic error of the same size as the effect.

\item \textbf{Turn edges into a rate with an uncertainty.} The analysis is
short and every line of it is a decision worth defending in a README.
\begin{pycode}
# witness/rate.py
import json
import numpy as np

v = np.load("samples.npy")
meta = json.load(open("meta.json"))
fs = meta["actual_rate"]

# Threshold midway between the two settled levels, not at a fixed voltage:
# the marker's high level depends on the series resistor and the HAT's input.
lo, hi = np.percentile(v, 5), np.percentile(v, 95)
thr = 0.5 * (lo + hi)

above = v > thr
cross = np.flatnonzero(np.diff(above.astype(np.int8)) != 0)
# Linear interpolation between the two samples that straddle the threshold.
frac = (thr - v[cross]) / (v[cross + 1] - v[cross])
t_edge = (cross + frac) / fs

n = np.arange(t_edge.size)
slope, intercept = np.polyfit(n, t_edge, 1)
resid = t_edge - (slope * n + intercept)

period = slope                      # one edge per sampling interrupt
rate = 1.0 / period
# Standard error of the slope of a straight-line fit.
se = resid.std(ddof=2) / (n.std() * np.sqrt(n.size))
print(f"edges           {t_edge.size}")
print(f"sample period   {period*1e6:.4f} us")
print(f"sample rate     {rate:.4f} Hz  +/- {rate*se/period:.4f} Hz")
print(f"jitter, rms     {resid.std()*1e9:.1f} ns")
print(f"jitter, peak    {np.abs(resid).max()*1e9:.1f} ns")
\end{pycode}
Report all four numbers. The rate without the jitter hides a build that is on
average right and individually wrong, which is exactly what the tick build with
a long conversion wait looks like.

\item \textbf{Build the gated frequency counter on the board.} The witness is
on the other side of the room and needs a Pi. The board can measure a periodic
signal by itself, and from here on it does. One 32-bit timer counts edges
arriving on its external trigger input; a second timer, clocked from the system
clock, opens and closes the gate.
\begin{ccode}
/* instr/freqcount.c: counts edges on TIM5_ETR over a gate from TIM7. */
void freq_start(uint32_t gate_ms)
{
    TIM5->PSC  = 0u;
    TIM5->ARR  = 0xFFFFFFFFu;
    TIM5->SMCR = TIM_SMCR_ECE;        /* external clock mode 2, on ETR */
    TIM5->CR1  = TIM_CR1_CEN;

    TIM7->PSC = 28000u - 1u;          /* 280 MHz -> 10 kHz */
    TIM7->ARR = (gate_ms * 10u) - 1u;
    TIM7->DIER = TIM_DIER_UIE;
    TIM7->CR1  = TIM_CR1_CEN;
}

static volatile uint32_t last, delta;
void TIM7_IRQHandler(void)
{
    TIM7->SR = 0u;
    uint32_t now = TIM5->CNT;
    delta = now - last;               /* unsigned, so the wrap is free */
    last  = now;
}
\end{ccode}
The resolution of this instrument is one count in the gate. On a 1 kHz signal
with a one second gate that is one part in a thousand, which is far too coarse
to see the effect this chapter is after. The upgrade is the classic one and is
worth the extra twenty lines: stop counting edges in a fixed window and instead
capture the system-clock time of the Nth edge. With a 32-bit capture timer at
280 MHz, a thousand periods are timed to about 3.6 ns, which is under four parts
in a billion of the interval. Resolution then stops being the limit and the
oscillator becomes the limit, which is the honest place for it to be.

\item \textbf{Test the analysis before you trust it.} The rate script is the
instrument's most fragile part, so it is tested on synthetic input with a known
answer, on the host, in continuous integration. This is the same discipline
chapter 12 turns into a rig.
\begin{pycode}
# witness/test_rate.py
import numpy as np
from rate import edges_and_fit       # refactored out of the script

def test_known_rate():
    fs, f, secs = 100000.0, 500.0, 5.0
    t = np.arange(int(fs * secs)) / fs
    v = 1.65 + 1.5 * np.sign(np.sin(2 * np.pi * f * t))
    v += np.random.default_rng(0).normal(0, 0.01, v.size)
    period, jitter = edges_and_fit(v, fs)
    assert abs(1.0 / period - 2 * f) < 0.05     # one edge per half period
    assert jitter < 2e-6
\end{pycode}
\end{steps}

\subsection*{Build, flash and debug}

\diagram{c06_timing}{One millisecond as three of the four parties see it. The witness's sample instants are its own and are never aligned to the board's; the interpolated crossing between two of them is what makes a rate out of a square wave.}

Three builds, one application, three binaries. The CMake target names carry the
variant so a flashed board can be identified from its console banner.

\begin{shellcode}
cmake -B build -DCMAKE_TOOLCHAIN_FILE=cmake/arm-none-eabi.cmake
cmake --build build --target acq_systick acq_timer acq_dma_double -j
probe-rs run --chip STM32H7A3ZITx build/acq_dma_double.elf
\end{shellcode}

On the Pi, the witness runs and the analysis follows it, so a run produces a
number and not a folder of samples somebody will interpret later.

\begin{shellcode}
ssh pi@witness 'cd witness && python scan.py && python rate.py' | tee run.log
\end{shellcode}

\begin{note}{When the rate is right and the jitter is not}
The usual cause is that the marker is toggled after work rather than before it,
so the edge carries the duration of whatever the interrupt did first. Move the
toggle to the first statement and measure again. If the jitter survives that,
the next suspects in order are a higher-priority interrupt that preempts the
sampling one, a wait state on the memory the buffer lives in, and the converter
sampling time rather than the trigger. The gated counter cannot separate those;
the cycle counter placed around each candidate can.
\end{note}

If the board stops responding to the debugger after a build that enables the
transfer engine, the recovery is the same as chapter 1's: hold the reset while
the flashing tool connects, or use the probe's connect-under-reset option. No
option bytes are written in this chapter, so there is nothing here that can put
the board beyond that recovery.

\subsection*{Verification and acceptance criteria}

\begin{itemize}
\item The witness reports a sample rate whose stated uncertainty does not
contain the difference between the three builds, or it reports that it does and
the chapter says the three are indistinguishable at this rate. Either result is
publishable; a result with no uncertainty is not.
\item The three builds compile from one unmodified application source, proven by
a checksum of that file being identical in all three build logs.
\item The cycle counter's self-test runs on a board that has never had a
debugger attached in that power cycle, and the result is recorded in the README
either way. If the trace counter does not run without a debugger, the timer
backend is used and the chapter says so.
\item The overrun count is zero over a ten minute run for the double-buffered
build, and is greater than zero when the processing is deliberately slowed past
the block budget. A counter that never increments has not been tested.
\item The rate analysis passes its synthetic test on the host, in continuous
integration, with no board attached.
\item The marker toggle's cost in cycles is measured and written in the
repository, because seven later chapters inherit it.
\end{itemize}

\subsection*{Variants}

\begin{table}[H]\centering\small
\begin{tabular}{p{24mm}p{26mm}p{44mm}p{22mm}p{20mm}}
\toprule
Axis & Variant & What changes & Cost & Built in full in \\
\midrule
Time and safety & SysTick & The core's own exception at 1 kHz. No timer is spent, and the conversion wait sits inside the exception & Jitter tied to whatever else preempts it & Here \\
Time and safety & General timer & The timer's trigger output starts the conversion with no software in the path, which is the first build whose rate is a hardware property & One timer & Here \\
Execution model & DMA double buffered & Two memory pointers and a current-target bit, so the application always has a whole half to itself & Cache maintenance becomes mandatory & Here \\
Execution model & DMA circular with half transfer & The same ping-pong with one pointer and a half-transfer interrupt. Available on every transfer engine here & The idle half's address cannot be changed & Chapter 4 \\
Execution model & Polled loop or an interrupt with a flag & Read the converter in the main loop, or let the end-of-conversion interrupt raise a flag the loop acts on & No fixed rate at all, or one sample of latency & Chapter 3 \\
Execution model & RTOS task & A task blocks on a notification from the transfer interrupt & A kernel and its tick & Chapter 20 \\
Time and safety & Low-power timer & The same 1 kHz from a timer that survives the deeper sleep states & Fewer features & Chapter 7 \\
Peripheral & Hardware oversampling & The converter averages in hardware before the transfer, so the rate falls and the resolution rises with no processor involvement & Latency and a decision about the shift & Chapter 18 \\
Intelligence & Edge host & Stream every sample to the Pi and do the work there & A framed link and its bandwidth & Chapter 9 and 11 \\
\bottomrule
\end{tabular}
\caption{Variants for chapter 6. The three built here are the three whose difference the witness can actually resolve; the rest are referenced because their difference shows up at rates this chapter does not reach.}
\end{table}

The comparison across the top three rows is the part this book has not found
published. SOURCE.md records as a gap that nobody has measured several execution
models solving one problem on one board with a stated method. This chapter
measures three of them, states the method, and leaves the remaining three to
chapters 3, 4 and 20 using the same witness and the same analysis script, which
is what makes the numbers comparable across the volume.

\subsection*{Pitfalls}

\begin{itemize}
\item Copying a converter trigger source number or a transfer request number
from a project for the better-known member of this family. The tables differ.
The failure is a converter that never starts, which reads as a dead transfer.
\item Placing the sample buffer in tightly coupled memory because it is fast.
The main transfer engines cannot reach it. The transfer appears to run.
\item Leaving the data cache enabled and reading the buffer without
invalidating it. The processor returns bytes that were correct one block ago.
\item Believing the requested scan rate on the witness rather than the rate it
granted. A rounded scan rate is a systematic error, not noise.
\item Toggling the marker after the interrupt's work instead of before it, then
reporting the resulting spread as the board's jitter.
\item Comparing two clocks that share an oscillator. If the witness is ever
driven from the board, or the board from the witness, the measurement proves
nothing.
\item Reporting an average rate with no uncertainty and no jitter, or assuming
the trace cycle counter runs without a debugger. The second is item 25 on the
confirm list precisely because the answer decides which instrument this book
uses.
\end{itemize}

\subsection*{Best practices applied}

\begin{itemize}
\item Every measurement names its instrument. The rate and the jitter come from
the MCC 118; the cycle costs come from this chapter's own wrapper; the cells
that neither has filled say so.
\item The instrument is tested before it is trusted. The analysis script runs
against synthetic input with a known answer, and the cycle counter tests itself
at startup rather than assuming a debug block is enabled.
\item The three builds share one application source, so the comparison measures
the acquisition strategy and nothing else, and the uncertainty reported with
every number is derived from the fit rather than asserted.
\item The overrun path is exercised deliberately, because an error counter that
has never incremented is a decoration.
\item Facts that are not settled are written as confirm-list items with a
reference, not as assertions with a plausible number.
\end{itemize}

\subsection*{Stretch goals}

\begin{itemize}
\item Replace the gated counter with the reciprocal counter sketched in step 9
and show the resolution improving by three orders of magnitude while the stated
accuracy does not move, because the oscillator has become the limit.
\item Run the witness against the board's low-speed crystal instead of its
sampling interrupt, by routing a divided crystal output to the marker pin, and
compare the two oscillators the board carries. Chapter 7 needs that number.
\item Raise the rate to 10 kHz and then 50 kHz and find where each of the three
builds stops meeting its budget. The tick build will be the first to stop and
the witness can say by how much.
\item Publish the analysis script and its synthetic test as a small standalone
package, since no permissively licensed tool in the pool turns a sampled square
wave into a rate with an uncertainty.
\end{itemize}

\subsection*{Roadmap and next steps}

Chapter 7 takes the timer built here into the sleep states, where the
peripheral bus clock stops and the general timer stops with it, which is why the
low-power timer exists and why that chapter rather than this one builds it.

For going deeper, the community embedded roadmap is the map worth following,
and its central advice, that projects settle questions reading leaves open, is
this volume's thesis as well. The written series behind the five-step
progression cited above is the best companion to this chapter specifically, and
reading it against RM0455 rather than against the part it was written for is
itself a useful exercise. For the scheduling mathematics that decides when a
sampling deadline is actually met rather than usually met, the university
specialisation on real-time embedded systems listed in Appendix H is the only
free course in the pool that teaches it properly; nothing else here does. The
vendor's own free training modules cover this family, though there is no course
specific to this part.

\subsection*{Portfolio evidence}

\begin{itemize}
\item A repository with three build targets, one application source, and a
README that opens with the architecture figure and states the measured rate with
its uncertainty in the first paragraph.
\item The run log from the witness, showing the granted scan rate, the edge
count, the rate, the standard error and both jitter figures.
\item A short written method, dated and written before the first run, saying
what would count as a negative result.
\item The cycle counter wrapper and the frequency counter as a small library
with their self-tests, reused unchanged by later chapters, which is better
evidence than a larger piece of code used once.
\end{itemize}

\subsection*{Sources}

Normative references:
\begin{itemize}
\item Reference manual RM0455, for the timer inventory, the trigger output
selection, the converter's external trigger table and the transfer request
table. Not RM0433.
\item The STM32H7A3xI datasheet, for the alternate-function table that decides
which pins carry the external trigger input and the capture inputs, and for the
converter's channel mapping.
\item The board user manual for MB1363, for which Zio pins remain free and
carry a timer channel.
\item The Cortex-M7 technical reference manual and the architecture reference
manual, for the trace block, the cycle counter and the conditions under which it
counts.
\item The MCC 118 user guide, for the single-channel and aggregate scan rates,
the input range and whether an external trigger exists.
\end{itemize}

Reusable implementations:
\begin{itemize}
\item Ian Ross's five-step timer and converter progression on a Cortex-M7
Nucleo, MIT, with the written series that explains each step.\newline
\url{https://github.com/ian-ross/stm32-timer-adc-dma}\newline
\url{https://www.skybluetrades.net/blog/2020/11/2020-11-26-stm32-timer-adc-dma-3/}
\item Majerle's transfer-driven serial examples, MIT, for the half and full
transfer pattern reused here.\newline
\url{https://github.com/MaJerle/stm32-usart-uart-dma-rx-tx}
\item GorgonMeducer's performance counter, Apache-2.0, which chooses the system
tick over the trace counter for exactly the reason step 2 tests.\newline
\url{https://github.com/GorgonMeducer/perf_counter}
\item The flashing and debugging tool whose target list includes this die.\newline
\url{https://github.com/probe-rs/probe-rs}
\end{itemize}
   % Sampling on a timer at exactly 1 kHz
\project{07}{Stop mode, RTC wake, and a battery number}{NUCLEO-H7A3ZI-Q with an nRF-PPK2}{Low-power modes, clock restore after wake, charge per cycle}

\begin{keyfacts}
\item[Board] NUCLEO-H7A3ZI-Q, with an nRF-PPK2 in the supply path
\item[Peripherals] PWR and RCC for the mode and the clock restore, the RTC on the fitted 32.768 kHz crystal, LPTIM1, EXTI for the wake line, three GPIO markers for the instrument's digital channels
\item[Toolchain] arm-none-eabi-gcc with CMake as in chapter 1; Python 3 on the host with the instrument's interface installed from its package index
\item[Operating system] Bare metal
\item[Difficulty] 5 of 5
\item[Effort] 4 evenings of about four hours, and the last one is entirely arithmetic
\item[Deliverable] A duty-cycled application that sleeps between wakes, restores its clock tree correctly, and reports charge and energy per cycle with a spread over at least ten cycles, together with an honest statement of what this part is and is not good for
\end{keyfacts}

\subsection*{Why this project}

This chapter has to begin with the finding, because the finding changes what
the chapter is allowed to claim. \textbf{This part is a poor low-power node.}
The literature on energy-budgeted sensor nodes runs on the ultra-low-power
families: the STM32L4 and STM32U5 lines, the nRF52 and nRF54 series, the
MSP430. Those parts idle in microamps. This one idles in milliamps, on a board
whose debug probe is powered from the same cable. A search of the published
work turns up no thesis and no paper that builds a low-power sensor node on this
family at all, which is a negative result worth recording rather than a gap to
apologise for.

So the chapter keeps its subject and changes its claim. It is not "here is a
low-power node". It is "here is how you account for energy, and here is what
racing to sleep looks like on a part chosen for speed". Both halves are
transferable. The accounting method is the same on any part, and it is the half
that most published low-power writing does badly: an average current with no
spread, taken over an unstated interval, on a board with unstated peripherals
enabled. The racing-to-sleep half is specific and interesting here precisely
because the part is fast: a processor that finishes its work in a tenth of the
time can afford a much worse sleep current before it loses, and the crossover is
a number this chapter can produce.

There are two pieces of good news, both verified rather than assumed. The
32.768 kHz crystal is fitted on this board. That is confirmed from two
independent machine-readable sources that name this exact board: the upstream
board description in the other operating system's tree lists a 32.768 kHz
crystal oscillator, and the MicroPython board definition selects that oscillator
as the real-time clock source. So the real-time clock is usable here and its
accuracy is claimable, which is what makes a ten second wake period a period
rather than an approximation. And the vendor's firmware package carries three
power examples for this exact board under
\texttt{Projects/}\allowbreak\texttt{NUCLEO-H7A3ZI-Q/}\allowbreak\texttt{Examples/PWR}: standby, standby with a clock
wake, and the deeper stop mode with a clock wake. Working register sequences
for this board exist and do not have to be derived from a sibling part.

\begin{note}{The clock tree does not come back on its own}
This is the trap specific to this family and it is the reason the chapter has a
whole step for it. On wake from the deeper sleep state the system clock has
reverted to the internal oscillator and the phase-locked loop is off. Nothing
restores it. The symptom is not a failure: the application runs, the wake
completes, and the measured current after the wake stays higher than it should
while the core runs slower than it should. Every project that reports
disappointing energy numbers on this family should check this first. The
restore also has an order: the voltage scale and the flash wait states come
before the frequency, exactly as in chapter 1.
\end{note}

\subsection*{Prior art and what to reuse}

\begin{table}[H]\centering\small
\begin{tabular}{p{32mm}p{46mm}p{40mm}p{20mm}}
\toprule
Source & What it gives & What it does not & Licence \\
\midrule
The vendor's three power examples for this exact board & Working entry and wake sequences for standby, standby with a clock wake, and the deeper stop mode with a clock wake, written against RM0455 & No measurement, no clock restore worth the name, no accounting, and nothing about what a cycle costs & Basic examples are BSD-3-Clause; read the header of each file, because the middleware around them is not \\
The instrument's Python interface & Scripted control of the power instrument: mode, supply voltage, start and stop, and the sample stream with its digital channels & It is a transport, not a method. It computes no charge and segments no cycles & GPL-2.0. Install from the package index. Do not fork it into a portfolio repository \\
Golioth's article on measuring current with this instrument & The procedure: how to place the instrument, which mode to choose, and how to read a trace & It is an introduction on a different board, with no uncertainty treatment and no per-cycle statistics & Article, read and cite \\
The MicroPython board definition for this board & Confirms the low-speed crystal is fitted and is used as the real-time clock source & It is an interpreter, and its power behaviour is not this chapter's subject & MIT \\
The ultra-low-power families' datasheets and application notes & The comparison to beat, in numbers their vendors publish and stand behind & They are different silicon. Their figures are the target, never this part's result & Vendor documentation \\
\bottomrule
\end{tabular}
\caption{Prior art for chapter 7. No peer-reviewed method exists for this instrument, so the measurement rigour below is original work rather than a reproduction.}
\end{table}

What is left to write is the rigour and the restore. The clock restore
function, the phase marking that lets the trace be segmented automatically, the
per-cycle charge computation with a spread over at least ten cycles, and the
separation of the measurement from the battery arithmetic are the author's, and
chapters 10 and 12 use all four unchanged.

\subsection*{Parts from the inventory}

\begin{table}[H]\centering\small
\begin{tabular}{p{40mm}p{66mm}p{40mm}}
\toprule
Part & Role & Interface \\
\midrule
NUCLEO-H7A3ZI-Q & The subject. Duty cycles between a sense phase and a sleep state & Micro USB to the host \\
nRF-PPK2 & Ammeter or source meter up to 1 A, with eight digital inputs time-aligned to the current trace & USB to the host, two leads to the board \\
Three jumper wires & Two for the supply path, one common ground for the digital channels & Jumper header \\
Three more jumper wires & Phase markers from free GPIO to the instrument's digital inputs & Zio header to the instrument \\
Host PC & Runs the capture script and the analysis & USB \\
\bottomrule
\end{tabular}
\caption{Inventory items used in chapter 7. The instrument measures at most 1 A, which is comfortable for this board but is the reason a cellular modem is never placed in the same path; chapter 11 says why.}
\end{table}

\subsection*{System architecture}

\diagram{c07_arch}{The instrument sits in the supply path and watches three marker pins at the same time. Everything the chapter claims comes out of one trace with four channels in it.}

The instrument is in series with whatever it is measuring, which is the first
thing that has to be settled and is not settled yet. Two topologies are
possible on this board. In ammeter mode the instrument is placed across the
jumper that carries the microcontroller's own supply current, and the board is
powered as usual. In source-meter mode the instrument supplies the board
itself. Which jumper carries that current, exactly what it isolates, whether
the probe's own current is excluded, and whether the instrument can power the
board through it at the required voltage are items 6 and 34 on the confirm
list. Until they are settled the chapter measures in ammeter mode and states in
every caption that the probe is still powered from USB and is not in the
measured path.

\subsection*{Peripheral configuration}

\begin{table}[H]\centering\small
\begin{tabular}{p{20mm}p{28mm}p{24mm}p{30mm}p{30mm}}
\toprule
Peripheral & Mode & Clock source & Pins and function & Interrupt and transfers \\
\midrule
PWR & Stop with the deepest stop voltage scale the part allows & None & None & None \\
RCC & Restore path after wake: voltage scale, flash latency, phase-locked loop, switch & 8 MHz bypass from the probe & None & None \\
RTC & Calendar and a periodic alarm & The fitted 32.768 kHz crystal & None & Alarm through the wake line \\
LPTIM1 & Periodic counter that continues in the stop state & The same crystal & None & Update through the wake line \\
EXTI & Wake line for the alarm and for the user button & None & PC13 for the button & Wake from stop \\
GPIO markers & Push-pull outputs, one per phase & Peripheral bus & Three free Zio pins, item 13 on the confirm list & None \\
SysTick & Stopped before sleeping, restarted after & Core clock & None & Tickless build only \\
\bottomrule
\end{tabular}
\caption{Peripheral configuration. What each low-power mode retains, which sources can wake from it and what the wake latency is are item 27 on the confirm list and are read out of RM0455's low-power section before any of them is printed as fact.}
\end{table}

\subsection*{Wiring}

\diagram{c07_wiring}{The instrument in ammeter mode, in series with the board's own supply through the measurement jumper, with the three marker pins on its digital inputs and one common ground. Which jumper this is remains item 6 on the confirm list.}

Three safety points. The instrument measures at most 1 A, so nothing that draws
a pulse larger than that goes in this path, which excludes every cellular modem
on this bench. The board's logic is 3.3 V and the instrument's digital inputs
are driven by the board, never the other way round; the logic threshold of
those inputs is item 33 on the confirm list and is read before the markers are
trusted. And the ground between the instrument and the board must be common or
the digital channels float, which produces a trace that looks plausible and
means nothing.

\subsection*{Memory and timing budget}

\diagram{c07_mem}{The clock tree at the instant of wake, next to the clock tree the application expects. The difference between them is the whole of the trap, and nothing in the silicon closes it.}

\begin{table}[H]\centering\small
\begin{tabular}{p{44mm}p{28mm}p{28mm}p{30mm}}
\toprule
Quantity & Budget & Measured & Margin \\
\midrule
Wake period & 10.000000 s from the fitted crystal & not measured & not measured \\
Run current at 280 MHz & no budget until the first measurement & not measured & not measured \\
Stop current & no budget until the first measurement & not measured & not measured \\
Wake to first instruction & under 1 ms & not measured & not measured \\
Wake to restored clock tree & under 2 ms & not measured & not measured \\
Charge per cycle & established here & not measured & not measured \\
Energy per cycle & established here & not measured & not measured \\
Spread over ten cycles & within 5 percent of the median & not measured & not measured \\
\bottomrule
\end{tabular}
\caption{The budget for chapter 7. The current rows carry no budget on purpose. A budget invented before the first measurement becomes the number the measurement is read against, and no figure in microamps is written anywhere in this book until an instrument has produced it.}
\end{table}

One row in that table is exact and it is worth saying why. The wake period is
ten seconds because the crystal is 32.768 kHz and the real-time clock's
prescalers divide it by exactly 32768, so a ten second alarm is ten seconds and
not ten seconds plus a drift term. That is only true because the crystal is
fitted, which is a confirmed fact about this board and not a family assumption.
On a board where the low-speed oscillator is the internal one, the same code
produces a period with a tolerance of several percent and every per-cycle figure
inherits it.

\subsection*{Firmware design (UML)}

\diagram{c07_uml}{The duty cycle as a state machine, with the state most projects on this family omit drawn in full. The path that skips the restore is drawn too, because it is the one that runs.}

The application is a loop with a sleep in it, and the only design decision that
matters is where the boundaries of a cycle are. This chapter puts them at the
wake, so a cycle is wake, restore, sense, prepare, sleep, and the next wake. The
marker pin that encloses the whole cycle is what lets the analysis segment the
trace without a human looking at it, and that is the difference between a
measurement you can repeat ten times and a screenshot.

\subsection*{Data flow (ASCII)}

\begin{asciiart}
  board                                            instrument                host
  +-------------------------------------+          +---------------+         +------------------+
  | RTC alarm  (32.768 kHz crystal)     |          | nRF-PPK2      |         | capture.py       |
  |   |                                 |          |               |         |   mode, voltage  |
  |   v                                 |   IDD    | ammeter in    |  USB    |   start, stop    |
  | wake: HSI, PLL off, slow            +--------->| series with   +-------->|                  |
  |   |                                 |          | the supply    |         | cycles.py        |
  |   v                                 |          |               |         |   segment by D0  |
  | restore: VOS, latency, PLL, switch  |          | D0 cycle      |<--------+   integrate I dt |
  |   |                                 |  D0 D1 D2| D1 sense      |         |   per-cycle Q    |
  |   v                                 +--------->| D2 restore    |         |   median, spread |
  | sense: one sample, one feature      |          |               |         |                  |
  |   |                                 |          | 4 channels,   |         | battery.py       |
  |   v                                 |          | one timebase  |         |   Q -> lifetime  |
  | prepare: markers low, tick stopped  |          +---------------+         |   assumptions    |
  |   |                                 |                                    +------------------+
  |   v                                 |
  | stop: SLEEPDEEP, WFI                |
  +-------------------------------------+
\end{asciiart}

\subsection*{Repository layout}

\begin{plaincode}
nucleo-h7a3-lowpower/
  CMakeLists.txt                  # three sleep strategies from one application
  cmake/arm-none-eabi.cmake       # unchanged from chapter 1
  src/main.c                      # the duty cycle, identical in all three
  src/sleep.h                     # the one interface all three satisfy
  src/sleep_rtc.c                 # wake on a real-time clock alarm
  src/sleep_lptim.c               # wake on the low-power timer
  src/sleep_tickless.c            # stop the tick, sleep to the next deadline
  src/clocks.c                    # restore_after_stop(), the chapter's core
  src/markers.c                   # three phase pins, shared with chapter 10
  meas/capture.py                 # instrument control, writes a raw trace
  meas/cycles.py                  # segment, integrate, per-cycle statistics
  meas/battery.py                 # the arithmetic, and what it assumes
  meas/test_cycles.py             # the analysis tested on a synthetic trace
  meas/requirements.txt           # the interface pinned, never vendored
  docs/method.md                  # written before the first capture
  README.md
\end{plaincode}

\subsection*{Steps}

\begin{steps}
\item \textbf{Install the instrument's interface without vendoring it.} The
Python interface to this instrument is GPL-2.0. That licence is entirely fine
to depend on and entirely wrong to copy into a portfolio repository that you
intend to publish under a permissive licence. Install it, pin it, and record
why in the README.
\begin{shellcode}
python -m venv .venv && . .venv/bin/activate
pip install ppk2-api
pip freeze > meas/requirements.txt
\end{shellcode}

\item \textbf{Prove the instrument before you believe it.} An instrument that
has not been checked against something known is an opinion. Put a resistor of
known value across the supply at a known voltage and confirm the reading is
where Ohm's law puts it, then record the deviation in the README as the
instrument's working uncertainty for the rest of this book.
\begin{pycode}
# meas/capture.py -- the API names follow the published examples of the
# version pinned in requirements.txt. Check them against that version.
from ppk2_api.ppk2_api import PPK2_API

ppk = PPK2_API("/dev/ttyACM0")
ppk.get_modifiers()                 # calibration constants from the device
ppk.use_ampere_meter()              # the board keeps its own supply
ppk.start_measuring()
raw = ppk.get_data()
current_uA, digital = ppk.get_samples(raw)
ppk.stop_measuring()
\end{pycode}
The digital channels come back alongside the current samples on one timebase,
which is the property this whole chapter is built on. How many channels, at what
logic threshold, at what sample rate, and whether an exported file carries them
are item 33 on the confirm list.

\item \textbf{Measure the run state first and write down what was enabled.} A
current figure with no statement of which peripherals were on is not a
measurement. Record the core frequency, the voltage scale, which peripheral
clocks are enabled, whether the caches are on, and whether the debugger is
attached. Attach that list to the number, permanently.

\item \textbf{Enter the stop state with a real-time clock alarm.} The domain
structure of this part differs from its better-known sibling: it has a CPU
domain and a smart-run domain where the other has three power domains, so the
bit names in the power control registers are not the same. Take the sequence
from RM0455's low-power section and from the vendor's example for this exact
board.
\begin{ccode}
/* sleep_rtc.c: arm the alarm, then stop. Register names from the CMSIS
   device header for this part; the deepest stop voltage scale the part
   allows is read out of RM0455, not assumed from a sibling. */
void sleep_until_next_alarm(uint32_t seconds)
{
    rtc_set_alarm_in(seconds);              /* alarm A, through the wake line */

    PWR->CR1 = (PWR->CR1 & ~PWR_CR1_SVOS_Msk) | PWR_SVOS_DEEPEST;
    SCB->SCR |= SCB_SCR_SLEEPDEEP_Msk;      /* stop, not plain sleep */

    marker_clear(MARK_CYCLE);               /* the trace's cycle boundary */
    __DSB();
    __WFI();                                /* execution resumes on the line
                                               after this one, at reduced
                                               speed, with PLL1 off */
    marker_set(MARK_CYCLE);
    SCB->SCR &= ~SCB_SCR_SLEEPDEEP_Msk;
}
\end{ccode}
Two details decide whether this works at all. Every wake source must be enabled
in the wake-line register as well as in its own peripheral, because the
peripheral interrupt alone does not leave the stop state. And any pending
interrupt at the moment of the wait instruction causes it to return
immediately, which presents as a board that never sleeps and is almost always
a flag left set from the previous cycle.

\item \textbf{Restore the clock tree, and measure the difference.} This is the
step the chapter exists for. Run one capture with the restore commented out and
one with it in, and put the two traces side by side.
\begin{ccode}
/* clocks.c: the order is voltage scale, then wait states, then frequency.
   Reversing it is the classic way to produce a board that runs at reset
   speed, or one that does not run. */
void restore_after_stop(void)
{
    marker_set(MARK_RESTORE);

    pwr_set_run_voltage_scale();            /* back to the scale 280 MHz needs */
    flash_set_latency_for_280mhz();         /* wait states BEFORE the frequency */

    RCC->CR |= RCC_CR_HSEON;                /* bypass input from the probe */
    while ((RCC->CR & RCC_CR_HSERDY) == 0u) { }
    RCC->CR |= RCC_CR_PLL1ON;
    while ((RCC->CR & RCC_CR_PLL1RDY) == 0u) { }

    MODIFY_REG(RCC->CFGR, RCC_CFGR_SW, RCC_CFGR_SW_PLL1);
    while ((RCC->CFGR & RCC_CFGR_SWS) != RCC_CFGR_SWS_PLL1) { }

    marker_clear(MARK_RESTORE);
}
\end{ccode}
The restore costs time and the marker measures it. That cost is what makes the
racing-to-sleep argument quantitative rather than rhetorical: a part that takes
longer to become fast again has a longer fixed overhead per cycle, and below
some wake period the overhead is the whole budget.

\item \textbf{Build the low-power timer variant.} The general timer of chapter 6
stops in the stop state because its bus clock stops. The low-power timer does
not, because it runs from the same crystal the real-time clock uses and its
logic continues. That is the entire reason it exists and it is why this chapter
rather than chapter 6 builds it.
\begin{ccode}
/* sleep_lptim.c: a periodic wake without the calendar. */
void lptim_start_periodic(uint32_t ticks)
{
    /* kernel clock select: the fitted 32.768 kHz crystal, not the internal
       low-speed oscillator. The selection field is in RCC's kernel clock
       configuration register for this part. */
    rcc_select_lptim1_kernel_clock_lse();

    LPTIM1->CR  = 0u;
    LPTIM1->IER = LPTIM_IER_ARRMIE;         /* autoreload match */
    LPTIM1->CR  = LPTIM_CR_ENABLE;
    LPTIM1->ARR = ticks - 1u;               /* write only while enabled */
    LPTIM1->CR |= LPTIM_CR_CNTSTRT;         /* continuous */
}
\end{ccode}
The autoreload register on this peripheral is written after the peripheral is
enabled and the write takes time to be accepted. Writing it before enabling, or
writing it again without waiting, is the standard first fault with this
peripheral and it presents as a period that is right sometimes.

\item \textbf{Build the tickless variant.} The millisecond tick is a source of
wakes that do nothing. In this build the tick is stopped before sleeping and
the elapsed time is reconstructed on wake from the counter that kept running.
\begin{ccode}
/* sleep_tickless.c: stop the tick, sleep, then account for the gap. */
void sleep_tickless(uint32_t ms)
{
    uint32_t before = lptim_count();
    SysTick->CTRL &= ~SysTick_CTRL_ENABLE_Msk;

    lptim_set_next_match(before + ms_to_ticks(ms));
    sleep_until_next_alarm_or_match();

    uint32_t slept = lptim_count() - before;        /* unsigned, wrap is free */
    tick_advance(ticks_to_ms(slept));               /* the clock did not stop */
    SysTick->CTRL |= SysTick_CTRL_ENABLE_Msk;
}
\end{ccode}
Two things are worth stating plainly. The reconstruction has a resolution of one
crystal tick, which is about 30.5 \textmu{}s, so a millisecond software clock
built this way is accurate but granular. And the same scheme on an operating
system is that system's tickless idle, which chapter 20 builds; doing it by hand
first is what makes that chapter's version legible rather than magical.

\item \textbf{Mark the phases and segment the trace automatically.} Three
markers: one enclosing the whole cycle, one on the sense phase, one on the
restore. The analysis never asks a human where a cycle begins.
\begin{pycode}
# meas/cycles.py
import numpy as np

def per_cycle_charge(current_uA, digital, fs, cycle_bit=0):
    """Charge in microcoulombs for each complete cycle in the trace."""
    d = (digital >> cycle_bit) & 1
    edges = np.flatnonzero(np.diff(d) != 0)
    rises = edges[d[edges + 1] == 1]
    out = []
    for a, b in zip(rises[:-1], rises[1:]):
        out.append(current_uA[a:b].sum() / fs)      # uA * s = uC
    return np.asarray(out)

def report(q_uC, volts=3.3):
    e_uJ = q_uC * volts
    return {
        "cycles":        int(q_uC.size),
        "charge_median": float(np.median(q_uC)),
        "charge_min":    float(q_uC.min()),
        "charge_max":    float(q_uC.max()),
        "energy_median": float(np.median(e_uJ)),
        "spread_pct":    float(100 * (q_uC.max() - q_uC.min())
                               / np.median(q_uC)),
    }
\end{pycode}
Report the median and the extremes over at least ten cycles, never a single
average current. An average current over a duty-cycled trace hides the thing
that decides a battery life, which is whether the cycles are alike.

\item \textbf{Test the analysis on a synthetic trace.} The segmentation is the
fragile part, so it is tested on the host with a trace whose answer is known by
construction, in continuous integration, with no instrument attached.
\begin{pycode}
# meas/test_cycles.py
import numpy as np
from cycles import per_cycle_charge

def test_known_charge():
    fs = 100000.0
    n_cycle = int(fs)                       # one second per cycle
    i = np.full(10 * n_cycle, 100.0)        # 100 uA flat, 10 cycles
    d = np.tile(np.r_[np.ones(n_cycle // 10), np.zeros(9 * n_cycle // 10)], 10)
    q = per_cycle_charge(i, d.astype(np.uint8), fs)
    assert q.size == 9                      # nine complete cycles in ten
    assert np.allclose(q, 100.0, rtol=1e-6) # 100 uA for 1 s is 100 uC
\end{pycode}

\item \textbf{Turn the measurement into a battery number and keep the two
apart.} The energy per cycle is a measurement. The lifetime is arithmetic over
that measurement and a set of assumptions, and the assumptions are where these
claims usually fail. State them next to the number, every time.
\begin{pycode}
# meas/battery.py
def lifetime_days(energy_per_cycle_uJ, cycles_per_day, cell_mAh, cell_V):
    """Arithmetic only. The measurement is the first argument; everything
    else is an assumption and is printed with the answer."""
    cell_J = cell_mAh * 1e-3 * 3600.0 * cell_V
    return cell_J / (energy_per_cycle_uJ * 1e-6 * cycles_per_day)
\end{pycode}
What that arithmetic does not model, and what the README says it does not model:
self-discharge, the capacity a coin cell actually delivers under pulsed load
rather than under its rated continuous load, temperature, the regulator's
efficiency at the currents involved, and the fact that this board is not the
product. Whether the backup domain is usable without a coin cell at all is item
10 on the confirm list. A battery number offered without that paragraph is
advertising rather than engineering.
\end{steps}

\subsection*{Build, flash and debug}

\diagram{c07_timing}{One cycle with the restore and one without, on the same axes, with the marker channel underneath. The shape of the difference is the point; the values on the vertical axis come from the instrument and are not drawn here because they have not been measured.}

\begin{shellcode}
cmake -B build -DCMAKE_TOOLCHAIN_FILE=cmake/arm-none-eabi.cmake
cmake --build build --target sleep_rtc sleep_lptim sleep_tickless -j
probe-rs run --chip STM32H7A3ZITx build/sleep_rtc.elf
python meas/capture.py --seconds 120 --out trace.npz
python meas/cycles.py trace.npz && python meas/battery.py
\end{shellcode}

\begin{note}{When the board sleeps and the debugger does not}
A debugger attached across a sleep state changes what you are measuring and
sometimes prevents the state entirely, because the debug block keeps clocks
running so that it can still see the core. Capture with the debugger detached
and the board reset by a power cycle, and say so in the caption. If the board
will not re-enter debug afterwards, connect under reset. If a build sleeps
immediately and never enumerates, hold the user button at reset and have the
first thing the application does be a check of that button, so there is always a
path back into a board that has decided to be quiet.
\end{note}

\subsection*{Verification and acceptance criteria}

\begin{itemize}
\item The instrument's reading against a known resistor at a known voltage is
recorded, with its deviation, before any board measurement is reported.
\item Charge and energy per cycle are reported as a median with a minimum and a
maximum over at least ten consecutive cycles. A single average current is not
an acceptable result in this chapter or in chapters 10 and 12.
\item The trace with the clock restore and the trace without it are captured
under otherwise identical conditions and both are published. The difference is
stated in microjoules per cycle and in restore milliseconds.
\item Every current figure carries the list of what was enabled when it was
taken: frequency, voltage scale, peripheral clocks, caches, and whether the
debugger was attached.
\item The wake period is checked against the real-time clock over at least an
hour and the drift is stated, which is claimable here because the crystal is
fitted.
\item The three sleep builds share one unmodified application source, and the
analysis passes its synthetic test on the host with no instrument attached.
\item The battery number is printed together with its assumptions in the same
output, so that the two cannot be separated by whoever quotes it.
\end{itemize}

\subsection*{Variants}

\begin{table}[H]\centering\small
\begin{tabular}{p{24mm}p{26mm}p{44mm}p{22mm}p{20mm}}
\toprule
Axis & Variant & What changes & Cost & Built in full in \\
\midrule
Time and safety & The real-time clock & A calendar alarm from the fitted crystal wakes the part. The period is exact and the wall-clock time survives the sleep & The calendar's own setup, and a backup domain question & Here \\
Time and safety & Low-power timer & A counter that continues in the stop state because it runs from the crystal rather than a bus clock. No calendar & Fewer features than a general timer & Here \\
Time and safety & Tickless idle & The millisecond tick is stopped before sleeping and the elapsed time is reconstructed from the counter that kept running & Granularity of one crystal tick & Here \\
Time and safety & General timer or SysTick & Both stop when their clock stops, so neither can wake the part from the stop state at all & Not usable here & Chapter 6 \\
Time and safety & Standby instead of stop & Deeper, and the vendor ships two examples of it for this board. Memory contents are not retained, so the application restarts & A restart is not a wake & Referenced here \\
Time and safety & Watchdog across the sleep & Whether the independent watchdog survives the stop state is item 28 on the confirm list and decides whether a sleeping node can be recovered automatically & A missed wake becomes a reset & Chapter 12 \\
Operating system & RTOS tickless idle & The same reconstruction, owned by the kernel, with the idle task deciding when it is safe & A kernel and its assumptions & Chapter 20 \\
Language & MicroPython & The interpreter exposes sleep states directly, which is the fastest way to see the current fall & No control of the restore & Chapter 1 \\
\bottomrule
\end{tabular}
\caption{Variants for chapter 7. The first three are built here because they are the three that still run when the bus clocks stop, which is the property that defines this chapter.}
\end{table}

The comparison worth making across those rows is not which sleep state is
deepest. It is which one leaves a usable timebase running, because a node that
sleeps perfectly and cannot tell how long it slept is not a node. The real-time
clock keeps a calendar, the low-power timer keeps a count, and the standby state
keeps neither in the general case. Chapter 10 takes the same three and asks a
different question, which is where inside a cycle the charge actually goes.

\subsection*{Pitfalls}

\begin{itemize}
\item Omitting the clock restore after wake and then reporting the energy. The
board works, so nothing points at the cause.
\item Restoring the frequency before the voltage scale and the flash wait
states, in either this chapter or chapter 1. The order is not decorative.
\item Taking a power measurement with the debugger attached and not saying so.
The debug block keeps clocks running and the number is not the node's.
\item Copying a power control register sequence from a project for the
better-known member of this family. This part has a different domain structure
and the bit names differ.
\item Enabling a wake source in its own peripheral but not in the wake-line
register. The interrupt fires and the part stays asleep.
\item Entering the wait instruction with a flag still pending from the previous
cycle, which returns immediately and presents as a board that never sleeps.
\item Writing the low-power timer's autoreload register before enabling the
peripheral, or without waiting for the write to be accepted.
\item Reporting an average current over a duty-cycled trace. It hides whether
the cycles are alike, which is the only thing a battery cares about.
\item Quoting a battery lifetime without the assumptions it rests on.
\end{itemize}

\subsection*{Best practices applied}

\begin{itemize}
\item The finding comes first. The chapter states that this part is a poor
low-power node before it states anything else, and every claim after that is
framed as method rather than as a result to be proud of.
\item Every measurement names its instrument, and every current figure carries
the configuration it was taken under.
\item A dependency's licence is stated before any code is written around it.
The instrument's interface is GPL-2.0, so it is installed and pinned and never
vendored.
\item Measurement and arithmetic are kept apart. The energy per cycle is
measured; the lifetime is computed, and the assumptions travel with it.
\item No figure in microamps appears anywhere until an instrument has produced
it, including in the budget table.
\item The spread is reported with the median, because a duty-cycled system that
is usually consistent and occasionally not is a different system from one that
is always consistent.
\end{itemize}

\subsection*{Stretch goals}

\begin{itemize}
\item Find the crossover. Sweep the sense-phase duration and the wake period and
find the point at which racing to sleep on a fast part stops paying, then state
it as a rule with its conditions. Nobody has published that curve for this
family.
\item Repeat the whole measurement on an ultra-low-power board if one ever
reaches this bench, using the same scripts, and publish the two together. The
method transfers even though the numbers do not.
\item Measure the restore time as a function of the phase-locked loop
configuration and find the cheapest configuration that still meets the
application's throughput requirement.
\item Add a fourth build that wakes on the user button rather than on time, and
show what an event-driven node costs when the events are rare.
\item Write the whole thing up as the measurement method it is, since no
peer-reviewed method for this instrument exists, and offer it for review rather
than only publishing numbers.
\end{itemize}

\subsection*{Roadmap and next steps}

Chapter 8 gives the sleeping node something to do between wakes, which is the
state machine this chapter's duty cycle is a skeleton of, and chapter 10 splits
the cycle measured here into its phases and asks where the charge goes.

For going deeper, the community embedded roadmap is the map worth following,
and its point that projects settle questions reading leaves open applies with
force here, because the published material on this specific question is thin.
The free course from a different silicon vendor on its own development kit is
the best free treatment of duty-cycled node design in the pool, and its
concepts transfer directly even though its silicon does not. For the standards a
reader will eventually meet on a product that claims a battery life, the
consumer device security baseline and the industrial component security series
are listed in Appendix H. Nothing in the pool teaches the measurement rigour
this chapter needs, which is why the chapter writes it.

\subsection*{Portfolio evidence}

\begin{itemize}
\item A repository with three sleep builds, one application source, and a
README that states the finding about this part in its first paragraph rather
than burying it.
\item Two published traces under otherwise identical conditions, with and
without the clock restore, and the difference stated in microjoules per cycle.
\item A per-cycle table over at least ten cycles with the median and both
extremes, and the configuration list attached to it.
\item A battery number printed together with the assumptions it rests on, in the
same output, so it cannot be quoted alone.
\end{itemize}

\subsection*{Sources}

Normative references:
\begin{itemize}
\item Reference manual RM0455, for the low-power modes, what each retains, the
wake sources, the domain structure and the power control register bit names.
Not RM0433, whose domain structure differs.
\item The STM32H7A3xI datasheet, for the voltage scales, the clock limits at
each scale and the supply configuration.
\item The board user manual for MB1363, for the measurement jumper, what it
isolates, and whether the probe's current is excluded from it.
\item The nRF-PPK2 user guide, for the two modes, the maximum current, the
digital input count and threshold, the sample rate and the export format.
\end{itemize}

Reusable implementations:
\begin{itemize}
\item The vendor's three power examples for this exact board, under
\texttt{Projects/}\allowbreak\texttt{NUCLEO-H7A3ZI-Q/}\allowbreak\texttt{Examples/PWR}: standby, standby with a clock
wake, and the deeper stop mode with a clock wake. Basic examples are
BSD-3-Clause; check each file's header.
\item The Python interface to the power instrument, GPL-2.0. Install from the
package index and do not vendor it.\newline
\url{https://github.com/IRNAS/ppk2-api-python}
\item Golioth's article on measuring current consumption with this
instrument, which is the method this chapter starts from.\newline
\url{https://blog.golioth.io/measuring-current-consumption-with-power-profiler-kit-ii/}
\item The MicroPython board definition for this exact board, MIT, which is the
second independent confirmation that the low-speed crystal is fitted.\newline
\url{https://github.com/micropython/micropython}
\end{itemize}
   % Stop mode, RTC wake, and a battery number

\part{The mioty node, everything except the radio}
\project{08}{The node's state machine: sense, feature, and a transmit that is a stub}{NUCLEO-H7A3ZI-Q, and nothing else}{Application state machine, bare metal}

\begin{keyfacts}
\item[Board] NUCLEO-H7A3ZI-Q alone. No shield, no radio, no instrument
\item[Peripherals] The three on-board LEDs as a state display, USART3 for the event log, the periodic wake of chapter 7, the user button as a forced event
\item[Toolchain] arm-none-eabi-gcc with CMake as in chapter 1, plus a host build of the same source with a C compiler and a test runner
\item[Operating system] Bare metal, run to completion, no dynamic allocation
\item[Difficulty] 3 of 5
\item[Effort] 3 evenings of about four hours
\item[Deliverable] A transition table whose every row is proven reachable by a test, a dispatcher that cannot block, and a transmit step that logs exactly the bytes a radio would have been handed
\end{keyfacts}

\subsection*{Why this project}

The seven chapters before this one produce bytes, samples, frames and energy.
None of them decides anything. This chapter builds the part of the node that
does: the control flow that says what happens next, in what order, and what
happens instead when a step does not complete. It is the smallest chapter in
the book by line count and the one that is hardest to get back if it is written
badly, because a control flow that has grown out of a set of boolean flags
cannot be tested, cannot be drawn, and cannot be explained to anybody who did
not write it.

The transmit step is a stub on purpose, and the purpose is not that the radio
is missing. It is that a node whose transmit is a stub can be run to completion
on a host, in continuous integration, with no hardware attached, and the frame
it would have sent can be compared byte for byte against the encoder that
produces it. That comparison is this chapter's acceptance criterion and it is
what makes chapter 9 and this chapter verify each other rather than merely
agreeing. The bench has no mioty radio and none is bought, so the stub is also
the honest representation of what exists.

\begin{note}{There is no prior art worth the name and that is the finding}
This is the only chapter in the volume whose technical work has no source to
start from, and it is worth saying plainly rather than padding a table. What
exists is either a framework whose licence excludes it from a permissively
published portfolio, or a thin wrapper around a switch statement that costs
more as a dependency than the code it replaces. Nobody publishes a sensor
node's application state machine, because the application layer is the one
layer that encodes the product, and a product is exactly what nobody gives
away. So the state machine here is the author's, and the chapter's contribution
is the discipline around it: a transition table that can be walked, a coverage
count that proves every row runs, and a stub that is verified against a real
encoder.
\end{note}

The third reason this chapter exists is that it is the fixed point of the
variants matrix. Every other chapter varies how something is executed: polled,
on interrupt, by transfer engine, in a task, in another language. This chapter
varies nothing. It is the application that each of those execution models
hosts, and it is written once so that chapters 11, 12, 16 and 20 can re-host it
without rewriting it. That is why the variants table below claims nothing as
built in full here. It is the only such table in the book, and the reason is
structural rather than an omission.

\subsection*{Prior art and what to reuse}

\begin{table}[H]\centering\small
\begin{tabular}{p{32mm}p{46mm}p{40mm}p{20mm}}
\toprule
Source & What it gives & What it does not & Licence \\
\midrule
The active-object framework and its free video course & A complete hierarchical state machine framework with modelling tools, and a fifty-six lesson course that walks one problem from a superloop to an event-driven architecture. It is the best teaching material on this subject anywhere & Its licence excludes it from a permissively published portfolio, and the course code is the most restrictive item in this book's pool. Read it, do not paste it & GPL-3.0 or commercial; course code AGPL-3.0 \\
Small permissive state machine libraries in general & A dispatcher, a table type and sometimes a code generator & None of them is a standard, and the dispatcher below is twelve lines. A dependency that replaces twelve lines costs more than it saves, and it has to be justified rather than assumed & Various \\
The vendor's example tree for this board & Working peripheral bring-up for everything this chapter drives & No application layer at all. Every example demonstrates one peripheral and then stops, which is precisely the gap & Mixed \\
Majerle's lightweight ring buffer & The event queue, unchanged from chapter 2, with its atomicity assumptions stated & Nothing about what an event means or who may post one & MIT \\
The host test runner and the fake-function framework & A test harness for the same source compiled for the host, and function fakes for the drivers & The fake framework's default history silently discards calls past ten arguments and seventeen calls, which is a trap worth knowing before a test passes for the wrong reason & MIT \\
\bottomrule
\end{tabular}
\caption{Prior art for chapter 8. Four of the five rows are infrastructure. The row that would carry the chapter does not exist, which is the finding rather than a gap to apologise for.}
\end{table}

What is left to write is the state machine, the dispatcher, the event queue
discipline, the coverage instrumentation and the stub. That is the whole
chapter. The reuse here is honest and small: the ring buffer from chapter 2 and
a test runner. Everything else is written, and the repository says so in its
first paragraph.

\subsection*{Parts from the inventory}

\begin{table}[H]\centering\small
\begin{tabular}{p{40mm}p{66mm}p{40mm}}
\toprule
Part & Role & Interface \\
\midrule
NUCLEO-H7A3ZI-Q & The whole chapter. Its three LEDs are the state display and its button is a forced event & Micro USB to the host \\
A USB data cable & Power, programming and the event log on one lead & Micro USB \\
Host PC & Builds the firmware, and builds the same state machine source natively to run the tests & Windows with the Linux subsystem, or a Raspberry Pi \\
\bottomrule
\end{tabular}
\caption{Inventory items used in chapter 8. Nothing is wired, nothing is bought, and no instrument is needed, which is unusual in this volume and is the point: this chapter's acceptance is a test result rather than a measurement.}
\end{table}

\subsection*{System architecture}

\diagram{c08_arch}{The state machine sits between the drivers below it and the codec beside it, and has no knowledge of either. The same source file is compiled twice: once for the board and once for the host, where the drivers are fakes and the test drives the events.}

Three properties of that arrangement are worth stating because they are what
make the chapter testable. The state machine calls no driver directly; it calls
action functions, and the actions are the only code that knows about hardware.
Events arrive only through the queue, never as function calls, so there is one
place where the order of events is decided. And no action blocks, so a
dispatch always completes and the loop always returns to the queue.

\subsection*{Peripheral configuration}

\begin{table}[H]\centering\small
\begin{tabular}{p{22mm}p{26mm}p{24mm}p{30mm}p{30mm}}
\toprule
Peripheral & Mode & Clock source & Pins and function & Interrupt and transfers \\
\midrule
GPIOB, GPIOE & Push-pull outputs, the state display & Peripheral bus & LD1 PB0, LD2 PE1, LD3 PB14, all confirmed & None \\
GPIOC & Input with the wake line & Peripheral bus & PC13, user button & Posts a forced event \\
USART3 & Asynchronous, 115200 8N1, the event log & Peripheral bus & Console pins, confirmed in the board manual & Interrupt driven as in chapter 3 \\
The periodic wake & The real-time clock alarm or the low-power timer of chapter 7 & The fitted 32.768 kHz crystal & None & Posts the tick event \\
The acquisition & Whichever of chapter 6's three builds is linked & As chapter 6 & As chapter 6 & Posts the block event \\
\bottomrule
\end{tabular}
\caption{Peripheral configuration. Everything here is already built in an earlier chapter; this chapter adds only the three LEDs as a display and the button as an event source.}
\end{table}

\subsection*{Wiring}

\diagram{c08_wiring}{Nothing is wired. The three LEDs are on the board and their pins are confirmed from two independent sources, which is what makes it safe to spend them on a state display.}

The state display is worth the three pins. Eight states fit in three LEDs
exactly, and a node whose current state is visible from across the room can be
diagnosed without a debugger, without a terminal and without stopping it. That
matters most in the chapters where the board is sealed inside a measurement
setup and stopping it would change the measurement. It is not a substitute for
the event log, and the chapter says so twice: the LEDs show where the node is,
the log shows how it got there.

\subsection*{Memory and timing budget}

\diagram{c08_mem}{Every byte this chapter owns, and where it lives. The transition table is constant and stays in flash; the context, the queue and the frame buffer are the whole of its static memory. There is no allocator in this node at all.}

\begin{table}[H]\centering\small
\begin{tabular}{p{44mm}p{28mm}p{28mm}p{30mm}}
\toprule
Quantity & Budget & Measured & Margin \\
\midrule
Transition table & under 1 kB in flash & not measured & not measured \\
Node context with frame[64] & 96 B of static memory & 96 B by construction & none needed \\
Event queue & 32 events of 1 B each & 32 B by construction & none needed \\
Coverage counters & 15 rows of 4 B & 60 B by construction & none needed \\
Dynamic allocation & zero bytes & zero by construction & not applicable \\
Worst dispatch, excluding actions & under 200 cycles & not measured & not measured \\
Worst action, the feature step & under 20000 cycles & not measured & not measured \\
Transition rows covered by test & all of them & not measured & not measured \\
\bottomrule
\end{tabular}
\caption{The budget for chapter 8. Three rows are exact because they are static allocations rather than measurements. The two cycle rows are measured with the wrapper built in chapter 6, and the last row is the chapter's acceptance criterion.}
\end{table}

\subsection*{Firmware design (UML)}

\diagram{c08_uml}{The state machine proper. Eight states, nine events, and every edge drawn, including the ones that are only taken when something does not complete. An edge that is not drawn here does not exist in the table either.}

The design rule that makes the rest of the chapter work is that the diagram and
the table are the same object. The table is the source of truth, the diagram is
generated from it or checked against it by hand at every change, and any
behaviour that is not a row in the table is a defect rather than a special
case. The most common way an embedded application becomes untestable is that
one condition is handled with an early return inside an action instead of as a
transition, and after five of those the diagram no longer describes the
program.

\subsection*{Data flow (ASCII)}

\begin{asciiart}
             posted by ISRs                          dispatched in the loop
  +---------------------------+            +-------------------------------------------+
  | tick     (RTC or LPTIM)   |            | for each event in the queue:              |
  | block    (acquisition)    |            |   find the row: (state, event, guard)     |
  | tx_ok / tx_fail (stub)    |  --------> |   run its action, which must not block    |
  | timeout  (guard timer)    |  ring      |   set the next state                      |
  | button   (EXTI PC13)      |  buffer    |   count the row as hit                    |
  | fault    (any action)     |  32 slots  |   show the state on LD1 LD2 LD3           |
  +---------------------------+            +-------------------------------------------+
                                                     |
                                                     v
          +------------------+   feature    +------------------+   frame   +-------------+
          | block of samples | -----------> | one scalar       | --------> | codec, ch 9 |
          +------------------+              +------------------+           +------+------+
                                                                                  |
                                                                                  v
                                          +-----------------------------------------------+
                                          | transmit stub: log the exact bytes, no radio   |
                                          | the host test compares them with the encoder   |
                                          +-----------------------------------------------+
\end{asciiart}

\subsection*{Repository layout}

\begin{plaincode}
nucleo-h7a3-node/
  CMakeLists.txt                  # two builds: the board, and the host tests
  cmake/arm-none-eabi.cmake       # unchanged from chapter 1
  node/node_sm.c                  # the table and the dispatcher, no hardware
  node/node_sm.h                  # states, events, the context type
  node/node_actions.c             # the actions, the only code that knows drivers
  node/node_actions.h             # the interface the host build fakes
  port/board/drivers.c            # real drivers, from chapters 3, 4, 6 and 7
  port/board/leds.c               # three LEDs as a three-bit state display
  port/host/fakes.c               # the same interface, faked, for the tests
  test/test_reachability.c        # every row of the table is exercised
  test/test_payload.c             # the stub's bytes against chapter 9's encoder
  test/events.txt                 # scripted event sequences, one per line
  tools/table_to_dot.py           # the diagram is generated from the table
  docs/states.md                  # the states and events, written before the code
  README.md
\end{plaincode}

\subsection*{Steps}

\begin{steps}
\item \textbf{Write down the states and the events before any code.} This is a
step, not advice. Eight states and nine events, each with one sentence saying
what it means and who may post it, in a file that is committed before
\texttt{node\_sm.c} exists. If a state cannot be described in one sentence it
is two states, and if an event has two meanings it is two events.

\item \textbf{Declare the table types.} A row is a from-state, an event, an
optional guard, an action and a to-state. Nothing else. Putting a timeout value
or a retry count in the row is the first step towards a table that describes
half the behaviour.
\begin{ccode}
/* node_sm.h */
typedef enum { S_INIT, S_IDLE, S_SENSE, S_FEATURE, S_ENCODE,
               S_TX, S_BACKOFF, S_FAULT, S_COUNT } node_state_t;

typedef enum { E_TICK, E_BLOCK, E_FEATURE_DONE, E_FRAME_READY,
               E_TX_OK, E_TX_FAIL, E_TIMEOUT, E_BUTTON, E_FAULT,
               E_COUNT } node_event_t;

typedef struct {
    node_state_t state;
    uint32_t     tx_attempts;
    uint32_t     cycles_ok;
    uint32_t     cycles_dropped;
    int32_t      feature;
    uint8_t      frame[64];
    uint8_t      frame_len;
} node_ctx_t;

typedef struct {
    node_state_t from;
    node_event_t on;
    bool  (*guard)(const node_ctx_t *);
    void  (*action)(node_ctx_t *);
    node_state_t to;
} node_row_t;
\end{ccode}

\item \textbf{Write the table.} It is the specification. Read it top to bottom
and you have read the node's behaviour, which is a property no set of boolean
flags ever has.
\begin{ccode}
/* node_sm.c */
static const node_row_t TABLE[] = {
  { S_INIT,    E_TICK,          NULL,        act_arm,       S_IDLE    },
  { S_IDLE,    E_TICK,          NULL,        act_start_acq, S_SENSE   },
  { S_IDLE,    E_BUTTON,        NULL,        act_force,     S_SENSE   },
  { S_SENSE,   E_BLOCK,         NULL,        act_feature,   S_FEATURE },
  { S_SENSE,   E_TIMEOUT,       NULL,        act_note_lost, S_IDLE    },
  { S_FEATURE, E_FEATURE_DONE,  NULL,        act_encode,    S_ENCODE  },
  { S_ENCODE,  E_FRAME_READY,   NULL,        act_tx_stub,   S_TX      },
  { S_TX,      E_TX_OK,         NULL,        act_done,      S_IDLE    },
  { S_TX,      E_TX_FAIL,       guard_retry, act_backoff,   S_BACKOFF },
  { S_TX,      E_TX_FAIL,       NULL,        act_drop,      S_IDLE    },
  { S_BACKOFF, E_TIMEOUT,       NULL,        act_tx_stub,   S_TX      },
  { S_BACKOFF, E_BUTTON,        NULL,        act_drop,      S_IDLE    },
  { S_FEATURE, E_FAULT,         NULL,        act_fault,     S_FAULT   },
  { S_ENCODE,  E_FAULT,         NULL,        act_fault,     S_FAULT   },
  { S_FAULT,   E_BUTTON,        NULL,        act_reset_ctx, S_INIT    },
};
\end{ccode}
The last three rows are where a flat table starts to cost something. A fault
can be raised from more than one state, and a flat table needs one row per
state that can raise it. Two are enough here; at six it would be worth the
hierarchical states the stretch goals point at, where the fault edge is drawn
once from a superstate. Say which of the two you chose and why, because the
reader will otherwise assume the rows were forgotten.

Two rows share a from-state and an event and are separated only by a guard.
That is deliberate and is the only place ordering matters: the first matching
row wins, so the guarded row comes first and the unguarded row is the fallback.
Write that rule as a comment above the table, because a reader who does not
know it will reorder the rows one day.

\item \textbf{Write the dispatcher.} Twelve lines. If it grows, something that
belongs in the table has been put here instead.
\begin{ccode}
/* node_sm.c */
static uint32_t hit[sizeof TABLE / sizeof TABLE[0]];

void node_dispatch(node_ctx_t *c, node_event_t ev)
{
    for (size_t i = 0; i < sizeof TABLE / sizeof TABLE[0]; i++) {
        const node_row_t *r = &TABLE[i];
        if (r->from != c->state || r->on != ev) continue;
        if (r->guard && !r->guard(c))           continue;
        if (r->action) r->action(c);            /* must not block */
        c->state = r->to;
        hit[i]++;                               /* the coverage counter */
        leds_show_state(c->state);
        log_transition(i, r->from, ev, r->to);
        return;
    }
    log_unhandled(c->state, ev);   /* not an error, but it is never silent */
}
\end{ccode}
An event with no matching row is logged and discarded. That is a deliberate
choice and the alternative is worse: a node that faults on an unexpected event
stops in the field for a reason nobody can reconstruct, while a node that
counts them can be asked afterwards how often it happened.

\item \textbf{Post events through one queue and nowhere else.} The queue is the
ring buffer of chapter 2, one byte per event, written by interrupt handlers and
read by the loop. Actions may post events too, and that is the only recursion
risk in the design: an action that dispatches directly instead of posting can
re-enter the dispatcher with a half-updated context.
\begin{ccode}
/* node_actions.c: actions post, they never dispatch. */
void node_post_from_isr(node_event_t ev)
{
    if (!rb_write(&evq, (uint8_t) ev, 1)) {
        overflow++;            /* counted, never silently discarded */
    }
}

int main(void)
{
    node_ctx_t ctx = { .state = S_INIT };
    for (;;) {
        uint8_t ev;
        while (rb_read(&evq, &ev, 1)) node_dispatch(&ctx, (node_event_t) ev);
        enter_idle();          /* chapter 7 decides how deeply */
    }
}
\end{ccode}

\item \textbf{Write the transmit stub so that it is worth verifying.} It logs
exactly the bytes a radio would have been handed, in order, and nothing else.
No summary, no pretty printing, no field names.
\begin{ccode}
/* node_actions.c */
void act_tx_stub(node_ctx_t *c)
{
    /* There is no radio on this bench. The stub is not a placeholder for
       one: it is the interface a radio would sit behind, and the bytes
       below are the contract chapter 9's encoder has to satisfy. */
    log_hex("TX", c->frame, c->frame_len);
    c->tx_attempts++;
    node_post(c->frame_len > 0 ? E_TX_OK : E_TX_FAIL);
}
\end{ccode}
The log line is machine readable on purpose, because the host test parses it.
A stub that prints something a human likes and a parser does not is a stub that
will be verified once and then trusted forever.

\item \textbf{Show the state on the three LEDs.} Eight states and three LEDs is
an exact fit, which is a small piece of luck worth taking.
\begin{ccode}
/* port/board/leds.c: LD1 PB0, LD2 PE1, LD3 PB14, all confirmed. */
void leds_show_state(node_state_t s)
{
    uint32_t v = (uint32_t) s;                 /* 0 to 7, three bits */
    GPIOB->BSRR = (v & 1u) ? (1u << 0)  : (1u << (0 + 16));
    GPIOE->BSRR = (v & 2u) ? (1u << 1)  : (1u << (1 + 16));
    GPIOB->BSRR = (v & 4u) ? (1u << 14) : (1u << (14 + 16));
}
\end{ccode}

\item \textbf{Prove every row is reachable.} This is the acceptance criterion
and it is a test, not an inspection. The same state machine source is compiled
for the host with faked actions, driven by scripted event sequences, and the
coverage array is checked at the end.
\begin{ccode}
/* test/test_reachability.c, host build, no hardware */
void test_every_row_is_reachable(void)
{
    node_ctx_t c = { .state = S_INIT };
    for (int line = 0; line < script_count(); line++) {
        node_event_t ev;
        node_state_t force;
        script_read(line, &force, &ev);
        c.state = force;               /* the script may place the machine */
        node_dispatch(&c, ev);
    }
    for (size_t i = 0; i < node_row_count(); i++) {
        TEST_ASSERT_MESSAGE(node_row_hits(i) > 0, node_row_name(i));
    }
}
\end{ccode}
Allowing the script to place the machine in a state is a compromise and it is
worth naming: it proves each row runs, not that each row is reachable from
reset. The stronger property is a breadth-first walk of the table from
\texttt{S\_INIT}, which finds rows no sequence of events can ever reach. Write
that walk as a second test; it is twenty lines and it finds real defects.

\item \textbf{Make the stub agree with the encoder byte for byte.} This is the
test that ties the two chapters together, and it fails loudly when either side
changes.
\begin{ccode}
/* test/test_payload.c */
void test_stub_matches_encoder(void)
{
    node_ctx_t c = { .state = S_ENCODE, .feature = -1234 };
    node_dispatch(&c, E_FRAME_READY);            /* runs act_tx_stub */

    uint8_t expect[64];
    size_t n = payload_encode(expect, sizeof expect, -1234);  /* chapter 9 */

    TEST_ASSERT_EQUAL_size_t(n, c.frame_len);
    TEST_ASSERT_EQUAL_HEX8_ARRAY(expect, c.frame, n);
}
\end{ccode}

\item \textbf{Decide what the node does when it cannot proceed.} The fault
state exists and it has exactly one way out, which is the button, and one other
way out which belongs to chapter 12: the watchdog. Write down which conditions
enter it and make sure each of them is one of the rows.
\end{steps}

\subsection*{Build, flash and debug}

\diagram{c08_timing}{The bare-metal dispatch timeline. One event is handled to completion before the next is taken, interrupts only post, and the queue depth is what absorbs a burst. Chapter 20 draws the schedule that replaces this one when a kernel is introduced.}

Two builds from one source tree, and the host build is the one that runs on
every commit.

\begin{shellcode}
cmake -B build-fw -G Ninja -DCMAKE_TOOLCHAIN_FILE=cmake/arm-none-eabi.cmake
cmake --build build-fw -j && cp build-fw/node.bin "$PROBE_DISK"/   # onto the probe disk

cmake -B hostbuild -DNODE_HOST_TESTS=ON
cmake --build hostbuild -j && ctest --test-dir hostbuild --output-on-failure
\end{shellcode}

The diagram is generated from the table rather than drawn beside it, so the two
cannot drift.

\begin{shellcode}
python tools/table_to_dot.py node/node_sm.c > docs/states.dot
dot -Tsvg docs/states.dot > docs/states.svg
\end{shellcode}

\begin{note}{When the node stops in a state and the log says nothing}
Read the LEDs first: three bits give you the state without stopping the board.
Then check, in this order, whether the event that should have moved it was ever
posted, whether the queue overflowed, and whether a row exists for that state
and event at all. The unhandled counter answers the third question directly,
which is why the dispatcher logs it rather than ignoring it. If the node is in
the fault state, the log line that preceded the entry names the action that
posted the fault event.
\end{note}

\subsection*{Verification and acceptance criteria}

\begin{itemize}
\item Every row of the transition table is exercised by the host test suite,
proven by the coverage counters and not by reading the code. A row with zero
hits stops the build.
\item A breadth-first walk of the table from the initial state reaches every
state, and any row it cannot reach is either removed or given a reachable path.
\item The bytes the transmit stub logs are identical to the bytes chapter 9's
encoder produces for the same input, compared as an array and not as a string.
\item No action blocks. This is checked by a test that asserts an upper bound
on the cycles spent inside each action, measured with chapter 6's cycle counter
on the board.
\item There is no dynamic allocation anywhere in the node. Proven by linking
without the allocator and observing that the link succeeds.
\item The event queue overflow counter is exercised deliberately by flooding
the queue, and is observed to increment rather than to lose events silently.
\item The three-bit LED display matches the state reported in the log over a
scripted run, checked once by hand and then trusted.
\end{itemize}

\subsection*{Variants}

\begin{table}[H]\centering\small
\begin{tabular}{p{24mm}p{26mm}p{44mm}p{22mm}p{20mm}}
\toprule
Axis & Variant & What changes & Cost & Built in full in \\
\midrule
Execution model & Polled loop or an interrupt with a flag & The events arrive differently. The table does not change at all, which is the property being demonstrated & Latency and jitter & Chapter 3 \\
Execution model & DMA double buffered & The block event carries a whole half buffer instead of one sample & Cache maintenance & Chapter 6 \\
Execution model & RTOS task set & The dispatcher becomes a task blocking on a queue, and the actions may then block, which changes the design rule & A kernel and its assumptions & Chapter 20 \\
Synchronisation & Queue or direct task notification & The ring buffer is replaced by a kernel object with the same single producer and single consumer contract & Five bytes per task & Chapter 20 \\
Language & C++ & The table becomes a constant array of structures with member function pointers, and the states become a scoped enumeration & A decision about exceptions & Chapter 9 \\
Intelligence and reach & Edge host over a framed link & The transmit stub is replaced by a framed link to a host, and the same bytes cross it & A link and its errors & Chapter 9 and 11 \\
Intelligence and reach & Classifier on the MCU & The feature action becomes a model evaluation and the frame carries a class rather than a scalar & Flash and cycles & Chapter 16 \\
Time and safety & Watchdog over the state machine & A software watchdog fed only from the dispatcher, so a stuck action is caught rather than a stuck loop & A crash record to keep & Chapter 12 \\
\bottomrule
\end{tabular}
\caption{Variants for chapter 8. This is the only variants table in the volume with nothing in the last column pointing here, and that is deliberate: this chapter is the application that all of those variants host, so it is written once and re-hosted rather than rebuilt.}
\end{table}

The thing to notice across those rows is how little the table changes. The
execution model rows change who posts an event and when; the operating system
row changes what the dispatcher blocks on; the language row changes the
syntax. None of them changes the states, the events or the edges between them,
and that is the return on writing the control flow as data rather than as
control flow.

\subsection*{Pitfalls}

\begin{itemize}
\item Keeping a state in a boolean. Two booleans are four states, three are
eight, and none of them appears in the diagram. This is the single most common
way an embedded application becomes untestable.
\item Handling a condition with an early return inside an action instead of as
a transition. After five of those the diagram no longer describes the program.
\item An action that blocks. It does not break anything visibly until the next
chapter adds a second source of events, and then it loses them.
\item An action that calls the dispatcher directly instead of posting an event.
The dispatcher re-enters with a half-updated context.
\item Reordering the table so that an unguarded row precedes a guarded row with
the same from-state and event. The guarded row becomes unreachable and nothing
reports it until the reachability test runs.
\item A stub that prints something readable instead of the exact bytes. It will
be verified once and trusted forever, and it will diverge.
\item Treating an unhandled event as a fault. A node that stops in the field
for an unexpected event cannot be diagnosed; one that counts them can.
\item Using the LED display as the only observability. Three bits say where the
node is and nothing about how it got there.
\end{itemize}

\subsection*{Best practices applied}

\begin{itemize}
\item The control flow is data. The table is the specification, the diagram is
generated from it, and behaviour that is not a row is a defect.
\item The acceptance criterion is a test result rather than an inspection, and
the test fails the build when a row is never reached.
\item Two chapters verify each other. The stub's bytes and the encoder's bytes
are compared as arrays, so a change on either side is caught immediately.
\item Nothing is allocated at run time and every buffer has a stated size, so
the memory figure in the budget table is exact rather than estimated.
\item Every discarded event is counted. There is no silent loss anywhere in
this chapter.
\item The layer boundary is real: the state machine knows no drivers, the
actions know no states, and the host build replaces one side without touching
the other.
\end{itemize}

\subsection*{Stretch goals}

\begin{itemize}
\item Write the breadth-first reachability walk and run it over a deliberately
broken table to confirm it finds the unreachable row.
\item Generate the transition table from the written specification rather than
the other way round, so that the document and the code cannot disagree.
\item Add hierarchical states, so that the fault entry is one edge from a
superstate rather than one edge from each state, and measure what that costs in
table size and in dispatcher complexity. The framework named above solves this
properly and is worth reading before deciding.
\item Record the last sixteen transitions in a small ring in memory that
survives a reset, so a node that faulted can say what it was doing. Whether the
backup memory survives a power cycle without a coin cell is item 10 on the
confirm list.
\item Run the same table under the kernel of chapter 20 without modification
and publish the diff, which should be confined to the dispatcher and the queue.
\end{itemize}

\subsection*{Roadmap and next steps}

Chapter 9 writes the encoder whose output this chapter's stub has been
comparing itself against, and closes the loop: the two chapters are accepted
together or not at all.

For going deeper, the free video course that walks one problem from a superloop
to an event-driven architecture across fifty-six lessons is the best companion
to this chapter anywhere, and its author's framework is the most serious
treatment of hierarchical state machines on small parts. Its code licence means
it is recommended here and never quoted, and Appendix H records that. The
community embedded roadmap makes the same point this chapter makes, that
projects settle questions reading leaves open, and the university specialisation
on real-time embedded systems listed there is where the scheduling mathematics
lives once the dispatcher becomes a task set in chapter 20.

\subsection*{Portfolio evidence}

\begin{itemize}
\item A repository whose README opens with the generated state diagram and
states, in one sentence, that the diagram is generated from the table rather
than drawn beside it.
\item A continuous integration run on a machine with no board attached, showing
the reachability test and the payload comparison passing.
\item The coverage output listing every transition row and its hit count, with
no zeros.
\item A short written specification of the states and events, committed before
the code, so that the order of work is visible in the history.
\end{itemize}

\subsection*{Sources}

Normative references:
\begin{itemize}
\item The board user manual for MB1363, for the LED and button pins and for the
console pins.
\item Reference manual RM0455, for the wake line the button event uses and for
the port configuration registers.
\item The C language standard for the ordering guarantees the dispatcher
relies on, and Arm's application note on memory barrier instructions for
M-profile, which is where chapter 2's queue gets its rules.
\end{itemize}

Reusable implementations:
\begin{itemize}
\item Majerle's lightweight ring buffer, MIT, reused unchanged as the event
queue.\newline
\url{https://github.com/MaJerle/lwrb}
\item The host test runner, MIT, and the fake-function framework, MIT, for the
host build.\newline
\url{https://github.com/ThrowTheSwitch/Unity}\newline
\url{https://github.com/meekrosoft/fff}
\item The active-object framework, GPL-3.0 or commercial, and its free video
course, whose code is AGPL-3.0. Both are recommended reading and neither is
quoted or vendored anywhere in this volume.\newline
\url{https://www.state-machine.com/products/qp}\newline
\url{https://www.state-machine.com/video-course}
\end{itemize}
   % Sense, feature, and a transmit that is a stub
\project{09}{The payload codec and its Python twin}{NUCLEO-H7A3ZI-Q with a host}{Bit packing, CBOR, round-trip property tests}

\begin{keyfacts}
\item[Board] NUCLEO-H7A3ZI-Q, with the host at the other end of the same USB cable
\item[Peripherals] USART3 for the framed link, and nothing else. The codec itself touches no peripheral at all
\item[Toolchain] arm-none-eabi-gcc with CMake as in chapter 1; a native compiler and Python 3 on the host; MicroPython on the board for the third variant
\item[Operating system] Bare metal on the board, anything on the host
\item[Difficulty] 3 of 5
\item[Effort] 3 evenings of about four hours, one of them on the C++ variant
\item[Deliverable] A 40-bit payload whose encoder is C and whose decoder is Python, a round-trip property test over random inputs that runs on the host in continuous integration, and a written justification for not using a standard serialisation format
\end{keyfacts}

\subsection*{Why this project}

Chapter 8 ends with a transmit step that logs a frame. This chapter writes the
code that produces that frame, and writes it twice: an encoder in C that runs on
the board and a decoder in Python that runs on the host. Two implementations of
one specification is normally a defect waiting to happen, and the whole of this
chapter is about turning that liability into the strongest test in the volume.
If the two agree over a hundred thousand random inputs, and both agree with a
set of vectors that a human worked out by hand, the specification is real.

The payload is bit-packed rather than serialised with a standard format, and
that decision has to be justified rather than assumed. Three good permissive
implementations of the obvious alternative exist and are named below. The
argument for packing bits is that this payload is 5 bytes and the same content
in a self-describing format is several times that, on a link where the payload
size is the design constraint. How many times depends entirely on the shape
chosen, and quoting one figure without naming the shape hides the decision:
measured over the golden vectors and 100000 random cases, a map with one-letter
keys averages 21.07 bytes, a positional array 11.07, and a map with full field
names 50.07. The array is the cheapest and gives up the only thing the format
was chosen for, which is being readable without the header. The argument against is everything else:
schema evolution, tooling, and the fact that a self-describing frame can be
read by somebody who does not have your header file. The chapter states both,
picks one, and says what would change the decision.

The third reason is that a codec is the one piece of firmware that can be
tested properly with no hardware at all. It has no timing, no peripherals and
no state. That makes it the right place to introduce property-based testing to
this volume: instead of a handful of examples, assert a property over random
inputs and run it on every commit. Chapter 12 turns that habit into a rig.

\begin{note}{Two implementations generated from one specification is not two implementations}
The field table below generates both the C header and the Python module, which
removes the possibility of the two drifting apart. It also removes their
independence, and an agreement between two outputs of one generator proves only
that the generator is self-consistent. That is why the golden vectors matter:
they are computed by hand from the written bit order, checked into the
repository, and no generator has ever touched them. The round-trip test catches
drift; the vectors catch a specification that was wrong from the start.
\end{note}

\subsection*{Prior art and what to reuse}

\begin{table}[H]\centering\small
\begin{tabular}{p{32mm}p{46mm}p{40mm}p{20mm}}
\toprule
Source & What it gives & What it does not & Licence \\
\midrule
NanoCBOR & The smallest of the three encoders, written for constrained nodes, and effectively public domain so it can be vendored without a thought & It is a serialisation format, so it carries type and key information this payload does not need. It does not pack bits & CC0 \\
zcbor & Schema-driven generation from a written description, used in production, and the closest of the three to the generate-both-sides idea this chapter uses & It generates for a self-describing format rather than for a fixed bit layout, and it brings a build step & Apache-2.0 \\
TinyCBOR & A minimalist alternative with a stable C API and a long history & The same objection as the other two, and no Python counterpart in the same repository & MIT \\
CrustyAuklet's bit packer & Exactly this chapter's idea done well: one format string shared between a C++ implementation and a Python counterpart, so the two sides cannot disagree about layout & It is C++14 and this book is C. That is a real obstacle and not a detail to paper over, and it is why the C++ variant below exists rather than a dependency & MIT \\
Python's standard library & \texttt{ctypes} to call the C encoder from the test, \texttt{struct} for the byte handling, and \texttt{random} with a recorded seed for reproducible cases & No shrinking, so a failing case is whatever random produced rather than the smallest failing case & Standard library \\
\bottomrule
\end{tabular}
\caption{Prior art for chapter 9. The bit packer is the closest thing to this chapter that exists and the language difference is the reason it is read rather than used; the C++ variant below is the honest response to that.}
\end{table}

What is left to write is the field specification, the generator, the C bit
writer, the Python bit reader, the golden vectors and the property test. The
three serialisation libraries are named, compared and then not used, and the
chapter says why in one paragraph rather than leaving the reader to guess that
they were not considered.

\subsection*{Parts from the inventory}

\begin{table}[H]\centering\small
\begin{tabular}{p{40mm}p{66mm}p{40mm}}
\toprule
Part & Role & Interface \\
\midrule
NUCLEO-H7A3ZI-Q & Runs the encoder, and in the third variant the decoder as well & Micro USB to the host \\
A USB data cable & Programming and the framed link on one lead & Micro USB \\
Host PC & Runs the decoder, the property test and the continuous integration job & Virtual COM port at 115200 \\
\bottomrule
\end{tabular}
\caption{Inventory items used in chapter 9. Nothing is wired and nothing is bought. The property test needs no board at all, which is the point of putting it in continuous integration.}
\end{table}

\subsection*{System architecture}

\diagram{c09_arch}{One specification, two generated headers, three consumers. The golden vectors sit outside the generator on purpose, because two outputs of one generator agreeing proves less than it looks.}

The generator is fifty lines and it earns its place by removing a class of
defect rather than by saving typing. A bit layout written twice will diverge,
and it will diverge silently, because both sides will still produce five bytes.
What it cannot do is check that the layout is the one the specification
intended, and that is what the hand-computed vectors are for.

\subsection*{Peripheral configuration}

\begin{table}[H]\centering\small
\begin{tabular}{p{22mm}p{26mm}p{24mm}p{30mm}p{30mm}}
\toprule
Peripheral & Mode & Clock source & Pins and function & Interrupt and transfers \\
\midrule
None, for the codec & The encoder is pure computation with no hardware dependency at all & None & None & None \\
USART3 & Asynchronous, 115200 8N1, carrying the framed link of the edge-host variant & Peripheral bus & Console pins, confirmed in the board manual & Circular transfer as in chapter 4 \\
The CRC unit & The frame checksum of chapter 5, reused unchanged & Peripheral bus & None & None \\
\bottomrule
\end{tabular}
\caption{Peripheral configuration. The first row is the interesting one: a codec that touches no peripheral is a codec that can be tested on a host, which is what makes the property test possible.}
\end{table}

\subsection*{Wiring}

\diagram{c09_wiring}{Nothing is wired. The edge-host variant uses the same cable the debugger is already on, because the probe presents a serial port alongside the debug interface.}

\subsection*{Memory and timing budget}

\diagram{c09_mem}{The 40 bits, drawn the way a reference manual draws a register, with the most significant bit leftmost and first on the wire. This is a payload layout rather than a peripheral register; chapters 5 and 17 draw the silicon registers.}

\begin{table}[H]\centering\small
\begin{tabular}{p{44mm}p{28mm}p{28mm}p{30mm}}
\toprule
Quantity & Budget & Measured & Margin \\
\midrule
Payload & 40 bits & 5 B by construction & none needed \\
Framed length on the wire & 9 B by construction & 9 B by construction & none needed \\
The same content as CBOR, map with short keys & about 17 B by arithmetic & 21.07 B mean & none needed \\
The same content as CBOR, positional array & not estimated & 11.07 B mean & none needed \\
The same content as CBOR, full field names & not estimated & 50.07 B mean & none needed \\
Encoder flash cost & under 512 B & 366 B & +146 B \\
Encoder cycles per frame & under 300 & 2\,955 and 3\,185 & OVER by about 10x \\
One frame at 115200 8N1 & 781 \textmu{}s by arithmetic & not measured & not measured \\
Random cases per test run & 100000 & 100000 by construction & none needed \\
\bottomrule
\end{tabular}
\caption{The budget for chapter 9. The two arithmetic rows are labelled as arithmetic rather than as measurements, because 9 bytes at ten bits each over 115200 is a calculation and the instrument that would confirm it is the one built in chapter 6. The flash and cycle rows were measured on the board on Saturday 3 October 2026.}
\end{table}

\textbf{The cycle budget was wrong by about a factor of ten, and the cause is one
function.} The budget said under 300 cycles per frame. The part says 2\,955, as an
upper bound that includes the measuring loop and a call that does not inline across
translation units.

The measurement is not the suspect. It is the two-point form with a warm-up that
chapters 1 and 2 arrived at, and an order of magnitude is far outside what the
harness could contribute: the harness is a loop and one call, tens of cycles, not
thousands. The encoder really is that expensive, and reading it says why.

\texttt{bw\_put} writes ONE BIT AT A TIME. For each of the forty bits it computes a
source shift, a destination byte index and a mask, then does a read-modify-write of
a byte in memory, under a branch on the bit's value. Forty iterations of that, five
calls to reach them, a \texttt{memset} first, and every instruction fetched from
flash with no instruction cache enabled. Seventy-odd cycles a bit is unremarkable
under those conditions; forty of them is the figure.

\textbf{Two figures are given because the encoder was measured twice and the
answers differ by 7.8 per cent.} The second build differed from the first by a
renamed variable and three extra \texttt{printf} lines. Nothing about the encoder
changed at all.

The new string literals grew \texttt{.rodata}, every byte of code after it moved,
and the hot loop landed differently against the 32 byte flash fetch lines. There is
no instruction cache enabled on this part in this volume yet, so every instruction
is fetched from flash and where a loop sits relative to a fetch boundary changes
what it costs.

This is the third time today that an unrelated edit moved a cycle figure. Chapter 1
watched its delay loop go from 9 cycles an iteration to 19 because four register
writes were added to the reset handler. It is not noise in the measurement, which
repeats to a few parts in a hundred thousand within one build; it is a real
property of the code as placed.

So the rule for every cycle figure in this volume: it belongs to a build, not to a
function, and a difference smaller than about ten per cent between two builds
cannot be attributed to the change under test. Chapter 2's comparison of ordering
barriers carries the same caveat and says so. Enabling the instruction cache, which
is chapter 19, is what would make these figures belong to the code instead.

Against a budget of 300 none of that matters: both measurements are about ten times
over and the conclusion is the same either way.

That is a defensible implementation and an indefensible budget. The bit-by-bit
writer is the clearest possible statement of the specification in
\texttt{docs/bitorder.md}, which is why it was written that way, and the codec's
correctness is the chapter's deliverable rather than its speed. What was wrong was
estimating 300 without counting what the loop does.

Whether it matters depends on the duty cycle, and that is the honest framing rather
than a reflex to optimise. At 64 MHz, 2\,955 cycles is 46 microseconds. One frame
per minute makes it irrelevant. One frame per millisecond makes it 4.6 per cent of
the processor. The node of chapters 8 to 12 sends on the order of one frame a
minute, so the encoder is not the thing to improve, and chapter 12's energy budget
is where that would show if it were.

A byte-oriented writer would be perhaps ten times faster and a great deal harder to
read against the specification. That trade is a stretch goal rather than a
correction, and it is recorded as one.


The nine bytes are worth spelling out because they are the argument for packing
bits. One byte of byte-stuffing overhead, five bytes of payload, two bytes of
checksum and one delimiter. Replace the five with seventeen and the frame is
twenty-one bytes, which is more than twice the air time on a link where air
time is the budget. On a link where it is not, the comparison inverts and the
self-describing format wins, and the chapter says so.

\subsection*{Firmware design (UML)}

\diagram{c09_uml}{The components and what each of them is allowed to know. The encoder knows the field table; the framing layer knows only that it was handed some bytes; the decoder knows the same field table and nothing about the board.}

The interface between the codec and everything else is two functions and a
buffer, and keeping it that narrow is what lets the same encoder be called from
the board, from a host test through a foreign function interface, and from the
C++ variant without change. Anything that widens it, an allocator, a logging
call, a dependency on the node context, costs the chapter its best property.

\subsection*{Data flow (ASCII)}

\begin{asciiart}
  codec/fields.py  (the specification, written once, read by everything)
        |
        +--> tools/gen_codec.py --> payload_fields.h   (C, for the board)
        |                       \-> payload_fields.py  (Python, for the host)
        |
        +--> docs/bitorder.md    (the written bit order: MSB first, bit 39 first on the wire)
                 |
                 v
        test/vectors.json   hand computed from the document, no generator touched it
                 |
   board                                          host
   +-----------------------------+                +-------------------------------+
   | node ctx (chapter 8)        |                | ctypes -> libpayload.so       |
   |   |                         |                |   encode(fields) -> 5 bytes   |
   |   v                         |                |        |                      |
   | payload_encode() -> 5 bytes |                |        v                      |
   |   |                         |   COBS+CRC16   | payload_decode() -> fields    |
   |   v                         |  9 bytes at    |        |                      |
   | lwpkt frame (chapter 5)     | ------------>  |        v                      |
   |   |                         |   115200 8N1   | assert fields == the input    |
   |   v                         |                | 100000 random cases, seeded   |
   | USART3                      |                +-------------------------------+
   +-----------------------------+
\end{asciiart}

\subsection*{Repository layout}

\begin{plaincode}
nucleo-h7a3-codec/
  CMakeLists.txt                  # board target, host shared library, tests
  codec/fields.py                 # the specification, and the only place it lives
  codec/payload.c                 # the bit writer and the bit reader in C
  codec/payload.h
  codec/payload.hpp               # the C++ variant, compile-time field widths
  codec/payload.py                # the Python twin: bit reader and writer
  tools/gen_codec.py              # emits payload_fields.h and payload_fields.py
  docs/bitorder.md                # the written bit order, the normative document
  test/vectors.json               # hand computed, never generated
  test/test_roundtrip.py          # the property test, ctypes into the C encoder
  test/test_vectors.py            # both sides against the hand-computed vectors
  test/test_cpp.cpp               # the C++ variant against the same vectors
  micropython/payload_mpy.py      # the same decoder under MicroPython on the board
  host/listen.py                  # the edge-host variant: frames in, fields out
  README.md
\end{plaincode}

\subsection*{Steps}

\begin{steps}
\item \textbf{Write the field table once, as data.} It is the specification and
everything else is generated from it or checked against it. Widths, signedness
and order, and nothing else: no encoding logic lives here.
\begin{pycode}
# codec/fields.py -- the only place the layout is written.
# Order is most significant bit first. Bit 39 is the first bit on the wire.
FIELDS = [
    # name        bits  signed  note
    ("version",    3,   False,  "format version, 0 to 7"),
    ("flags",      4,   False,  "bit 0 retry, bit 1 forced, bits 2 and 3 spare"),
    ("sequence",   9,   False,  "wraps at 512"),
    ("feature",   18,   True,   "two's complement, scaled by 1/16"),
    ("battery",    6,   False,  "units of 50 mV above 2.00 V"),
]
TOTAL_BITS = sum(f[1] for f in FIELDS)     # 40
assert TOTAL_BITS % 8 == 0, "the payload must be a whole number of bytes"
\end{pycode}

\item \textbf{Write the bit order down in prose and treat that document as
normative.} This is the step that is always skipped and it is the one that costs
a week when a second implementation appears. Three sentences are enough: the
most significant bit of the first field is the most significant bit of byte
zero; fields follow in table order with no padding between them; a signed field
is two's complement in its stated width. Everything else in this chapter is
checked against those three sentences.

\item \textbf{Write the bit writer in C.} It is twenty lines and it is worth
writing rather than depending on, because every dependency that could replace it
brings either a different language or a different format.
\begin{ccode}
/* codec/payload.c: MSB first, no padding between fields. */
#include <stdint.h>
#include <string.h>
#include "payload.h"

typedef struct { uint8_t *buf; size_t cap; size_t bit; } bw_t;

static int bw_put(bw_t *w, uint32_t value, unsigned width)
{
    if (w->bit + width > w->cap * 8u) return -1;
    for (unsigned i = 0; i < width; i++) {
        unsigned src = width - 1u - i;                 /* MSB of the field */
        unsigned dst = w->bit + i;
        uint8_t  bit = (uint8_t) ((value >> src) & 1u);
        uint8_t  mask = (uint8_t) (0x80u >> (dst & 7u));
        if (bit) w->buf[dst >> 3] |= mask;
        else     w->buf[dst >> 3] = (uint8_t) (w->buf[dst >> 3] & ~mask);
    }
    w->bit += width;
    return 0;
}

size_t payload_encode(uint8_t *out, size_t cap, const payload_t *p)
{
    if (cap < PAYLOAD_BYTES) return 0;
    memset(out, 0, PAYLOAD_BYTES);
    bw_t w = { out, cap, 0 };
    if (bw_put(&w, p->version,  PAYLOAD_W_VERSION))  return 0;
    if (bw_put(&w, p->flags,    PAYLOAD_W_FLAGS))    return 0;
    if (bw_put(&w, p->sequence, PAYLOAD_W_SEQUENCE)) return 0;
    if (bw_put(&w, (uint32_t) p->feature & PAYLOAD_M_FEATURE,
                                PAYLOAD_W_FEATURE))  return 0;
    if (bw_put(&w, p->battery,  PAYLOAD_W_BATTERY))  return 0;
    return PAYLOAD_BYTES;
}
\end{ccode}
The signed field is masked to its width before it is written, which is the line
that decides whether the codec is portable. A negative \texttt{int32\_t} shifted
right is implementation defined in C, so the value is converted to unsigned and
masked rather than shifted arithmetically.

\item \textbf{Write the reader, in C and in Python, and sign-extend explicitly.}
Sign extension is where a bit-packed codec goes wrong, because it works for
every positive test value.
\begin{pycode}
# codec/payload.py
from payload_fields import FIELDS, TOTAL_BITS

def _bits(data: bytes, start: int, width: int) -> int:
    v = 0
    for i in range(width):
        b = start + i
        v = (v << 1) | ((data[b >> 3] >> (7 - (b & 7))) & 1)
    return v

def decode(data: bytes) -> dict:
    if len(data) * 8 != TOTAL_BITS:
        raise ValueError(f"expected {TOTAL_BITS // 8} bytes, got {len(data)}")
    out, pos = {}, 0
    for name, width, signed, _ in FIELDS:
        v = _bits(data, pos, width)
        if signed and (v >> (width - 1)) & 1:
            v -= 1 << width                      # two's complement, explicit
        out[name] = v
        pos += width
    return out
\end{pycode}

\item \textbf{Compute the golden vectors by hand and check them in.} Take three
cases, work the bits out on paper from the written bit order, and write the
expected bytes into a file. Do this before running either implementation, so
that neither can talk you into its own answer.
\begin{plaincode}
test/vectors.json
[
  {"fields": {"version": 1, "flags": 0, "sequence": 0,
              "feature": 0, "battery": 0},
   "bytes": "2000000000"},
  {"fields": {"version": 7, "flags": 15, "sequence": 511,
              "feature": 131071, "battery": 63},
   "bytes": "FFFF7FFFFF"},
  {"fields": {"version": 7, "flags": 15, "sequence": 511,
              "feature": -1, "battery": 63},
   "bytes": "FFFFFFFFFF"},
  {"fields": {"version": 1, "flags": 2, "sequence": 5,
              "feature": -1, "battery": 40},
   "bytes": "2405FFFFE8"}
]
\end{plaincode}
The second and third vectors look like one case and are two, which is the error
this step exists to prevent and which an earlier printing of this chapter made.
The largest positive value of an eighteen-bit two's complement field is 131071,
which is a zero followed by seventeen ones, not eighteen ones. That zero is bit
seventeen of the stream, so byte two is \texttt{7F} and not \texttt{FF}. Five
bytes of \texttt{FF} is also a real case and it decodes to a feature of
\(-1\). Carry both.

The maximum case catches a width error, the all-ones case catches reading 131071
as eighteen set bits, and the negative case catches sign extension. The last
vector is the one to work out slowly; if your implementation disagrees with it,
do not change the vector.

\item \textbf{Run the C encoder from the test through a foreign function
interface.} The host build produces a shared library from exactly the same
source the board links, and the test calls into it. Nothing is reimplemented for
the test.
\begin{pycode}
# test/test_roundtrip.py
import ctypes, random, json
from payload import decode
from payload_fields import FIELDS

lib = ctypes.CDLL("./build/libpayload.so")

class Payload(ctypes.Structure):
    _fields_ = [("version", ctypes.c_uint32), ("flags", ctypes.c_uint32),
                ("sequence", ctypes.c_uint32), ("feature", ctypes.c_int32),
                ("battery", ctypes.c_uint32)]

def _random_case(rng):
    out = {}
    for name, width, signed, _ in FIELDS:
        lo = -(1 << (width - 1)) if signed else 0
        hi = (1 << (width - 1)) - 1 if signed else (1 << width) - 1
        out[name] = rng.randint(lo, hi)
    return out

def test_roundtrip_100k():
    rng = random.Random(20260920)          # the seed is recorded on purpose
    buf = (ctypes.c_uint8 * 8)()
    for _ in range(100000):
        case = _random_case(rng)
        n = lib.payload_encode(buf, 8, ctypes.byref(Payload(**case)))
        assert n == 5, f"encoder returned {n} for {case}"
        assert decode(bytes(buf[:5])) == case, case
\end{pycode}
The seed is fixed and recorded so that a failure is reproducible. A test whose
random input changes on every run finds more defects and reports them in a form
nobody can act on, so the compromise here is a fixed seed plus a second job that
runs with a random one and prints the seed it used.

\item \textbf{Build the C++ variant.} This is the language variant this chapter
owns, and the reason to build it is a real one: the field widths can be checked
at compile time instead of at run time, and the offsets stop being written by
hand.
\begin{cppcode}
// codec/payload.hpp -- C++17, exceptions and RTTI both disabled.
#include <array>
#include <cstdint>
#include <cstddef>

template <std::size_t Width, bool Signed>
struct Field { static constexpr std::size_t width = Width;
               static constexpr bool is_signed = Signed; };

using Version  = Field<3,  false>;
using Flags    = Field<4,  false>;
using Sequence = Field<9,  false>;
using Feature  = Field<18, true>;
using Battery  = Field<6,  false>;

constexpr std::size_t total_bits =
    Version::width + Flags::width + Sequence::width
  + Feature::width + Battery::width;
static_assert(total_bits % 8 == 0, "payload must be a whole number of bytes");
static_assert(total_bits == 40, "layout changed: regenerate the vectors");
\end{cppcode}
Two static assertions replace two runtime checks and one class of review
comment. What it costs is the argument the reader actually wants: exceptions
cost roughly 10 to 30 kB of unwinder tables, which on 2 MB of flash is a choice
rather than a necessity, and this build disables them anyway because nothing
here throws. Compile both variants, run both against the same vectors, and put
the two size outputs in the README rather than asserting which is smaller.

\item \textbf{Run the same decoder on the board under MicroPython.} The board
definition for this exact part exists in the interpreter's main branch and
prebuilt firmware is published, so this variant costs a drag-and-drop and no
build at all. It proves something worth proving: one decoder source serving the
host and the target.
\begin{shellcode}
cp firmware.hex /media/$USER/NODE_H7A3ZI/       # the probe's disk, chapter 1
mpremote cp codec/payload.py :payload.py
mpremote cp payload_fields.py :payload_fields.py
mpremote exec "import payload; print(payload.decode(bytes.fromhex('2405FFFFE8')))"
\end{shellcode}
Two constraints apply and they are worth stating before somebody discovers them
at the prompt. The decoder must avoid anything outside the interpreter's
subset, which for this module means no f-string formatting in the hot path and
no type annotations that require a module the firmware does not carry. And the
board definition's declared feature list is empty, so nothing here should
promise networking or USB host on that firmware without building it.

\item \textbf{Build the edge-host variant over the framed link.} This is the
second variant this chapter owns. The frame is the one from chapter 5: byte
stuffing, a 16-bit checksum from the hardware unit, and a delimiter.
\begin{pycode}
# host/listen.py -- frames in, fields out, on the virtual COM port
import serial
from cobs import decode_cobs          # chapter 5's Python side
from crc16 import crc16               # the same polynomial the unit uses
from payload import decode

port = serial.Serial("/dev/ttyACM0", 115200, timeout=1)
buf = bytearray()
while True:
    b = port.read(1)
    if not b:
        continue
    if b == b"\x00":                  # the delimiter
        frame = decode_cobs(bytes(buf)); buf.clear()
        body, got = frame[:-2], int.from_bytes(frame[-2:], "big")
        if crc16(body) != got:
            print("checksum mismatch, frame discarded"); continue
        print(decode(body))
    else:
        buf += b
\end{pycode}
The host is the same Python module the property test uses, which is the point:
the edge host is not a new implementation, it is the existing decoder given a
different source of bytes.

\item \textbf{Write the comparison with the serialisation format and commit
to a decision.} Encode the same fields with one of the three libraries named
above, measure the result, and put both numbers in the README with the
reasoning. State what would change the decision: a link where air time is not
the constraint, a payload that has to be read by a party who does not have the
header, or a schema that will change more than once a year.
\end{steps}

\subsection*{Build, flash and debug}

\diagram{c09_timing}{Nine bytes on the wire at 115200, and what the host is doing while they arrive. The frame duration is arithmetic rather than a measurement, and chapter 6 built the instrument that would confirm it.}

\begin{shellcode}
python tools/gen_codec.py codec/fields.py --out-c codec --out-py codec
cmake -B build -DPAYLOAD_HOST=ON && cmake --build build -j   # libpayload.so
python -m pytest test -q                                     # vectors, then 100k cases
cmake -B fw -DCMAKE_TOOLCHAIN_FILE=cmake/arm-none-eabi.cmake
cmake --build fw -j && probe-rs run --chip STM32H7A3ZITx fw/codec_demo.elf
\end{shellcode}

\begin{note}{When the two sides disagree by one byte}
The order of suspicion, and it is almost always the first: bit order inside a
field, where one side writes the most significant bit first and the other writes
the least. Then byte order of a multi-byte field that happens to straddle a
boundary. Then sign extension, which is invisible until a negative value
appears, which is why the third golden vector is negative. Then a width that was
changed in the specification and not regenerated on one side, which the static
assertion in the C++ header catches at compile time and the C build does not.
Print both byte strings in hexadecimal before reasoning about them; a
disagreement that looks like a field error is often a length error.
\end{note}

\subsection*{Verification and acceptance criteria}

\begin{itemize}
\item The round-trip property holds over 100000 random cases with a recorded
seed, run on the host in continuous integration with no board attached, and a
second job runs the same test with a random seed and prints it.
\item Both implementations reproduce the hand-computed vectors exactly,
including the all-ones case and the negative case.
\item The bytes the encoder produces are identical to the bytes chapter 8's
transmit stub logs for the same input, which is that chapter's acceptance
criterion and this one's.
\item The C++ variant produces byte-identical output to the C variant for every
vector and for the whole random run, and both size outputs are published rather
than one being asserted as smaller.
\item The Python decoder runs unmodified under the interpreter on the board and
returns the same fields for the vector bytes.
\item The framed link delivers a frame end to end and a deliberately corrupted
frame is discarded by the checksum rather than decoded into plausible rubbish.
\item The generator is run in continuous integration and the build fails if the
generated files differ from the committed ones, so a specification change cannot
be merged without the regeneration.
\end{itemize}

\subsection*{Variants}

\begin{table}[H]\centering\small
\begin{tabular}{p{24mm}p{26mm}p{44mm}p{22mm}p{20mm}}
\toprule
Axis & Variant & What changes & Cost & Built in full in \\
\midrule
Language & C++ & Field widths become template parameters and the layout is checked by static assertion at compile time rather than by a test at run time & A decision about exceptions and a second toolchain configuration & Here \\
Language & MicroPython on the target & The same Python decoder file runs on the board, so one source serves the host and the target & Determinism and direct control & Here, and chapter 1 as the five-minute path \\
Intelligence and reach & Edge host over a framed link & The frame crosses the virtual COM port to a host that runs the same decoder module the tests use & A link and its error handling & Here, and chapter 11 for the cellular case \\
Peripheral substitution & CBOR instead of bit packing & Self-describing, readable without the header, about three times the size for this payload & Air time, and a dependency & Named here, not built \\
Language & Host in Python & The decoder is Python here because the twin has to be readable; a different host language would be a second twin and a second liability & Another implementation to keep honest & Chapter 3 \\
Execution model & DMA circular receive on the host link & The frame arrives without the processor polling, which is what chapter 4 built & Cache maintenance & Chapter 4 \\
Intelligence and reach & Classifier output instead of a scalar & The feature field becomes a class index and a confidence, which changes the widths and nothing else in this chapter & Flash and cycles & Chapter 16 \\
\bottomrule
\end{tabular}
\caption{Variants for chapter 9. Three are built in full here because the matrix assigns them here; the serialisation format is named and compared and deliberately not built, which is a decision rather than an omission.}
\end{table}

The three built variants make one point between them. A specification written as
data, with the layout generated from it, survives being re-implemented in a
second language on the same machine, in a third language on a different machine,
and across a link. What does not survive that treatment is a layout that lives
only in the offsets of one encoder, and that is the failure mode the whole
chapter is arranged to prevent.

\subsection*{Pitfalls}

\begin{itemize}
\item Not writing the bit order down. Every other pitfall here is a consequence
of this one.
\item Shifting a signed integer right to extract a field. The result is
implementation defined in C. Convert to unsigned and mask.
\item Forgetting to sign-extend on the way back. It works for every positive
test value, which is why one golden vector is negative.
\item Testing only with values that are easy to read in hexadecimal. The
all-ones case and the boundary case are the ones that find width errors.
\item Letting the generated files drift from the specification. Run the
generator in continuous integration and fail the build on a difference.
\item Believing that two generated implementations agreeing proves the layout is
right. It proves the generator is self-consistent. The hand-computed vectors are
what prove the layout.
\item Choosing a serialisation format without measuring the size, or choosing
bit packing without measuring what the format would have cost.
\item Using a random seed that is not recorded. A failure nobody can reproduce
is a failure nobody will act on.
\end{itemize}

\subsection*{Best practices applied}

\begin{itemize}
\item The specification is data, written once, and both implementations are
generated from it. The written bit order is normative and is committed before
either implementation.
\item The independence the generator removes is restored deliberately, with
vectors computed by hand and checked in.
\item Every dependency names its licence before any code is written around it,
and the one that fits best is read rather than used because it is in a different
language, which is said plainly.
\item The test runs on the host with no hardware, on every commit, and the seed
that produced any failure is in the output.
\item The alternative that was not chosen is measured rather than dismissed, and
the conditions that would reverse the decision are written down.
\item Two chapters verify each other. The bytes here and the bytes chapter 8
logs are compared as arrays.
\end{itemize}

\subsection*{Stretch goals}

\begin{itemize}
\item Replace the seeded random generator with a property-based testing library
that shrinks a failing case to the smallest one that still fails. Hypothesis is
the usual choice for Python; read its licence before depending on it, because
this volume's prior-art pool does not record it.
\item Generate the decoder for a third language from the same field table and
run it against the same vectors, which costs an evening and tests the generator
rather than the codec.
\item Add a format version negotiation: the three version bits exist for that
reason, and a decoder that refuses an unknown version cleanly is a small piece of
work with a large payoff.
\item Measure the encoder in cycles with chapter 6's counter and compare the C,
the C++ and a handwritten table-free version, then publish all three rather than
the fastest.
\item Extend the round-trip test across the framed link itself, so that the
property is checked end to end on hardware in the rig chapter 12 builds.
\end{itemize}

\subsection*{Roadmap and next steps}

Chapter 10 takes the frame this chapter produces and asks what it cost to make
it, by splitting the duty cycle measured in chapter 7 into phases with marker
pins. Chapter 11 replaces the stub and the virtual COM port with a modem and a
command set, and the frame crosses the same interface unchanged.

For going deeper, the community embedded roadmap is the map worth following,
and the free open textbook that grew out of a university course is the closest
thing to a curriculum for the chapters after this one. On the serialisation
question specifically, the schema-driven generator named above is worth reading
even if the decision goes the other way, because its approach of generating from
a written description is the same argument this chapter makes in a smaller
space. Appendix H lists the courses; the one that walks a single problem from a
polled loop to an event-driven architecture is the best free companion to the
node chapters, though its code licence means it is recommended and never quoted.

\subsection*{Portfolio evidence}

\begin{itemize}
\item A repository whose README opens with the bit-layout figure and the
sentence that the layout is generated from one specification and checked against
vectors no generator produced.
\item A continuous integration run on a machine with no board attached, showing
the vectors passing, the round trip passing over 100000 cases, and the generated
files matching the committed ones.
\item Two size outputs, C and C++, published side by side with no claim about
which is better until the numbers are there.
\item A terminal recording of the same decoder file returning the same fields on
the host and on the board under the interpreter.
\end{itemize}

\subsection*{Sources}

Normative references:
\begin{itemize}
\item The written bit order document in this repository, which is normative for
both implementations and for the vectors.
\item The C language standard, for the rules on shifting signed integers and on
conversion to unsigned that the encoder depends on.
\item The board user manual for MB1363, for the virtual COM port pins used by
the edge-host variant.
\item Reference manual RM0455, for the checksum unit's configuration that
chapter 5 established and this chapter reuses unchanged.
\end{itemize}

Reusable implementations:
\begin{itemize}
\item NanoCBOR, CC0, the smallest of the three encoders and effectively public
domain.\newline
\url{https://github.com/bergzand/NanoCBOR}
\item zcbor, Apache-2.0, schema-driven and used in production.\newline
\url{https://github.com/NordicSemiconductor/zcbor}
\item TinyCBOR, MIT, the minimalist alternative.\newline
\url{https://github.com/intel/tinycbor}
\item CrustyAuklet's bit packer, MIT, which shares one format string between a
C++ implementation and a Python counterpart. It is C++14 where this book is C,
which is why it is read rather than used.\newline
\url{https://github.com/CrustyAuklet/bitpacker}
\item Majerle's lightweight packet library, MIT, which supplies the framing this
chapter's edge-host variant rides on.\newline
\url{https://github.com/MaJerle/lwpkt}
\end{itemize}
   % The payload codec and its Python twin
\project{10}{Where the energy goes: wake, sense, compute, send}{NUCLEO-H7A3ZI-Q with the nRF-PPK2}{Marker pins, per-phase charge accounting}

\begin{keyfacts}
\item[Board] NUCLEO-H7A3ZI-Q, with the nRF-PPK2 as the ammeter and as this book's logic recorder
\item[Peripherals] RTC for the wake, GPIO marker pins on free Zio pins, I2C for the sense phase, USART3 held silent during a capture
\item[Toolchain] arm-none-eabi-gcc with CMake on the target; Python on the host for capture, decode and the ledger
\item[Operating system] Bare metal
\item[Difficulty] 4 of 5
\item[Effort] 3 evenings of about four hours
\item[Deliverable] A per-phase charge ledger for one duty cycle, with the marker trace that proves the split, whose phases sum to the measured cycle total within five percent
\end{keyfacts}

\subsection*{Why this project}

Chapter 7 ends with one number: the charge drawn by one complete duty cycle,
measured with the nRF-PPK2 in ammeter mode. That number answers the question a
customer asks, which is how long a battery lasts. It answers no question an
engineer can act on. If the cycle costs more than the budget allows, the number
does not say whether the cost is in the wake-up transient, in the sensor read,
in the computation, or in the transmit stub that is standing in for a radio.
Halving the computation is pointless if the wake-up transient owns most of the
charge, and that is a thing you cannot know from a single figure.

This chapter turns the single number into a ledger. The method is small and it
is the whole chapter: the firmware raises a pattern on a few spare output pins
at each phase boundary, the instrument records those pins on its digital inputs
time-aligned with the current it is already sampling, and a host script cuts the
current trace at the marker edges and integrates each piece separately. The
acceptance test is then arithmetic rather than opinion. The phases must sum to
the whole within five percent, and if they do not, either the split is wrong or
there is charge being drawn in a phase nobody named.

There is a second reason this chapter exists on this bench in particular. There
is no oscilloscope, no logic analyser and no Joulescope here, confirmed on
Sunday 20 September 2026. The instrument's digital inputs are the only
time-aligned logic record available, and they are sampled by the same instrument
that samples the current, which is exactly the property a phase split needs. A
separate logic capture on a separate clock would have to be aligned after the
fact, and alignment error would appear as charge attributed to the wrong phase.
Using the ammeter's own inputs removes that class of problem, or rather moves it
inside the instrument where it can be characterised once.

\begin{note}{The sum is the test}
A per-phase breakdown that does not reconcile against the whole-cycle figure is
a plausible-looking picture and nothing more. Measure the cycle total first, as
chapter 7 does, before any markers exist. Then split it. The reconciliation
residual is the honest error bar on every per-phase number in the ledger, and it
belongs in the report next to them.
\end{note}

\subsection*{Prior art and what to reuse}

\begin{table}[H]\centering\small
\begin{tabular}{p{34mm}p{46mm}p{38mm}p{20mm}}
\toprule
Source & What it gives & What it does not & Licence \\
\midrule
Chapter 7 of this volume & The measurement set-up, the instrument in ammeter mode, the whole-cycle charge figure, and the clock-restore trap that keeps current high after a wake & One number per cycle and no way to attribute it & This volume \\
The Python interface to the power instrument & Scripted capture, start and stop, and export from the host so a run is repeatable rather than a screenshot & It is GPL-2.0. Install it from its package index and do not fork it into a portfolio repository & GPL-2.0 \\
Golioth's method article on this instrument & The practical set-up: source mode against ammeter mode, what the two leads do, and where the ground reference has to be & No phase splitting. It produces the single number that chapter 7 already has & Article, read and cite \\
Published empirical work on cellular module energy & The presentation model this chapter copies: a trace divided into named states with a charge for each & A different part, a different modem and a host that is not this board & Paper \\
Sensor-node lifetime estimation literature & The arithmetic from charge per cycle to battery months, including self-discharge and the duty-cycle algebra & It assumes a per-phase split exists and never says how to obtain one & Papers \\
\bottomrule
\end{tabular}
\caption{Prior art for chapter 10. The instrument's host library is installed and never vendored, because its licence is the one category in this book that a portfolio repository must keep at arm's length.}
\end{table}

What is left to write is the marker protocol, the decoder and the
reconciliation. No peer-reviewed method exists for this instrument, which is
recorded in this volume's provenance notes as a gap rather than an oversight, so
the rigour here is original work and has to be stated in enough detail that
somebody else can repeat it. That means writing down the marker encoding, the
cost of raising a marker, the assumed alignment between the digital inputs and
the current samples, and what happens to charge drawn while no marker is
asserted.

The phase splitting by marker pin is the author's contribution. Everything else
in the chapter is borrowed and credited: the instrument, its host library, the
set-up method and the presentation format.

\subsection*{Parts from the inventory}

\begin{table}[H]\centering\small
\begin{tabular}{p{40mm}p{66mm}p{40mm}}
\toprule
Part & Role & Interface \\
\midrule
NUCLEO-H7A3ZI-Q & The node under measurement. Its own probe is a load, so it is considered in the budget rather than ignored & Micro USB for flashing, then the question of what powers the target during a capture \\
nRF-PPK2 & Ammeter for the current, and the logic recorder for the marker pins on its digital inputs & USB to the host; measurement leads to the board; digital inputs to the marker pins \\
Jumper wires & Marker pins to the instrument's digital inputs, and the ground reference they share & Female headers on both sides, which is a constraint and not a detail \\
X-NUCLEO-IKS4A1 & Optional, to make the sense phase real rather than a delay. One shield at a time, never two & Arduino header \\
Host PC & Capture, decode, integrate, report & USB \\
\bottomrule
\end{tabular}
\caption{Inventory items used in chapter 10. Nothing is bought and nothing is soldered. Whether male-to-male jumper wires exist in the drawer is an open question on the bench list, and every header here is female.}
\end{table}

\subsection*{System architecture}

\diagram{c10_arch}{One instrument samples two things at once. The current path and the marker path share the instrument's timebase, which is the property that makes the phase split trustworthy rather than approximate.}

The firmware does not know it is being measured. It raises and lowers output
pins at points it would have to mark anyway for any kind of tracing, and those
pins happen to be wired to an instrument. That separation matters: the same
firmware runs when the instrument is not connected, and the marker writes stay
in the production build because removing them would change the timing of the
thing being characterised.

\subsection*{Peripheral configuration}

\begin{table}[H]\centering\small
\begin{tabular}{p{22mm}p{26mm}p{24mm}p{30mm}p{30mm}}
\toprule
Peripheral & Mode & Clock source & Pins and function & Interrupt and transfers \\
\midrule
RTC & Wake timer, periodic & LSE at 32.768 kHz, fitted on this board and confirmed & None & Wake-up interrupt from Stop mode \\
GPIO markers & Push-pull output, written through the set and reset register & Peripheral bus & Free Zio pins. Which pins are free alongside a fitted shield is on the confirm list & None \\
I2C1 & Fast mode, sense phase & Peripheral bus & Arduino D14 and D15 by convention, checked in the board manual & Polled here; DMA in chapter 4 \\
USART3 & Console, held silent during a capture & Peripheral bus & Confirm in the board manual & None during a capture \\
DWT cycle counter & Free-running & Core clock & None & None. Cross-check only \\
\bottomrule
\end{tabular}
\caption{Peripheral configuration for chapter 10. The console is configured and then not used during a capture, because a line of text at 115200 is charge that belongs to no phase and would appear in the reconciliation residual.}
\end{table}

The cycle counter earns its place as an independent check on the phase
boundaries. It costs nothing to read, it runs on a completely different
mechanism from the instrument, and if the phase durations it reports disagree
with the marker edge times the instrument recorded, one of the two is wrong and
it is worth knowing which before trusting either. Whether the counter works
with no debugger attached is on the confirm list for this silicon, and if it
turns out to need one, the cross-check is lost and the marker edges stand alone.

\subsection*{Wiring}

\diagram{c10_wiring}{The measurement wiring. Three open questions are drawn as questions: which jumper carries the current measurement and what it isolates, how many digital inputs the instrument really has and at what threshold, and whether the instrument can supply the board at all in source mode.}

\begin{note}{Three limits that are not negotiable}
The instrument measures at most 1 A, so nothing that draws more goes through it.
The board's logic is 3.3 V and nothing above that reaches a pin. The cellular
modems of chapter 11 need 2 A peaks from their own supply and are never fed from
the board or through this instrument, which is why the send phase here is a stub
and the real modem waits for the next chapter.
\end{note}

\subsection*{Memory and timing budget}

\diagram{c10_mem}{The ledger, drawn as a shape rather than as numbers. Every quantity in it is unmeasured until the bench produces it, and the figure exists to fix the arithmetic and the reconciliation band before any number is written into it.}

\begin{table}[H]\centering\small
\begin{tabular}{p{44mm}p{28mm}p{28mm}p{30mm}}
\toprule
Quantity & Budget & Measured & Margin \\
\midrule
Charge per cycle, whole & from chapter 7 & not measured & not measured \\
Wake phase charge & 15 percent of cycle & not measured & not measured \\
Sense phase charge & 25 percent of cycle & not measured & not measured \\
Compute phase charge & 10 percent of cycle & not measured & not measured \\
Send phase charge & 40 percent of cycle & not measured & not measured \\
Sleep phase charge & 10 percent of cycle & not measured & not measured \\
Reconciliation residual & below 5 percent & not measured & not measured \\
Cost of one marker write & below 0.1 percent of cycle & not measured & not measured \\
Marker code in flash & 256 B & not measured & not measured \\
Ledger record in RAM & 128 B & not measured & not measured \\
\bottomrule
\end{tabular}
\caption{The budget for chapter 10. The percentages in the budget column are the hypothesis to be tested, written down before the measurement so that agreement is evidence and disagreement is a finding. Nothing enters the measured column until the instrument has produced it.}
\end{table}

Writing the budget down first is not a formality. A per-phase measurement is
easy to read backwards: once a trace is on screen it is very natural to decide
that the tall part must be the transmit and the flat part must be the sleep,
and to stop asking. Stating the expected split in advance makes a surprise
visible as a surprise.

\subsection*{Firmware design (UML)}

\diagram{c10_uml}{The duty cycle as a state machine, with the marker code each state asserts. The code is raised on entry and never cleared inside a state, so every instant of the cycle carries exactly one code and charge cannot fall between two phases.}

The design rule that makes the arithmetic work is that the marker code changes
exactly once per transition and never inside a state. If a marker is cleared at
the end of a phase and the next is raised at the start of the next, the interval
between the two writes belongs to no phase, and on a part running at 280 MHz
that interval is short but it is not zero. Writing the whole code in one store
to the set and reset register removes the gap entirely: the old code and the new
code change in the same cycle, and there is no instant where the pins show
neither.

An encoding decision follows from that. A one-hot code with one pin per phase is
the easiest to read on a screen and the easiest to decode, and it costs one pin
per phase. A binary code on three pins covers eight phases and frees pins for
other uses, at the cost of a decoder that has to handle the transient if the
pins are not written atomically. Since the set and reset register does write them
atomically, the binary code carries no real risk here, and the choice comes down
to how many pins the fitted shield leaves free. That is an open item on the
bench list and it is decided at the bench, not in this text.

\subsection*{Data flow (ASCII)}

\begin{asciiart}
  board                                    instrument                      host
  +-------------------------------+        +--------------------+   +-----------------------+
  | RTC wake  -> marker code 001  |        |                    |   | capture script        |
  | sense     -> marker code 010  | pins   | 8 digital inputs   |   |   start, stop, export |
  | compute   -> marker code 011  +------->|   same timebase    |   |        |              |
  | send stub -> marker code 100  |        |        +           |   |        v              |
  | sleep     -> marker code 000  |        | current sampler    |-->| decode marker edges   |
  |                               |        |   ammeter mode     |USB|        |              |
  | VDD, GND                      | leads  |        ^           |   |        v              |
  |   +---------------------------+------->|--------+           |   | integrate per phase   |
  +-------------------------------+        +--------------------+   |        |              |
                                                                    |        v              |
                                                                    | reconcile vs total    |
                                                                    | ledger.json           |
                                                                    +-----------------------+
\end{asciiart}

\subsection*{Repository layout}

\begin{plaincode}
nucleo-h7a3-energy-ledger/
  CMakeLists.txt
  ld/stm32h7a3zi.ld              # from chapter 1, unchanged
  src/phase.h                    # the marker protocol, one header, no code
  src/phase.c                    # one store to BSRR, nothing else
  src/cycle.c                    # the duty cycle state machine of chapter 8
  src/sense.c                    # shield read, or a calibrated delay if bare
  src/main.c
  host/capture.py                # drives the instrument, writes a raw export
  host/decode.py                 # marker edges -> phase intervals
  host/ledger.py                 # integrate, reconcile, emit ledger.json
  host/report.py                 # the figure in this chapter, regenerated
  host/requirements.txt          # the instrument library, installed not vendored
  docs/method.md                 # the alignment assumption, written down
  ledgers/                       # one ledger.json per run, committed
  README.md
\end{plaincode}

The ledger directory is committed deliberately. A run that is not kept is a
demonstration; a run that is kept next to the firmware revision that produced it
is a measurement, and chapter 12 turns the committed ledgers into a regression
test that fails a build when the charge moves.

\subsection*{Steps}

\begin{steps}
\item \textbf{Reproduce the whole-cycle number before splitting anything.} Run
chapter 7's measurement again on the firmware you are about to instrument, and
write the figure down. If it does not reproduce within a few percent of what
chapter 7 recorded, stop here: something in the set-up has moved, and a phase
split built on an unstable baseline will attribute the instability to whichever
phase happens to be running when it occurs.

\item \textbf{Write the marker protocol as a header before any code.} The
protocol is the contract between the firmware and the decoder, and it is short
enough to read in one screen.
\begin{ccode}
/* phase.h: the marker protocol. Three pins, binary code, one store per change.
 * Code 000 means sleep and is also the reset state, so an unpowered or
 * halted board reads as sleep rather than as an undefined phase.          */
#ifndef PHASE_H
#define PHASE_H

typedef enum {
    PHASE_SLEEP   = 0u,   /* 000 */
    PHASE_WAKE    = 1u,   /* 001  clock restore included, see chapter 7 */
    PHASE_SENSE   = 2u,   /* 010 */
    PHASE_COMPUTE = 3u,   /* 011 */
    PHASE_SEND    = 4u,   /* 100  a stub until chapter 11 */
    PHASE_FAULT   = 7u    /* 111  raised by the fault handler, never cleared */
} phase_t;

void phase_init(void);           /* clock the port, configure the three pins */
void phase_set(phase_t p);       /* exactly one store to BSRR                */

#endif
\end{ccode}
The fault code is worth the pin it costs. A trace that ends in a code nobody
raises deliberately tells you the run is invalid without anyone reading a log.

\item \textbf{Make a marker cost one store.} The set and reset register takes a
set mask in its low half and a reset mask in its high half, so setting some pins
and clearing others is a single write with no read, no modify and no interrupt
window in the middle.
\begin{ccode}
/* The three marker pins are adjacent, at PHASE_PIN0 and up, on one port.
 * Adjacency is a convenience, not a requirement; a scattered set works with
 * a lookup table instead of a shift.                                      */
#define PHASE_PIN0   4u
#define PHASE_MASK   (7u << PHASE_PIN0)

void phase_set(phase_t p)
{
    uint32_t set   = ((uint32_t) p << PHASE_PIN0) & PHASE_MASK;
    uint32_t clear = (~set) & PHASE_MASK;
    PHASE_PORT->BSRR = set | (clear << 16);   /* one store, atomic on the pins */
}
\end{ccode}
Two properties follow and both matter for the arithmetic. The write is atomic
with respect to interrupts, so a marker change cannot be interrupted halfway and
leave the pins showing a code that was never a phase. And the port clock must
already be running: a pin written before its clock is enabled does nothing and
reports no error, which is the same trap as chapter 1 and it is worth checking
first when the markers are simply absent from a capture.

\item \textbf{Bound the cost of a marker and put the bound in the report.} The
markers are inside the thing being measured, so their cost is part of the
answer. Measure it the only way the bench allows: run the whole cycle with the
markers compiled in and again with \texttt{phase\_set} reduced to an empty
function, and compare the whole-cycle charge from the instrument. The difference
is an upper bound on the total marker cost, because it also contains whatever
the compiler did differently. State the bound; do not claim the markers are
free.

\item \textbf{Wire the digital inputs and settle three questions at the bench.}
The inventory records eight digital inputs on this instrument. How many there
are, what logic threshold they use, at what rate they are sampled, whether the
export carries them alongside the current samples, and how tightly they are
aligned with the current timebase are all open items on the confirm list. None
of them is written as fact in this chapter. Settle them by capturing a known
square wave from a timer output on a marker pin, at a period you have set
yourself, and reading it back out of the export.
\begin{shellcode}
python host/capture.py --seconds 5 --out runs/align_check.raw
python host/decode.py runs/align_check.raw --expect-period-ms 10
\end{shellcode}
If the decoded period differs from the programmed period by more than the
instrument's own sample interval, the alignment assumption is not safe and the
phase boundaries inherit that error.

\item \textbf{Capture one duty cycle with room either side.} A capture that
starts inside the cycle produces a first phase that is truncated and a ledger
that silently under-counts it. Start the capture in sleep, wait for at least two
complete cycles, and stop in sleep, then let the decoder find whole cycles by
looking for the code transitions rather than by trusting the capture window.

\item \textbf{Decode the marker edges into phase intervals.} The decoder is the
piece worth writing carefully, because every later number depends on it. Treat a
code change as an instantaneous boundary, reject any interval shorter than the
instrument's sample interval as a decoding artefact rather than as a phase, and
count how many it rejected so the count can go in the report.
\begin{pycode}
def intervals(samples, codes, dt):
    """Yield (code, t_start, t_end) for each run of a constant marker code."""
    out, start, cur = [], 0.0, codes[0]
    for i in range(1, len(codes)):
        if codes[i] != cur:
            out.append((cur, start, i * dt))
            start, cur = i * dt, codes[i]
    out.append((cur, start, len(codes) * dt))
    return [iv for iv in out if (iv[2] - iv[1]) >= dt]
\end{pycode}

\item \textbf{Integrate charge per phase, then reconcile.} Charge is the
integral of current over time, which on a uniformly sampled trace is a sum times
the sample interval. Integrate each phase separately, sum the phases, and
compare against the integral of the whole cycle computed independently from the
same samples without reference to the markers at all.
\begin{pycode}
def ledger(current_a, codes, dt):
    per = {}
    for code, t0, t1 in intervals(current_a, codes, dt):
        i0, i1 = int(t0 / dt), int(t1 / dt)
        per[code] = per.get(code, 0.0) + sum(current_a[i0:i1]) * dt
    total = sum(current_a) * dt              # computed without the markers
    split = sum(per.values())
    residual = abs(total - split) / total if total else 1.0
    return per, total, residual              # residual is the acceptance test
\end{pycode}
The two totals are computed by different routes on purpose. If the residual were
defined as the difference between the phase sum and itself it would always be
zero and the test would prove nothing.

\item \textbf{Attribute the residual rather than hiding it.} A residual above
five percent has a small number of usual causes, and they are distinguishable.
Charge drawn in an interval the decoder rejected as too short shows up as a
rejection count above zero. Charge drawn while the fault code is asserted shows
up as a phase nobody expected. Charge drawn by something that is not the target
at all, such as the on-board probe, shows up as a floor under every phase
including sleep, which is the reason the measurement jumper and what exactly it
isolates is on the confirm list. Report which of these it was.

\item \textbf{Run the experiment the ledger is for.} With a per-phase split in
hand, reorder the cycle and measure again. The classic question is whether it
pays to finish the computation quickly at full clock and return to sleep, or to
run slower and longer, and the answer depends entirely on the ratio between the
compute phase charge and the fixed cost of the wake transient. On a part chosen
for speed the expectation is that racing to sleep wins, but the point of the
chapter is that the expectation is now testable.
\begin{shellcode}
python host/capture.py --seconds 12 --out runs/race.raw
python host/ledger.py runs/race.raw --baseline ledgers/baseline.json
\end{shellcode}
\end{steps}

\subsection*{Build, flash and debug}

\diagram{c10_timing}{One duty cycle as the instrument records it. The current trace and the three marker lanes share one timebase, and the phase boundaries are the marker edges rather than features picked out of the current by eye.}

The build is chapter 1's build with two files added, so there is nothing new to
say about it. What is new is that the board is being measured while it runs,
which changes how it is flashed and how it is recovered.

\begin{shellcode}
cmake -B build-fw -G Ninja -DCMAKE_TOOLCHAIN_FILE=cmake/arm-none-eabi.cmake
cmake --build build-fw -j
cp build-fw/firmware.bin "$PROBE_DISK"/   # onto the probe disk
\end{shellcode}

Flash first, then set up the measurement, then reset the board and capture. A
debugger attached during a capture changes the answer: the core is not allowed
to enter the deeper low-power modes with a debug connection active in the usual
configuration, and the probe itself draws current that may or may not be inside
the measurement path depending on the jumper. Detach, reset, capture.

\begin{note}{When every phase looks the same}
In order of likelihood: the port clock for the marker pins was never enabled, so
the pins never move and the decoder sees one long phase; the digital input
ground and the board ground are not common, so the inputs float and read as
noise rather than as a code; the pins chosen collide with a fitted shield, which
is why the free-pin question is on the confirm list; the capture window is
shorter than one cycle; or the instrument is in source mode when it should be in
ammeter mode, in which case the current trace is of the instrument supplying the
board rather than of the board drawing. Check them in that order.
\end{note}

To recover a board that stops responding after a low-power experiment, hold the
reset line low while the flashing tool connects, which is the connect-under-reset
path chapter 7 also needs. A part that has entered a deep low-power mode with no
wake source configured is not damaged and is not permanently inaccessible; it is
simply not listening, and a reset is enough. Nothing in this chapter touches
option bytes, and nothing in this book does until the question of whether a bad
option-byte write can leave this board unrecoverable without external tooling
has been answered.

\subsection*{Verification and acceptance criteria}

\begin{itemize}
\item The per-phase charges sum to the independently computed cycle total within
five percent, on ten consecutive cycles from one capture, with the residual for
each cycle printed rather than averaged away.
\item The whole-cycle total reproduces chapter 7's figure on the same firmware
within the repeatability that chapter 7 itself established, so the ledger is
known to be splitting the same quantity chapter 7 measured.
\item The marker cost bound is stated as a number with its method, obtained by
comparing whole-cycle charge with markers compiled in against a build where the
marker function is empty.
\item The decoder reports its rejection count, and a run with any rejected
interval is reported as such rather than quietly accepted.
\item The phase durations from the marker edges agree with the durations the
cycle counter reports from inside the firmware, which is an independent check on
the instrument's alignment. If the cycle counter turns out to need a debugger on
this part, this criterion is dropped and the chapter says so.
\item A ledger file is emitted for every run, carries the firmware revision, and
is committed next to it.
\item Every number in the ledger names the instrument that produced it. Nothing
is written in the measured column of the budget table until it has been
measured.
\end{itemize}

\subsection*{Variants}

\begin{table}[H]\centering\small
\begin{tabular}{p{24mm}p{26mm}p{44mm}p{22mm}p{20mm}}
\toprule
Axis & Variant & What changes & Cost & Built in full in \\
\midrule
Time and safety & RTC wake against low-power timer & The wake source changes the shape and the cost of the wake phase, which the ledger now resolves separately from everything else & Different wake latency & Chapter 7 \\
Execution model & DMA for the sense phase & The core sleeps through the transfer instead of polling, which moves charge from the compute phase to the sense phase rather than removing it & Setup cost per transfer & Chapter 4 \\
Execution model & An RTOS task set & Phases become tasks and the marker writes move into the task bodies, with the idle task raising the sleep code & Tick cost during sleep & Chapter 20 \\
Intelligence and reach & A real transmit instead of the stub & The send phase stops being a delay and starts being a modem, on its own supply and outside this instrument's 1 A limit & A second supply and a second measurement path & Chapter 11 \\
Intelligence and reach & Compute on the host instead of the node & The compute phase disappears and the send phase grows, which the ledger can price directly & More airtime & Chapter 16 \\
Language & Host tooling in Python & The capture, decode and ledger scripts here & None & Chapter 3 \\
Peripheral substitution & Markers on timer output channels & A timer channel can raise a marker without the core, which removes the marker cost from the phase being measured & More configuration & Chapter 17 \\
\bottomrule
\end{tabular}
\caption{Variants for chapter 10. This chapter builds no variant in full, which is deliberate: it contributes a measurement method that every one of these variants can then be priced with, and the book's matrix assigns each variant to the chapter where it teaches most.}
\end{table}

\subsection*{Pitfalls}

\begin{itemize}
\item Clearing a marker at the end of a phase and raising the next at the start
of the next. The interval between the two writes belongs to no phase and its
charge appears in the residual. One store per transition removes the gap.
\item Trusting the digital inputs to be aligned with the current samples without
checking. The alignment is on the confirm list for a reason, and an alignment
error appears as charge attributed to the neighbouring phase, which is the most
plausible-looking kind of wrong answer.
\item Leaving the console enabled during a capture. A single line of text at
115200 is charge drawn outside any phase and it will not be obvious in the trace.
\item Measuring with the debugger attached. The low-power behaviour differs and
the probe's own consumption may be inside the measurement path.
\item Reading the split off the current trace by eye and using the markers only
as decoration. If the boundaries come from the shape of the current, the method
is circular and the reconciliation test no longer means anything.
\item Averaging the residual across cycles. A method that reconciles on average
and not per cycle is hiding a real effect, usually a phase whose duration varies.
\item Choosing marker pins that collide with a fitted shield. Which Zio pins with
spare general-purpose function remain free alongside a shield is an open bench
question and it is answered before the pins are chosen, not after.
\end{itemize}

\subsection*{Best practices applied}

\begin{itemize}
\item The budget is written before the measurement, so that agreement is
evidence and disagreement is a finding rather than an embarrassment.
\item Two independent routes compute the quantity that the acceptance test
compares. A test that compares a number with itself is not a test.
\item The instrument's own timebase carries both signals, which removes a whole
class of alignment error rather than correcting for it afterwards.
\item The cost of the instrumentation is bounded and reported, because the
markers are inside the system being characterised.
\item Open questions are written as questions. The digital input count, the
threshold, the alignment, the measurement jumper and the free Zio pins are all
marked as items for the bench, and none of them is asserted.
\item The dependency with the restrictive licence is installed from its package
index and is not copied into the repository.
\end{itemize}

\subsection*{Stretch goals}

\begin{itemize}
\item Emit the ledger as a machine-readable record with the firmware revision,
the instrument serial number and the capture parameters, and have chapter 12's
rig compare it against the committed baseline automatically.
\item Add a sixth phase for the clock restore inside the wake phase, which
chapter 7 identifies as the expensive part of waking this family, and find out
what fraction of the wake charge it owns.
\item Reproduce the whole measurement with the shield fitted and the sense phase
reading a real sensor at a real rate, and compare against the delay-based stub,
which prices the sensor rather than the firmware around it.
\item Sweep the core frequency across the voltage scales the datasheet allows
and plot charge per cycle against frequency. The shape of that curve is the
racing-to-sleep question in one figure and no published result covers this part.
\item Extend the decoder to cope with a marker pin that is shared with a shield
signal by ignoring intervals shorter than a stated threshold, and report how
often that happens rather than assuming it does not.
\end{itemize}

\subsection*{Roadmap and next steps}

Chapter 11 replaces the send stub with a real modem, which is where most of the
charge in a node like this eventually goes, and it has to do so on a separate
supply because the modem's peak demand is outside this instrument's range.
Chapter 12 takes the ledger produced here and makes it a gate: a build whose
charge per cycle has moved by five percent turns the build red without anybody
at the bench.

For the wider path, the community embedded engineering roadmap distinguishes
firmware work from embedded Linux work and from hardware work and says which to
weight at which stage, and the accompanying essay on bridging the gap between
reading and building makes the same argument this volume makes, which is that
projects teach what reading does not. Its central advice is the reason this
chapter insists on a measurement rather than a description. Appendix H lists the
courses; for the energy material specifically, there is no course on this bench's
instrument, and that absence is itself recorded in this volume's provenance notes
as one of the gaps the work fills.

\subsection*{Portfolio evidence}

\begin{itemize}
\item A public repository containing the firmware, the host tooling, the method
document that states the alignment assumption, and one committed ledger per run.
\item A figure showing one duty cycle with the current trace and the marker lanes
on one timebase, generated by the report script from a committed run rather than
captured from a screen.
\item The reconciliation table: ten cycles, the phase sum, the independent total
and the residual for each, with nothing averaged.
\item A short written finding on the racing-to-sleep question for this part, with
the method stated well enough that somebody with the same instrument could
disagree with it.
\end{itemize}

\subsection*{Sources}

Normative references:
\begin{itemize}
\item Reference manual RM0455, for the low-power modes, what each retains, the
wake sources and the set and reset register behaviour on this part. Not RM0433,
which documents its sibling.
\item The board user manual for MB1363, for the current measurement jumper, what
it isolates, and which header pins are free alongside a fitted shield.
\item The STM32H7A3xI datasheet, for the supply currents at each voltage scale
and for the clock limits that bound the racing-to-sleep experiment.
\item The instrument's own user guide, for the digital input count, the logic
threshold, the sample rate and the export format.
\end{itemize}

Reusable implementations:
\begin{itemize}
\item The Python interface to the power instrument, GPL-2.0. Install it from its
package index; do not fork it into a portfolio repository.\newline
\url{https://github.com/IRNAS/ppk2-api-python}
\item Golioth's method article on measuring current consumption with this
instrument, which is the practical set-up this chapter starts from.\newline
\url{https://blog.golioth.io/measuring-current-consumption-with-power-profiler-kit-ii/}
\item The vendor's own power examples for this exact board, under
\texttt{Projects/}\allowbreak\texttt{NUCLEO-H7A3ZI-Q/}\allowbreak\texttt{Examples/PWR/}, which are the reference for
entering and leaving the low-power modes correctly.
\item The peer-reviewed empirical study of narrow-band module energy whose
per-state presentation format this chapter copies, noting that its host is a
single-board computer and not this part.\newline
\url{https://jtde.telsoc.org/index.php/jtde/article/view/955}
\item The community embedded engineering roadmap, for where this chapter sits in
a longer path.\newline
\url{https://github.com/m3y54m/Embedded-Engineering-Roadmap}
\end{itemize}
   % Where the energy goes: wake, sense, compute, send
\project{11}{An AT engine that never blocks}{NUCLEO-H7A3ZI-Q with the Waveshare SIM7020E}{Asynchronous command queue, unsolicited results, timeouts, backoff}

\begin{keyfacts}
\item[Board] NUCLEO-H7A3ZI-Q, with the Waveshare SIM7020E narrow-band modem on its own supply
\item[Peripherals] A USART with circular DMA receive and idle-line detection, a general timer for the timeout wheel, GPIO for the module's power control line
\item[Toolchain] arm-none-eabi-gcc with CMake on the target; a Python fault injector on the host
\item[Operating system] Bare metal, which is the whole point of the chapter
\item[Difficulty] 4 of 5
\item[Effort] 4 evenings of about four hours
\item[Deliverable] A command engine whose main loop never waits, which survives one hundred injected faults with every one ending in a defined state, and which stops in a named state when no SIM card is fitted
\end{keyfacts}

\subsection*{Why this project}

The usual driver for a modem of this kind is a chain of functions that each
write a command and then sit in a loop waiting for a reply. It works on a bench
with a good signal and it fails everywhere else, because everything the module
does that is not an immediate reply looks to that code like a stopped system. A
reply that arrives after the code gave up is discarded, and then answers the
next command. A message the module volunteers on its own, which is ordinary
behaviour and not an error, is read as the reply to whatever command happens to
be outstanding. A network registration that takes tens of seconds blocks the
sensor loop for tens of seconds. None of this is exotic; it is the normal life
of a cellular module, and a blocking driver has no vocabulary for it.

This chapter builds the alternative, which is small enough to hold in your head:
a queue of command records, one command in flight at a time, a per-command
timeout taken from the module's own documentation rather than guessed, a table
of unsolicited result codes dispatched independently of whatever command is
running, and an exponential backoff with a cap for the cases where retrying is
the right answer. The main loop calls one function that returns immediately, and
every path through the engine ends somewhere named.

The reason this is a chapter rather than a dependency is a gap in the prior art
that is worth stating precisely. Majerle's cellular library is MIT licensed, is
well engineered, explicitly supports the SIM7070G that is also on this bench, and
even bundles a messaging client. It requires an operating system and has no
bare-metal path. That is exactly the gap this chapter fills, and it is also the
reason chapter 20 exists: once there is a kernel, that library becomes the right
answer and this engine becomes the thing you understand it with. The same
library does not list the SIM7600E-H or the SIM7020E at all, and the silicon
vendor's own cellular middleware covers none of the three modems on this bench,
targeting a different set of module vendors entirely. There is no maintained
permissive component that does this job for this module on bare metal.

\begin{note}{The narrow-band module has no point-to-point protocol}
On the SIM7600E-H there is a choice: run a point-to-point link and put a network
stack on the microcontroller, or drive the module by commands. The SIM7020E
removes the choice. It offers no point-to-point protocol, so every byte of
payload is handed over inside a command and every result comes back as text.
That is not a restriction to design around; it decides the architecture, and it
is why this is necessarily a command-set chapter.
\end{note}

\subsection*{Prior art and what to reuse}

\begin{table}[H]\centering\small
\begin{tabular}{p{34mm}p{46mm}p{38mm}p{20mm}}
\toprule
Source & What it gives & What it does not & Licence \\
\midrule
Majerle's cellular library & A complete, maintained engine with a command queue, unsolicited result handling and a bundled messaging client, and it explicitly supports the SIM7070G that is also on this bench & It requires an operating system and has no bare-metal path, and it does not list the SIM7600E-H or the SIM7020E & MIT \\
The silicon vendor's cellular middleware & Nothing applicable here & It covers none of the three modems on this bench, targeting a different set of module vendors & Vendor licence \\
Majerle's host-side parser for the radio coprocessor & The clearest readable model of how a command queue, a line classifier and an unsolicited dispatcher fit together & It speaks a different command set, so the structure transfers and the commands do not & MIT \\
A peer-reviewed empirical study of this exact module & Realistic expectations for signal level, registration time and failure rate, which is what the retry and backoff policy has to be sized against & Its host is a single-board computer running a general-purpose operating system, not a bare-metal microcontroller & Paper \\
Chapters 2 and 4 of this volume & The byte path: a ring buffer and a circular DMA receive with idle-line detection, which is what makes a non-blocking engine possible at all & No framing of replies and no notion of a command & This volume \\
\bottomrule
\end{tabular}
\caption{Prior art for chapter 11. The one library that would do this job needs a kernel, which is the gap this chapter fills and the argument chapter 20 answers from the other side.}
\end{table}

What is left to write is the engine itself: the command record, the queue, the
timeout wheel, the line classifier, the unsolicited dispatcher, the backoff
policy and the fault injector that proves the whole thing. The byte path below
it is chapters 2 and 4, unchanged. The structural model is borrowed from the
host-side parser named above, which is the clearest published statement of how
these pieces fit together even though it speaks to different hardware.

One more thing is borrowed and it is the most important: the module's own AT
command manual states a maximum response time for each command. Those numbers
are the only defensible source for a per-command timeout. A timeout invented at
the keyboard is a number that will be wrong in the field in one direction or the
other, and the direction that fails quietly is the one that is too long.

\subsection*{Parts from the inventory}

\begin{table}[H]\centering\small
\begin{tabular}{p{40mm}p{66mm}p{40mm}}
\toprule
Part & Role & Interface \\
\midrule
NUCLEO-H7A3ZI-Q & Runs the engine & Micro USB to the host for flashing and the console \\
Waveshare SIM7020E & The narrow-band module under test. No point-to-point protocol, so command set only & Serial at 3.3 V logic, plus a power control line \\
Joy-it SBC-POW-BB or a USB supply & The module's own supply. It is never fed from the board & Barrel or USB, with the ground made common to the board \\
Renkforce USB/TTL cable & Talking to the module directly from the host, before any firmware exists, and again when the firmware disagrees with the module & Serial at 3.3 V. Disconnect the red 5 V lead \\
Waveshare SIM7070G & The comparison module. The library named above supports it, which makes it the control case & Serial \\
Host PC & Terminal, then the fault injector & USB \\
\bottomrule
\end{tabular}
\caption{Inventory items used in chapter 11. A SIM card is not in the inventory list, and the acceptance test is written so that its absence is a pass rather than an obstacle.}
\end{table}

\subsection*{System architecture}

\diagram{c11_arch}{The engine sits between an application that never waits and a byte path that never parses. Each layer has one job, and the two most common designs for this problem both come from collapsing two of these layers into one.}

Two collapses are worth naming because both are common and both are expensive.
Parsing inside the receive interrupt puts an unbounded amount of work at high
priority and makes every later timing measurement meaningless. Letting the
application call a send-and-wait function puts the modem's latency into the
sensor loop, which is the problem this chapter exists to remove. The layering in
the figure is not ceremony; each boundary is there because of a failure that
happens without it.

\subsection*{Peripheral configuration}

\begin{table}[H]\centering\small
\begin{tabular}{p{22mm}p{26mm}p{24mm}p{30mm}p{30mm}}
\toprule
Peripheral & Mode & Clock source & Pins and function & Interrupt and transfers \\
\midrule
USART to modem & Asynchronous, 115200 8N1 to start, module default confirmed first & Peripheral bus & A free USART on the Zio header. Which instance reaches Arduino D0 and D1 is on the confirm list & Circular DMA receive plus idle line; DMA transmit \\
USART3 console & Asynchronous, 115200 & Peripheral bus & Confirm in the board manual & Interrupt driven, chapter 3 \\
General timer & Up counter, 1 kHz tick & Peripheral bus & None & Update interrupt drives the timeout wheel \\
GPIO power control & Push-pull output & Peripheral bus & One free Zio pin & None \\
GPIO status input & Input with pull, optional & Peripheral bus & One free Zio pin & None. Polled \\
\bottomrule
\end{tabular}
\caption{Peripheral configuration for chapter 11. The baud rate is confirmed against the module before it is written into a header, because a module that has been left in a different rate by a previous experiment answers nothing and looks like a dead module.}
\end{table}

The timeout wheel does not need a dedicated timer. A millisecond tick from any
source works, and on a part with fifteen 16-bit timers and two 32-bit ones there
is no shortage. What it does need is to be independent of the command path, so
that a command which never completes still has something counting on its behalf.

\subsection*{Wiring}

\diagram{c11_wiring}{Two supplies, one ground. The module's peak demand is the reason for the first, and the fact that serial signalling is referenced to ground is the reason for the second. The cross between transmit and receive is the single most common wiring fault on this bench.}

\begin{note}{The module takes 2 A peaks and never takes them from the board}
A cellular module draws short, high current peaks when its transmitter turns on.
The figure on this bench is 2 A, which is more than the board can supply and
more than the measurement instrument of chapters 7 and 10 will pass. The module
gets its own supply, the grounds are made common so the serial signalling has a
reference, and the power instrument stays out of this path entirely. If the
module resets whenever it tries to register, the supply is the first suspect and
usually the only one.
\end{note}

A second safety point applies to the USB to serial cable used for talking to the
module directly. Disconnect its red 5 V lead. The module has its own supply and
a second one arriving down the serial cable is at best redundant.

\subsection*{Memory and timing budget}

\diagram{c11_mem}{Where the engine's bytes live. Everything is statically allocated, because an engine whose failure mode is exhaustion of a heap has replaced one unbounded behaviour with another.}

\begin{table}[H]\centering\small
\begin{tabular}{p{44mm}p{28mm}p{28mm}p{30mm}}
\toprule
Quantity & Budget & Measured & Margin \\
\midrule
Engine code in flash & 6 kB & not measured & not measured \\
Command queue, 8 slots & 512 B & not measured & not measured \\
Receive ring & 1 kB & not measured & not measured \\
Line assembly buffer & 256 B & not measured & not measured \\
Unsolicited result table & 256 B & not measured & not measured \\
Stack used by the engine & 512 B & not measured & not measured \\
Time in the engine per tick & below 20 \textmu{}s & not measured & not measured \\
Longest path with interrupts off & below 5 \textmu{}s & not measured & not measured \\
Commands faulted without a hang & 100 of 100 & not measured & not measured \\
\bottomrule
\end{tabular}
\caption{The budget for chapter 11. The two timing rows are what "never blocks" means as a number, and they are measured with the cycle counter rather than asserted. Nothing enters the measured column until the bench produces it.}
\end{table}

The receive ring is sized for the longest burst the module can produce between
two visits to the engine, not for the longest reply. Those are different
quantities and confusing them is how a ring that looks generous overflows. The
longest reply matters for the line assembly buffer instead, and a line longer
than that buffer is truncated deliberately and reported as a truncated line
rather than silently joined to the next one.

\subsection*{Firmware design (UML)}

\diagram{c11_uml}{The engine as a state machine. Every arrow out of a waiting state is either a reply, a timeout or an unsolicited message, and there is no arrow that leads nowhere. The two stopped states are outcomes and not faults.}

The state machine has one property that carries most of its value: the states
where the engine waits are the states where the application is free. In the
waiting states the engine holds a deadline and nothing else, and the function
the main loop calls returns immediately after checking whether the deadline has
passed and whether a complete line has arrived. Nothing in the engine ever spins.

Unsolicited messages are handled outside the state machine on purpose. A message
the module volunteers is not a reply and must not be treated as one, even when
one is outstanding. The line classifier decides which it is before the state
machine sees it, using the simple and reliable rule that a final result code
ends a command and anything matching an entry in the unsolicited table is
dispatched to that entry's handler no matter what state the engine is in.

\subsection*{Data flow (ASCII)}

\begin{asciiart}
  application           AT engine                     byte path             module
  +---------------+    +-------------------------+   +----------------+    +----------+
  | sensor loop   |    | queue: 8 command slots  |   | DMA tx         |    |          |
  |   |           |    |   +---------------+     |-->| USART tx       |--->|          |
  |   | at_submit |--->|   | cmd, timeout, |     |   +----------------+    | SIM7020E |
  |   |           |    |   | retries, done |     |   +----------------+    |          |
  |   | at_poll   |--->|   +---------------+     |   | DMA rx, circular|<---|          |
  |   |  returns  |    | one in flight at a time |<--| idle line        |   +----------+
  |   |  at once  |    |   deadline from the     |   | ring buffer      |       |
  |   v           |    |   module's own manual   |   +----------------+        | own
  | done callback |<---+-------------------------+          |                  | 2 A
  +---------------+    | line classifier              <-----+                  | supply
                       |   final result -> completes the command in flight
                       |   intermediate  -> appended to that command's result
                       |   unsolicited   -> dispatched from the table, any state
                       +-------------------------------------------------------+
\end{asciiart}

\subsection*{Repository layout}

\begin{plaincode}
nucleo-h7a3-at-engine/
  CMakeLists.txt
  src/at_engine.c              # queue, state machine, timeout wheel
  src/at_engine.h              # the whole public surface, three functions
  src/at_lines.c               # line assembly and classification
  src/at_urc.c                 # the unsolicited table and its dispatch
  src/at_backoff.c             # exponential with a cap and jitter
  src/modem_sim7020e.c         # the command set for this module, nothing else
  src/uart_dma.c               # from chapter 4, unchanged
  src/ringbuf.c                # from chapter 2, unchanged
  src/main.c
  test/test_engine.c           # host build, no hardware, chapter 12 runs it
  test/faults.py               # the hundred injected faults
  test/transcripts/            # recorded module replies, good and damaged
  docs/timeouts.md             # every timeout with the manual page it came from
  README.md
\end{plaincode}

The timeouts document is not decoration. It is the file a reviewer opens first,
because every number in it should be traceable to a page of the module's own
manual, and the ones that are not should be marked as measured on the bench with
the instrument that measured them.

\subsection*{Steps}

\begin{steps}
\item \textbf{Talk to the module from the host before writing any firmware.}
Power it from its own supply, connect the serial cable with the red 5 V lead
disconnected, and confirm it answers. This settles the baud rate, the line
ending the module expects and whether echo is on, all of which otherwise become
firmware bugs.
\begin{shellcode}
python -m serial.tools.miniterm /dev/ttyUSB0 115200 --eol CRLF
# then, one at a time:  AT   ATE0   AT+CMEE=2   AT+CPIN?   AT+CSQ
\end{shellcode}
Turning echo off with \texttt{ATE0} and turning verbose errors on with
\texttt{AT+CMEE=2} are the first two commands the engine will send, for the same
reasons you want them here: echo doubles every line the parser has to consider,
and a numeric error code with no text is a support request waiting to happen.

\item \textbf{Build the byte path and never parse inside it.} The receive side is
chapter 4's circular DMA with idle-line detection, feeding chapter 2's ring
buffer. The interrupt does one thing: it records how far the transfer counter
has moved and sets a flag. Line assembly happens in the engine, at whatever
priority the main loop runs at.

\item \textbf{Define the command record before writing the engine.} The record is
the contract, and most design questions answer themselves once it exists.
\begin{ccode}
typedef enum { AT_OK, AT_ERROR, AT_TIMEOUT, AT_ABANDONED } at_result_t;

typedef void (*at_done_fn)(at_result_t r, const char *last_line, void *ctx);

typedef struct {
    const char  *text;         /* without the line ending; the engine adds it */
    uint32_t     timeout_ms;   /* from the module's manual, never invented    */
    uint8_t      max_retries;  /* 0 means send once and report what happened  */
    uint8_t      flags;        /* AT_F_NO_ECHO_CHECK, AT_F_BINARY_TAIL, ...   */
    at_done_fn   done;         /* called once, in every outcome, including    */
    void        *ctx;          /* timeout and abandonment                     */
} at_cmd_t;
\end{ccode}
The completion callback is called exactly once for every command submitted, in
every outcome. That single rule removes a whole family of leaks in which the
caller waits for a callback that a failure path never reaches.

\item \textbf{Write the engine as a function that returns immediately.} Three
functions are the entire public surface: submit, poll and cancel.
\begin{ccode}
/* Called from the main loop as often as convenient. Never waits, never spins,
 * and performs at most one state transition per call so the worst-case time
 * through it is bounded and measurable.                                     */
void at_poll(at_engine_t *e, uint32_t now_ms)
{
    const char *line;
    while ((line = at_next_line(e)) != NULL) {     /* bounded: ring is finite */
        if (at_urc_dispatch(e, line))              /* unsolicited, any state  */
            continue;
        at_feed_reply(e, line);                    /* belongs to the command  */
    }
    if (e->state == AT_ST_WAITING && (int32_t)(now_ms - e->deadline_ms) >= 0)
        at_on_timeout(e, now_ms);
    if (e->state == AT_ST_IDLE)
        at_start_next(e, now_ms);
}
\end{ccode}
The deadline comparison is written as a signed difference on purpose, so that it
stays correct when the millisecond counter wraps. A comparison written the
obvious way is correct for forty-nine days and then is not.

\item \textbf{Classify every line before the state machine sees it.} Three
classes and nothing else: a final result code ends the command in flight, an
unsolicited code is dispatched from the table, and anything else is an
intermediate line belonging to the command in flight.
\begin{ccode}
/* Final result codes, per the standard command set. The module's own manual
 * adds module-specific ones and they go in this table, not in the code.    */
static const char *const FINAL_OK[]  = { "OK", "CONNECT", NULL };
static const char *const FINAL_BAD[] = { "ERROR", "+CME ERROR:", "+CMS ERROR:",
                                         "NO CARRIER", "ABORTED", NULL };

at_class_t at_classify(const char *line)
{
    if (match_any(line, FINAL_OK))  return AT_CLASS_FINAL_OK;
    if (match_any(line, FINAL_BAD)) return AT_CLASS_FINAL_BAD;
    if (at_urc_lookup(line))        return AT_CLASS_URC;
    return AT_CLASS_INTERMEDIATE;
}
\end{ccode}
An empty line is not a class. Modules pad replies with blank lines and a
classifier that treats one as an intermediate result will attach it to whatever
command is running. Skip them in the line assembler, before classification.

\item \textbf{Give the unsolicited codes a table and a handler each.} Registration
status, signal quality and incoming payload notifications all arrive without
being asked for, and some of them arrive in the middle of a command. The table
is data, the handlers are code, and neither touches the state machine.
\begin{ccode}
static const at_urc_t URC_TABLE[] = {
    { "+CEREG:",  on_registration },   /* network registration changed      */
    { "+CSQ:",    on_signal       },   /* only unsolicited if enabled       */
    { "RDY",      on_module_ready },   /* the module restarted underneath us */
    { NULL,       NULL }
};
\end{ccode}
The module-restart notification deserves its own handler even if it does nothing
but set a flag, because everything the engine believes about the module's state
is wrong after one, and the flag is what tells the application to reconfigure
rather than to carry on.

\item \textbf{Take every timeout from the module's manual and write down where.}
Some commands answer in milliseconds and some take tens of seconds. Network
registration is the long one, and the published empirical study of this exact
module is the source for what to expect in practice, with the caveat that its
host is a single-board computer and not this board.
\begin{ccode}
/* docs/timeouts.md carries the manual page for every one of these.          */
static const at_cmd_t CMD_ECHO_OFF = { "ATE0",     300,  1, 0, on_done, NULL };
static const at_cmd_t CMD_VERBOSE  = { "AT+CMEE=2", 300, 1, 0, on_done, NULL };
static const at_cmd_t CMD_SIM      = { "AT+CPIN?",  5000, 0, 0, on_sim,  NULL };
static const at_cmd_t CMD_SIGNAL   = { "AT+CSQ",    2000, 2, 0, on_csq,  NULL };
static const at_cmd_t CMD_REGSTAT  = { "AT+CEREG?", 5000, 2, 0, on_reg,  NULL };
\end{ccode}
The data-plane commands are module-specific and are not reproduced here from
memory. They come from the module's own AT command manual, which is the
normative reference for this chapter, and they go in one file so that porting to
the SIM7070G touches one file and not the engine.

\item \textbf{Back off, with a cap and with jitter.} Retrying immediately after a
failure that was caused by the network being busy is the one strategy guaranteed
to make things worse. Double the interval, cap it, and scatter it.
\begin{ccode}
uint32_t at_backoff_ms(uint8_t attempt, uint32_t base_ms, uint32_t cap_ms)
{
    uint32_t d = base_ms;
    for (uint8_t i = 0; i < attempt && d < cap_ms; i++) d <<= 1;
    if (d > cap_ms) d = cap_ms;
    return d - (d >> 3) + (at_rand_small() % (d >> 2 ? d >> 2 : 1));
}
\end{ccode}
The jitter matters more than it looks. Without it a fleet of nodes that all lose
the network at the same time all return at the same time.

\item \textbf{Inject a hundred faults and require a defined state from every
one.} The injector sits on the host, between the module and the board, or
replaces the module entirely with a recorded transcript. Four classes, and they
are the four that actually occur.
\begin{pycode}
FAULTS = [
    ("no_reply",      lambda c: b""),                     # module says nothing
    ("truncated",     lambda c: c.reply[: len(c.reply) // 2]),
    ("urc_midway",    lambda c: c.reply[:3] + b"\r\n+CEREG: 0\r\n" + c.reply[3:]),
    ("late_reply",    lambda c: (c.delay(c.timeout_ms + 500), c.reply)[1]),
]
# 25 of each, across the command set, with the engine's state logged after each.
\end{pycode}
The acceptance rule is not that the engine succeeds. It is that every one of the
hundred ends in one of the named results, that the completion callback fires
exactly once in each, and that none of them leaves the engine in a state from
which the next command cannot start.

\item \textbf{Run it with no SIM card and call the stop a pass.} With no card
fitted, the card status query returns an error and the engine has nowhere to go.
The correct behaviour is to stop in a named state, report it once and not retry
forever, because no amount of backoff produces a SIM card.
\begin{shellcode}
python test/faults.py --port /dev/ttyUSB0 --count 100 --report faults.json
python test/faults.py --scenario no_sim --expect-state AT_ST_STOPPED_NO_SIM
\end{shellcode}
\end{steps}

\subsection*{Build, flash and debug}

\diagram{c11_timing}{Two exchanges on the serial link. The first carries an unsolicited message arriving in the middle of a command, which a blocking driver would read as the reply. The second shows a reply that arrives after its deadline and is discarded rather than being attached to the next command.}

\begin{shellcode}
cmake -B build-fw -G Ninja -DCMAKE_TOOLCHAIN_FILE=cmake/arm-none-eabi.cmake
cmake --build build-fw -j
cp build-fw/firmware.bin "$PROBE_DISK"/   # onto the probe disk
\end{shellcode}

Debugging this chapter is mostly reading the link rather than the code. Keep a
transcript of every byte in both directions, with a timestamp, and write it to
the console on a second port so that reading it does not change the timing of
the thing being read. A transcript with timestamps answers almost every question
this chapter raises, and a transcript without them answers none of the
interesting ones.

\begin{note}{When the module answers the host but not the board}
In order of likelihood: transmit and receive are not crossed, which is the most
common wiring fault on this bench and produces exactly this symptom; the grounds
are not common, so the signalling has no reference and reads as noise; the
module is at a different baud rate because a previous experiment changed it; the
line ending the module expects is not the one being sent; echo is on and the
parser is treating the echoed command as the reply; or the module is resetting
on every transmit because it is being fed from the board rather than from its
own supply. Check them in that order and check the supply first if the symptom
comes and goes.
\end{note}

To recover from a module that has been left in an unknown configuration, use the
power control line rather than the command set. A module that will not answer at
any rate is easier to restart than to interrogate, and the restart notification
it emits when it comes back is already in the unsolicited table.

\subsection*{Verification and acceptance criteria}

\begin{itemize}
\item One hundred injected faults, twenty-five of each of the four classes,
every one ending in a defined result, with the completion callback fired exactly
once per submitted command and none of them leaving the engine unable to start
the next command.
\item No run reaches a state from which no transition is possible other than the
two named stopped states, and those two are reported as outcomes rather than
counted as failures.
\item With no SIM card fitted the engine stops in the named no-card state,
reports once, and does not retry, which is recorded as a pass.
\item The worst-case time spent inside the poll function, measured with the
cycle counter, stays below the budget with the longest reply the module can
produce present in the ring.
\item The longest interval with interrupts disabled stays below the budget,
measured the same way, because an engine that never blocks the main loop but
disables interrupts for a millisecond has moved the problem rather than solved
it.
\item Every timeout in the firmware appears in the timeouts document with either
the manual page it came from or the bench measurement and instrument that
produced it.
\item A late reply arriving after its deadline is discarded and is not attached
to the following command, proven by a transcript from the fault injector rather
than by inspection of the code.
\item The engine compiles and its tests run on the host with no hardware
present, which is what chapter 12 needs from it.
\end{itemize}

\subsection*{Variants}

\begin{table}[H]\centering\small
\begin{tabular}{p{24mm}p{26mm}p{44mm}p{22mm}p{20mm}}
\toprule
Axis & Variant & What changes & Cost & Built in full in \\
\midrule
Intelligence and reach & Edge host over a framed link & The board speaks a framed protocol to a Raspberry Pi and the Pi owns the modem, which removes the command set from the firmware entirely and is the right answer whenever a host is present & A host must exist and stay powered & Here and chapter 9 \\
Operating system & The maintained library on a kernel & With an operating system underneath, the MIT cellular library named above becomes the correct choice and this engine becomes the thing that explains it & A kernel, and its memory & Chapter 20 \\
Peripheral substitution & The SIM7070G instead & The same engine, a different command file, and a module the maintained library does support, which makes it the control case for this chapter & Different network and different manual & Chapter 17 \\
Peripheral substitution & The SIM7600E-H instead & A module that does offer a point-to-point protocol, so a network stack on the board becomes possible and the command set stops being the only route & A network stack, and far more code & Chapter 17 \\
Execution model & DMA receive with idle line & The byte path this chapter stands on & Setup complexity & Chapter 4 \\
Synchronisation & A queue between application and engine & With a kernel, the submit call becomes a queue send and the poll loop becomes a task & Kernel overhead & Chapter 20 \\
Language & Host tooling in Python & The fault injector here & None & Chapter 3 \\
\bottomrule
\end{tabular}
\caption{Variants for chapter 11. The edge-host row is the one that most often wins in practice, and this chapter builds it in full alongside chapter 9 precisely so that the comparison is honest rather than rhetorical.}
\end{table}

\subsection*{Pitfalls}

\begin{itemize}
\item Parsing inside the receive interrupt. It puts unbounded work at high
priority, and every timing number the chapter produces afterwards is about the
parser rather than about the system.
\item Treating an unsolicited message as the reply to the command in flight.
This is the failure that a bench with a good signal never reproduces and that a
field deployment reproduces on the first day.
\item Attaching a late reply to the next command. Once this happens the engine is
one command out of step and stays that way, and every subsequent result is
wrong in a way that looks like a module fault.
\item Comparing deadlines with a plain unsigned comparison. It is correct until
the millisecond counter wraps, which is far enough away to be missed and close
enough to happen.
\item Inventing timeouts. The module's manual states a maximum response time per
command and that is the number. A guessed timeout that is too long turns a
failure into a stopped system.
\item Retrying without backoff and without jitter, which converts one node's bad
minute into a fleet's bad hour.
\item Feeding the module from the board. The 2 A peak is real, and the symptom is
a module that resets whenever it tries to register, which looks like a firmware
fault and is not one.
\item Leaving echo on and writing the parser around it. Turn it off in the
second command the engine sends and the classifier gets simpler.
\item Allocating command records on a heap. An engine whose failure mode is heap
exhaustion has replaced one unbounded behaviour with another.
\end{itemize}

\subsection*{Best practices applied}

\begin{itemize}
\item The prior art is named precisely, including the one library that would
have done the job and the exact reason it cannot be used here. A chapter that
reimplements a maintained library without saying why has failed.
\item Every allocation is static and every buffer has a stated size and a stated
behaviour when it is full.
\item The completion callback fires exactly once per command in every outcome,
which is the single rule that keeps the caller's state machine honest.
\item Every timeout is traceable to a document or to a measurement with its
instrument named.
\item Failure is tested by injection rather than by waiting for it, and the
acceptance criterion is a defined state rather than a success.
\item A stop is allowed to be a pass. Engineering that cannot distinguish "this
cannot work and here is why" from "this is broken" produces systems that retry
forever.
\end{itemize}

\subsection*{Stretch goals}

\begin{itemize}
\item Record real transcripts from both modules on the bench and replay them
into the host build of the engine, so the test suite runs with no hardware and
still exercises real module behaviour.
\item Add a watchdog on the module rather than on the processor: if no line of
any kind has arrived for longer than the longest expected quiet period, restart
the module through its power control line and report it.
\item Price the engine in energy using chapter 10's ledger. The send phase stops
being a stub, and the registration interval becomes the most expensive decision
in the node.
\item Port the command file to the SIM7070G and run the same fault suite against
the maintained library on a kernel, which is a comparison nobody has published
for these modules.
\item Add a second queue at a higher priority for commands that must overtake a
long-running one, and measure what that does to the worst case through the poll
function.
\end{itemize}

\subsection*{Roadmap and next steps}

Chapter 12 takes this engine's host build and makes it a unit test suite that
runs on every commit, and then takes the board itself and makes it a
hardware-in-the-loop job. The fault injector written here is the most valuable
thing this chapter hands forward, because a test that can produce a late reply
on demand is worth more than any amount of careful reading.

For the wider path, the community embedded engineering roadmap places command
set work and connectivity where this chapter sits, between the peripheral
drivers below it and the system design above it, and the accompanying essay on
closing the gap between reading and building makes the case for exactly this
kind of exercise. The free course from a different silicon vendor on its own
connectivity software is the best structured material on the layer above this
one, and its concepts transfer even though its hardware does not. Appendix H
lists the courses. There is no course and no textbook covering this module on a
bare-metal microcontroller, which is why the chapter states its method in enough
detail to be repeated.

\subsection*{Portfolio evidence}

\begin{itemize}
\item A public repository with the engine, the module command file kept
separate from it, the host test build, the fault injector and the recorded
transcripts.
\item The fault report: one hundred injected faults, the class of each, the
result each produced, and the engine state after each, as a table rather than a
narrative.
\item The timeouts document, every entry traceable to a manual page or to a
bench measurement with its instrument named.
\item A transcript showing an unsolicited message arriving in the middle of a
command and the command still completing correctly, which is the clearest one
page summary of what the chapter is for.
\end{itemize}

\subsection*{Sources}

Normative references:
\begin{itemize}
\item The SIM7020 series AT command manual, for the command set, the final
result codes, the unsolicited result codes and the maximum response time of
every command used here. This is the normative reference for the chapter.
\item The SIM7020 series hardware design document, for the supply requirement,
the peak current and the power control line timing.
\item The 3GPP command set specification for terminal equipment, for the
standard commands and result codes that are not module specific.
\item Reference manual RM0455, for the USART, its idle-line detection and the
transfer engine that serves it. Not RM0433, which documents its sibling.
\item The board user manual for MB1363, for which USART instance reaches the
Arduino header pins.
\end{itemize}

Reusable implementations:
\begin{itemize}
\item Majerle's cellular library, MIT, which supports the SIM7070G on this
bench, bundles a messaging client, and requires an operating system.\newline
\url{https://github.com/MaJerle/lwcell}
\item Majerle's host-side parser for the radio coprocessor's command set, MIT,
whose queue and unsolicited dispatch structure this chapter follows.\newline
\url{https://github.com/MaJerle/lwesp}
\item Majerle's transfer-driven USART examples, MIT, which are the byte path
chapter 4 builds on and this chapter stands on.\newline
\url{https://github.com/MaJerle/stm32-usart-uart-dma-rx-tx}
\item The peer-reviewed empirical study of this exact module, for realistic
signal and registration expectations, noting its host is a single-board
computer.\newline
\url{https://jtde.telsoc.org/index.php/jtde/article/view/955}
\item A transport-agnostic messaging client, MIT, for the layer above this one
once a byte stream exists.\newline
\url{https://github.com/FreeRTOS/coreMQTT}
\end{itemize}
   % An AT engine that never blocks
\project{12}{Energy as a regression test, and the rig that runs it}{NUCLEO-H7A3ZI-Q with a Raspberry Pi and the nRF-PPK2}{Host unit tests, size gate, hardware in the loop, watchdogs}

\begin{keyfacts}
\item[Board] NUCLEO-H7A3ZI-Q under test, a Raspberry Pi as the self-hosted runner, the nRF-PPK2 as the instrument
\item[Peripherals] The independent and the window watchdog, the debug freeze bits, and whichever retained memory can hold a crash record
\item[Toolchain] arm-none-eabi-gcc with CMake for the target; a C test framework with a faking header on the host; a test runner on the Pi
\item[Operating system] Bare metal on the target; Linux on the runner
\item[Difficulty] 4 of 5
\item[Effort] 4 evenings of about four hours
\item[Deliverable] A pipeline in which a deliberate five percent charge regression turns the build red with nobody at the bench, and a two-layer watchdog that leaves a readable record when the firmware stops
\end{keyfacts}

\subsection*{Why this project}

Chapter 10 produces a per-phase charge ledger. It is a good measurement and it
has the same problem every good manual measurement has, which is that it is
taken once, by the person who cared, on the day they cared. Six commits later
somebody adds a retry to the sense path and the charge per cycle grows by eight
percent, and nobody finds out until the field trial. A measurement that is not
checked automatically is a demonstration. This chapter turns the ledger into a
gate.

There are three gates and they are deliberately in increasing order of cost.
Host unit tests run in seconds on a cloud runner with no hardware at all, and
they catch the majority of ordinary defects in the codec of chapter 9, the
engine of chapter 11 and the ledger arithmetic of chapter 10. A size gate reads
the map file and fails a build whose flash or static memory has grown past its
stated limit, which costs nothing and catches the class of change that quietly
consumes the margin. Only then does anything touch hardware: a self-hosted
runner on a Raspberry Pi downloads the artifact the cloud built, flashes the
board, runs the measurement and compares the ledger against the committed
baseline. Every published source on this subject agrees on that architecture,
and the reason is simple: cloud runners are cheap and plentiful and have no
board attached, and the machine with the board attached should do as little as
possible.

The second half of the chapter is watchdogs, and they belong here rather than in
a chapter of their own because a rig that resets boards unattended is exactly
the thing that needs to tell a stopped firmware from a slow one, and to leave
evidence either way. A watchdog that simply resets the part produces a rig that
reports intermittent failures and no information. A watchdog designed properly
produces a rig that reports which task stopped responding and when, which is a
different quality of result.

\begin{note}{A test nobody runs is documentation with a worse licence}
The acceptance criterion for this chapter is stated in terms of absence: the
regression is caught with nobody at the bench. Every design decision below
follows from that. If a step needs a person to plug something in, to press
reset, or to decide whether a number looks reasonable, it is not part of the
gate and it has to be replaced by something that is.
\end{note}

\subsection*{Prior art and what to reuse}

\begin{table}[H]\centering\small
\begin{tabular}{p{32mm}p{46mm}p{40mm}p{20mm}}
\toprule
Source & What it gives & What it does not & Licence \\
\midrule
Unity & A C test framework small enough to run on the target as well as the host & No mocking and no build system, so something has to supply both & MIT \\
Ceedling & Project generation, test runners and mock generation on top of Unity & It needs Ruby. On a Windows host with a Linux subsystem that is a real friction point, worth saying before a reader spends an evening on it & MIT \\
CppUTest & Memory-leak detection between setup and teardown, the one capability the others do not have & It is C++ underneath, so the build is heavier and the messages are further from the C you wrote & BSD \\
fff & Fake functions in a single header, the lightest fit for a C project & Its default call history silently drops calls past ten arguments and past seventeen calls, so a test asserting on the twentieth call passes for the wrong reason & MIT \\
Memfault's framework comparison & A direct, current comparison written by people who ship firmware & It compares frameworks, not rigs & Article \\
Golioth's series on automated hardware testing & The rig: a self-hosted runner, fixtures that power and flash a board, hardware as a test resource & Their boards and their cloud service, so the fixtures transfer and the specifics do not & Article \\
Memfault's watchdog article & The two-layer design, plus a task-liveness bitmask & No code for this part and no per-task timeout policy & Article \\
Ganssle's essay on watchdogs & The clearest statement of why a naive watchdog gives false confidence & It predates this silicon and does not discuss debug freeze & Article \\
The other operating system's task watchdog & The only shipping reference implementation of per-task timeouts & The common kernel ships no task-watchdog abstraction at all & Apache-2.0 \\
\bottomrule
\end{tabular}
\caption{Prior art for chapter 12. The frameworks and the rig are prior art; the architecture that joins them and the energy gate on top are not.}
\end{table}

What is left to write is the energy gate itself, the fixtures that make the
instrument a test resource rather than an instrument, and the watchdog layer.
Two observations from the sources deserve to be carried forward rather than
rediscovered. The first is the fff call-history limit: it is a default, it is
documented, and it is the sort of default that produces a green test for the
wrong reason, so the project either raises the limit explicitly in its
configuration header or asserts on call counts in a way that does not depend on
history depth. The second is that nobody has published a comparison between the
per-task watchdog abstraction of the other operating system and the absence of
one in the common kernel. That comparison is a small piece of original work and
it fits inside this chapter.

\subsection*{Parts from the inventory}

\begin{table}[H]\centering\small
\begin{tabular}{p{40mm}p{66mm}p{40mm}}
\toprule
Part & Role & Interface \\
\midrule
NUCLEO-H7A3ZI-Q & The board under test. Its own probe is the programmer, so no separate debug hardware is needed on the rig & USB to the Raspberry Pi \\
Raspberry Pi 4 & The self-hosted runner. It downloads the artifact, flashes, runs the tests and reports & Ethernet or wireless to the cloud; USB to the board and to the instrument \\
nRF-PPK2 & The instrument, driven from the runner rather than by hand & USB to the Pi; measurement leads to the board \\
Powered USB hub & The board, the instrument and the Pi all draw from one supply otherwise, and a runner that resets itself under load is not a runner & USB \\
Jumper wires & The marker pins of chapter 10, so the energy gate can report per phase rather than per cycle & Female headers both ends \\
Host PC & Writing the tests. It is not part of the rig, which is the point & USB \\
\bottomrule
\end{tabular}
\caption{Inventory items used in chapter 12. Nothing is bought and nothing is soldered, so the rig has no relay and cannot cut power to the board. That constraint shapes the recovery design below.}
\end{table}

\subsection*{System architecture}

\diagram{c12_arch}{Three gates in increasing order of cost. The cloud builds and publishes an artifact; the machine with the board attached downloads it and does as little as possible. Every published source on this subject arrives at the same shape.}

The split is not about saving money on compute. It is about failure modes. A
cloud runner that fails tells you the build is broken. A self-hosted runner that
fails might mean the firmware is broken, or the board has come unplugged, or the
instrument has enumerated under a different device name, or somebody has taken
the shield off. Keeping those two kinds of failure on separate machines means
the first kind is never confused with the second, and it means the expensive,
fragile half of the pipeline runs only against a binary that has already passed
everything cheap.

\subsection*{Peripheral configuration}

\begin{table}[H]\centering\small
\begin{tabular}{p{22mm}p{26mm}p{24mm}p{30mm}p{30mm}}
\toprule
Peripheral & Mode & Clock source & Pins and function & Interrupt and transfers \\
\midrule
Independent watchdog & Free running, reset on expiry & An internal low-speed oscillator, which is what makes it independent of the rest of the clock tree & None & Reset only. No interrupt \\
Window watchdog & Windowed, early refresh also resets & Peripheral bus, so it stops when that stops & None & Early wake-up interrupt, used to write the record \\
Debug freeze control & Freeze both watchdogs while halted & Debug clock. Confirm whether that clock must be enabled before the freeze write takes effect & None & None \\
Retained memory & Crash record across a reset & Backup domain. Whether it is usable without a coin cell is on the confirm list & None & None \\
GPIO markers & As chapter 10 & Peripheral bus & Free Zio pins & None \\
USART3 console & Test transcript to the runner & Peripheral bus & Confirm in the board manual & Interrupt driven \\
\bottomrule
\end{tabular}
\caption{Peripheral configuration for chapter 12. The two watchdogs are complementary rather than alternative: one is independent of the clock tree and cannot be talked out of a reset, and the other notices a refresh that arrives too early as well as one that arrives too late.}
\end{table}

The difference between the two watchdogs is worth stating plainly because it
decides which one carries which job. The independent watchdog runs from its own
oscillator, so it survives almost anything that goes wrong with the main clock
tree, and that independence is exactly why it cannot tell you anything before it
resets the part. The window watchdog runs from the peripheral bus and refuses a
refresh that arrives too early as well as one that arrives too late, which
catches a loop running away and refreshing in a tight cycle, and it offers an
interrupt shortly before it expires. That interrupt is where the crash record
gets written.

\subsection*{Wiring}

\diagram{c12_wiring}{The rig as it actually sits on the bench. There is no relay and no soldering iron here, so power cannot be cut under program control, and the recovery path has to work with what the probe and the reset line already provide.}

\begin{note}{The rig cannot cut power and that changes the recovery design}
With no relay in the inventory, a board that has stopped in a way a reset does
not clear will stay stopped until somebody unplugs it. The design answer is to
make that case impossible rather than recoverable: the independent watchdog is
configured in the very first firmware the rig ever flashes, the runner uses the
connect-under-reset path rather than assuming the target is responsive, and the
job reports a distinct outcome for "the board did not answer" so it is never
confused with a test failure.
\end{note}

\subsection*{Memory and timing budget}

\diagram{c12_mem}{The size gate. Both limits are stated in the repository and read by a script from the map file, so growth is refused at the point it happens rather than noticed when something no longer fits.}

\begin{table}[H]\centering\small
\begin{tabular}{p{44mm}p{28mm}p{28mm}p{30mm}}
\toprule
Quantity & Budget & Measured & Margin \\
\midrule
Flash limit for the node image & 256 kB & not measured & not measured \\
Static memory limit & 64 kB & not measured & not measured \\
Crash record size & 128 B & not measured & not measured \\
Host unit test suite, wall time & below 60 s & not measured & not measured \\
Size gate, wall time & below 2 s & not measured & not measured \\
Hardware job, wall time & below 10 min & not measured & not measured \\
Charge per cycle, tolerance & 5 percent of baseline & not measured & not measured \\
Independent watchdog period & 4 s & not measured & not measured \\
Software watchdog lead time & 250 ms before the hardware one & not measured & not measured \\
\bottomrule
\end{tabular}
\caption{The budget for chapter 12. The limits are deliberately far below what the part provides, because a gate set at the hardware limit fires only when it is already too late to be useful. Nothing enters the measured column until the bench produces it.}
\end{table}

\subsection*{Firmware design (UML)}

\diagram{c12_uml}{One job, end to end. The artifact crosses from the cloud to the bench exactly once, and the Raspberry Pi never compiles anything.}

The sequence has one property worth defending: the Raspberry Pi does not build.
If it did, the thing being tested would be a build produced by a machine nobody
else has, with whatever toolchain version happened to be installed on it, and a
green result would mean less than it appears to. Downloading the artifact makes
the tested binary and the published binary the same bytes.

The watchdog design that runs inside that binary has two layers. A software
watchdog, driven from a timer, expires slightly before the hardware one and its
handler writes a record: which activities had checked in, which had not, and the
program counter it was interrupted at. Then it stops refreshing, and the
hardware watchdog resets the part a short time later. The record survives only
if it is written somewhere the reset does not clear, and whether the backup
domain on this board is usable without a coin cell is on the confirm list, so
the fallback is a region of ordinary static memory that the startup code checks
before it zeroes anything.

\subsection*{Data flow (ASCII)}

\begin{asciiart}
  commit                cloud runner                    artifact store        self-hosted Pi
  +--------+   push   +------------------------+      +---------------+     +------------------+
  | git    |--------->| 1. build firmware.elf  |      |               |     | 4. download      |
  +--------+          | 2. host unit tests     |      | firmware.elf  |     | 5. flash target  |
                      |    codec, engine, ledger|---->| firmware.map  |---->| 6. run pytest    |
                      | 3. size gate from .map |      | sizes.json    |     | 7. drive PPK2    |
                      +------------------------+      +---------------+     | 8. ledger.json   |
                          |  red here means                                 | 9. compare       |
                          |  the code is wrong                              +------------------+
                          v                                                        |
                      +------------------------+                                   v
                      | no board involved      |                       red here means the code
                      | seconds, not minutes   |                       is wrong OR the bench is
                      +------------------------+                       and the report says which
\end{asciiart}

\subsection*{Repository layout}

\begin{plaincode}
nucleo-h7a3-node/
  CMakeLists.txt
  src/                           # the node firmware of chapters 8 to 11
  src/wdt_hw.c                   # IWDG and WWDG setup, debug freeze
  src/wdt_soft.c                 # the task-liveness bitmask and the early handler
  src/crash_record.c             # what survives a reset, and where it lives
  test/host/CMakeLists.txt       # the host build: no target toolchain
  test/host/test_codec.c         # chapter 9
  test/host/test_at_engine.c     # chapter 11, with the recorded transcripts
  test/host/test_ledger.c        # chapter 10 arithmetic
  test/host/fff_config.h         # the call-history depth raised, deliberately
  tools/size_gate.py             # reads the map file, fails on growth
  rig/conftest.py                # fixtures: board, serial, instrument
  rig/test_energy.py             # the gate: ledger against baseline
  rig/test_watchdog.py           # deliberate stop, record recovered
  rig/runner_node/               # the same runner in Node.js, built in full here
  rig/runner_java/               # and in Java, for the same reason
  baselines/ledger_baseline.json # committed, and changed only on purpose
  .github/workflows/build.yml    # cloud: build, test, size gate, publish
  .github/workflows/bench.yml    # self-hosted: download, flash, measure
  README.md
\end{plaincode}

\subsection*{Steps}

\begin{steps}
\item \textbf{Split the firmware so most of it builds on the host.} This is the
step that decides whether the rest of the chapter is easy or impossible. Every
module that contains logic worth testing gets a hardware-free interface, and the
register access lives behind it. The codec of chapter 9 and the ledger
arithmetic of chapter 10 are already pure. The engine of chapter 11 needs one
seam, which is the function that hands bytes to the transport.
\begin{ccode}
/* The seam. On the target this writes to the transfer engine; on the host the
 * test provides it and records what was written.                           */
void at_port_write(const uint8_t *buf, size_t len);
uint32_t at_port_now_ms(void);
\end{ccode}
Two functions is the right size for a seam. A seam with twenty functions is a
second driver that also has to be maintained.

\item \textbf{Choose the framework and record the reasons.} Unity for the
assertions, a faking header for the seams, and a decision about the generator.
\begin{shellcode}
git submodule add https://github.com/ThrowTheSwitch/Unity third_party/unity
git submodule add https://github.com/meekrosoft/fff  third_party/fff
# Ceedling generates all of this and more, and needs Ruby. On a Windows host
# with a Linux subsystem that is an extra runtime to install and keep working.
\end{shellcode}
The generator is worth it on a project with many modules and is not worth it on
a project with four. State which case you are in rather than following a
default.

\item \textbf{Configure the faking header before writing the first fake.} Its
default history depth is the trap most worth documenting in this chapter.
\begin{ccode}
/* test/host/fff_config.h
 * The default history records ten arguments and seventeen calls and then
 * silently drops the rest. A test that asserts on the twentieth call passes
 * because the history is empty, not because the behaviour is right.        */
#define FFF_ARG_HISTORY_LEN   32
#define FFF_CALL_HISTORY_LEN  256
#include "fff.h"
\end{ccode}
If a suite needs leak detection between setup and teardown as well, that is the
one thing the C++ framework in the table does and the others do not, and it is a
reason to run a second small suite rather than to move everything.

\item \textbf{Write the host tests that would have caught real defects.} Not
coverage for its own sake: three suites, each aimed at a failure that has
actually happened. Round-trip property tests for the codec, the four fault
classes for the engine replayed from recorded transcripts, and the
reconciliation arithmetic for the ledger with a hand-computed example.
\begin{shellcode}
cmake -S test/host -B build-host && cmake --build build-host
ctest --test-dir build-host --output-on-failure
\end{shellcode}

\item \textbf{Add the size gate and make it read the map file.} The size output
alone is not enough, because it reports totals and the interesting question is
which section grew.
\begin{pycode}
LIMITS = {"flash": 256 * 1024, "sram": 64 * 1024}

def gate(map_path, elf_path):
    used = read_sections(elf_path)              # text + rodata + data, and bss
    over = {k: (v, LIMITS[k]) for k, v in used.items() if v > LIMITS[k]}
    for k, (v, lim) in sorted(over.items()):
        print(f"FAIL {k}: {v} bytes, limit {lim}, over by {v - lim}")
    print_biggest_symbols(map_path, n=15)       # always, not only on failure
    return 1 if over else 0
\end{pycode}
Printing the fifteen largest symbols on every run, not only on failures, is what
makes the gate useful rather than annoying. The reader who sees the list on a
passing build learns where the flash went.

\item \textbf{Build the two watchdog layers, and configure the hardware one
first.} This is the variant this chapter owns, so both watchdogs are built here
in full rather than named.
\begin{ccode}
/* Layer one: the independent watchdog. Its own oscillator, so it survives a
 * clock tree that has gone wrong. Period and prescaler limits come from
 * RM0455 for this part, not from a sibling.                                */
void wdt_hw_init(uint32_t period_ms)
{
    IWDG->KR  = 0x5555u;                 /* unlock the prescaler and reload  */
    IWDG->PR  = IWDG_PRESCALER_FOR(period_ms);
    IWDG->RLR = IWDG_RELOAD_FOR(period_ms);
    IWDG->KR  = 0xCCCCu;                 /* start; it cannot be stopped now  */
    while (IWDG->SR != 0u) { }           /* wait for the registers to settle */
}

/* Layer two: the window watchdog, refreshed only from the liveness check, and
 * with its early interrupt enabled so the record can be written.           */
void wdt_window_init(void)
{
    WWDG->CFR |= WWDG_CFR_EWI;           /* early wake-up interrupt          */
    WWDG->CR   = WWDG_CR_WDGA | WWDG_COUNTER_START;
}
\end{ccode}
The independent watchdog cannot be stopped once started, which is the property
that makes it trustworthy and the property that makes the debug freeze bits
necessary.

\item \textbf{Freeze the watchdogs while halted, and check that the freeze took
effect.} Without this, every breakpoint resets the part, and the resulting
half-hour of confusion is a rite of passage nobody needs.
\begin{ccode}
/* The documented trap: on this family the debug block has its own clock, and
 * a write to the freeze register before that clock is enabled is accepted
 * and does nothing. Enable, write, then read back and assert.              */
void wdt_debug_freeze(void)
{
    RCC_ENABLE_DEBUG_CLOCK();                    /* first, always            */
    DBGMCU->APB1FZR1 |= DBGMCU_IWDG_STOP | DBGMCU_WWDG_STOP;
    if ((DBGMCU->APB1FZR1 & DBGMCU_IWDG_STOP) == 0u)
        panic("debug freeze write did not take effect");
}
\end{ccode}
Confirm the register name, the bit positions and the clock-enable requirement in
RM0455 for this part before trusting any of it. The requirement that the debug
clock be enabled first is documented and is a silent failure when it is missed,
which is exactly why the read-back is in the code rather than in a comment.

\item \textbf{Give every periodic activity a liveness bit and one honest
question.} The software layer keeps a bitmask. Each activity sets its bit when
it completes a cycle, and the checker refreshes the hardware watchdog only when
every registered bit is set, then clears them all.
\begin{ccode}
static volatile uint32_t live_mask;        /* set by the activities          */
static uint32_t          live_expected;    /* set at registration            */

void wdt_task_alive(uint8_t id) { live_mask |= (1u << id); }

void wdt_check(void)                        /* from the timer, not from main  */
{
    if ((live_mask & live_expected) == live_expected) {
        live_mask = 0u;
        IWDG->KR = 0xAAAAu;                 /* refresh, and only here        */
        WWDG->CR = WWDG_COUNTER_START;
    }
    /* otherwise: do nothing. The hardware layer is the one that acts.       */
}
\end{ccode}
The question this design does not answer, and which nothing in the published
material answers well, is what timeout to give an activity that is correctly
blocked for a long time. A sensor waiting on a watermark interrupt at a low rate
is not stopped, and a network attach that legitimately takes tens of seconds is
not stopped either. The honest options are a per-activity deadline rather than
one global period, which is what the other operating system's task watchdog
provides and the common kernel does not provide at all, or an explicit
suspend-and-resume around known long waits. Chapter 20 builds the kernel side;
this chapter records that the comparison between the two has not been published.

\item \textbf{Make the instrument a test fixture.} The rig's test code should
never reason about USB device names. Put that in a fixture, make the fixture
fail loudly when the hardware is absent, and give that failure its own outcome
so the report distinguishes it from a test failure.
\begin{pycode}
@pytest.fixture(scope="session")
def board(request):
    port = request.config.getoption("--port")
    if not flash(elf="artifacts/firmware.elf", port=port, connect_under_reset=True):
        pytest.exit("BENCH: target did not answer", returncode=3)
    return Board(port)

@pytest.fixture(scope="session")
def meter():
    ppk = PPK2Instrument.open()          # GPL-2.0 library, installed not vendored
    ppk.ammeter_mode()
    yield ppk
    ppk.close()
\end{pycode}
Exiting with a distinct code when the bench is at fault is the single most
useful thing in the whole rig. It is the difference between a pipeline people
trust and a pipeline people rerun until it goes green.

\item \textbf{Close the loop with the energy gate, and prove it by breaking it on
purpose.} The gate compares the ledger produced by this run against the
committed baseline and fails outside the tolerance.
\begin{pycode}
def test_charge_per_cycle(board, meter):
    led = run_ledger(board, meter, cycles=10)
    base = json.load(open("baselines/ledger_baseline.json"))
    drift = (led["total_c"] - base["total_c"]) / base["total_c"]
    assert led["residual"] < 0.05, f"phases do not reconcile: {led['residual']:.3f}"
    assert abs(drift) < 0.05, f"charge per cycle moved by {drift:+.1%}"
\end{pycode}
Then prove the gate works by adding a commit that wastes charge deliberately, a
spin loop of a known length in the compute phase sized to cost about five
percent, and watching the pipeline turn red with nobody present. A gate that has
never failed is not known to work.
\end{steps}

\subsection*{Build, flash and debug}

\diagram{c12_timing}{The two watchdog layers on one timeline. The software layer expires first and its handler writes the record; the hardware layer resets the part a short time later. The interval between them is the whole reason the record exists.}

\begin{shellcode}
# on the runner, which never compiles:
gh run download --name firmware --dir artifacts
pytest rig/ --port /dev/ttyACM0 --junitxml=report.xml
\end{shellcode}

The runner in the repository is written twice, in Node.js and in Java, because
this book's variant matrix assigns the host-language axis to this chapter. Both
do the same four things: fetch the artifact, invoke the flashing tool, run the
test suite and post the report. The comparison is not about which language is
better. It is about what each one costs to install on a Raspberry Pi and how
each behaves when a subprocess stops responding, which is the only failure the
runner itself has to handle well.

\begin{note}{When the pipeline is red and the code is fine}
In order of likelihood: the board enumerated under a different device name after
a reboot, which is why the fixture takes the port as an option and fails with
its own exit code; the artifact is stale because the download step succeeded
against an older run; the instrument is in the wrong mode, so the ledger is of
the instrument supplying the board rather than of the board drawing; the shield
was left fitted from another chapter and its quiescent current is inside the
measurement; or a previous job left the target halted at a breakpoint and the
flashing step needs the connect-under-reset path. Each has a distinct signature
in the report if the fixtures are written to produce one.
\end{note}

To recover a board the rig cannot reach, the order is: connect under reset,
then mass erase, then reflash the last known good artifact. Nothing here touches
option bytes, and nothing in this book does until the question of whether a bad
option-byte write can leave this board unrecoverable without external tooling
has been answered.

\subsection*{Verification and acceptance criteria}

\begin{itemize}
\item A commit that raises charge per cycle by about five percent turns the
build red with nobody at the bench, and the failure message names the quantity,
the baseline and the drift rather than saying an assertion failed.
\item The per-phase reconciliation from chapter 10 runs as part of the same job,
and a run whose residual exceeds five percent fails as a measurement fault
rather than passing with a plausible-looking ledger.
\item A commit that grows the image past the stated flash or static memory limit
fails in under two seconds on a machine with no board attached, and the report
lists the fifteen largest symbols whether it passed or failed.
\item The faking header's call history is configured explicitly, and a test
exists that would fail if the default depth were restored, so the trap cannot
return quietly.
\item A deliberately stopped activity produces a crash record that survives the
reset and is read back by the rig, naming which liveness bits were missing.
\item The debug freeze configuration is read back and asserted at startup, so a
freeze write that did not take effect reports itself instead of producing
breakpoints that reset the part.
\item A bench fault, meaning an absent board or an absent instrument, exits with
its own code and is reported as a bench fault and never as a test failure.
\item Both host runners, in Node.js and in Java, drive the same job to the same
report, and the installation cost of each on the Raspberry Pi is recorded.
\end{itemize}

\subsection*{Variants}

\begin{table}[H]\centering\small
\begin{tabular}{p{24mm}p{26mm}p{44mm}p{22mm}p{20mm}}
\toprule
Axis & Variant & What changes & Cost & Built in full in \\
\midrule
Time and safety & Independent watchdog against window watchdog & One runs from its own oscillator and cannot be stopped, the other notices an early refresh and offers an interrupt before it expires. The design here uses both, for different jobs & Two peripherals to configure and a freeze question to settle & Here \\
Language & Host runner in Node.js or Java & The same four steps written twice, compared on installation cost on a Raspberry Pi and on behaviour when a subprocess stops responding & Two runtimes to keep working & Here \\
Language & Host tooling in Python & The test suite and the size gate & None & Chapter 3 \\
Synchronisation & Event group and stream buffer & With a kernel, the liveness bitmask becomes an event group and the test transcript becomes a stream buffer, which removes the hand-written mask entirely & A kernel & Chapter 20 \\
Operating system & Per-task watchdog from the other operating system & A per-activity deadline instead of one global period, which is the only shipping answer to the long-blocked activity question & A different operating system & Chapter 20 \\
Execution model & Run the unit tests on the target as well & The assertion framework is small enough to run on the part, which catches the defects that only appear with real register behaviour & Flash, and a slower loop & Chapter 17 \\
\bottomrule
\end{tabular}
\caption{Variants for chapter 12. The two rows marked here are the ones the book's matrix assigns to this chapter, and both are built rather than described.}
\end{table}

\subsection*{Pitfalls}

\begin{itemize}
\item Building the firmware on the machine that tests it. The tested binary and
the published binary then differ, and a green result means less than it looks.
\item Accepting the faking header's default call history. Past ten arguments and
seventeen calls it drops the rest, and the test that asserts on the twentieth
call passes because the history is empty.
\item Writing the debug freeze register before the debug clock is enabled. The
write is accepted and does nothing, which is a documented silent failure, and
the symptom is a part that resets at every breakpoint.
\item Refreshing the watchdog from a timer interrupt with no liveness condition.
The part then stays alive while doing nothing useful, which is the naive design
Ganssle's essay is about.
\item Giving a correctly blocked activity the same deadline as a fast one. A
sensor waiting on a low-rate watermark interrupt is not stopped, and a watchdog
that resets it is worse than no watchdog.
\item Reporting a missing board as a test failure. The pipeline then loses the
trust that makes it useful, and people rerun it until it passes.
\item Committing a new energy baseline to make a red build go green. The
baseline changes in its own commit, with a reason, or the gate means nothing.
\end{itemize}

\subsection*{Best practices applied}

\begin{itemize}
\item The cheap gates run first and on a machine with no hardware, so a failing
build costs seconds and a bench visit is the last resort rather than the first.
\item The tested artifact is the published artifact, byte for byte.
\item Bench faults have their own outcome and are never reported as test
failures.
\item A trap in a dependency is configured out explicitly and a test protects
the configuration, rather than being written in a comment that survives one
refactor.
\item The watchdog is two layers, so a stop produces a record and not only a
reset.
\item A register write whose failure is silent is read back and asserted.
\item The gate is proven by making it fail on purpose. A gate that has never
fired is not known to work, and every limit and tolerance lives in a file whose
changes are reviewed.
\end{itemize}

\subsection*{Stretch goals}

\begin{itemize}
\item Publish the ledger for every commit and plot charge per cycle over the
history of the project. The interesting result is rarely the regression that
fired the gate; it is the slow growth underneath it.
\item Write the comparison nobody has published: the per-task watchdog
abstraction of the other operating system against the absence of one in the
common kernel, with the long-blocked activity as the test case, using the
kernel build from chapter 20.
\item Extend the size gate into a symbol-level differential report, so a pull
request shows which functions grew rather than only the total.
\item Add a second board to the rig and make the suite run on both, which is
where fixture design stops being theoretical.
\item Run the same host suite under the C++ framework for its leak detection
between setup and teardown, and report what it finds that the lighter framework
did not.
\end{itemize}

\subsection*{Roadmap and next steps}

Chapter 13 takes the rig for granted and starts producing the thing it will
measure: a sensor logger that never polls. From here on, every chapter that
produces a number can add it to this pipeline, and the ones that do are the ones
whose claims survive contact with the next six commits.

For the wider path, the community embedded engineering roadmap puts build
systems, testing and continuous integration in the band a firmware engineer
reaches after the peripheral work and before the system design, which is exactly
where this chapter sits, and the accompanying essay on closing the gap between
reading and building argues that a working pipeline teaches more than any amount
of material about pipelines. Appendix H lists the courses. For the testing
material specifically, the comparison article named in the sources is the best
current starting point, and the hardware rig series is the only published
treatment of the self-hosted half that is written by people with boards on their
desks.

\subsection*{Portfolio evidence}

\begin{itemize}
\item A public repository whose pipeline runs on every commit, with a badge that
reflects a gate rather than a compile.
\item The commit that deliberately wastes about five percent of the charge, and
the red build it produced, linked from the README as the proof the gate works.
\item The size report from a passing build, showing the fifteen largest symbols,
which is the artifact reviewers find most immediately readable.
\item A recovered crash record from a deliberately stopped activity, naming the
liveness bits that were missing.
\item A short written comparison of the two host runners, on installation cost
on a Raspberry Pi and on behaviour when a subprocess stops responding.
\end{itemize}

\subsection*{Sources}

Normative references:
\begin{itemize}
\item Reference manual RM0455, for the independent and window watchdog register
maps, the prescaler and reload limits, the debug freeze bits, whether the
independent watchdog survives the low-power modes, and whether the debug clock
must be enabled before a freeze write takes effect. Not RM0433, which documents
its sibling.
\item The board user manual for MB1363, for whether the backup domain is usable
without a coin cell, which decides where the crash record can live.
\item The STM32H7A3xI datasheet, for the internal low-speed oscillator's
frequency tolerance, which is what the watchdog period budget has to absorb.
\end{itemize}

Reusable implementations:
\begin{itemize}
\item Unity, MIT, the C assertion framework.\newline
\url{https://github.com/ThrowTheSwitch/Unity}
\item Ceedling, MIT, which generates the scaffolding and needs Ruby.\newline
\url{https://github.com/ThrowTheSwitch/Ceedling}
\item CppUTest, BSD, for memory-leak detection between setup and teardown.\newline
\url{https://github.com/cpputest/cpputest}
\item fff, MIT, header-only fakes, whose default call history depth is the trap
this chapter configures out.\newline
\url{https://github.com/meekrosoft/fff}
\item Memfault's comparison of the C unit test and mocking frameworks.\newline
\url{https://interrupt.memfault.com/blog/unit-test-mocking}
\item Golioth's series on automated hardware testing with a self-hosted
runner.\newline
\url{https://blog.golioth.io/automated-hardware-testing-using-pytest/}
\item Memfault's article on watchdog best practice, which is the source for the
two-layer design and the task-liveness bitmask.\newline
\url{https://interrupt.memfault.com/blog/firmware-watchdog-best-practices}
\item Ganssle's essay on why naive watchdogs give false confidence.\newline
\url{https://www.ganssle.com/watchdogs.htm}
\item The other operating system's task watchdog, Apache-2.0, the only shipping
reference implementation of per-task timeouts.\newline
\url{https://docs.zephyrproject.org/latest/services/task_wdt/index.html}
\end{itemize}
   % Energy as a regression test, and the rig that runs it

\part{Sensors on the shield}
\project{13}{The IKS4A1: FIFO, watermark, interrupt}{NUCLEO-H7A3ZI-Q with the X-NUCLEO-IKS4A1}{Sensor FIFO, watermark interrupt, the three bus topologies}

\begin{keyfacts}
\item[Board] NUCLEO-H7A3ZI-Q with the X-NUCLEO-IKS4A1 fitted, and no other shield
\item[Peripherals] I2C in fast mode, one external interrupt line for the watermark, a timer for the log timebase, USART3 for the transcript
\item[Toolchain] arm-none-eabi-gcc with CMake, with the register-level driver collection vendored under its permissive licence
\item[Operating system] Bare metal. The kernel version of the same exercise is chapter 20
\item[Difficulty] 4 of 5
\item[Effort] 4 evenings of about four hours
\item[Deliverable] A logger that never polls a sensor, driven entirely by the sensor's own FIFO watermark interrupt, proven by a ten-minute run at a stated rate with zero dropped samples counted on the board
\end{keyfacts}

\subsection*{Why this project}

Reading a sensor in a loop is the default and it is the wrong default for
anything that wants to sleep. At a hundred samples a second, polling means the
processor wakes a hundred times a second, and almost all of that work is bus
transactions that return nothing new. Every sensor on this shield has a
first-in first-out buffer and at least one interrupt line, and using them changes
the question from "how often does the processor wake" to "how deep is the
buffer", which is a question with a much better answer. The chapter's deliverable
is a logger with no polling anywhere in it: the sensor fills its own buffer, tells
the processor when the buffer reaches a level the firmware chose, and the
processor drains it in one burst and goes back to sleep.

The second reason this chapter exists is that the shield is unusually rich and
the richness is in the wiring rather than in the sensors. It carries two inertial
units, each with a programmable engine of a different kind, an accelerometer, a
magnetometer, a pressure sensor, a humidity and temperature sensor, a second
temperature sensor, and a detachable electrode add-on for the electrostatic
sensing channel. What makes it interesting is that those parts can be arranged
on the bus in three different ways, and the arrangement is set by jumpers before
any code runs. Choosing the arrangement is a design decision with real
consequences for the interrupt routing and for what the firmware can batch
together, and it is the part of this exercise that no driver library does for
you.

The prior art is exactly the right amount of incomplete. The register-level
driver collection covers every sensor on this shield, is permissively licensed,
and sits behind a porting shim of two functions. It gives no board layer, no
orchestration of the bus arrangements and no build system. That is the correct
division: the register maps and the buffer tag decoding are tedious and already
done, and the design decisions are left where they belong.

\begin{note}{One shield at a time and never two}
The three shields in this inventory fit the same header and collide on bus
addresses and on the 3.3 V budget. Only one is fitted at a time. This is a
standing rule in the sibling lab set and it is repeated here because the failure
it prevents is subtle: two devices answering at one address produce readings that
are plausible, stable and wrong.
\end{note}

\subsection*{Prior art and what to reuse}

\begin{table}[H]\centering\small
\begin{tabular}{p{32mm}p{46mm}p{40mm}p{20mm}}
\toprule
Source & What it gives & What it does not & Licence \\
\midrule
The register-level driver collection & Every sensor on this shield behind a porting shim of two functions: register maps, the buffer tag decoding, the unit conversions & No board layer, no orchestration of the bus arrangements, no interrupt routing and no build system & BSD-3-Clause \\
The vendor's expansion package & Complete worked examples for this exact shield, and a sensor-fusion middleware & The fusion middleware ships as pre-compiled objects under a restrictive licence and must not be vendored. The examples are tied to the vendor's development environment & Mixed, flagged \\
The upstream shield description in the other operating system's tree & The three bus arrangements set out more clearly than in the vendor's own manual, with the jumper positions named for each & It is a device tree overlay rather than C, so the information transfers and the code does not & Apache-2.0 \\
The shield user manual UM3239 & The normative reference: the sensor complement, the bus addresses and the jumper map & It is a manual. The arrangement it describes least clearly is the one the overlay above describes best & Vendor document \\
Chapters 3 and 10 of this volume & The interrupt discipline that keeps work out of the handler, and the energy ledger that prices batching against polling & Nothing sensor specific & This volume \\
\bottomrule
\end{tabular}
\caption{Prior art for chapter 13. The permissive driver collection is vendored with its copyright headers intact; the fusion middleware is read about and never copied, because pre-compiled objects under a licence tied to one vendor's silicon do not belong in a portfolio repository.}
\end{table}

What is left to write is the board layer, and that is more than it sounds. It
means choosing the bus arrangement and recording why, writing the porting shim
onto this part's I2C peripheral, configuring the buffer depth and watermark for
a stated sample rate, routing the sensor's interrupt to a free pin on this
board, draining the buffer in one burst without holding the bus longer than the
next watermark allows, decoding the tags, and counting drops on the board rather
than inferring them on the host.

\subsection*{Parts from the inventory}

\begin{table}[H]\centering\small
\begin{tabular}{p{40mm}p{66mm}p{40mm}}
\toprule
Part & Role & Interface \\
\midrule
NUCLEO-H7A3ZI-Q & Runs the logger & Micro USB for flashing and the transcript \\
X-NUCLEO-IKS4A1 & The shield. Two inertial units with programmable engines, an accelerometer, a magnetometer, a pressure sensor, a humidity and temperature sensor, a second temperature sensor & Arduino header, I2C, plus interrupt lines \\
The shield's electrode add-on & Detachable, for the electrostatic sensing channel. Not used in the acceptance run, and named so the reader knows it is there & Connector on the shield \\
One jumper wire & The sensor interrupt line to a free header pin, if the shield's own routing does not reach a usable one & Female both ends \\
nRF-PPK2 & Optional, to price batching against polling with chapter 10's ledger & Leads to the board \\
Host PC & Receives the transcript and checks it & USB \\
\bottomrule
\end{tabular}
\caption{Inventory items used in chapter 13. Nothing is bought and nothing is soldered. The other two shields stay in the drawer for the whole chapter.}
\end{table}

\subsection*{System architecture}

\diagram{c13_arch}{Four layers and one seam. Everything below the shim is vendored and permissively licensed; everything above it is this chapter. The bus arrangement is not a layer, it is a decision that changes what the layers above can do.}

The seam is two functions, read and write, and it is worth noticing how much
that buys. Every sensor on the shield, and every sensor in the collection that
is not on this shield, is reachable through the same two functions, so the
board layer is written once and the sensor set becomes a configuration question
rather than a porting question. The cost is that the shim has to be correct
about repeated starts and about register auto-increment, which are the two
places where a bus implementation that works for a single-byte read quietly
fails for a burst.

\subsection*{Peripheral configuration}

\begin{table}[H]\centering\small
\begin{tabular}{p{22mm}p{26mm}p{24mm}p{30mm}p{30mm}}
\toprule
Peripheral & Mode & Clock source & Pins and function & Interrupt and transfers \\
\midrule
I2C & Fast mode, 400 kHz, repeated start & Peripheral bus & Arduino D14 and D15 by the usual convention for this board family, checked in the board manual before use & Polled for setup; DMA for the burst drain \\
External interrupt & Rising edge, one line & Peripheral bus & The sensor's first interrupt line, on whichever free header pin the shield routes it to & The handler sets a flag and returns \\
Timer & Up counter, 1 kHz & Peripheral bus & None & Timebase for the transcript \\
USART3 & Asynchronous, 115200 & Peripheral bus & Confirm in the board manual & Interrupt driven, chapter 3 \\
DWT cycle counter & Free running & Core clock & None & Measures the drain, if it works with no debugger attached \\
\bottomrule
\end{tabular}
\caption{Peripheral configuration for chapter 13. Two entries are deliberately not written as facts: which pins the bus reaches and which header pin the interrupt line arrives on are checked against the board manual and UM3239 respectively.}
\end{table}

The bus addresses are not reprinted in this chapter. They are in UM3239 and in
the sibling lab set, which is authoritative for them, and reproducing a table of
hexadecimal addresses from memory is the single easiest way to cost somebody an
evening. Read them from the manual, put them in one header, and have the
identification step below prove that every one of them is right before anything
else runs.

\subsection*{Wiring}

\diagram{c13_wiring}{The shield fits the header and the only decision is the jumper block. Three arrangements are possible and the figure names what each one does to the bus and to the interrupt routing.}

The three arrangements are the reason this chapter has a wiring figure at all.
In the first, every sensor sits directly on the bus and the processor addresses
each one. In the second and the third, one of the two inertial units acts as the
hub and reads the others itself, presenting their data inside its own buffer,
and which of the two units takes that role is the difference between the second
arrangement and the third. The hub arrangements reduce the number of
transactions the processor performs and increase the amount of configuration the
firmware has to get right, and they change which device's interrupt line
matters.

\subsection*{Memory and timing budget}

\diagram{c13_mem}{The sensor's buffer, which is the silicon this chapter is really about. Each entry carries a tag that says which source it came from, which is what makes one buffer able to hold several data streams and what makes the drain a decode rather than a copy.}

\begin{table}[H]\centering\small
\begin{tabular}{p{44mm}p{28mm}p{28mm}p{30mm}}
\toprule
Quantity & Budget & Measured & Margin \\
\midrule
Sample rate, inertial & 104 Hz & not measured & not measured \\
Watermark level & half the buffer depth & not measured & not measured \\
Time to fill to watermark & from the rate and the level & not measured & not measured \\
Time to drain one burst & below one tenth of the fill time & not measured & not measured \\
Landing buffer in RAM & 2 kB & not measured & not measured \\
Time in the interrupt handler & below 3 \textmu{}s & not measured & not measured \\
Dropped samples in ten minutes & zero & not measured & not measured \\
Bus occupancy & below 20 percent & not measured & not measured \\
\bottomrule
\end{tabular}
\caption{The budget for chapter 13. The buffer depth is deliberately absent: it is in the sensor datasheet and it is not assumed from a sibling part, because the family shares register names across parts with different depths. Nothing enters the measured column until the bench produces it.}
\end{table}

The one ratio that decides whether this design works is the drain time against
the fill time. If a burst takes a tenth of the time the buffer needs to refill
to the watermark, there is an order of magnitude of margin for an interrupt that
is late, for a bus that is shared and for a drain that is occasionally preempted.
If it takes half, the design works on the bench and produces intermittent drops
in the field, which is the worst of the available outcomes because it looks like
a sensor fault.

\subsection*{Firmware design (UML)}

\diagram{c13_uml}{One watermark event, end to end. The handler sets a flag and returns; everything that takes time happens in the main loop, where it can be interrupted by the next watermark without losing anything.}

The discipline here is the one chapter 3 establishes and it matters more with a
bus in the picture. The interrupt handler must not start a bus transaction. A
burst read at 400 kHz takes a long time by interrupt standards, and a handler
that performs one blocks everything at that priority for the duration, including
the next watermark. Setting a flag and returning costs a few cycles and moves
the work to a place where it can be measured, bounded and preempted.

The drop counter is the second design decision worth defending. The sensor
reports its own overrun condition, and the firmware also checks the tag sequence
as it decodes, so a gap is detected two ways. Both counters live on the board
and are printed by the board, because a host that counts the lines it received
can only prove that the link worked.

\subsection*{Data flow (ASCII)}

\begin{asciiart}
  sensor                              MCU                                     host
  +---------------------------+       +---------------------------------+     +-------------+
  | sample at the batch rate  |       |                                 |     |             |
  |   |                       |       | EXTI handler: set flag, return  |     | transcript  |
  |   v                       | INT   |   |                             |     |  reader     |
  | FIFO: tag + 6 data bytes  |------>|   v                             |     |             |
  |   |   |   |   |           |       | main loop: one burst read       |     |  checks the |
  |   v   v   v   v           | I2C   |   |                             |UART |  board's own|
  | watermark reached --------|<----->|   v                             |---->|  counters   |
  |                           | burst | decode tags -> per source rings |     |             |
  | overrun flag if not drained       |   |                             |     +-------------+
  +---------------------------+       |   v                             |
                                      | drop counter: overrun + tag gap |
                                      +---------------------------------+
\end{asciiart}

\subsection*{Repository layout}

\begin{plaincode}
nucleo-h7a3-iks4a1-fifo/
  CMakeLists.txt
  third_party/stmems/            # vendored, BSD-3-Clause, headers kept intact
  src/bus_shim.c                 # the two functions the drivers need
  src/shield_iks4a1.c            # the board layer: addresses, topology, reset
  src/topology.h                 # which arrangement, and why, in a comment
  src/fifo_drain.c               # burst read, tag decode, per source rings
  src/dropcount.c                # overrun flag plus tag sequence check
  src/exti.c                     # one line, one flag
  src/main.c
  tools/check_transcript.py      # host side: it checks, it does not count
  docs/topology.md               # the three arrangements and the jumper map
  docs/addresses.h.txt           # copied from UM3239 by hand and checked twice
  README.md
\end{plaincode}

\subsection*{Steps}

\begin{steps}
\item \textbf{Choose the bus arrangement before writing code and write down
why.} There are three and they are not interchangeable. Everything on one bus is
the simplest to debug and the most transactions. Either inertial unit as the hub
reduces transactions and moves configuration into the hub device. The jumper map
is in UM3239 and the arrangements are set out more clearly in the upstream
shield description in the other operating system's tree, which is worth reading
first even though it is not C.
\begin{plaincode}
docs/topology.md
  chosen: all sensors directly on the bus
  why:    the acceptance run needs one interrupt source and the simplest
          possible failure mode; the hub arrangements are chapter 16's
          question, where the transaction count is the thing being measured
  jumpers: per UM3239, recorded here after being read off the silkscreen
\end{plaincode}

\item \textbf{Vendor the driver collection and write the shim.} Two functions,
and they are the only place this chapter touches the bus peripheral directly.
\begin{ccode}
/* The porting shim. Every driver in the collection reaches the bus through
 * exactly these two, which is why the board layer is written once.        */
int32_t bus_write(void *handle, uint8_t reg, const uint8_t *data, uint16_t len)
{
    const uint8_t addr = *(const uint8_t *) handle;
    return i2c_mem_write(addr, reg, data, len) ? 0 : -1;
}

int32_t bus_read(void *handle, uint8_t reg, uint8_t *data, uint16_t len)
{
    const uint8_t addr = *(const uint8_t *) handle;
    return i2c_mem_read(addr, reg, data, len) ? 0 : -1;   /* repeated start */
}
\end{ccode}
The read must use a repeated start rather than a stop between the register
address and the data phase, and the sensors must be left with register
auto-increment enabled, or a burst read returns the same register many times.
Both failures produce data that looks like a stuck sensor rather than like a bus
fault.

\item \textbf{Prove every address before anything else.} Each device has an
identification register with a fixed value. Read all of them at startup and stop
with a named error if any disagrees, which turns the most common shield fault
into one line of output instead of an afternoon.
\begin{ccode}
static const struct { const char *name; uint8_t addr; uint8_t id; } ROLL[] = {
    /* addresses and identity values are transcribed from UM3239 and the
     * sensor datasheets; nothing here is remembered rather than read      */
    { "imu.fusion",   ADDR_IMU_FUSION,   ID_IMU_FUSION   },
    { "imu.ispu",     ADDR_IMU_ISPU,     ID_IMU_ISPU     },
    { "accel",        ADDR_ACCEL,        ID_ACCEL        },
    { "magnetometer", ADDR_MAG,          ID_MAG          },
    { "pressure",     ADDR_PRESS,        ID_PRESS        },
    { "temperature",  ADDR_TEMP,         ID_TEMP         },
};
\end{ccode}
The humidity and temperature part on this shield does not follow the same
identification convention as the others, so it is checked with its own command
rather than being forced into this table.

\item \textbf{Configure the buffer and the watermark, in that order.} Set the
batching rate for each source that should appear in the buffer, then the buffer
mode, then the watermark level. A watermark set before the batching rate is
applied against a buffer nothing is being written into.
\begin{ccode}
/* Continuous mode: when full, the oldest entry is replaced and the overrun
 * flag is raised. The alternative stops writing, which hides the overrun
 * and produces a buffer that silently stops advancing.                    */
imu_fifo_set_batch_rate(&ctx, BATCH_104_HZ_ACCEL | BATCH_104_HZ_GYRO);
imu_fifo_set_mode(&ctx, FIFO_MODE_CONTINUOUS);
imu_fifo_set_watermark(&ctx, WATERMARK_WORDS);     /* half the depth        */
imu_int1_route(&ctx, INT1_FIFO_THRESHOLD);
\end{ccode}
Choosing continuous mode rather than the stop-when-full mode is a deliberate
trade: continuous mode loses the oldest samples and tells you it did, and the
other mode keeps the oldest samples and stops, which is harder to notice.

\item \textbf{Route the interrupt and keep the handler empty.} One line, one
edge, one flag.
\begin{ccode}
volatile uint32_t fifo_ready;        /* the only thing the handler writes   */

void EXTI_IMU_IRQHandler(void)
{
    EXTI_CLEAR_PENDING(IMU_INT_LINE);
    fifo_ready++;                    /* counted, not set, so a missed drain */
}                                    /* is visible rather than absorbed     */
\end{ccode}
Incrementing rather than setting is a small thing that pays for itself. If the
counter is ever above one when the main loop looks at it, a watermark arrived
while the previous drain was still running, which is exactly the condition the
fill-to-drain ratio is supposed to prevent.

\item \textbf{Drain in one burst and decode the tags.} The buffer's level
register says how many entries are waiting. Read them all in one transaction and
decode afterwards, because a transaction per entry converts a cheap drain into
an expensive one.
\begin{ccode}
void fifo_drain(void)
{
    uint16_t words = imu_fifo_level(&ctx);
    if (words > LANDING_WORDS) words = LANDING_WORDS;   /* bounded, always  */
    bus_read(&ctx, FIFO_DATA_OUT, landing, words * WORD_BYTES);

    for (uint16_t i = 0; i < words; i++) {
        const uint8_t *w = &landing[i * WORD_BYTES];
        switch (w[0] >> TAG_SHIFT) {                     /* the source tag  */
        case TAG_ACCEL: ring_push(&acc_ring, w + 1); break;
        case TAG_GYRO:  ring_push(&gyr_ring, w + 1); break;
        default:        unknown_tag++;                  break;
        }
    }
}
\end{ccode}
The landing buffer is bounded and the level is clamped to it. An unbounded read
driven by a register value is a buffer overflow waiting for a bus fault to
trigger it.

\item \textbf{Count drops on the board, two ways.} The sensor's overrun flag is
the first. The tag sequence is the second: each source in the buffer advances in
a known order, so a gap in the sequence is a drop even when the overrun flag was
missed.
\begin{ccode}
void dropcount_update(const uint8_t *word)
{
    static uint8_t expect_cnt;
    uint8_t cnt = (word[0] >> CNT_SHIFT) & CNT_MASK;
    if (cnt != expect_cnt) drops += (uint8_t)(cnt - expect_cnt);
    expect_cnt = (uint8_t)(cnt + 1u);
}
\end{ccode}
Both counters are printed by the board. A host that counts the lines it received
proves the serial link worked and nothing else, which is the failure this whole
step exists to avoid.

\item \textbf{Size the watermark from the ratio and not from a guess.} Fill time
is the watermark level divided by the batching rate. Drain time is the
transaction length at 400 kHz plus the decode. Measure the drain with the cycle
counter, compute the fill from numbers you set, and require an order of
magnitude between them.
\begin{shellcode}
# printed by the board at startup, from its own configuration:
# rate 104 Hz  watermark 128 words  fill 1230 ms  drain 000 us (cycle counter)
\end{shellcode}
If the ratio is not there, lower the watermark rather than raising the bus
speed. A lower watermark costs more wake-ups and is easy to price with chapter
10's ledger; a faster bus costs signal margin and is not.

\item \textbf{Run it for ten minutes and let the board report.} The acceptance
criterion is zero drops over a ten-minute run at a stated rate, counted on the
board by both mechanisms.
\begin{shellcode}
python tools/check_transcript.py --minutes 10 --expect-drops 0 \
    --port /dev/ttyACM0 --out runs/iks4a1_10min.json
\end{shellcode}
The host script checks the board's own counters and the monotonicity of the
timebase. It does not count samples, because then a dropped serial line would
look like a dropped sample.

\item \textbf{Price the design against polling.} Run the same logger twice, once
in this form and once with the buffer disabled and the sensor polled at the
sample rate, and measure both with chapter 10's per-phase ledger. This is the
number that justifies the whole chapter, and it is a number this book can produce
and most writing on the subject cannot.
\end{steps}

\subsection*{Build, flash and debug}

\diagram{c13_timing}{Fill against drain on one timebase. The interrupt line rises at the watermark, the burst is one transaction, and the margin in the figure is the ratio the budget table requires.}

\begin{shellcode}
cmake -B build-fw -G Ninja -DCMAKE_TOOLCHAIN_FILE=cmake/arm-none-eabi.cmake
cmake --build build-fw -j
cp build-fw/firmware.bin "$PROBE_DISK"/   # onto the probe disk
\end{shellcode}

\begin{note}{When the bus answers but the data is wrong}
In order of likelihood: register auto-increment is not enabled, so a burst read
returns the first register repeated; the read uses a stop instead of a repeated
start, so the sensor forgets which register was addressed; the shield is in a
different bus arrangement from the one the firmware assumes, so a device that
should be behind the hub is being addressed directly and answers nothing; two
devices share an address because another shield is also fitted, which is the
rule this chapter opens with; or the watermark was configured before the
batching rate and applies to a buffer nothing writes into. Check the
identification roll call first, because it fails cleanly for the first four.
\end{note}

To recover a shield that has stopped answering, reset the sensors through their
software reset bit rather than power cycling the board, and confirm the reset
completed by reading the identification registers again. Nothing in this chapter
touches option bytes, and nothing in this book does until the question of
whether a bad option-byte write can leave this board unrecoverable without
external tooling has been answered.

\subsection*{Verification and acceptance criteria}

\begin{itemize}
\item A ten-minute run at a stated sample rate with zero dropped samples,
counted on the board by both the overrun flag and the tag sequence check, and
printed by the board rather than inferred by the host.
\item No polling anywhere in the logger. Proven by a build in which the only
path that reads the buffer is reached from the interrupt flag, and by a
transcript whose wake count matches the watermark count.
\item The interrupt flag is never above one when the main loop reads it, over
the whole ten minutes, which is the direct evidence that the fill-to-drain ratio
holds.
\item The drain time is measured with the cycle counter and is at most one tenth
of the computed fill time, with both numbers printed by the board at startup.
\item The identification roll call passes for every device on the chosen bus
arrangement, and a deliberately wrong address in the table produces a named
error rather than silence.
\item The chosen bus arrangement is recorded with its jumper positions and the
reason it was chosen, and the firmware refuses to run if the devices it expects
to see are not the ones that answer.
\item A comparison run against a polled version of the same logger, measured
with chapter 10's ledger, gives a charge per sample for each.
\end{itemize}

\subsection*{Variants}

\begin{table}[H]\centering\small
\begin{tabular}{p{24mm}p{26mm}p{44mm}p{22mm}p{20mm}}
\toprule
Axis & Variant & What changes & Cost & Built in full in \\
\midrule
Operating system & The other operating system & It has an in-tree description of this shield and of this board, so the three bus arrangements become a configuration choice rather than firmware. The comparison is cheap and worth making & A different build system and a kernel & Chapter 20 \\
Execution model & An RTOS task set & The drain becomes a task woken by the interrupt, which removes the flag and makes the wake latency measurable as a distribution & Kernel overhead & Chapter 20 \\
Execution model & DMA for the burst read & The drain stops occupying the core, which matters as the watermark grows & Setup per transfer, and a cache question & Chapters 4 and 19 \\
Peripheral substitution & One inertial unit as the hub & The processor addresses one device and that device reads the others, which reduces transactions and moves configuration into the sensor & More configuration to get right & Chapter 16 \\
Intelligence and reach & The classifier inside the sensor & Both inertial units carry a programmable engine, of different kinds, so the decision can be made before the data ever crosses the bus & A separate toolchain per engine & Chapter 16 \\
Synchronisation & A queue between handler and drain & With a kernel the flag becomes a notification and the rings become queues & Kernel overhead & Chapter 20 \\
Language & Host tooling in Python & The transcript checker here & None & Chapter 3 \\
\bottomrule
\end{tabular}
\caption{Variants for chapter 13. This chapter builds none of them in full, which is deliberate: it establishes the buffer and watermark discipline that chapters 16 and 20 then vary, and the book's matrix assigns each variant to the chapter where it teaches most.}
\end{table}

\subsection*{Pitfalls}

\begin{itemize}
\item Starting a bus transaction inside the interrupt handler. A burst at 400
kHz is long by interrupt standards and it blocks the next watermark.
\item Leaving register auto-increment disabled. The burst returns one register
repeated and the data looks like a stuck sensor.
\item Using a stop instead of a repeated start between the register address and
the data phase. It works for a single byte and fails for a burst.
\item Configuring the watermark before the batching rate, so the level applies
to a buffer nothing is writing into.
\item Choosing the stop-when-full buffer mode without meaning to. It keeps the
oldest samples and stops advancing, which is much harder to notice than losing
the oldest and being told.
\item Reading the buffer level register and trusting it without bounding it
against the landing buffer.
\item Counting samples on the host. A dropped serial line then looks like a
dropped sample, and the number that matters is the one the board counts.
\item Assuming the buffer depth from a sibling part. The family shares register
names across parts with different depths, and this volume's standing rule about
inheriting from siblings applies to sensors as well as to the processor.
\item Fitting a second shield. Two devices at one address give readings that are
plausible, stable and wrong.
\end{itemize}

\subsection*{Best practices applied}

\begin{itemize}
\item The vendored dependency keeps its copyright headers and its licence file,
and the licence-flagged middleware is named and not copied.
\item The seam between vendored code and this chapter's code is two functions,
so the boundary is obvious to a reader and cheap to test.
\item Every address and identity value is transcribed from the manual and
checked by the firmware at startup rather than trusted.
\item The design decision that cannot be seen in the code, which is the bus
arrangement, is written down next to the code with its reason.
\item The evidence for the acceptance criterion is produced by the board, which
is the only party that can distinguish a lost sample from a lost line of text.
\item The sizing rule is a ratio with an order of magnitude of margin, stated
before the measurement.
\end{itemize}

\subsection*{Stretch goals}

\begin{itemize}
\item Repeat the acceptance run in each of the three bus arrangements and report
transactions per second, processor wake-ups per second and charge per sample for
each. Nothing published compares the three on one board.
\item Add the pressure and humidity sources to the same buffer through the hub
arrangement and show that the tag decode absorbs them without a second drain
path.
\item Attach the electrode add-on and log the electrostatic sensing channel
alongside the inertial data, which is the one source on this shield that almost
nothing is written about.
\item Move the drain to a transfer engine and measure what that does to the time
the core spends awake, which is the thread chapters 4 and 19 pick up.
\item Push the sample rate until drops appear, and report the rate at which the
fill-to-drain ratio stops holding. A design whose limit is known is worth more
than one that has only been seen to work.
\end{itemize}

\subsection*{Roadmap and next steps}

Chapter 14 fits the other inertial shield and asks a different question, which
is what happens when two full-scale ranges are needed at once. Chapter 16 takes
the bus arrangements from here and turns them into the measurement that decides
where a classifier should run, and chapter 20 rebuilds this same logger on a
kernel, where the shield already has an upstream description and the three
arrangements become configuration.

For the wider path, the community embedded engineering roadmap places sensor
integration and bus work in the same band as the peripheral drivers of the early
chapters and recommends exactly this kind of exercise over further reading, and
the accompanying essay on closing the gap between reading and building makes the
same case. Appendix H lists the courses. For in-sensor processing specifically,
independent writing is almost non-existent, which is recorded in this volume's
provenance notes as a gap rather than an oversight, and it is the reason the
stretch goal above is worth doing.

\subsection*{Portfolio evidence}

\begin{itemize}
\item A public repository with the board layer, the shim, the vendored
collection with its headers intact, and a topology document that states which
arrangement was chosen and why.
\item The ten-minute transcript with the board's own drop counters at zero, and
the startup line showing the computed fill time and the measured drain time.
\item A short comparison of charge per sample between the buffered logger and a
polled one, measured with the instrument named.
\item The identification roll call output, including one deliberately broken run
showing the named error, which is the clearest evidence that the checks are real.
\end{itemize}

\subsection*{Sources}

Normative references:
\begin{itemize}
\item The shield user manual UM3239, for the sensor complement, the bus
addresses and the jumper map that sets the three bus arrangements.
\item The datasheet of each sensor used, for the buffer depth, the tag format,
the batching rates and the identification register value. The depth is not
assumed from a sibling part.
\item Reference manual RM0455, for this part's I2C peripheral, its fast-mode
timing registers and the external interrupt controller. Not RM0433, which
documents its sibling.
\item The board user manual for MB1363, for the Arduino header mapping and for
which pins remain free while this shield is fitted.
\end{itemize}

Reusable implementations:
\begin{itemize}
\item The register-level driver collection covering every sensor on this shield,
BSD-3-Clause, vendored with its copyright headers intact.\newline
\url{https://github.com/STMicroelectronics/STMems_Standard_C_drivers}
\item The upstream shield description in the other operating system's tree,
Apache-2.0, which sets out the three bus arrangements more clearly than the
vendor's own manual.\newline
\url{https://docs.zephyrproject.org/latest/boards/shields/x_nucleo_iks4a1/doc/index.html}
\item Ready-made classifier configurations for the in-sensor engine,
BSD-3-Clause, for the stretch goal and for chapter 16.\newline
\url{https://github.com/STMicroelectronics/st-mems-machine-learning-core}
\item The in-sensor programmable core examples, which declare no licence at
their root and so are checked file by file before anything is reused.\newline
\url{https://github.com/STMicroelectronics/st-mems-ispu}
\item The upstream board description for this exact board in the other operating
system, Apache-2.0, which makes the chapter 20 comparison cheap.\newline
\url{https://docs.zephyrproject.org/latest/boards/st/nucleo_h7a3zi_q/doc/index.html}
\end{itemize}
   % The IKS4A1: FIFO, watermark, interrupt
\project{14}{The IKS5A1: low and high g at once, wide pressure}{NUCLEO-H7A3ZI-Q with X-NUCLEO-IKS5A1}{Simultaneous acceleration ranges, dual full-scale barometer}

\begin{keyfacts}
\item[Board] NUCLEO-H7A3ZI-Q with X-NUCLEO-IKS5A1 on the Arduino header, and no second shield
\item[Peripherals] I2C1 for every sensor, two sensor interrupt lines on free Zio pins, TIM2 as the single timebase, USART3 for the console
\item[Toolchain] arm-none-eabi-gcc with CMake, the vendor's register-level sensor drivers vendored under their own licence
\item[Operating system] Bare metal
\item[Difficulty] 4 of 5
\item[Effort] 4 evenings of about four hours
\item[Deliverable] A shock and altitude logger that records a fine-resolution acceleration window and a wide-range pressure trace around an event the fine channel cannot see, because the coarse channel and the fine channel are both running the whole time
\end{keyfacts}

\subsection*{Why this project}

Every other sensor chapter in this book picks a range. This one refuses to. The
part that gives this shield its character reports a low-range and a high-range
acceleration measurement at the same time, from the same die, on the same
timebase, with no switching between them. That single property changes what a
logger can promise. A logger with one channel either resolves the quiet
behaviour and saturates on the event, or survives the event and cannot see the
quiet behaviour. A logger with both channels running does neither, and the
record it writes is one record rather than two recordings that have to be
aligned afterwards.

The arithmetic is worth doing before any code. Take a sixteen-bit two's
complement output, which is what these parts report, and assume the full scale
is symmetric. At a full scale of 16 g there are 32768 counts on each side, so
one count is about 0.49 mg. At a full scale of 256 g one count is about 7.8 mg.
A vibration signature at 10 mg is a little over twenty counts on the fine
channel and barely more than one count on the coarse one. That is the whole
argument for simultaneous ranges in two numbers, and it is the reason the
chapter exists. Confirm both full-scale lists and the output word format in the
part's datasheet before you trust the constants; the shape of the argument does
not change, but the constants might.

The second property is the barometer, which has two full scales rather than one:
a narrow scale for atmospheric work, where the finer sensitivity earns you
altitude resolution, and a wide scale that keeps reporting where the narrow one
stops. Pressure is the cheapest altitude estimate available and the only one on
this bench that costs no radio, and a sealed enclosure or a short water column
leaves the atmospheric range entirely. The chapter treats the range selection as
a runtime decision with a cost attached to it rather than as a configuration
constant set once at bring-up.

\begin{note}{What is actually fitted to this shield}
The X-NUCLEO-IKS5A1 carries the ISM6HG256X, an inertial measurement unit with
simultaneous low-g and high-g acceleration detection; the ISM330IS, an inertial
unit with the in-sensor programmable core; the IIS2DULPX, an ultra-low-power
accelerometer; the IIS2MDC, a magnetometer; and the ILPS22QS, a barometer with
two selectable full scales of 1260 hPa and 4060 hPa. Four independent sources
agree on that list, the vendor's own product page among them, and UM3550 is the
reference. Two parts named in older bench notes are not on this board. The
LSM6DSV32X is a consumer part on an industrial shield and reaches these boards
only through an adapter kit. The LSM6DSO16IS is the other shield's
programmable-core part, and since the two are siblings with the same core and
different qualification, it is the likely source of the confusion. Chapter 16
uses that other part deliberately.
\end{note}

Confirming that on the silkscreen takes one minute and is Step 1. Read the part
markings, run an address scan, and compare the result against UM3550. Then stop
looking at the shield and spend the rest of the chapter on the measurement,
which is where the engineering is.

\subsection*{Prior art and what to reuse}

\begin{table}[H]\centering\small
\begin{tabular}{p{34mm}p{46mm}p{38mm}p{20mm}}
\toprule
Source & What it gives & What it does not & Licence \\
\midrule
The vendor's register-level driver collection & A complete register map and a typed function per setting for every part on this shield, behind a two-function porting shim & No board layer, no bus arbitration, no ordering between two sensors, no build system, and no policy about what to record & BSD-3-Clause \\
The shield user manual UM3550 & The fitted parts, the addresses, the jumper options and which interrupt line reaches which header pin & Nothing about how to run two acceleration ranges as one measurement & Vendor document \\
The other operating system's in-tree shield definitions & The clearest published statement of how these shields route a bus and an interrupt, for the sibling shield & Whether a definition exists in that tree for this shield is checked at bring-up rather than assumed & Apache-2.0 \\
The vendor's expansion package & Working examples for the same silicon, and a sensor-fusion library & Its fusion component ships as pre-compiled objects under a licence tied to this vendor's silicon, so it is read and never vendored & Mixed, flagged \\
Chapter 13 of this book & The FIFO, the watermark interrupt and the three bus topologies, already built and measured & Nothing about a second range or a second full scale, which is this chapter & This book \\
\bottomrule
\end{tabular}
\caption{Prior art for chapter 14. The driver collection is permissive and is vendored with its copyright headers intact; the fusion middleware is not, and the difference is stated in the repository README rather than discovered by a reader later.}
\end{table}

What is left to write is the part that no source supplies: the policy. The
drivers will set any register you name. Nothing in them decides which channel
arms the recording, how long a pre-trigger window is, when the barometer changes
scale, what a record contains, or what happens when two interrupt lines assert
within the same microsecond. That policy is the chapter, and it is also the
thing a reader can carry to a different sensor next year.

\subsection*{Parts from the inventory}

\begin{table}[H]\centering\small
\begin{tabular}{p{40mm}p{66mm}p{40mm}}
\toprule
Part & Role & Interface \\
\midrule
NUCLEO-H7A3ZI-Q & Bus master, event engine, record writer, console & Micro USB to the host \\
X-NUCLEO-IKS5A1 & The five sensors, fitted to the Arduino header. The only shield on the board & I2C1 plus two interrupt lines \\
A USB data cable & Power, programming and console & Micro USB \\
nRF-PPK2 & Current, and its digital inputs used as the bench's logic recorder for the two interrupt lines & IDD jumper and two digital inputs \\
Host PC & Build, flash, terminal and the record decoder & Windows with the Linux subsystem, or a Raspberry Pi \\
\bottomrule
\end{tabular}
\caption{Inventory items used in chapter 14. Nothing is bought and nothing is soldered. One shield is on the board at a time, because two shields collide on addresses and on the 3V3 budget.}
\end{table}

\begin{note}{One shield at a time}
The three shields on this bench are not stackable in practice even though they
are stackable mechanically. Several of the parts across them answer on the same
addresses, and an address that two devices acknowledge is worse than an address
that nobody answers, because the bus reports success and the bytes are a
blend of two devices. The 3V3 rail on the board is the second reason: it is
sized for the board and one shield. This rule comes from the kit labs in the
sibling volume and is treated as authoritative here rather than re-derived.
\end{note}

\subsection*{System architecture}

\diagram{c14_arch}{One bus, five sensors, two interrupt lines and one timebase. The two acceleration channels come from the same part and therefore share a clock, which is the property the whole chapter is built on.}

The shape is conventional and the interesting part is what is not in it. There
is no fusion library, no orientation estimate and no filtering on the board
beyond a scale conversion. The firmware's job is to decide what to keep and to
keep it with an honest timestamp. Everything else can be done on the host from
the record, and anything done on the board that could have been done on the host
costs energy in the place where energy is scarce.

\subsection*{Peripheral configuration}

\begin{table}[H]\centering\small
\begin{tabular}{p{22mm}p{26mm}p{24mm}p{30mm}p{30mm}}
\toprule
Peripheral & Mode & Clock source & Pins and function & Interrupt and transfers \\
\midrule
I2C1 & Fast mode, 400 kHz & Peripheral bus, source selected in RCC & D14 and D15 on the Arduino header, which are PB9 and PB8 by the Nucleo-144 convention. Confirm in MB1363 & Interrupt-driven transfers. No DMA in this chapter \\
GPIO input & Rising edge, external interrupt & Peripheral bus & Low-g watermark line from the ISM6HG256X & EXTI, pre-emption priority 6 \\
GPIO input & Rising edge, external interrupt & Peripheral bus & High-g event line from the same part & EXTI, pre-emption priority 5 \\
TIM2 & Free running, 32-bit, 1 MHz tick & Peripheral bus & None & Read by both handlers. No interrupt \\
GPIO output & Push-pull marker pins & Peripheral bus & Two free Zio pins into the PPK2 digital inputs & None \\
USART3 & Asynchronous, 115200 8N1 & Peripheral bus & The console pins of chapter 1 & Polled \\
\bottomrule
\end{tabular}
\caption{Peripheral configuration. The high-g line is given the higher pre-emption priority because it is rare and time critical, while the watermark line is frequent and tolerant of a few microseconds of latency. The Arduino header mapping is convention until it is read in MB1363.}
\end{table}

The bus runs at 400 kHz rather than 100 kHz because the FIFO drain is the
chapter's throughput constraint and the difference is a factor of four in the
time the handler occupies. Three things decide whether that rate is actually
available. The pull-up resistors on the shield set the rise time, and they are
the shield's and not yours, so the practical ceiling is a property of the board
you are holding. The peripheral's timing register has to be programmed from the
peripheral clock that RCC is actually delivering, which on this part is not the
same branch as on its better-known sibling and is read in RM0455. And any device
on the bus may stretch the clock, so a transfer that completes in a computed
time on paper does not necessarily complete in that time on the bench. The
chapter computes the expected drain time, then measures it, and the difference
between the two is a result rather than an inconvenience.

TIM2 is the single timebase for both channels and for the barometer, and it is
32 bits wide, which at a 1 MHz tick wraps in about 71.6 minutes. Every record
carries the raw tick and the host does the conversion. Timestamping in the
handler rather than in the main loop is the whole point: a timestamp taken after
the record has been assembled measures when the firmware got around to it, not
when the sensor asserted.

\subsection*{Wiring}

\diagram{c14_wiring}{The shield sits on the Arduino header and nothing is wired by hand except the two marker pins into the power instrument. The interrupt lines reach header pins through the shield itself; which pins they reach is read from UM3550 rather than assumed.}

Two safety points before the shield goes on. Everything here is 3.3 V logic and
no 5 V signal reaches a GPIO on this board. The PPK2 measures up to 1 A, which
is far above anything this pair draws, but its digital inputs have their own
threshold and are read as logic only; they are not an analogue instrument. There
is no soldering iron on this bench, so any option on this shield that needs a
solder bridge moved is out of scope and the chapter says which configuration it
is using instead.

\subsection*{Memory and timing budget}

\diagram{c14_mem}{The record and the ring. The pre-trigger ring holds raw low-g samples so that the window before an event survives the event; the record written to the log holds both channels and the pressure trace with one timebase across all three.}

\begin{table}[H]\centering\small
\begin{tabular}{p{44mm}p{28mm}p{28mm}p{30mm}}
\toprule
Quantity & Budget & Measured & Margin \\
\midrule
Pre-trigger ring, 2048 low-g samples & 12 kB of SRAM & not measured & not measured \\
One event record, 500 ms window & about 5 kB & not measured & not measured \\
Firmware image including five drivers & 48 kB of flash & not measured & not measured \\
High-g line to first stored sample & 200 \textmu{}s & not measured & not measured \\
Barometer full-scale change to first trustworthy sample & 50 ms & not measured & not measured \\
Charge per stored event & not budgeted & not measured & not measured \\
\bottomrule
\end{tabular}
\caption{The budget table. The ring figure is arithmetic: 2048 samples of six bytes is 12288 bytes. The window figure is arithmetic too: 500 ms at 1667 Hz is 834 samples of six bytes, about 5 kB. Everything in the measured column stays blank until the DWT cycle counter or the PPK2 has produced it.}
\end{table}

The interesting budget line is the last timing row. Changing the barometer's
full scale is not free: the sensitivity constant changes, and the first
conversions after the change are taken under the new configuration but not
necessarily settled. Treating the change as instantaneous produces a pressure
trace with a step in it that looks exactly like a real event, which is the worst
kind of error because it is plausible. The chapter measures the settling
interval with the data-ready line into a PPK2 digital input and then encodes the
answer as a blanking period in the firmware.

\subsection*{Firmware design (UML)}

\diagram{c14_uml}{The event sequence. The low-g channel fills a ring continuously; the high-g line is what decides that a stretch of that ring is worth keeping. The barometer is asked for a wide-scale trace only inside the event window.}

The design is an armed recorder, not a sampler. In the quiet state the low-g
channel streams into a circular buffer through the FIFO and the watermark
interrupt built in Chapter 13, and nothing is stored. The high-g line asserts on
a threshold crossing configured in the sensor rather than tested in firmware,
which matters because a threshold tested in firmware needs the samples to have
arrived first, and by then the leading edge of the event is already inside the
ring rather than ahead of it. On assertion the firmware marks the ring position,
starts a post-trigger countdown, changes the barometer to its wide scale, waits
out the blanking period, and then assembles one record from the ring, the high-g
samples and the pressure trace.

Two ordering rules keep it honest. First, the ring is never rewound: the
pre-trigger window is whatever the ring already held, and if the ring is shorter
than the requested window the record says so in a flag rather than silently
returning a short window. Second, a second high-g assertion inside an open
window extends the window rather than starting a second record, because two
overlapping records of one physical event are harder to interpret than one long
one.

\subsection*{Data flow (ASCII)}

\begin{asciiart}
  X-NUCLEO-IKS5A1                       STM32H7A3                         host
  +---------------------------+         +-----------------------------+   +-------------+
  | ISM6HG256X                |  I2C1   | i2c1_xfer (interrupt)       |   | decoder.py  |
  |   low-g path  -> FIFO     +-------->|   |                         |   |   |         |
  |   high-g path -> threshold|         |   v                         |   |   v         |
  |                           |  INT1   | ring[2048] 6 B per sample   |   | CSV + plot  |
  |   INT1 watermark  --------+-------->|   |                         |   |             |
  |   INT2 high-g     --------+-------->|   +--> arm: mark, countdown  |   |             |
  +---------------------------+  INT2   |        |                    |   |             |
  | ILPS22QS 1260 / 4060 hPa  |  I2C1   |        v                    |   |             |
  |   data-ready      --------+-------->| record assembler            |   |             |
  +---------------------------+         |   |                         |   |             |
                                        |   v                         |USART3           |
  PPK2 digital in D0 <-- marker A       | framed record ---------------+-->|             |
  PPK2 digital in D1 <-- marker B       +-----------------------------+   +-------------+
\end{asciiart}

\subsection*{Repository layout}

\begin{plaincode}
nucleo-h7a3-iks5a1-shocklog/
  CMakeLists.txt
  cmake/arm-none-eabi.cmake
  third_party/stmems/           # vendored, BSD-3-Clause, headers untouched
    ism6hg256x_reg.[ch]
    ilps22qs_reg.[ch]
    iis2mdc_reg.[ch]
    LICENSE                     # copied verbatim, not summarised
  src/bus_i2c1.c                # the two-function porting shim
  src/shock_cfg.c               # ranges, rates, thresholds, in one place
  src/ring.c                    # pre-trigger ring, from chapter 2
  src/recorder.c                # the policy: arm, extend, assemble
  src/baro_range.c              # full-scale selection and its blanking period
  src/timebase.c                # TIM2 at 1 MHz, read from both handlers
  src/main.c
  tools/decode_record.py        # the host twin of the record format
  tools/scan_i2c.py             # address scan, run once at bring-up
  docs/PROVENANCE.md            # what is vendored, under which licence
  README.md
\end{plaincode}

\subsection*{Steps}

\begin{steps}
\item \textbf{Confirm the shield before writing a line of driver code.} Read the
part markings on the silkscreen and compare them against UM3550. You are looking
for ISM6HG256X, ISM330IS, IIS2DULPX, IIS2MDC and ILPS22QS. If you find an
LSM6DSV32X or an LSM6DSO16IS you are holding a different board, and every
address and register in this chapter is wrong for it. Then confirm electrically
with a scan, because a marking you misread and an address that answers are two
independent pieces of evidence.
\begin{shellcode}
python tools/scan_i2c.py --port /dev/ttyACM0
# prints every address that acknowledges, one line each
\end{shellcode}
Write the addresses the scan returns into a header with a comment naming UM3550
as the source. This book does not print sensor addresses it has not read on the
bench, and neither should your repository.

\item \textbf{Vendor the register-level drivers and write the shim.} The
collection is BSD-3-Clause, so it can be copied into the repository with its
copyright headers intact. Record that in \texttt{docs/PROVENANCE.md} on the same
day you copy it, not later. Each driver needs exactly two functions from you.
\begin{ccode}
/* The porting shim. One instance per sensor, differing only in address. */
typedef struct { uint16_t addr; } bus_ctx_t;

int32_t stmems_write(void *h, uint8_t reg, const uint8_t *buf, uint16_t n)
{
    const bus_ctx_t *c = (const bus_ctx_t *) h;
    return i2c1_write_reg(c->addr, reg, buf, n) ? 0 : -1;
}

int32_t stmems_read(void *h, uint8_t reg, uint8_t *buf, uint16_t n)
{
    const bus_ctx_t *c = (const bus_ctx_t *) h;
    return i2c1_read_reg(c->addr, reg, buf, n) ? 0 : -1;
}
\end{ccode}
The driver's own configuration function names are taken from its header when you
compile against it. This book does not print an API it has not compiled, which
is why the configuration below goes through a wrapper of your own.

\item \textbf{Configure both acceleration channels and prove they are both
live.} This is the step the chapter exists for. One structure holds both, and
neither is a mode you switch into.
\begin{ccode}
typedef struct {
    uint16_t lo_odr_hz;    /* fine channel rate, e.g. 1667 */
    uint16_t lo_fs_g;      /* fine channel full scale, e.g. 16 */
    uint16_t hi_fs_g;      /* coarse channel full scale, from the datasheet */
    uint16_t hi_thresh_mg; /* threshold evaluated in the sensor, not here */
    uint16_t fifo_wtm;     /* watermark in samples, see chapter 13 */
} shock_cfg_t;

/* Acceptance for this step: with the board held still, the fine channel
   reports about 1 g on one axis with counts of noise, and the coarse
   channel reports the same 1 g with far fewer counts of resolution.
   Both at the same time, from one part, on one timestamp. */
\end{ccode}
If only one channel ever reports, you have configured a range rather than two
channels, and the rest of the chapter does not work. Fix it here.

\item \textbf{Bring up the FIFO and the watermark line.} The mechanism is
Chapter 13's and is not repeated: the sensor fills its own FIFO, asserts a line
at the watermark, and the handler drains a known number of samples rather than
polling a status bit. What is new here is where the samples go, which is a
circular pre-trigger ring rather than a consumer.
\begin{ccode}
void EXTI_wtm_handler(void)
{
    const uint32_t t = timebase_ticks();      /* TIM2, before any bus work */
    uint8_t raw[FIFO_WTM * 6];
    (void) i2c1_read_fifo(ACC_ADDR, raw, sizeof raw);
    ring_push_block(&pretrig, raw, FIFO_WTM, t);
    marker_pulse(MARKER_A);                   /* PPK2 digital input D0 */
}
\end{ccode}

\item \textbf{Arm the recorder on the high-g line.} The threshold lives in the
sensor. The handler does almost nothing, because the handler runs at the moment
that matters and everything it does is latency in the record.
\begin{ccode}
volatile uint32_t g_event_tick;
volatile uint8_t  g_event_pending;

void EXTI_highg_handler(void)
{
    g_event_tick = timebase_ticks();
    g_event_pending = 1;                      /* the main loop assembles */
    marker_pulse(MARKER_B);                   /* PPK2 digital input D1 */
}
\end{ccode}
Measure the interval from the marker pulse to the first byte of the assembled
record with the DWT cycle counter, and cross-check the two marker pulses on the
PPK2 digital inputs, which are time-aligned with its current trace. That
cross-check is what this bench has instead of a logic analyser, and it is enough
for intervals of this length.

\item \textbf{Change the barometer's full scale and find out what that costs.}
Select the narrow 1260 hPa scale for the quiet state, where the finer
sensitivity is what gives altitude any resolution, and the wide 4060 hPa scale
inside an event window. Then measure the cost of the change rather than assuming
it is free.
\begin{ccode}
/* Blanking: samples taken inside this window are tagged, not discarded,
   so the host can see what the firmware chose to distrust. */
#define BARO_BLANK_US  50000u      /* placeholder until measured */

void baro_set_range(baro_range_t r)
{
    baro_apply_range(r);                       /* wrapper over the driver */
    g_baro_valid_from = timebase_ticks() + BARO_BLANK_US;
}
\end{ccode}
To measure it: hold the pressure steady, change the scale, and record every
sample with its timestamp. The interval from the change to the point where the
reading stops moving is the blanking period. Replace the placeholder with the
measured figure and say in the README which instrument produced it.

\item \textbf{Define the record and write its host twin at the same time.} One
record, one timebase, both channels, explicit flags for everything the firmware
distrusted. Write the decoder in the same sitting, because a format with no
second implementation is a format with no specification.
\begin{ccode}
typedef struct __attribute__((packed)) {
    uint32_t tick_event;      /* TIM2 ticks at the high-g assertion */
    uint16_t n_pre;           /* fine-channel samples before the event */
    uint16_t n_post;          /* fine-channel samples after it */
    uint16_t n_baro;          /* pressure samples in the window */
    uint8_t  lo_fs_g;
    uint8_t  hi_fs_g;
    uint8_t  baro_range;      /* 0 = 1260 hPa, 1 = 4060 hPa */
    uint8_t  flags;           /* bit0 short window, bit1 extended, bit2 baro blanked */
} record_hdr_t;
\end{ccode}

\item \textbf{Produce a controlled event and check the record against it.} No
instrument on this bench generates a calibrated shock, so do not pretend to one.
A repeatable event is enough: let the board and shield, still attached to their
cable, settle onto a foam pad from a fixed low height, ten times. The record
should show a quiet fine channel, a coarse-channel excursion at the moment the
marker pin pulsed, and a pressure trace that is unremarkable, which is itself
the result. Then repeat with the assembly inside a sealed bag that you squeeze,
which moves the pressure and leaves the acceleration alone, and check that the
two channels are independent in the record.

\item \textbf{Decode and plot on the host.} The decoder reads the framed record
from the console, converts counts to engineering units using the scale codes in
the header rather than a constant compiled into the script, and derives altitude
from pressure with the standard relation. Keeping the scale in the record is
what lets a record recorded under one configuration be read years later.
\begin{pycode}
# altitude from pressure, the standard barometric relation
h_m = 44330.0 * (1.0 - (p_hpa / p0_hpa) ** 0.1903)
\end{pycode}

\item \textbf{Find the rate at which the bus stops keeping up.} The fine channel
rate is the one number in the configuration that has a hard ceiling, and the
ceiling is not in the sensor. Compute the expected drain time first, so that you
know what you are looking for.
\begin{shellcode}
# 128 samples of 6 bytes is 768 bytes, plus address and register overhead.
# At 400 kHz each byte costs 9 bit times, so 768 x 9 / 400000 = 17.3 ms.
# Computed from the bus rate. Not measured.
\end{shellcode}
Then raise the rate in steps, and at each step record two numbers with the DWT
cycle counter: the time spent inside the watermark handler, and the fraction of
the sample period that represents. When the handler occupies more than about a
quarter of the period, the design has stopped being a logger and has become a
bus benchmark, and the answer is either a larger watermark, a faster bus, or the
transfer engine of Chapter 4. Record where that point is, because it is the
number a reader of your repository will want and nobody has published it for
this shield.
\end{steps}

\subsection*{Build, flash and debug}

\diagram{c14_timing}{The two interrupt lines and the window they define. The pre-trigger portion is already in the ring when the high-g line asserts, which is the only way to record the leading edge of an event the fine channel cannot detect in time.}

\begin{shellcode}
cmake -B build-fw -G Ninja -DCMAKE_TOOLCHAIN_FILE=cmake/arm-none-eabi.cmake
cmake --build build-fw -j
cp build-fw/shocklog.bin "$PROBE_DISK"/   # onto the probe disk
python tools/decode_record.py --port /dev/ttyACM0 --out events/
\end{shellcode}

Recovering a board that will not answer follows Chapter 1: the probe presents a
disk, and copying a known-good binary onto it needs no tools at all. The failure
specific to this chapter is different, and it is a silent one.

\begin{note}{When the high-g channel reports nothing}
In order of likelihood: the second channel was never enabled, so the part is
reporting one range twice; the threshold is set in the wrong units, since a
threshold expressed in milli-g and a threshold expressed in full-scale counts
differ by whatever the full scale happens to be; the interrupt line is routed to
a header pin the firmware is not watching, which UM3550 settles in a minute; the
line is configured for the wrong edge; and the pull configuration on the input
pin opposes the shield's own. Check them in that order, and check the first one
by reading both channels back with the board held still, where a working pair
reports the same acceleration at visibly different resolutions.
\end{note}

\subsection*{Verification and acceptance criteria}

\begin{itemize}
\item With the board held still, both acceleration channels report the same
value within their respective resolutions, simultaneously, from one part, with
one timestamp. This is the chapter's central claim and it is checked first.
\item The pre-trigger window in a stored record contains samples whose
timestamps are earlier than the high-g assertion tick, proving the ring held the
leading edge rather than the firmware reconstructing it.
\item A second high-g assertion inside an open window extends that window and
sets the extension flag. Two overlapping records for one physical event never
appear.
\item The barometer's blanking period is a measured number in the README with
its instrument named, not the placeholder constant, and samples inside it are
tagged rather than discarded.
\item The host decoder reads the record using only the scale codes carried in
the header, proven by decoding a record captured under a different full-scale
configuration.
\item A pressure event with no acceleration and an acceleration event with no
pressure change both produce records in which the other channel is unremarkable.
\item Every measured figure in the repository names the DWT cycle counter, TIM2
or the PPK2. Nothing is asserted from a feeling about how long something took.
\end{itemize}

\subsection*{Variants}

\begin{table}[H]\centering\small
\begin{tabular}{p{24mm}p{26mm}p{44mm}p{22mm}p{20mm}}
\toprule
Axis & Variant & What changes & Cost & Built in full in \\
\midrule
Execution model & Polled loop & Poll the status bit instead of watching the watermark line. Simpler, and it loses the leading edge of exactly the events this chapter is about & The pre-trigger window becomes a fiction & Chapter 3 \\
Execution model & DMA into the ring & The FIFO drain becomes a transfer the core does not supervise, which matters once the fine channel rate rises & Cache maintenance on the ring buffer & Chapter 4 and Chapter 19 \\
Time and safety & RTC and a low-power timer & The quiet state stops being a busy loop and becomes a wait, which is what makes a logger a logger & Wake latency added to the pre-trigger reasoning & Chapter 7 \\
Intelligence and reach & Classifier in the sensor & The ISM330IS on this same shield carries the programmable core, so the arming decision could be made in the sensor rather than by a threshold & A toolchain and a licence audit & Chapter 16 \\
Operating system & An RTOS task set & The drain, the assembler and the console become three tasks with stated priorities and a measured latency distribution & Static allocation and a scheduler to justify & Chapter 20 \\
Peripheral substitution & A different bus instance & The same five sensors read through another I2C instance, which changes the clock domain and therefore what survives a low-power mode & Pin availability while a shield is fitted & Chapter 17 \\
\bottomrule
\end{tabular}
\caption{Variants for chapter 14. Nothing in this table is built in full here; this chapter's contribution is the simultaneous-range policy, and each variant is named so a reader knows where the full treatment lives.}
\end{table}

\subsection*{Pitfalls}

\begin{itemize}
\item Trusting an older bench note about which parts are on this shield. Two
parts that appear in such notes are not fitted, and one of them is a sibling of
a part that is, which makes the mistake survive a casual check.
\item Fitting two shields at once. Addresses collide and the 3V3 budget is
shared. One shield is on the board at a time.
\item Treating the two acceleration channels as a range selection. If your
configuration code has a branch that chooses between them, you have not
configured this part, you have configured a different one.
\item Setting a threshold in the wrong units. Milli-g and full-scale counts are
not the same number and the error scales with the full scale, so it is small at
16 g and large at 256 g.
\item Timestamping after the bus transaction. The handler reads the timebase
first, before any I2C work, or the timestamp measures the bus rather than the
event.
\item Assuming the barometer's full-scale change is instantaneous. The step it
produces in an untagged trace looks exactly like a real pressure event.
\item Compiling the scale constants into the host decoder. The record carries
its own scale codes so that it can be read later, by someone else, on a
different configuration.
\item Reading register offsets from a tutorial written for a different member of
this sensor family. The sibling parts share a core and differ in qualification,
addresses and available registers.
\end{itemize}

\subsection*{Best practices applied}

\begin{itemize}
\item The prior art is named and its licence is stated before any code is
written around it, and the vendored copy keeps its copyright headers and its
licence file verbatim rather than in summary.
\item Every measurement names its instrument, and a measurement that has not
been taken is written as not measured rather than estimated.
\item Facts about the shield carry their evidence: the part list is confirmed
from four independent sources and then again on the silkscreen and on the bus,
and the addresses are read rather than printed from memory.
\item The record format has two implementations from the first day, one on the
board and one on the host, which is the cheapest specification there is.
\item What the firmware distrusts is flagged in the data rather than removed
from it, so that a later reader can disagree with the firmware.
\end{itemize}

\subsection*{Stretch goals}

\begin{itemize}
\item Add the IIS2DULPX as a third acceleration channel and compare its noise
floor against the fine channel of the main part at the same rate. The
ultra-low-power part is not a smaller version of the other one, and the
comparison is a genuine result because nobody has published it for this shield.
\item Use the magnetometer to record orientation at the moment of the event, so
that a record says which way up the assembly was. Keep the estimate on the host
using a permissive orientation filter rather than putting fusion on the board.
\item Store records in flash rather than streaming them, which raises the
write-granularity and error-correction questions this family has and which are
on the confirm list rather than settled.
\item Drive the arming decision from the in-sensor core on the ISM330IS and
compare the latency and the charge against the threshold path built here. That
comparison is Chapter 16's subject and this shield can host half of it.
\end{itemize}

\subsection*{Roadmap and next steps}

Chapter 15 changes the sensor and keeps the discipline: a time-of-flight array
whose output is sixty-four values that nobody wants, reduced on the
microcontroller to a decision that somebody does. Chapter 16 then takes the
arming decision built here and moves it to three different places to find out
what each one costs.

For going wider, the community embedded engineering roadmap distinguishes
firmware from embedded Linux from hardware and is explicit about which to weight
at which stage, and its central claim, that projects outrank reading, is this
book's claim too. The accompanying essay on bridging the gap between the two is
the short version. For the sensor-intelligence direction this chapter points at,
the free open textbook that grew out of a university course is the closest thing
to a canonical curriculum, and the free-to-audit course on embedded machine
learning covers the same ground with exercises. Appendix H lists them with their
licences, which matters here: one widely recommended video course releases its
code under terms that make it something to recommend and never to quote.

\subsection*{Portfolio evidence}

\begin{itemize}
\item A public repository whose README opens with the architecture figure and
states in its first paragraph that this shield carries the ISM6HG256X and not
the part older notes name, with the four sources listed.
\item A plot of one event showing both acceleration channels on one time axis,
with the high-g assertion marked, and the pre-trigger portion visibly earlier
than the assertion.
\item A PPK2 capture in which the two marker pins and the current trace are
time-aligned, which is what this bench has instead of a logic analyser and is
worth saying so in the caption.
\item The provenance document listing every vendored file, its licence and the
date it was copied.
\end{itemize}

\subsection*{Sources}

Normative references:
\begin{itemize}
\item UM3550, the user manual for the X-NUCLEO-IKS5A1, for the fitted parts, the
addresses, the jumper options and the interrupt routing to header pins.
\item The datasheets of the ISM6HG256X, ISM330IS, IIS2DULPX, IIS2MDC and
ILPS22QS, for the full-scale lists, the output word format, the output data
rates and the interrupt behaviour.
\item Reference manual RM0455 for I2C1 and the external interrupt controller on
this part. Not RM0433, which documents its better-known sibling.
\item The board user manual for MB1363, for the Arduino header mapping, which
pins D14 and D15 reach, and which Zio pins stay free while a shield is fitted.
\end{itemize}

Reusable implementations:
\begin{itemize}
\item The vendor's register-level driver collection, BSD-3-Clause, covering
every sensor on this shield behind a two-function porting shim.\newline
\url{https://github.com/STMicroelectronics/STMems_Standard_C_drivers}
\item The other operating system's in-tree shield definitions, Apache-2.0, which
document the bus topologies of the sibling shield more clearly than the vendor's
own manual.\newline
\url{https://docs.zephyrproject.org/latest/boards/shields/x_nucleo_iks4a1/doc/index.html}
\item The in-sensor processor examples for the programmable core on the
ISM330IS. The repository declares no licence at its root, so files are checked
individually and nothing is vendored on faith.\newline
\url{https://github.com/STMicroelectronics/st-mems-ispu}
\item The current home of the well-known orientation filter, MIT, for the
host-side orientation stretch goal.\newline
\url{https://github.com/xioTechnologies/Fusion}
\item The Python interface to the power instrument, GPL-2.0. Install it from its
package index; do not fork it into a portfolio repository.\newline
\url{https://github.com/IRNAS/ppk2-api-python}
\end{itemize}
   % The IKS5A1: low and high g at once, wide pressure
\project{15}{Eight by eight time of flight, decided on the MCU}{NUCLEO-H7A3ZI-Q with X-NUCLEO-53L8A1}{Sensor firmware upload, zone reduction, hysteresis}

\begin{keyfacts}
\item[Board] NUCLEO-H7A3ZI-Q with X-NUCLEO-53L8A1 on the Arduino header, and no second shield
\item[Peripherals] I2C1 in fast mode, the sensor's interrupt line on a header pin, the independent watchdog, TIM2 as the timebase, USART3 for the console
\item[Toolchain] arm-none-eabi-gcc with CMake, and a driver taken from the redistributable republication rather than from behind a click-through
\item[Operating system] Bare metal
\item[Difficulty] 4 of 5
\item[Effort] 4 evenings of about four hours
\item[Deliverable] A ranging node that uploads an 84 kilobyte firmware image to its sensor at every power-on, reduces sixty-four zones to one status byte with hysteresis, and sends the host that byte and nothing else
\end{keyfacts}

\subsection*{Why this project}

This sensor reports an eight by eight grid of distances. Sixty-four zones, each
with a distance, a target status and a measure of how much the sensor trusts
itself, arriving several times a second. Almost nobody wants sixty-four numbers.
What a system wants is a decision: the way is clear, something is approaching,
something is present. Sending sixty-four values to a host so that the host can
compute one is the design mistake this chapter exists to avoid, and it is a
common one because the driver makes the sixty-four values so easy to obtain.

The second reason is a property of the part that changes the shape of the whole
firmware. The sensor has volatile memory only. It holds no firmware of its own
between power cycles, so the microcontroller stores an image of about 84
kilobytes in its own flash and pushes it across the bus at every power-on. At
400 kHz with nine bit times per byte that transfer is
$86016 \times 9 / 400000$, which is about 1.94 seconds. That figure is computed
from the bus rate and is not measured. It is also not a detail: it means the
node is blind for seconds after reset, it means the watchdog cannot be armed
with a short window across that period, it means 84 kilobytes of the
microcontroller's 2 megabytes belong to a peripheral, and it means a supply
disturbance that resets the sensor without resetting the microcontroller must be
detected and answered rather than ignored.

The third reason is that the decision belongs on the microcontroller and this
chapter argues for that placement explicitly. Chapter 16 tests the argument by
building the same decision in three places and measuring each. Here the choice
is made on the grounds that a status byte crossing a link costs one transaction
where a frame costs many, and that a policy small enough to fit in a few hundred
lines is a policy a customer can be shown.

\begin{note}{What the sensor is actually running}
The image uploaded here is firmware for a programmable core inside the sensor,
not a configuration blob. The core belongs to a processor family that came out
of this vendor's many-core research work, and there is a good independent
write-up of that lineage listed in the sources. Reading it is not required to
finish the chapter, but it answers the question that the upload raises: this is
a small computer being loaded, which is why the image is this large and why the
part needs seconds rather than milliseconds before it can range.
\end{note}

\subsection*{Prior art and what to reuse}

\begin{table}[H]\centering\small
\begin{tabular}{p{34mm}p{46mm}p{38mm}p{20mm}}
\toprule
Source & What it gives & What it does not & Licence \\
\midrule
The vendor's Arduino organisation, republishing the ultra-lite driver & The complete driver including the firmware image, under a licence that can be redistributed. This is the path this chapter takes & It arrives wrapped in a library for a different environment, so the wrapper is stripped and the copyright headers are kept & BSD-3-Clause \\
The same driver from behind the click-through on the vendor's own site & The same code, and the licence a reader has to accept in a browser & No public repository, so no way to pin a version or show provenance in a build & Click-through \\
The vendor's firmware package for this silicon & Board integration examples for the same shield & A per-file licence mix. It is audited file by file before anything is copied, and this chapter copies nothing from it & Mixed, flagged \\
UM3120, the shield user manual & The bus selection, the interrupt routing, the header mapping and the power arrangement & Nothing about what to do with sixty-four numbers & Vendor document \\
The independent write-up on the sensor's internal core & Why the image is this size and why it is volatile & Nothing you need to compile & Reading only \\
\bottomrule
\end{tabular}
\caption{Prior art for chapter 15. The licence column is the reason this chapter exists in the form it does: one of these two identical drivers can be published in a portfolio repository with its provenance shown and one cannot.}
\end{table}

What is left to write is the port and the policy. The port means taking the
driver out of its wrapper, writing the four platform functions it needs against
this board's I2C driver, placing the 84 kilobyte image where the linker can see
it, and recording in \texttt{docs/PROVENANCE.md} which upstream commit it came
from and under which licence. The policy is everything that turns a frame into a
byte: which zones are believed, how many have to agree, and what stops the
output changing state every time a person shifts their weight.

\subsection*{Parts from the inventory}

\begin{table}[H]\centering\small
\begin{tabular}{p{40mm}p{66mm}p{40mm}}
\toprule
Part & Role & Interface \\
\midrule
NUCLEO-H7A3ZI-Q & Holds the sensor's firmware image, uploads it, reduces frames, drives the console & Micro USB to the host \\
X-NUCLEO-53L8A1 & The ranging sensor on the Arduino header. The only shield on the board & I2C1 plus an interrupt line \\
A USB data cable & Power, programming and console & Micro USB \\
nRF-PPK2 & The current trace that shows the upload as a plateau, plus one digital input as a marker & IDD jumper and one digital input \\
Host PC & Build, flash, and a terminal that receives one byte per decision & Windows with the Linux subsystem, or a Raspberry Pi \\
\bottomrule
\end{tabular}
\caption{Inventory items used in chapter 15. Nothing is bought and nothing is soldered, which is why the bus selection question below has to be answered before the chapter is planned rather than after.}
\end{table}

\begin{note}{The one question that has to be answered before this chapter starts}
Whether this shield selects its bus with a jumper or with a solder bridge is on
this volume's confirm list and is settled by reading UM3120 and by looking at
the board. There is no soldering iron on this bench. If the selection needs a
bridge moved, this chapter is I2C only, the faster bus is described and not
built, and every timing figure in it is an I2C figure. That is the assumption
the text below is written under, and the paragraph after the wiring figure says
what changes if the answer turns out to be a jumper.
\end{note}

\subsection*{System architecture}

\diagram{c15_arch}{The image lives in the microcontroller's flash and is pushed into the sensor's volatile memory at every power-on. After that, frames come back and only a byte leaves the board.}

The asymmetry in that figure is the chapter. Eighty-four kilobytes travel one
way, once per power cycle, and one byte travels the other way, several times a
second, for as long as the node runs. A design that got this backwards, pushing
frames to a host that makes the decision, would move a kilobyte per frame across
a link for the life of the product.

The arithmetic makes the case better than the argument does. A result block of
about a kilobyte at fifteen frames per second is roughly 15 kB per second
leaving the board, continuously. A status byte at the same rate is 15 bytes per
second. That is a factor of a thousand in link traffic, and on any link that
costs energy per byte, which is every link in Part II of this book, it is a
factor of a thousand in the part of the energy budget the radio owns. The
reduction itself costs a few hundred microseconds of a core that is otherwise
idle between frames. The trade is not close, and the only reason it is worth
stating at all is that the opposite design is the one the driver's interface
gently encourages.

\begin{note}{The one number that decides this chapter's architecture}
Everything above follows from the sensor having no non-volatile memory of its
own. If it held its firmware, the board would boot in milliseconds, the flash
would be free, the watchdog ordering would be ordinary and the re-upload state
would not exist. One packaging decision inside a part the size of a grain of
rice determines four separate properties of the firmware around it. That is
worth noticing, because the next sensor you choose will have its own version of
this and it will not be in the marketing material either.
\end{note}

\subsection*{Peripheral configuration}

\begin{table}[H]\centering\small
\begin{tabular}{p{22mm}p{26mm}p{24mm}p{30mm}p{30mm}}
\toprule
Peripheral & Mode & Clock source & Pins and function & Interrupt and transfers \\
\midrule
I2C1 & Fast mode, 400 kHz & Peripheral bus, source selected in RCC & D14 and D15 on the Arduino header. Confirm in MB1363 & Interrupt-driven. The upload is the case for DMA and Chapter 4 says why \\
GPIO input & Falling edge, external interrupt & Peripheral bus & The sensor's data-ready line, routed by the shield & EXTI, pre-emption priority 6 \\
GPIO output & Push-pull, sensor reset & Peripheral bus & The shield's reset line if it is routed & None \\
IWDG & Enabled after the upload completes & Low-speed internal oscillator & None & Reset source \\
TIM2 & Free running, 32-bit, 1 MHz tick & Peripheral bus & None & Read by the handler \\
GPIO output & Marker pin & Peripheral bus & One free Zio pin into a PPK2 digital input & None \\
USART3 & Asynchronous, 115200 8N1 & Peripheral bus & The console pins of chapter 1 & Polled \\
\bottomrule
\end{tabular}
\caption{Peripheral configuration. The watchdog row is the one that catches people: a watchdog armed before the upload with a window shorter than two seconds produces a board that resets forever and looks like a hardware fault.}
\end{table}

\subsection*{Wiring}

\diagram{c15_wiring}{The shield on the header, one marker wire to the instrument, and the bus selection that has to be read rather than assumed. Nothing else is placed by hand.}

If the bus selection turns out to be a jumper rather than a bridge, the faster
serial bus becomes available and the upload time falls by roughly the ratio of
the two bit rates, which would turn a two-second blind period into a fraction of
a second. That is the single most valuable configuration change available in
this chapter, which is why the question is worth an hour with UM3120 before any
code is written. Everything else in the chapter is unchanged by the answer: the
same driver, the same platform functions, the same policy.

Safety, as everywhere in this book: 3.3 V logic only, no 5 V into a GPIO, and
the instrument's digital inputs are logic inputs rather than an analogue
instrument. This is also the chapter where the absence of an oscilloscope costs
least. The event being measured lasts about two seconds, and the instrument's
current trace resolves it directly as a plateau with a beginning and an end.

\subsection*{Memory and timing budget}

\diagram{c15_mem}{Where the image lives before it is uploaded and where it lives afterwards. The microcontroller's flash holds a peripheral's firmware for the whole life of the product, and the sensor holds it only until the supply is interrupted.}

\begin{table}[H]\centering\small
\begin{tabular}{p{44mm}p{28mm}p{28mm}p{30mm}}
\toprule
Quantity & Budget & Measured & Margin \\
\midrule
Sensor firmware image in flash & 84 kB of 2 MB & not measured & not measured \\
Upload time at 400 kHz & 1.94 s computed & not measured & not measured \\
Result block per frame & about 1 kB & not measured & not measured \\
Frame period at 15 Hz & 66.7 ms & not measured & not measured \\
Reduction of 64 zones to a byte & 200 \textmu{}s & not measured & not measured \\
Watchdog window after arming & 1 s & not measured & not measured \\
Charge per decision & not budgeted & not measured & not measured \\
\bottomrule
\end{tabular}
\caption{The budget table. The upload figure is arithmetic from the bus rate and the image size and is marked computed rather than measured, which is a different claim. The result block figure depends on which outputs the driver's configuration header enables and is read there rather than guessed.}
\end{table}

Two of those rows deserve a sentence. The flash row is 4.2 percent of the
part's program memory spent on a peripheral's firmware, which on a 2 megabyte
part is affordable and on a smaller one would not be, and that is a real
architectural consequence of choosing this sensor. The watchdog row is a
constraint rather than a measurement: the window is chosen after the reduction
time is known, not before, and the watchdog is armed after the upload rather
than across it.

\subsection*{Firmware design (UML)}

\diagram{c15_uml}{The node's state machine, including the hysteresis that stops the output changing state whenever a target sits on a threshold. The re-upload path exists because a supply disturbance can reset the sensor without resetting the board.}

The machine has four operating states and one that most implementations omit.
From reset the node is cold: nothing has been uploaded and no frame is
meaningful. Uploading is a single long operation with the watchdog not yet
armed. Ranging is the steady state, and inside it the reported decision moves
between clear, approaching and present according to a policy with two distinct
thresholds and a confirmation count. The state most implementations omit is the
one entered when the sensor stops answering or answers with an implausible
identity register: the node returns to uploading rather than reporting stale
distances, and the status byte says so while it does.

Hysteresis is not an optional refinement here. A single threshold on a distance
that hovers near it produces an output that alternates at the frame rate, which
is worse than useless because it looks like activity. Two thresholds and a
confirmation count turn that into a decision that changes when something
actually changed. The numbers in the figure, entering at 800 mm and leaving at
950 mm with three agreeing frames out of five, are a starting point to be
measured against a real target rather than constants to be trusted.

The confirmation is counted in frames rather than in milliseconds, and that is
deliberate. A count of frames degrades correctly: if the sensor is running more
slowly than configured because the chosen resolution costs more than expected, a
three of five rule still means three independent pieces of evidence, where a two
hundred millisecond rule might mean one. It also makes the policy testable on
the host, because a recording of frames can be replayed through the same C code
with no clock involved at all. The cost is that decision latency becomes a
function of the frame rate, which then has to be stated in the interface
documentation rather than discovered by whoever integrates the node.

\subsection*{Data flow (ASCII)}

\begin{asciiart}
  power-on
     |
     v
  +----------------------------------------+          +--------------------------+
  | STM32H7A3 flash                        |          | VL53L8CX                 |
  |   vl53l8cx_fw[86016] const, read only  |--I2C1--->| volatile program memory  |
  |   application text and data            |  1.94 s  |   no non-volatile store  |
  +----------------------------------------+          +------------+-------------+
                                                                   | data ready
  +----------------------------------------+                       v
  | frame handler                          |<------ I2C1 ----- result block, ~1 kB
  |   validity:   status code and sigma    |
  |   reduce:     nearest, occupancy, col  |
  |   hysteresis: 800 / 950 mm, 3 of 5     |
  +-------------------+--------------------+
                      | one byte
                      v
  +----------------------------------------+          +--------------------------+
  | USART3 framed link                     |--------->| host: state and count    |
  +----------------------------------------+          +--------------------------+
\end{asciiart}

\subsection*{Repository layout}

\begin{plaincode}
nucleo-h7a3-53l8a1-zones/
  CMakeLists.txt
  cmake/arm-none-eabi.cmake
  third_party/vl53l8cx/        # from the redistributable republication
    vl53l8cx_api.[ch]          # wrapper stripped, headers untouched
    vl53l8cx_buffers.h         # the 84 kB image, as a const array
    LICENSE                    # BSD-3-Clause, copied verbatim
  src/platform_i2c1.c          # the four functions the driver asks for
  src/bringup.c                # identity check, upload, arm the watchdog
  src/zones.c                  # validity, reduction, hysteresis
  src/status_byte.c            # the one byte that leaves the board
  src/timebase.c
  src/main.c
  tools/status_decode.py       # the host twin of the byte
  tools/frame_dump.py          # diagnostic mode only, 64 values
  docs/PROVENANCE.md           # upstream commit, licence, date copied
  README.md
\end{plaincode}

\subsection*{Steps}

\begin{steps}
\item \textbf{Read UM3120 and settle the bus question before anything else.}
Find the bus selection, whether it is a jumper or a solder bridge, which signals
the faster bus would use, where the data-ready line lands on the header and
whether the reset line is routed. Write the answers into the README as a table
with the manual named as the source. If the selection is a bridge, write one
sentence saying the chapter is I2C only and why, and move on without regret.

\item \textbf{Take the driver from the path that can be republished.} Both
copies are the same code. Only one of them can be shown in a portfolio
repository with a pinned version.
\begin{shellcode}
git clone https://github.com/stm32duino/VL53L8CX vendor-src
# keep: the driver core, the platform header, the firmware image array
# remove: the wrapper for the other environment and its examples
# keep untouched: every copyright header and the LICENSE file
\end{shellcode}
Record the upstream commit hash, the licence and the date in
\texttt{docs/PROVENANCE.md} in the same sitting. A provenance file written later
is a provenance file written from memory.

\item \textbf{Write the four platform functions.} The driver asks for a small,
well-defined amount from you, which is what makes it portable and what makes
this step short.
\begin{ccode}
/* Multi-byte register access, a delay, and a timebase. Nothing else. */
uint8_t VL53L8CX_RdMulti(VL53L8CX_Platform *p, uint16_t reg,
                         uint8_t *val, uint32_t size);
uint8_t VL53L8CX_WrMulti(VL53L8CX_Platform *p, uint16_t reg,
                         uint8_t *val, uint32_t size);
uint8_t VL53L8CX_WaitMs(VL53L8CX_Platform *p, uint32_t ms);
uint8_t VL53L8CX_GetTickCount(VL53L8CX_Platform *p, uint32_t *ms);
\end{ccode}
The register address is sixteen bits wide rather than eight, which is the first
thing to get right in the I2C driver and the first thing to check if the
identity register reads as zero.

\item \textbf{Place the image and prove the linker put it where you meant.} The
array is const and belongs in flash. If it ends up in initialised data it costs
84 kilobytes of flash and 84 kilobytes of SRAM, and the startup code copies it
on every reset for no reason.
\begin{shellcode}
arm-none-eabi-size -A build-fw/zones.elf
arm-none-eabi-nm --size-sort -S build-fw/zones.elf | tail -5
\end{shellcode}
Expect the largest symbol in the image to be the firmware array, in a read-only
section. If it is not, the const qualifier or the section attribute is wrong.

\item \textbf{Upload, and measure what the upload costs.} Pulse the marker pin
before and after, and capture the current trace on the instrument. The upload
appears as a plateau with a clear beginning and end, and the instrument's
digital input puts the marker pulses on the same time axis as the current. Read
the elapsed time from the trace and the elapsed cycles from the DWT counter, and
put both in the README with the instruments named.
\begin{ccode}
marker_set(1);
const uint32_t t0 = timebase_ticks();
status = vl53l8cx_init(&dev);          /* this is where the 84 kB goes */
const uint32_t t1 = timebase_ticks();
marker_set(0);
printf("upload %lu ticks at 1 MHz\r\n", (unsigned long)(t1 - t0));
\end{ccode}
Compare the measured figure against the 1.94 seconds computed above. If the
measured value is substantially larger, the difference is bus overhead per
transaction, clock stretching by the sensor, or a per-chunk delay in the driver,
and finding out which is a better exercise than accepting the number.

\item \textbf{Arm the watchdog after the upload and not before.} The independent
watchdog on this part runs from its own oscillator and does not care what the
core is doing. A window shorter than the upload produces a board that resets
during boot, forever, with no message, which reads as a hardware fault and is
not one.
\begin{ccode}
/* Order matters. Chapter 12 builds the watchdog design in full; this
   chapter only has to avoid the one mistake that this sensor invites. */
bringup_upload_firmware();      /* seconds, watchdog not yet running */
iwdg_start(IWDG_WINDOW_MS);     /* only now */
\end{ccode}
The alternative, kicking the watchdog from inside the upload loop, is defensible
and is also how a watchdog quietly stops protecting anything. If you take it,
say in a comment which specific failure it still catches.

\item \textbf{Configure eight by eight and read one frame.} Choose the
resolution, the ranging frequency and the target order, then read a single frame
and print all sixty-four values once, as a diagnostic, to prove the geometry is
what you think it is. Hold a flat surface at a known distance, at an angle, and
check that the grid tilts the way the geometry says it should. Then never print
sixty-four values again outside diagnostic mode.

\item \textbf{Decide which zones to believe.} Every zone carries a target status
and a sigma, and a zone that reports a distance with a status the driver's
header describes as unreliable is not data. This filter is the difference
between a reduction that works in a corridor and one that works only on a bench.
\begin{ccode}
static bool zone_valid(const frame_t *f, unsigned z)
{
    if (!status_is_trusted(f->status[z])) return false;   /* from the header */
    if (f->sigma_mm[z] > SIGMA_LIMIT_MM)  return false;   /* sensor's own doubt */
    if (f->dist_mm[z] < DIST_MIN_MM)      return false;   /* cover glass */
    return true;
}
\end{ccode}

\item \textbf{Reduce, then apply hysteresis, then pack the byte.} The reduction
is deliberately simple and deliberately explainable, which is the argument of
this book's intelligence chapters in miniature.
\begin{ccode}
typedef struct {
    uint16_t nearest_mm;     /* smallest valid distance in the frame  */
    uint8_t  n_valid;        /* how many of the 64 zones were believed */
    uint8_t  column;         /* occupancy centroid, 0..7               */
} reduced_t;

/* Two thresholds, never one. Three agreeing frames out of five, which
   at 15 Hz is a decision in about 200 ms. */
#define ENTER_MM  800u
#define LEAVE_MM  950u
#define CONFIRM_N 3u
#define CONFIRM_M 5u

uint8_t status_byte(const reduced_t *r, zone_state_t *st)
{
    zone_state_update(st, r->nearest_mm, ENTER_MM, LEAVE_MM,
                      CONFIRM_N, CONFIRM_M);
    return (uint8_t)((st->ready   ? 0x80u : 0u)
                   | ((st->state  & 0x03u) << 5)
                   | ((occupancy_bucket(r->n_valid) & 0x07u) << 2)
                   |  (st->confidence & 0x03u));
}
\end{ccode}

\item \textbf{Write the host twin and prove the link carries a byte.} The
decoder is ten lines and it is the specification of the interface. Then run the
node for an hour against a moving target and count state changes. A policy that
produces hundreds of changes in an hour of ordinary activity has thresholds that
are too close together, and the fix is a measurement rather than an opinion.
\begin{pycode}
def decode(b: int) -> dict:
    return {
        "ready":      bool(b & 0x80),
        "state":      ["clear", "approaching", "present", "leaving"][(b >> 5) & 3],
        "occupancy":  (b >> 2) & 7,
        "confidence": b & 3,
    }
\end{pycode}
\end{steps}

\subsection*{Build, flash and debug}

\diagram{c15_timing}{Power-on to first decision. The blind period is the upload, the watchdog is armed only after it, and the steady state that follows is short frames with a short reduction inside each one.}

\begin{shellcode}
cmake -B build-fw -G Ninja -DCMAKE_TOOLCHAIN_FILE=cmake/arm-none-eabi.cmake
cmake --build build-fw -j
cp build-fw/zones.bin "$PROBE_DISK"/   # onto the probe disk
python tools/status_decode.py --port /dev/ttyACM0
\end{shellcode}

Recovering a board that will not answer is Chapter 1's material: the probe
presents a disk and a known-good binary can be copied onto it with no tools. The
failure specific to this chapter is a board that resets every two seconds, and
the cause is almost always the watchdog ordering above.

\begin{note}{When the sensor answers once and then stops}
In order of likelihood: the firmware was never uploaded, because the identity
register was read before the upload and looked plausible; the register address
was treated as eight bits when it is sixteen; the supply dipped far enough to
reset the sensor and not the board, which is exactly the case the re-upload
state exists for; the data-ready line is on a different header pin than assumed
and the firmware is waiting on a pin nobody drives; and the frame rate was set
above what the configured resolution supports, so frames are reported as not
ready rather than as an error. Read the identity register again after the upload
rather than before, which distinguishes the first case from the rest in one
line.
\end{note}

\subsection*{Verification and acceptance criteria}

\begin{itemize}
\item The host receives a status byte and never sixty-four values, except in an
explicitly named diagnostic mode that is off by default. This is the chapter's
acceptance criterion and it is checked by reading the link, not the source.
\item The upload time is a measured number in the README, with the instrument
named, next to the computed figure it is being compared against.
\item The firmware image appears in a read-only section in the size output, and
the largest symbol in the image is that array.
\item A deliberate supply interruption to the sensor alone, with the board still
running, produces a re-upload and a status byte with the ready bit clear during
it, and never a stale distance presented as current.
\item A target held steady at the threshold distance for one minute produces at
most one state change. A target moved across the threshold and back produces
exactly two.
\item The watchdog is armed after the upload, proven by a test build whose
window is deliberately shorter than the upload and which still boots.
\item The provenance document names the upstream commit, the licence and the
date, and the build fails if the vendored tree has been modified without the
document being updated.
\end{itemize}

\subsection*{Variants}

\begin{table}[H]\centering\small
\begin{tabular}{p{24mm}p{26mm}p{44mm}p{22mm}p{20mm}}
\toprule
Axis & Variant & What changes & Cost & Built in full in \\
\midrule
Execution model & DMA for the upload & An 84 kilobyte transfer is the strongest argument for the transfer engine in this book, and it frees the core for the whole blind period & Cache maintenance on a buffer in main memory & Chapter 4 and Chapter 19 \\
Peripheral substitution & The faster serial bus & The same driver and the same platform functions, with the upload time falling by the ratio of the bit rates. Available only if the selection is a jumper & Pins, and the answer to the confirm-list question & Chapter 17 \\
Time and safety & Window watchdog instead of the independent one & A watchdog that also objects to being kicked too early, which catches a different set of faults & More thought about the steady-state period & Chapter 12 \\
Intelligence and reach & A classifier on the frame & The reduction becomes a learned model over sixty-four zones rather than a threshold policy, which is a better fit for gesture than for occupancy & Explainability, and a training set & Chapter 16 \\
Operating system & An RTOS task set & The upload becomes a task that blocks, rather than a function that occupies the whole boot & A scheduler to justify and a stack to size & Chapter 20 \\
Execution model & Polled data ready & Poll the status instead of watching the line. Simpler, and it wastes the bus it is polling & Bus bandwidth and charge & Chapter 3 \\
\bottomrule
\end{tabular}
\caption{Variants for chapter 15. Nothing here is built in full in this chapter; its own contribution is the upload discipline and the reduction policy, and each variant names the chapter that owns it.}
\end{table}

\subsection*{Pitfalls}

\begin{itemize}
\item Arming the watchdog before the upload. The board resets forever and the
symptom looks like hardware.
\item Treating the sixteen-bit register address as eight bits. Everything reads
as zero and the sensor looks absent.
\item Letting the firmware image land in initialised data. It costs the flash
twice and the startup code copies 84 kilobytes on every reset.
\item Reading the identity register before the upload, deciding the sensor is
present, and never noticing that the upload did not happen.
\item Sending sixty-four values to the host because the driver makes them
convenient. The interface is the byte.
\item One threshold instead of two. The output then changes at the frame rate
whenever a target sits near the boundary.
\item Believing every zone. A zone with an untrusted status code and a wide
sigma is not a distance, and including it moves the nearest-distance result
without any target moving.
\item Taking the driver from behind the click-through and publishing it. Same
code, different licence, and the repository is the thing that cannot be
published.
\item Assuming a per-frame time from the configured frequency. The configured
frequency is what the sensor attempts at the chosen resolution, and the confirm
list has the resolution and rate combinations on it for a reason.
\end{itemize}

\subsection*{Best practices applied}

\begin{itemize}
\item The dependency's licence decided which copy of identical code is used, and
that decision is documented in the repository rather than being invisible.
\item Provenance is written on the day the code is copied, with an upstream
commit hash, and the build checks that the vendored tree and the document agree.
\item A computed figure and a measured figure are labelled differently and
printed next to each other, so that the reader can see the model and the
measurement disagree.
\item The interface is the smallest thing that carries the decision, and the
diagnostic mode that violates that is named, defaulted off, and never used to
justify the design.
\item A failure mode that most implementations ignore, the sensor losing its
firmware without the board noticing, is a state in the machine rather than a
paragraph in a README.
\end{itemize}

\subsection*{Stretch goals}

\begin{itemize}
\item Compress the firmware image in flash and decompress it during the upload.
The transfer is bus-bound rather than core-bound, so the decompression may be
free in wall-clock terms, and measuring whether it is would be a genuinely new
result for this part.
\item Keep the sensor powered across a microcontroller reset and skip the
upload, which turns a two-second boot into a fast one at the cost of a more
complicated power design and a harder verification story.
\item Use the occupancy centroid to report direction of travel rather than
presence, which costs nothing extra on the bus and turns a status byte into
something a lighting controller can act on.
\item Run the reduction on frames captured to the host and replay them through
the same C code compiled for the host, which is Chapter 12's testing approach
applied to a policy that would otherwise only ever be tested by waving a hand.
\end{itemize}

\subsection*{Roadmap and next steps}

Chapter 16 takes the decision built here and moves it: the same classification in
the sensor, on the microcontroller and on the host, with latency, memory and
charge measured for each, which is the measurement that decides whether this
chapter's placement argument is right.

For going wider, the community embedded engineering roadmap separates firmware
from embedded Linux from hardware and says which to weight at which stage, and
its accompanying essay on bridging the gap is the short version of the same
argument. The free open textbook that grew out of a university course is the
closest thing to a canonical curriculum for the sensor-intelligence direction,
and the free-to-audit course on embedded machine learning covers the same ground
with exercises. For the state-machine discipline this chapter leans on, the free
course that walks one problem from a polled loop to an event-driven architecture
across fifty-six lessons is the best companion in this volume, with the standing
caution that its code licence makes it something to recommend and never to
quote. Appendix H lists all of them with their licences.

\subsection*{Portfolio evidence}

\begin{itemize}
\item A public repository containing a working port of a permissively licensed
driver, with a provenance document naming the upstream commit, the licence and
the date, and a README that explains why this copy and not the other one.
\item An instrument capture showing the upload as a current plateau with the
marker pulses on the same time axis, annotated with both the computed and the
measured duration.
\item A one-hour recording of status bytes against ordinary activity, with the
state-change count, which is the evidence that the hysteresis was chosen rather
than guessed.
\item The size output showing where the 84 kilobyte array landed.
\end{itemize}

\subsection*{Sources}

Normative references:
\begin{itemize}
\item UM3120, the user manual for the X-NUCLEO-53L8A1, for the bus selection,
the interrupt routing, the reset line and the header mapping.
\item The VL53L8CX datasheet and its user manual, for the resolution and rate
combinations, the target status codes and the result block contents.
\item Reference manual RM0455 for I2C1 and the independent watchdog on this
part. Not RM0433, which documents its better-known sibling.
\item The board user manual for MB1363, for which Arduino header pins reach
which microcontroller pins.
\end{itemize}

Reusable implementations:
\begin{itemize}
\item The vendor's Arduino organisation, which republishes the driver under
BSD-3-Clause. This is the copy that can be ported and published.\newline
\url{https://github.com/stm32duino/VL53L8CX}
\item An independent write-up of the processor lineage inside these sensors,
which explains why the image is this size and why it is volatile.\newline
\url{https://the6p4c.github.io/2022/06/13/stxp70-sthorm-p2012.html}
\item The vendor's register-level driver collection, BSD-3-Clause, which is the
model for the two-function porting shim used across the shield chapters.\newline
\url{https://github.com/STMicroelectronics/STMems_Standard_C_drivers}
\item Memfault's article on watchdog design, for the two-layer approach this
chapter's ordering rule is a special case of.\newline
\url{https://interrupt.memfault.com/blog/firmware-watchdog-best-practices}
\item The Python interface to the power instrument, GPL-2.0. Install it from its
package index; do not fork it into a portfolio repository.\newline
\url{https://github.com/IRNAS/ppk2-api-python}
\end{itemize}
   % Eight by eight time of flight, decided on the MCU
\project{16}{Where the classifier runs: sensor, MCU or host}{NUCLEO-H7A3ZI-Q with X-NUCLEO-IKS4A1 and a Raspberry Pi}{In-sensor engine against on-MCU model against host, measured}

\begin{keyfacts}
\item[Board] NUCLEO-H7A3ZI-Q with X-NUCLEO-IKS4A1, and a Raspberry Pi as the host. One shield only
\item[Peripherals] I2C1 to the shield, the FIFO watermark line, USART3 as the framed link to the host, TIM2 and the cycle counter as instruments, one marker pin
\item[Toolchain] arm-none-eabi-gcc with CMake, a permissive decision-tree generator, the signal-processing library, and the in-sensor toolchain for the third leg
\item[Operating system] Bare metal
\item[Difficulty] 5 of 5
\item[Effort] 6 evenings of about four hours
\item[Deliverable] One labelled dataset, one feature definition, one trained model, and the same classification running in three places with latency, memory and charge per decision measured for each
\end{keyfacts}

\subsection*{Why this project}

Every discussion of intelligence on small systems eventually reaches the
question of where the decision should be computed, and almost every answer is an
assertion. The three candidate placements are the sensor, which now contains a
programmable core of its own; the microcontroller, which has 1312 kB of declared
static memory and a signal-processing instruction set; and a host at the other
end of a link. Each
placement has an obvious argument in its favour and no published measurement
that compares all three on one task, on one set of hardware, with a stated
method. That gap is the reason this chapter exists.

The method is deliberately narrow so that the comparison means something. One
labelled dataset, collected once. One feature definition, written once. One
model, trained once. Then the same model is compiled three ways and placed in
three locations, and three quantities are measured for each: latency from the
end of a data window to a usable class, memory occupied, and charge consumed per
decision. Anything that differs between the legs other than the placement is a
confound, and the design of this chapter is mostly the work of removing
confounds.

There is one recent preprint that measures two of the three placements and
reaches a counter-intuitive conclusion: the host wins, because the cost is
dominated by the bus read rather than by the classification. Reading a register
over this bus is three transactions, not one, and if the decision is made in the
sensor the microcontroller still has to wake up and perform them. That is a
genuine insight and it is worth testing. It is also unreviewed, its authors are
unaffiliated on the page, and its host is a far heavier machine than anything on
this bench. Treat it as the hypothesis, not the answer. A host that is itself a
microcontroller, with a wake-up measured in microseconds rather than in tens of
milliseconds, could plausibly reverse the ranking, and naming that is part of
this chapter's contribution.

\begin{note}{What this chapter has and has not produced}
The result does not exist yet. What follows is the method that produces it, and
every measured column in every table in this chapter stays blank until an
instrument has filled it. That is the standing rule of this book and it applies
with particular force here, because the whole value of the chapter is that its
numbers are new. A chapter that filled them in from plausible expectation would
be worth less than no chapter at all.
\end{note}

\subsection*{Prior art and what to reuse}

\begin{table}[H]\centering\small
\begin{tabular}{p{34mm}p{46mm}p{38mm}p{20mm}}
\toprule
Source & What it gives & What it does not & Licence \\
\midrule
The in-sensor processor examples & Working programs for the programmable core in the shield's inertial part, including prebuilt outputs so a demonstration runs before the toolchain is installed & It declares no licence at its root, so files are checked one at a time and nothing is vendored on faith & None at root \\
The ready-made classifier configurations for the in-sensor engine & Complete worked configurations for the other in-sensor engine, and the clearest statement of what that engine can express & It replaced a predecessor that was withdrawn in 2025, so a great deal of existing writing points at addresses that no longer resolve & BSD-3-Clause \\
The in-sensor state machine collection & The same for the state-machine engine, which is the right tool when the decision is a sequence rather than a classification & Nothing about features or training & BSD-3-Clause \\
A permissive decision-tree and anomaly generator for small targets & C output for trees, forests and naive Bayes from about 2 kB of flash, which is the model this chapter compiles three ways & No feature extraction and no data collection & MIT \\
The signal-processing library & The transforms and statistics the feature stage is built from, with a Python wrapper whose API mirrors the C API so both sides run the same test vectors & No model and no training & Apache-2.0 \\
The permissive neural runtime & A clean path for the neural comparison leg, with optimised kernels that do help on this core through its signal-processing instructions & It is the comparison, not the baseline, and its speed-up figures in vendor material belong to later cores & Apache-2.0 \\
One recent preprint on placement & The hypothesis, and a measurement of two of the three placements & Review, affiliation, and a host anywhere near this size & Reading only \\
\bottomrule
\end{tabular}
\caption{Prior art for chapter 16. The repository with no licence at its root is the one most likely to end up copied by accident, because it is the one that ships ready-to-run outputs.}
\end{table}

What is left to write is the comparison itself, which is the part nobody has
published, plus the two pieces of discipline that make it a comparison rather
than three demonstrations: a feature definition that produces bit-identical
values in all three implementations, and a measurement protocol that states what
is inside and outside each boundary.

\begin{note}{Writing that points at addresses which no longer resolve}
Both in-sensor repositories replaced predecessors that were discontinued during
2025. A large amount of tutorial material, including material that is otherwise
good, still links to the old locations. When a search result for this engine
looks authoritative and its links do not resolve, that is the reason, and the
current repositories are named in the sources at the end of this chapter.
\end{note}

\subsection*{Parts from the inventory}

\begin{table}[H]\centering\small
\begin{tabular}{p{40mm}p{66mm}p{40mm}}
\toprule
Part & Role & Interface \\
\midrule
NUCLEO-H7A3ZI-Q & The middle placement, and the bus master and link driver for the other two & Micro USB to the host \\
X-NUCLEO-IKS4A1 & The sensor, including the inertial part with the programmable core that hosts the first placement. The only shield on the board & I2C1 plus the watermark line \\
Raspberry Pi 4 & The third placement, at the far end of a framed serial link & USB serial, or the header serial pins \\
Renkforce USB/TTL cable & The link between the board and the host when the header pins are used. The red 5 V lead is disconnected & 3.3 V serial \\
nRF-PPK2 & Charge per decision for the node. It measures up to 1 A and the host is outside its range and outside this measurement & IDD jumper and one digital input \\
Host PC & Data collection, training, and the build for all three legs & Windows with the Linux subsystem \\
\bottomrule
\end{tabular}
\caption{Inventory items used in chapter 16. Nothing is bought. The single-board computer is the host under test and is deliberately not the machine that trains the model, so that the training environment is not accidentally part of the comparison.}
\end{table}

\subsection*{System architecture}

\diagram{c16_arch}{The same model in three places, and what crosses each boundary. The boundaries are where the cost is, which is the preprint's insight and the thing this chapter is built to measure.}

Read that figure by the arrows rather than by the boxes. In the first leg a
class index crosses the bus. In the second a window of samples crosses the bus
and nothing crosses the link. In the third a window of samples crosses the bus
and then crosses the link as well. The work performed is nearly identical in all
three; what differs is how much data moves and how many times something has to
wake up to move it.

\subsection*{Peripheral configuration}

\begin{table}[H]\centering\small
\begin{tabular}{p{22mm}p{26mm}p{24mm}p{30mm}p{30mm}}
\toprule
Peripheral & Mode & Clock source & Pins and function & Interrupt and transfers \\
\midrule
I2C1 & Fast mode, 400 kHz & Peripheral bus, source selected in RCC & D14 and D15 on the Arduino header. Confirm in MB1363 & Interrupt-driven. DMA is a variant and Chapter 4 owns it \\
GPIO input & Rising edge, external interrupt & Peripheral bus & FIFO watermark line, or the in-sensor result line in the first leg & EXTI, pre-emption priority 6 \\
USART3 & Asynchronous, 115200 8N1 & Peripheral bus & The framed link to the single-board computer & Polled here, DMA in Chapter 4 \\
TIM2 & Free running, 32-bit, 1 MHz & Peripheral bus & None & The timestamp on every window \\
DWT cycle counter & Free running & Core clock & None & The instrument for the compute stages \\
GPIO output & Marker pin, one pulse per stage & Peripheral bus & One free Zio pin into a PPK2 digital input & None \\
\bottomrule
\end{tabular}
\caption{Peripheral configuration. The marker pin is the instrument that makes the three legs comparable: it separates wake, read, compute and send inside one current trace, which is the method Chapter 10 builds and this chapter applies.}
\end{table}

\subsection*{Wiring}

\diagram{c16_wiring}{The shield on the board, the instrument across the board's own supply, and the host at the end of a serial link. What the instrument can and cannot see is drawn deliberately, because the boundary of the measurement is part of the result.}

The safety points are the ones this bench always has, plus one that is specific
here. Everything is 3.3 V logic. The red 5 V lead of the serial cable is
disconnected, because the board and the single-board computer both have their
own supplies and joining them is how a ground loop or a back-powered board
happens. The instrument measures up to 1 A, which is ample for the board and
shield and nowhere near enough for the single-board computer, so the host's own
energy is not measured in this chapter and is named as excluded rather than
quietly omitted.

\subsection*{Memory and timing budget}

\diagram{c16_mem}{Footprint in the three placements, drawn to make one point: on this part memory is not the binding constraint. The reference anomaly model from the standard suite is about 270 kB and fits into this part several times over. In the sensor it does not fit at all.}

\begin{table}[H]\centering\small
\begin{tabular}{p{44mm}p{28mm}p{28mm}p{30mm}}
\toprule
Quantity & Budget & Measured & Margin \\
\midrule
Feature stage, flash on the MCU & 6 kB & not measured & not measured \\
Decision tree, flash on the MCU & 4 kB & not measured & not measured \\
Neural comparison model, flash & 64 kB & not measured & not measured \\
Window buffer, SRAM & 8 kB & not measured & not measured \\
In-sensor program memory & read from the application note & not measured & not measured \\
Latency, sensor placement & 5 ms & not measured & not measured \\
Latency, MCU placement & 5 ms & not measured & not measured \\
Latency, host placement & 30 ms & not measured & not measured \\
Charge per decision, node & not budgeted & not measured & not measured \\
Charge per decision, host & outside the instrument & not measured & excluded by design \\
\bottomrule
\end{tabular}
\caption{The budget table. The last row is the honest statement of what this bench cannot do, and it is also the row that stops this chapter from claiming to have settled the preprint's question in the general case.}
\end{table}

\begin{note}{Memory is almost never the binding constraint on this part}
This part declares 1312 kB of static memory across five regions, plus a tightly
coupled instruction memory that neither of the machine-readable sources for this
die states, and 2048 kB of flash. The commonly repeated figure of about 1.4
megabytes is reached only by adding that instruction memory in, which is why
this book states the declared total and names the addition separately. Either
way the part sits at the large end of the envelope that small-system machine
learning literature assumes. Most of that literature was written for parts an order of
magnitude smaller, and its central anxiety, fitting the model, does not apply
here. What binds instead is throughput and energy: one core, no accelerator, and
no vector extension. This core does have the signal-processing instruction
extension, so the optimised neural kernels genuinely help, through that path
rather than through vectors. Every impressive vendor speed-up figure for those
kernels belongs to a later core than this one, and a chapter that implied
otherwise would be repeating the error this whole volume is about.
\end{note}

\subsection*{Firmware design (UML)}

\diagram{c16_uml}{One feature definition, one trained model, three implementations. The dashed arrows are claims of conformance; the gate underneath is what turns a claim into evidence, and it compares outputs rather than source.}

The design is a compiler problem more than a firmware problem. The feature
definition exists once, as a specification and as a set of test vectors. Three
implementations claim to satisfy it: one in the in-sensor toolchain, one in C
against the signal-processing library, and one in Python on the host. The tree
exists once as trained coefficients and is emitted three times. Before any
measurement is taken, all three implementations are fed the same recorded window
and must produce the same class, on every window in the test set. Until that
holds, the three legs are three different experiments and their numbers cannot
be compared.

That equivalence is easier to state than to obtain, and the usual cause of
failure is numeric rather than logical. The in-sensor leg may work in a fixed
point format the other two do not use. The host may compute in double precision
where the microcontroller computes in single. The fix is to define the feature
values in the format the most constrained implementation can produce and to
require the other two to match it, rather than the other way around, which is
the opposite of what is comfortable and the only version that works.

\subsection*{Data flow (ASCII)}

\begin{asciiart}
  collection, once               training, once                 deployment, three ways
  +------------------------+     +------------------------+     +------------------------+
  | shield to board        |     | features in Python,    |     | A: features and tree   |
  |   raw labelled windows +---->|   mirroring the C API  |     |    in the sensor core  |
  |   stored on the host   |     | tree fitted once,      +---->| B: features and tree   |
  |                        |     | exported three ways    |     |    on the MCU          |
  |                        |     |                        |     | C: window over the     |
  +------------------------+     +------------------------+     |    link, tree on host  |
                                                                +------------------------+
                                              |
                                              v
  +--------------------------------------------------------------------------------------+
  | equivalence gate: all three produce the same features and the                        |
  |   same class on every window of the held-out set                                     |
  +--------------------------------------------------------------------------------------+
                                              |
                                              v
  +--------------------------------------------------------------------------------------+
  | measurement: DWT cycles, TIM2 timestamps, marker pin into the                        |
  |   instrument. Latency as median, p99 and maximum; flash and                          |
  |   SRAM; charge per decision. The host's own energy is outside                        |
  |   the instrument and is named as outside the comparison.                             |
  +--------------------------------------------------------------------------------------+
\end{asciiart}

\subsection*{Repository layout}

\begin{plaincode}
nucleo-h7a3-classifier-placement/
  CMakeLists.txt
  dataset/
    raw/                        # labelled windows, one file per session
    protocol.md                 # how each class was produced, in detail
  features/
    features.md                 # the definition, and the fixed-point format
    features.c  features.h      # the C implementation, CMSIS-DSP
    features.py                 # the Python implementation, same test vectors
    vectors/                    # inputs and expected outputs, shared by all
  model/
    train.py                    # fit once, export three ways
    tree_mcu.c                  # generated, permissive generator
    tree_host.py
    ispu/                       # the in-sensor program and its build
  src/leg_a_sensor.c            # read a class index, nothing else
  src/leg_b_mcu.c               # window, features, tree, on the board
  src/leg_c_host.c              # window, frame, send, wait for the class
  src/measure.c                 # DWT, TIM2, marker pin
  tools/equivalence.py          # the gate: three implementations, one answer
  tools/collect.py  tools/report.py
  docs/PROVENANCE.md            # every dependency, its licence, the date
  README.md
\end{plaincode}

\subsection*{Steps}

\begin{steps}
\item \textbf{Choose a task the bench can actually produce and write the
protocol down.} Four classes of movement of the board and shield on a table:
still, sliding, tapped, and lifted and rotated. Write
\texttt{dataset/protocol.md} describing exactly how each class is produced,
including how long each session runs and who produced it. Hand-generated data
limits what can be claimed about accuracy, and the protocol is what makes the
limitation visible rather than hidden. The accuracy of the classifier is not
this chapter's result; the placement comparison is.

\item \textbf{Collect the dataset once, and never collect it again.} Stream raw
windows over the link with timestamps and labels, and store them on the host.
Every later leg is fed from these files, including the legs that run on the
board, because a leg fed from live movement is being fed different data from the
others and the comparison is void.
\begin{shellcode}
python tools/collect.py --port /dev/ttyACM0 --label still --seconds 120
python tools/collect.py --port /dev/ttyACM0 --label sliding --seconds 120
python tools/collect.py --port /dev/ttyACM0 --label tapped --seconds 120
python tools/collect.py --port /dev/ttyACM0 --label rotated --seconds 120
\end{shellcode}

\item \textbf{Define the features in the most constrained format first.} Decide
the window length, the overlap and the feature list, and decide the numeric
format that the in-sensor implementation can produce. Then write that down as
the definition and make the other two implementations match it.
\begin{ccode}
/* One window: 128 samples of three axes at 104 Hz, about 1.23 s.
   Features are computed in Q15 because that is what the most
   constrained of the three implementations can produce. */
#define WIN_N      128u
#define WIN_AXES     3u
typedef struct {
    q15_t mean[WIN_AXES];
    q15_t rms[WIN_AXES];
    q15_t p2p[WIN_AXES];          /* peak to peak */
    q15_t band[WIN_AXES][4];      /* energy in four bands, from the FFT */
} features_t;
\end{ccode}

\item \textbf{Build the equivalence gate before building any leg.} This is the
step that turns three demonstrations into one experiment, and doing it later
never works because by then each implementation has its own quiet assumption.
\begin{pycode}
# Every implementation reads the same vectors and must agree exactly.
for name, impl in (("c", run_c_features), ("py", run_py_features),
                   ("ispu", run_ispu_features)):
    for vec in load_vectors("features/vectors"):
        assert impl(vec.input) == vec.expected, (name, vec.id)
\end{pycode}

\item \textbf{Train once and export three times.} The tree is fitted on the host
with the Python feature implementation, and the same fitted coefficients are
emitted as C for the microcontroller, as a program for the in-sensor core, and
as a Python object for the host leg. A tree with a bounded depth is chosen
deliberately, because bounded depth is bounded latency, and bounded latency is
half of what makes a classifier acceptable in a control loop.

\item \textbf{Leg A, in the sensor.} Load the in-sensor program, configure the
part to run it on its own data, and have the microcontroller do nothing but wait
for the line and read the result. Start from the prebuilt outputs in the
examples repository, which run before the toolchain is installed and prove the
path works, then replace them with your own. Audit the licence of every file you
keep, because the repository declares none at its root.
\begin{ccode}
void EXTI_result_handler(void)          /* leg A: the whole of the work */
{
    marker_pulse(MARKER_WAKE);
    const uint32_t t = timebase_ticks();
    uint8_t klass;
    i2c1_read_reg(IMU_ADDR, ISPU_RESULT_REG, &klass, 1);  /* three transactions */
    decision_post(t, klass);
    marker_pulse(MARKER_DONE);
}
\end{ccode}
Note what the comment records. A one-byte register read is an address write, a
repeated start and a data read. That is the cost the preprint says dominates,
and here it is, in the leg that was supposed to avoid work.

\item \textbf{Leg B, on the microcontroller.} Drain the FIFO on the watermark
line into a window, compute the features with the signal-processing library, and
evaluate the generated tree. Measure the two compute stages separately, because
the feature stage and the model are different sizes of problem and a single
number hides which one matters.
\begin{ccode}
const uint32_t c0 = DWT->CYCCNT;
features_t f;  features_compute(&win, &f);
const uint32_t c1 = DWT->CYCCNT;
const uint8_t k = tree_predict((const q15_t *) &f);
const uint32_t c2 = DWT->CYCCNT;
printf("feat %lu cyc, tree %lu cyc\r\n",
       (unsigned long)(c1 - c0), (unsigned long)(c2 - c1));
\end{ccode}

\item \textbf{Leg C, on the host.} Send the window over the framed link from
Chapter 5 and wait for the class to come back. Measure the round trip from the
board's side, which is the only side that matters for the node's latency, and
record the host's own processing time separately so that the two are not
confused.
\begin{pycode}
# On the single-board computer. Same features, same tree, different machine.
win = codec.decode(frame)
f = features.compute(win)          # the Python implementation, gated above
k = tree.predict(f)
port.write(codec.encode_class(k))
\end{pycode}

\item \textbf{Add the neural comparison, and hold it to the same gate.} Train a
small network on the same features, run it on the microcontroller under the
permissive runtime with the optimised kernels enabled, and compare it against
the tree on three axes: accuracy, kilobytes, and the spread of its latency. The
thesis worth testing is that the transform plus the statistical model wins on
accuracy per kilobyte, is deterministic in latency, and can be explained to a
customer. One peer-reviewed paper argues exactly that, and this is the chapter
that can check it on this part.

\item \textbf{Measure all three legs with one protocol and publish the
protocol.} For each leg: one hundred decisions, the marker pin pulsed at wake
and at done, the current trace captured, and the cycle counts recorded. Report
median, ninety-ninth percentile and maximum rather than a mean, because the
maximum is what a control loop has to tolerate. State which boundaries are
inside each measurement, and state that the host's own energy is outside all of
them.
\begin{shellcode}
python tools/report.py --leg a --n 100 --out results/leg_a.json
python tools/report.py --leg b --n 100 --out results/leg_b.json
python tools/report.py --leg c --n 100 --out results/leg_c.json
python tools/report.py --compare results/ --out results/README.md
\end{shellcode}
\end{steps}

\subsection*{Build, flash and debug}

\diagram{c16_timing}{One decision in each of the three placements, on one time axis. The shaded stages are what the marker pin separates inside the current trace. Every interval is a budget until an instrument has replaced it.}

\begin{shellcode}
cmake -B build-fw -G Ninja -DCMAKE_TOOLCHAIN_FILE=cmake/arm-none-eabi.cmake -DLEG=B
cmake --build build-fw -j
cp build-fw/placement.bin "$PROBE_DISK"/   # onto the probe disk
python tools/equivalence.py        # must pass before any result is recorded
\end{shellcode}

Recovering a board that will not answer is Chapter 1's material. The failure
specific to this chapter is subtler and does not present as a failure at all:
three legs that produce slightly different classes and a comparison that is
quietly meaningless. The equivalence gate is the only defence, and it belongs in
the build rather than in a habit.

\begin{note}{When the three legs disagree}
In order of likelihood: the window boundaries differ, because one leg starts its
window at a FIFO watermark and another at a sample count; the numeric format
differs, and a rounding rule that is invisible in one implementation changes a
feature at the last bit; the sample rate is not what was configured, because the
sensor's actual rate and its nominal rate differ by a few tenths of a percent;
one implementation normalises and another does not; and the tree was exported
from a different fit than the one deployed. Check them by feeding a recorded
window through all three and comparing the features before the class, which
localises the disagreement to a stage instead of to a leg.
\end{note}

\subsection*{Verification and acceptance criteria}

\begin{itemize}
\item The equivalence gate passes: all three implementations produce identical
features and identical classes on every window of the held-out set. The build
fails if it does not.
\item Every number in the results directory names the instrument that produced
it, and every number that has not been produced is absent rather than estimated.
\item Latency is reported as median, ninety-ninth percentile and maximum over at
least one hundred decisions per leg, never as a mean.
\item Memory is reported from the size output for the microcontroller legs and
from the in-sensor toolchain's own output for the sensor leg, with the two
clearly not being the same kind of number.
\item Charge per decision for the node is reported for all three legs, and the
exclusion of the host's own energy is stated in the same table rather than in a
footnote elsewhere.
\item The neural comparison is reported on accuracy, kilobytes and latency
spread, and the conclusion is stated even if it contradicts the thesis this
chapter set out to test.
\item The dataset protocol is published with the dataset, so that a reader can
see exactly how limited the hand-generated data is.
\end{itemize}

\subsection*{Variants}

\begin{table}[H]\centering\small
\begin{tabular}{p{24mm}p{26mm}p{44mm}p{22mm}p{20mm}}
\toprule
Axis & Variant & What changes & Cost & Built in full in \\
\midrule
Intelligence and reach & Classifier in the sensor & The decision is made on the sensor's own core and the microcontroller reads a class index. The bus read it still has to perform is the thing under test & A toolchain, and a repository with no licence at its root & Here \\
Intelligence and reach & Model on the MCU & Raw windows cross the bus, the features and the tree run on the board, and nothing crosses the link & Core time and flash, neither of which is scarce here & Here \\
Intelligence and reach & Model on the host & The window crosses the link and the class comes back. The node becomes a sensor and a modem & Link energy, and a host that must be awake & Here. The framed link itself is Chapter 9 \\
Intelligence and reach & Local only & No host at all, which is the baseline the first two legs already are & Nothing & Chapters 8 and 13 \\
Execution model & DMA into the window & The FIFO drain stops occupying the core, which matters most in the leg that has the most work to do & Cache maintenance on the window buffer & Chapter 4 and Chapter 19 \\
Operating system & An RTOS task set & The three stages become tasks with stated priorities and a measured latency distribution, which is a better home for the ninety-ninth percentile question & A scheduler to justify & Chapter 20 \\
Language & Host in Python & The host leg as written. A compiled host would change the host's numbers and not the node's & Host-side latency & Chapter 3 \\
\bottomrule
\end{tabular}
\caption{Variants for chapter 16. The three placements are this chapter's own, which is why the whole comparison lives here rather than being spread over three chapters where nothing could be compared.}
\end{table}

\subsection*{Pitfalls}

\begin{itemize}
\item Comparing three legs that were fed different data. A leg driven by live
movement is not comparable with a leg driven by a recording, however carefully
the movement is repeated.
\item Letting the numeric formats differ. The equivalence has to be defined in
the format the most constrained implementation can produce, which is not the
format that is most convenient to train in.
\item Reporting a mean latency. The interesting quantity is the tail, and the
mean hides exactly the behaviour that decides whether a classifier can sit
inside a control loop.
\item Assuming the in-sensor leg does no work on the microcontroller. It still
wakes, still performs a multi-transaction bus read, and that cost is the
preprint's whole argument.
\item Vendoring from a repository that declares no licence at its root because
its prebuilt outputs made the demonstration easy. Read it, run it, and check
each file before keeping any of it.
\item Quoting a vendor speed-up figure for the optimised neural kernels. Those
figures belong to later cores with a vector extension this part does not have.
\item Treating the popular training service as a supported path. It treats this
part as a porting target rather than a supported board: no ready-made firmware,
no data forwarder, and an exported kit that declares no single machine-readable
licence, with one optimising component proprietary and another sold only at the
top tier. The permissive runtime is the clean path and this chapter uses it.
\item Presenting the host's energy as unmeasured when it is in fact
unmeasurable on this bench. Those are different statements and the second one
has to be said out loud.
\item Claiming the comparison settles the question in general. It settles it for
this task, this data, this model and these three machines, which is already more
than anyone has published.
\end{itemize}

\subsection*{Best practices applied}

\begin{itemize}
\item The experiment is designed to remove confounds before it is designed to
produce numbers, and the equivalence gate is in the build rather than in the
author's discipline.
\item Every dependency's licence is stated before code is written around it, and
the one dependency with no licence at its root is named as such in the same
sentence that recommends reading it.
\item The boundary of the measurement is published with the measurement. What
the instrument cannot see is part of the result.
\item The claim is scoped to what was measured, and the configuration that could
overturn it, a microcontroller host with a short wake-up, is named rather than
left for a reader to discover.
\item Tail statistics rather than means, which is the reporting method the
serious real-time measurement literature uses and almost no vendor material
does.
\end{itemize}

\subsection*{Stretch goals}

\begin{itemize}
\item Add a fourth leg with a microcontroller as the host, using one of the
radio coprocessor boards on this bench over a serial link. Its wake-up is orders
of magnitude shorter than the single-board computer's, and it is the
configuration most likely to reverse the preprint's ranking. That result would
be new and it is within reach of this bench.
\item Run the standard small-system machine learning benchmark suite on this
part. No published result exists for this core: the suite's reference baseline
is a smaller core, the recent vendor submissions are a different core and an
accelerated preview part, and the vendor's own harness for running it has been
dormant since 2024.
\item Repeat the comparison with the state-machine engine rather than the
classifier engine for the in-sensor leg, which suits a decision that is a
sequence rather than a snapshot and would show whether the placement ranking
depends on the kind of decision.
\item Replace the hand-generated dataset with a mechanically repeatable one,
which is the single change that would turn the accuracy numbers from
illustrative into citable.
\end{itemize}

\subsection*{Roadmap and next steps}

Chapter 17 changes the subject from where code runs to how much a layer of
abstraction costs, and it uses the same instruments: the cycle counter, the size
output and the marker pin. Chapter 18 builds the transform stage this chapter
depends on, with a numerical acceptance test, and Chapter 20 gives the three
stages somewhere sensible to live when they have to share a processor with
anything else.

For going deeper on the machine-learning side, the free open textbook that grew
out of a university course and an online certificate is the closest thing to a
canonical curriculum for this material, and it is the right next step after this
chapter rather than before it. The free-to-audit course on embedded machine
learning, taught by the author of several widely used video series, covers the
same ground with exercises. The community embedded engineering roadmap is the
map for everything around it, and its central claim, that projects outrank
reading, is the reason this chapter is a build rather than a survey. Appendix H
lists them with their licences, including the two whose terms make them
something to recommend and never to quote.

\subsection*{Portfolio evidence}

\begin{itemize}
\item A public repository containing a labelled dataset, its collection
protocol, one feature definition with shared test vectors, and three
implementations that provably agree.
\item A results directory with latency, memory and charge for three placements,
each number naming its instrument, and a stated boundary for what the instrument
could not see.
\item A short written result that says which placement won on this task and by
how much, what the preprint predicted, and where the two agree and disagree.
That is a new result and it should be presented as one.
\item The equivalence gate running in the build, which is the piece of
engineering judgement in this chapter that a reviewer will recognise fastest.
\end{itemize}

\subsection*{Sources}

Normative references:
\begin{itemize}
\item UM3239, the user manual for the X-NUCLEO-IKS4A1, for the sensor
complement, the addresses and the three bus topologies.
\item The datasheet and the application note for the inertial part that carries
the programmable core, for the core's program memory size, its instruction set
and how a program is loaded.
\item Reference manual RM0455 for I2C1, USART3, the cycle counter and the timer
on this part. Not RM0433, which documents its better-known sibling.
\item The board user manual for MB1363, for the Arduino header mapping and the
free Zio pins that carry the marker signals.
\end{itemize}

Reusable implementations:
\begin{itemize}
\item The in-sensor processor examples, which ship prebuilt outputs so a
demonstration runs before the toolchain is installed. No licence is declared at
the repository root, so every file is checked individually.\newline
\url{https://github.com/STMicroelectronics/st-mems-ispu}
\item Ready-made configurations for the in-sensor classifier engine,
BSD-3-Clause. This repository replaced one that was discontinued in 2025.\newline
\url{https://github.com/STMicroelectronics/st-mems-machine-learning-core}
\item The in-sensor state machine collection, BSD-3-Clause, with the same
history of replacing a discontinued predecessor.\newline
\url{https://github.com/STMicroelectronics/st-mems-finite-state-machine}
\item The permissive generator for trees, forests, naive Bayes and anomaly
detection on small targets, MIT, from about 2 kB of flash.\newline
\url{https://github.com/emlearn/emlearn}
\item The signal-processing library, Apache-2.0, whose Python wrapper mirrors the
C API so that both sides run the same test vectors.\newline
\url{https://github.com/ARM-software/CMSIS-DSP}
\item The optimised neural kernels, Apache-2.0, which help on this core through
its signal-processing instructions rather than through a vector extension it
does not have.\newline
\url{https://github.com/ARM-software/CMSIS-NN}
\item The permissive neural runtime, Apache-2.0, which is the clean path for the
comparison leg.\newline
\url{https://github.com/tensorflow/tflite-micro}
\item The peer-reviewed argument that classical models plus good features
outperform neural networks per kilobyte on small targets.\newline
\url{https://arxiv.org/pdf/2107.09448}
\item The unreviewed preprint that measures two of the three placements and
reaches the counter-intuitive result this chapter treats as its hypothesis.\newline
\url{https://arxiv.org/abs/2606.00524}
\item The free open textbook that is the closest thing to a canonical curriculum
for this material.\newline
\url{https://mlsysbook.ai/}
\end{itemize}
   % Where the classifier runs: sensor, MCU or host

\part{Inside the Cortex-M7}
\project{17}{The same driver twice: vendor layer against registers}{NUCLEO-H7A3ZI-Q with one shield}{Abstraction layers, code size and cycles, and the Rust variant}

\begin{keyfacts}
\item[Board] NUCLEO-H7A3ZI-Q with one shield on the Arduino header. Either sensor shield serves, because what is measured is the transaction and not the sensor
\item[Peripherals] I2C1 for the sensor read, GPIO for the toggle measurement, the DWT cycle counter, TIM2 and SysTick as the substituted timebases, USART3 for the console
\item[Toolchain] arm-none-eabi-gcc with CMake for the two C implementations, and a Rust cross toolchain for the third and fourth
\item[Operating system] Bare metal, plus one real-time framework that schedules on the interrupt controller
\item[Difficulty] 4 of 5
\item[Effort] 5 evenings of about four hours
\item[Deliverable] One sensor read implemented four ways that return byte-identical results, with code size and cycle counts measured for each and the method published alongside the numbers
\end{keyfacts}

\subsection*{Why this project}

Ask whether a vendor's hardware abstraction layer is worth its cost and you will
get a confident answer within seconds and a measurement within never. This is
the weakest-evidenced topic in this book. There is a great deal of opinion, one
piece of reproducible bench work, and no well-argued case from the
vendor-layer side at all, which is itself remarkable: the layer is used by an
enormous number of projects and nobody has published a defence of it with
numbers. The chapter says that plainly and then supplies the measurement.

The one measured comparison worth the name is an engineer's bench work on a
Cortex-M7 part from this family running at 550 MHz, which found the vendor's pin
toggle call running at a fraction of the speed of a direct register write. That
is a real result and it is cited here with the caveat this whole volume exists
for: it was taken on a different member of the family, with a different
reference manual, a different clock tree and a different bus structure. The
shape of the finding transfers; the numbers do not.

The reason to do this on a sensor read rather than on a pin toggle is that the
answer changes, and the change is the finding. A pin toggle is two instructions
of real work, so any overhead at all dominates it. A six-byte register read over
a bus at 400 kHz is about nine bytes of traffic at nine bit times each, which is
81 bit times, which is about 202 microseconds, which at 280 MHz is roughly
56,700 core cycles of waiting. A call path costing a few hundred cycles is a
fraction of one percent of that. Both statements are true at once, and a reader
who has only met the first one will make bad decisions with it. The useful
quantity is not the overhead but the ratio of overhead to work, and this chapter
measures that ratio on two operations that sit at opposite ends of it.

\begin{note}{What this chapter is allowed to claim}
That the two implementations return identical bytes, and what each one cost in
flash and in cycles on this part at a stated clock, a stated optimisation level
and a stated cache state. Not that one layer is better. Not that the result
transfers to another part. Not that the ratio found here holds for a different
operation. The value of a measurement is in its boundaries, and a chapter on
this subject that overreached would be adding to the problem it is about.
\end{note}

\subsection*{Prior art and what to reuse}

\begin{table}[H]\centering\small
\begin{tabular}{p{34mm}p{46mm}p{38mm}p{20mm}}
\toprule
Source & What it gives & What it does not & Licence \\
\midrule
One engineer's bench measurement of pin toggling on this silicon family & The only reproducible comparison of the vendor call against a direct register write, with a method described well enough to repeat & It was taken on a different member of this family at a different clock, so its numbers are the shape of the answer and not the answer & Reading only \\
The vendor's firmware package for this board & Parallel example trees at both abstraction levels for this exact board, which is the raw material this chapter measures & Its examples are under a licence tied to this vendor's silicon, so they are read and measured rather than copied into a portfolio & Mixed, flagged \\
The Rust hardware crate for this family & A feature that names this part, routing it through its twin's peripheral module with the correct reference manual selected. That is the right engineering call and it is explained here rather than hidden & No example for this exact board, and a different set of assumptions about ownership and initialisation & 0BSD \\
The Rust peripheral crate collection & A generated peripheral module for this part, with typed register access that is itself an abstraction worth measuring & Nothing above the register level & Apache-2.0 or MIT \\
The async framework & It names this exact package and pin count, which is a stronger statement of support than most projects make & No worked example for this part at the time of writing, only for its twin. That gap is worth filling & Apache-2.0 or MIT \\
The real-time framework & Scheduling built on the interrupt controller, which is an architecture-level mechanism, so it works here without anything part-specific & No peripheral support of its own, by design & Apache-2.0 or MIT \\
The flashing and debugging tool & A target manifest carrying this exact die & Nothing about abstraction layers & Apache-2.0 or MIT \\
\bottomrule
\end{tabular}
\caption{Prior art for chapter 17. The gap in this table is the interesting part: there is no entry for a well-argued case in favour of the vendor layer, because there is not one to cite.}
\end{table}

What is left to write is the measurement and its method, plus the two
implementations that have to agree before the measurement means anything. The
method is the deliverable. Anyone can produce a number; producing a number with
the clock, the optimisation level, the cache state, the repetition count and the
reported statistic all stated is what makes it usable by somebody else, and it
is exactly what the existing writing on this subject does not do.

\subsection*{Parts from the inventory}

\begin{table}[H]\centering\small
\begin{tabular}{p{40mm}p{66mm}p{40mm}}
\toprule
Part & Role & Interface \\
\midrule
NUCLEO-H7A3ZI-Q & The whole chapter. Its probe is also the flashing path for both toolchains & Micro USB to the host \\
One sensor shield & The target of the read. Either the four-sensor or the five-sensor shield serves, because the operation under test is a bus transaction & I2C1 on the Arduino header \\
A USB data cable & Power, programming and console & Micro USB \\
nRF-PPK2 & The marker pin recorder, and a charge figure per operation if you want one & IDD jumper and one digital input \\
Host PC & Two toolchains, two build systems, one comparison script & Windows with the Linux subsystem \\
\bottomrule
\end{tabular}
\caption{Inventory items used in chapter 17. Nothing is bought and nothing is soldered. One shield is on the board at a time.}
\end{table}

\subsection*{System architecture}

\diagram{c17_arch}{Four implementations of one operation, one equality harness and one set of instruments. The harness is what makes the four comparable; the instruments are what make the comparison a measurement rather than an impression.}

The figure has one property worth pointing out. The two C implementations and
the two Rust implementations form the same ladder with different names. In C the
rungs are the vendor layer, the low-level layer and raw register access. In Rust
they are the hardware abstraction crate, the generated peripheral crate and a
raw volatile write. The peripheral crate is not the bottom of the ladder even
though it looks like registers, because it is a typed, ownership-checked view of
them and that view has a cost that can be measured like any other.

\subsection*{Peripheral configuration}

\begin{table}[H]\centering\small
\begin{tabular}{p{22mm}p{26mm}p{24mm}p{30mm}p{30mm}}
\toprule
Peripheral & Mode & Clock source & Pins and function & Interrupt and transfers \\
\midrule
I2C1 & Fast mode, 400 kHz & Peripheral bus, source selected in RCC & D14 and D15 on the Arduino header. Confirm in MB1363 & Blocking in all four implementations, so that only the layer differs \\
GPIO output & Push-pull, maximum speed & Peripheral bus & One free Zio pin, toggled for the second operation & None \\
GPIO output & Push-pull marker & Peripheral bus & One free Zio pin into a PPK2 digital input & None \\
DWT cycle counter & Free running & Core clock & None & The primary instrument \\
TIM2 & Free running, 32-bit, 1 MHz & Peripheral bus & None & The first substituted timebase \\
SysTick & Free running, reload at maximum & Core clock or its divided form & None & The second substituted timebase \\
USART3 & Asynchronous, 115200 8N1 & Peripheral bus & The console pins of chapter 1 & Polled \\
\bottomrule
\end{tabular}
\caption{Peripheral configuration. Every implementation uses the blocking form of the transfer, because comparing a blocking vendor call against an interrupt-driven register implementation would measure the execution model rather than the abstraction.}
\end{table}

\subsection*{Wiring}

\diagram{c17_wiring}{The shield supplies a device to talk to and the toggled pin supplies the other end of the ratio. One marker wire reaches the instrument. Nothing else is placed by hand.}

3.3 V logic only, as everywhere in this book, and no 5 V signal into a pin. The
toggled pin drives nothing but the instrument's digital input, so there is no
load to argue about. There is no oscilloscope on this bench, which matters more
here than in most chapters: a toggle at these rates cannot be observed
externally with what is available, so the cycle counter is the instrument and
the instrument's own cost is subtracted from every reading. The substitute for a
scope is stated rather than implied, and the subtraction is in the harness.

\subsection*{Memory and timing budget}

\diagram{c17_mem}{The two registers a pin toggle can go through. One store to the set and reset register changes a pin with no read; a read, modify and write of the output data register is three accesses and has a window in it. The difference is atomicity first and speed second, and the layer you call decides which one you get.}

\begin{table}[H]\centering\small
\begin{tabular}{p{44mm}p{28mm}p{28mm}p{30mm}}
\toprule
Quantity & Budget & Measured & Margin \\
\midrule
Flash, vendor-layer read & 4 kB & not measured & not measured \\
Flash, register-level read & 1 kB & not measured & not measured \\
Cycles, vendor-layer read & 57,000 & not measured & not measured \\
Cycles, register-level read & 56,800 & not measured & not measured \\
Cycles, vendor-layer toggle & 40 & not measured & not measured \\
Cycles, register-level toggle & 4 & not measured & not measured \\
Bus time for the read at 400 kHz & 202 \textmu{}s computed & not measured & not measured \\
Cost of reading the cycle counter & subtracted & not measured & not measured \\
\bottomrule
\end{tabular}
\caption{The budget table. The two read rows are deliberately close together and the two toggle rows are deliberately far apart, because that is the prediction this chapter exists to test. The bus figure is arithmetic from nine bytes at nine bit times and is marked computed rather than measured.}
\end{table}

Three things have to be stated with every cycle count on this part or the number
is not reusable. The core clock, because a cycle is not a time. The optimisation
level, because a call that is inlined at one level and not at another is a
different program. And the cache state, because this core has instruction and
data caches and a first run through cold code can cost many times what the
tenth run costs. The harness reports the first run and the median of the
remainder separately, and Chapter 19 explains why the two differ.

\subsection*{Firmware design (UML)}

\diagram{c17_uml}{One read, two call paths. The right-hand path is what the left-hand path does on your behalf, and reading the two side by side is the fastest way to see what the abstraction is actually for.}

What the vendor layer does that the register implementation does not is the
substance of the chapter, and it is not nothing. It checks that the peripheral
is in a state where the call is legal. It takes a lock, so that two callers
cannot interleave. It computes a deadline from a millisecond tick and abandons
the transfer if the deadline passes, which turns a stuck bus into an error code
rather than a stopped program. It inspects the error flags afterwards and maps
them onto a small set of return values. And it presents the same signature
whether the transfer is blocking, interrupt-driven or performed by the transfer
engine, which is worth a great deal when a design changes its mind.

The register implementation does none of that by default, which is why the
honest form of this chapter's question is not whether the abstraction is worth
its cost but which of those five services you actually need. A register-level
implementation that adds a timeout and error mapping, because it needs them,
converges on the vendor layer and its measured cost converges too. That
convergence is a result, and it is the part of this subject that opinion pieces
consistently miss.

\subsection*{Data flow (ASCII)}

\begin{asciiart}
  one operation, defined once
     |
     +--> impl 1: C, vendor layer ----------+
     +--> impl 2: C, registers -------------+
     +--> impl 3: Rust, peripheral crate ---+
     +--> impl 4: Rust, hardware crate -----+
                                            |
                                            v
  +-----------------------------------------------------------------------------------+
  | equality harness on the host: the four byte strings must be identical             |
  |   before any number anywhere in this chapter is recorded                          |
  +-----------------------------------------------------------------------------------+
                                            |
                                            v
  +-----------------------------------------------------------------------------------+
  | instruments: DWT cycles, first run and median reported                            |
  |   separately; arm-none-eabi-size per object; TIM2 and                             |
  |   SysTick as the substituted timebases; marker pin to PPK2                        |
  +-----------------------------------------------------------------------------------+
                                            |
                                            v
  +-----------------------------------------------------------------------------------+
  | one table: core clock, optimisation level, cache state, repetition                |
  |   count, median and maximum, in every row                                         |
  +-----------------------------------------------------------------------------------+
\end{asciiart}

\subsection*{Repository layout}

\begin{plaincode}
nucleo-h7a3-two-layers/
  c/
    CMakeLists.txt
    src/op_vendor.c          # the sensor read through the vendor layer
    src/op_registers.c       # the same read, written against the registers
    src/op_toggle.c          # both toggles, the other end of the ratio
    src/measure.c            # DWT, and the subtraction of its own cost
    src/timebase_dwt.c       # peripheral substitution, implementation 1
    src/timebase_tim2.c      # peripheral substitution, implementation 2
    src/timebase_systick.c   # peripheral substitution, implementation 3
  rust/
    Cargo.toml               # the crate features are the load-bearing part
    src/bin/pac.rs           # the read through the generated peripheral crate
    src/bin/hal.rs           # the read through the hardware crate
    src/bin/embassy.rs       # the worked example this part did not have
    memory.x                 # checked against RM0455, not inherited
  tools/
    equality.py              # four byte strings, one assertion
    collect.py  report.py    # the table, with its conditions in it
  docs/METHOD.md             # clock, optimisation, cache state, N, statistic
  README.md
\end{plaincode}

\subsection*{Steps}

\begin{steps}
\item \textbf{Define the operation before writing any implementation.} Write it
down in \texttt{docs/METHOD.md}: read six bytes from a named register of the
sensor at a named address, over I2C1 at 400 kHz, blocking, with the result
printed as twelve hexadecimal characters. Ambiguity here is what makes most
published comparisons unusable, because two implementations that do slightly
different work are not two implementations of one operation.

\item \textbf{Write the equality harness first.} It is ten lines and it is what
stops the chapter from comparing two different programs.
\begin{pycode}
results = {name: read_line(port, name) for name in
           ("vendor", "registers", "pac", "hal")}
assert len(set(results.values())) == 1, results
print("identical:", next(iter(results.values())))
\end{pycode}

\item \textbf{Implementation one, the vendor layer.} Use the blocking memory
read from the vendor's layer, configured by its own initialisation structure.
Keep it idiomatic: the point is to measure the layer as it is actually used, not
a stripped version of it that nobody writes.
\begin{ccode}
/* The vendor's own call. Note what it takes that the register version
   does not: a handle carrying state and a lock, and a timeout. */
if (HAL_I2C_Mem_Read(&hi2c1, DEV_ADDR << 1, REG_ADDR,
                     I2C_MEMADD_SIZE_8BIT, buf, 6, TIMEOUT_MS) != HAL_OK) {
    return -1;
}
\end{ccode}

\item \textbf{Implementation two, the registers.} Write the transfer against the
peripheral's own registers. One register on this peripheral carries the target
address, the direction, the byte count, the start condition and the automatic
end condition together, which is the single clearest example in the part of what
an abstraction is hiding. Read the field positions in RM0455 before writing
them; this book does not print a register map it has not checked, and the
reference manual for this part is not the one most search results return.
\begin{ccode}
/* Phase one: write the register address. Phase two: repeated start and
   read six bytes. The transfer-configuration register carries the target
   address, the direction, the count, the start and the automatic end. */
i2c_write_cr2(DEV_ADDR, 1, I2C_WRITE, I2C_SOFTEND);
i2c_put_byte(REG_ADDR);
i2c_wait_tc();                       /* transfer complete, not stop */
i2c_write_cr2(DEV_ADDR, 6, I2C_READ, I2C_AUTOEND);
for (int i = 0; i < 6; i++) buf[i] = i2c_get_byte();
i2c_wait_stop_and_clear();
\end{ccode}
Everything the vendor layer does and this does not is now visible as an absence:
no state check, no lock, no deadline, no error mapping. Add them back one at a
time later and watch the cost converge.

\item \textbf{Measure code size, per object rather than per image.} An image
figure is dominated by whatever else is linked in. The honest comparison is the
size of the objects that implement the operation, plus whatever they drag in.
\begin{shellcode}
arm-none-eabi-size -t build/CMakeFiles/*.dir/src/op_vendor.c.obj
arm-none-eabi-size -t build/CMakeFiles/*.dir/src/op_registers.c.obj
arm-none-eabi-nm --size-sort -S build/two_layers.elf | tail -20
\end{shellcode}

\item \textbf{Measure cycles, and subtract the instrument.} Reading the cycle
counter costs cycles. Measure an empty interval first and subtract it from every
reading, then report the first run and the median of the rest separately.
\begin{ccode}
static uint32_t cyc_overhead;

void measure_calibrate(void)
{
    uint32_t best = UINT32_MAX;
    for (int i = 0; i < 64; i++) {
        const uint32_t a = DWT->CYCCNT;
        const uint32_t b = DWT->CYCCNT;
        if (b - a < best) best = b - a;
    }
    cyc_overhead = best;              /* subtracted from every reading */
}

uint32_t measure_run(void (*op)(void))
{
    const uint32_t a = DWT->CYCCNT;
    op();
    const uint32_t b = DWT->CYCCNT;
    return (b - a) - cyc_overhead;
}
\end{ccode}
Whether the cycle counter works on this part with no debugger attached is on
this volume's confirm list. Check it before relying on it, because if it needs a
probe the measurement is only available on a bench and not in the field.

\item \textbf{Do the same on a pin toggle and compute the ratio.} This is the
step that produces the chapter's finding. The toggle is where the published
comparison found a large difference, and the read is where the bus dominates.
Report both, and report overhead as a fraction of the operation rather than as a
number of cycles.
\begin{ccode}
static inline void toggle_registers(void)   /* one store, no read */
{
    GPIOx->BSRR = PIN_MASK;                  /* set   */
    GPIOx->BSRR = PIN_MASK << 16;            /* reset */
}
\end{ccode}
The register figure above is what that line is doing. A read, modify and write
of the output register would produce the same waveform and would not be atomic,
which is a correctness difference and not only a speed one.

\item \textbf{Substitute the peripheral, and keep the job the same.} This
chapter owns that axis for the whole book, so it is built here in full. The job
is to produce an elapsed interval; three peripherals can do it.
\begin{ccode}
/* One interface, three peripherals behind it. Selected at build time so
   that nothing is paid for at run time. */
typedef struct {
    void     (*start)(void);
    uint32_t (*ticks)(void);     /* free-running count */
    uint32_t hz;                 /* what a tick is worth */
    uint32_t width_bits;         /* where it wraps */
} timebase_t;

extern const timebase_t TB_DWT;      /* core clock, 32 bits, no prescale  */
extern const timebase_t TB_TIM2;     /* 1 MHz, 32 bits, survives more     */
extern const timebase_t TB_SYSTICK;  /* core clock, 24 bits, counts down  */
\end{ccode}
Then compare them on resolution, wrap period, the cost of a read, and
availability. At 280 MHz the cycle counter wraps in about 15.3 seconds, a 32-bit
timer at 1 MHz wraps in about 71.6 minutes, and a 24-bit down counter at the
core clock wraps in about 60 milliseconds, which makes the third one unusable
for anything longer than an interrupt handler and perfectly good inside one.

\item \textbf{The Rust variant, in full.} The crate features are the part of
this that matters and they are verified: the hardware crate has a feature naming
this part which routes to its twin's peripheral module and selects the right
reference manual, and the async framework names this exact package and pin
count.
\begin{plaincode}
[dependencies]
cortex-m      = "0.7"
cortex-m-rt   = "0.7"
# The feature names are the verified part of this file. Take the versions
# from your own lockfile rather than from a page in a book.
stm32h7xx-hal = { version = "0.16", features = ["stm32h7a3", "rt"] }
\end{plaincode}
That routing is worth a paragraph in your README rather than a shrug. The crate
is not pretending this part is its twin; it is stating that the peripheral
module and the reference manual selection of the twin are the correct ones for
this die, which is precisely the distinction this whole volume is about, made
correctly, in a build file.
\begin{rustcode}
// The same six bytes, through the generated peripheral crate and through
// the hardware crate. Two rungs of the same ladder, in one language.
let dp = pac::Peripherals::take().unwrap();
let mut buf = [0u8; 6];

// hardware crate: ownership, typed pins, a configured bus
let mut i2c = dp.I2C1.i2c((scl, sda), 400.kHz(), ccdr.peripheral.I2C1,
                          &ccdr.clocks);
i2c.write_read(DEV_ADDR, &[REG_ADDR], &mut buf).unwrap();
\end{rustcode}
Build and flash with the same probe the C build uses, which is the practical
reason this variant is cheap to add on this bench.
\begin{shellcode}
rustup target add thumbv7em-none-eabihf
cargo build --release --bin hal
probe-rs run --chip STM32H7A3ZITxQ target/thumbv7em-none-eabihf/release/hal
\end{shellcode}

\item \textbf{Write the worked example this part did not have.} The async
framework names this exact package and pin count and had, at the time of
writing, a worked example only for its twin. Writing one for this part is a
small piece of genuinely new work, it is exactly the kind of contribution that
makes a portfolio repository useful to somebody else, and it is worth offering
upstream. The real-time framework needs no such work, because it schedules on
the interrupt controller, which is an architecture-level mechanism rather than a
part-specific one, and that is worth one sentence in the README so a reader
understands why one of the two needed porting and the other did not.
\end{steps}

\subsection*{Build, flash and debug}

\diagram{c17_timing}{What the two toggles look like against the core clock, and where the cycles go in the two reads. The toggle is where the layer costs a multiple; the read is where it costs a rounding error. Both are true and the chapter reports both.}

\begin{shellcode}
cmake -B c/build -DCMAKE_TOOLCHAIN_FILE=cmake/arm-none-eabi.cmake -DOPT=O2
cmake --build c/build -j
probe-rs run --chip STM32H7A3ZITxQ c/build/two_layers.elf
python tools/equality.py --port /dev/ttyACM0
python tools/report.py --out docs/RESULTS.md
\end{shellcode}

Recovering a board that will not answer is Chapter 1's material. The failure
specific to this chapter is a comparison that looks fine and is not.

\begin{note}{When the two implementations return different bytes}
In order of likelihood: the address was shifted in one implementation and not in
the other, because the vendor layer expects the seven-bit address in the upper
bits of an eight-bit field and a register implementation usually does not; the
register implementation stopped the transfer after phase one instead of
producing a repeated start; the automatic end condition was configured in the
phase where it should not be, so the sensor saw two transactions rather than
one; the byte count was written to the wrong field; or the sensor auto-increments
its register pointer in one case and not in the other because a configuration
bit differs. Compare the first byte only, then two bytes, then six. The point at
which they diverge names the fault.
\end{note}

\subsection*{Verification and acceptance criteria}

\begin{itemize}
\item All four implementations return byte-identical results, proven by the
harness rather than by inspection, and the harness runs in the build.
\item The cost of the abstraction is a number with its conditions attached: core
clock, optimisation level, cache state, repetition count, and median together
with maximum. A number without those is not published.
\item The instrument's own cost is measured and subtracted, and the subtraction
is visible in the source rather than folded into a constant.
\item The first run and the median of the remainder are reported separately, and
the difference between them is acknowledged as a cache effect rather than
averaged away.
\item Code size is reported per object with what it drags in, not as a
difference between two image sizes.
\item The ratio of overhead to work is reported for both operations, and the
chapter's conclusion is stated as a function of that ratio rather than as a
verdict on the layer.
\item The timebase substitution is demonstrated on all three peripherals, with
the wrap period of each stated, and the cycle counter's dependence on a probe
resolved rather than assumed.
\item The Rust build flashes with the same probe as the C build and the
peripheral module routing is explained in the README in one paragraph.
\end{itemize}

\subsection*{Variants}

\begin{table}[H]\centering\small
\begin{tabular}{p{24mm}p{26mm}p{44mm}p{22mm}p{20mm}}
\toprule
Axis & Variant & What changes & Cost & Built in full in \\
\midrule
Language & Rust & The same read at two rungs of the ladder, with the hardware crate routing this part through its twin's peripheral module and selecting the right reference manual & A second toolchain, and a different model of ownership to learn & Here \\
Peripheral substitution & One job on three peripherals & The elapsed-interval job moved between the cycle counter, a 32-bit timer and the system tick, with resolution and wrap period compared & Each has a different wrap period and a different dependence on the debugger & Here \\
Peripheral substitution & The same read on another bus instance & The identical transfer through a different peripheral instance, which on this family can mean a different clock domain and therefore different behaviour in low-power modes & Pin availability while a shield is fitted, which is on the confirm list & Here, gated on the header map \\
Language & C++ & The same read behind a template that resolves the layer at compile time, with no run-time cost and a different set of arguments about readability & Toolchain settings to justify & Chapter 9 \\
Execution model & Interrupt and DMA forms & The vendor layer presents the same signature for all three, which is one of the services being paid for & A different measurement entirely & Chapters 3 and 4 \\
Operating system & An RTOS task & The lock the vendor layer takes starts to matter once two tasks can call the same peripheral & Priority inheritance and a mutex to size & Chapter 20 \\
\bottomrule
\end{tabular}
\caption{Variants for chapter 17. Two axes are built in full here, which is why this chapter is longer on implementation and shorter on narrative than its neighbours.}
\end{table}

\subsection*{Pitfalls}

\begin{itemize}
\item Comparing a blocking call against an interrupt-driven implementation. That
measures the execution model and attributes the result to the abstraction.
\item Quoting a cycle count without the clock, the optimisation level and the
cache state. The number is then unusable by anybody, including you, in six
months.
\item Measuring only the pin toggle. It is the operation where overhead
dominates, and generalising from it is how this subject acquired its confident
opinions.
\item Forgetting the address shift. The vendor layer and a register
implementation disagree about where the seven-bit address sits, and the symptom
is a device that does not acknowledge.
\item Reporting a mean over runs. The first run through cold code is a different
measurement from the tenth, and averaging them produces a number that describes
neither.
\item Taking a linker script or a memory layout from the better-known member of
this family for the Rust build. The same rule applies in every language:
\texttt{memory.x} is checked against RM0455.
\item Treating the generated peripheral crate as the bottom of the ladder. It is
a typed view of registers with its own cost, and measuring it as though it were
raw access hides a rung.
\item Concluding that the vendor layer is expensive. What was measured is that
it is expensive relative to an operation with almost no work in it, which is a
much narrower statement and the only one the data supports.
\end{itemize}

\subsection*{Best practices applied}

\begin{itemize}
\item The equality harness exists before either implementation, so the
comparison cannot quietly become a comparison of two different operations.
\item The instrument's own cost is measured and subtracted rather than assumed
negligible.
\item Every published number carries its conditions in the same table, which is
the single thing the existing writing on this subject does not do.
\item The absence of prior art is stated as a finding rather than filled in with
confidence.
\item A dependency's engineering decisions are explained rather than worked
around: the crate that routes this part through its twin's peripheral module is
doing the right thing, and the README says why.
\item The conclusion is scoped to the ratio that was measured, and the operation
where the opposite conclusion holds is published in the same table.
\end{itemize}

\subsection*{Stretch goals}

\begin{itemize}
\item Add the services back to the register implementation one at a time,
namely the state check, the lock, the deadline and the error mapping, and plot
cost against services. The point at which the two lines meet is the real answer
to this chapter's question and nobody has published that curve.
\item Repeat the toggle measurement with the instruction cache disabled and with
the code placed in tightly coupled memory, which separates the abstraction's
cost from the memory system's. That is Chapter 19's subject approached from this
side.
\item Build the same four implementations at three optimisation levels and
publish the nine-cell table. Most of the disagreement in this field is people
comparing different cells of it.
\item Write the async framework example for this part properly, with a
peripheral that is not a pin, and offer it upstream. It is a small contribution
and it closes a real gap.
\item Measure charge per operation rather than cycles per operation with the
power instrument, which is the quantity that actually matters in Part II and
which nobody reports for abstraction layers at all.
\end{itemize}

\subsection*{Roadmap and next steps}

Chapter 18 takes the same instruments to a different question, namely what a
transform costs in two numeric formats with a numerical acceptance test, and
Chapter 19 explains the cache effect that this chapter's first run and median
are separated by. Chapter 20 is where the lock that the vendor layer takes stops
being decorative.

For going deeper, the community embedded engineering roadmap separates the
firmware path from the embedded Linux path and is explicit about weighting, and
the accompanying essay on bridging the gap makes the same argument in a page.
For the Rust direction this chapter opens, the standard progression is the
embedded working group's own book followed by the async framework's examples,
and the one caution worth repeating is the one this chapter is built on: check
that an example names your part and not its family. Appendix H lists the
courses; the free one whose code is under the most permissive licence in the
pool is the safest to build an exercise on, although its examples are not on
this board.

\subsection*{Portfolio evidence}

\begin{itemize}
\item A public repository containing four implementations of one operation that
provably return the same bytes, in two languages, with one build command each.
\item A results table whose every row carries its clock, optimisation level,
cache state, repetition count and statistic, which is the artefact this subject
currently lacks entirely.
\item A short written conclusion stating the ratio of overhead to work for two
operations at opposite ends of it, and saying plainly which of the two the
popular opinion is derived from.
\item A worked example for this exact part in a framework that had one only for
its twin, with the peripheral module routing explained.
\end{itemize}

\subsection*{Sources}

Normative references:
\begin{itemize}
\item Reference manual RM0455, for the I2C peripheral's transfer-configuration
register, the GPIO set and reset register, the cycle counter and the timers on
this part. Not RM0433, which documents its better-known sibling.
\item The STM32H7A3xI datasheet, for the clock limits that the cycle counts are
divided by.
\item The board user manual for MB1363, for the Arduino header mapping and the
free Zio pins used for the toggled pin and the marker.
\item The Cortex-M7 technical reference manual, for the cycle counter and the
cache behaviour that separates a first run from a median.
\end{itemize}

Reusable implementations:
\begin{itemize}
\item The only reproducible bench comparison of the vendor toggle call against a
direct register write, taken on a different member of this family at 550 MHz.\newline
\url{https://metebalci.com/blog/stm32h7-gpio-toggling/}
\item The Rust hardware crate for this family, 0BSD, whose feature for this part
routes to its twin's peripheral module and selects the right manual.\newline
\url{https://github.com/stm32-rs/stm32h7xx-hal}
\item The generated peripheral crates for this silicon, which carry a module for
this part.\newline
\url{https://github.com/stm32-rs/stm32-rs}
\item The async framework, Apache-2.0 or MIT, which names this exact package and
pin count and had no worked example for this part at the time of writing.\newline
\url{https://github.com/embassy-rs/embassy}
\item The real-time framework, which works here because it schedules on the
interrupt controller rather than on anything part-specific.\newline
\url{https://github.com/rtic-rs/rtic}
\item The flashing and debugging tool whose target manifest carries this exact
die.\newline
\url{https://github.com/probe-rs/probe-rs}
\item Oblaukhov's CMake package layer for this silicon family, MIT, which is the
practical way to build the vendor-layer implementation without its IDE.\newline
\url{https://github.com/ObKo/stm32-cmake}
\end{itemize}
   % The same driver twice: vendor layer against registers
\project{18}{Transforms and filters with a numerical acceptance test}{NUCLEO-H7A3ZI-Q}{CMSIS-DSP, float32 against Q15, cycles, memory placement}

\begin{keyfacts}
\item[Board] NUCLEO-H7A3ZI-Q, and nothing else. The signal is a stored vector, so the result is reproducible on any desk
\item[Peripherals] RCC and PWR for 280 MHz, the floating-point unit, the cycle counter in the debug block, the two L1 caches, USART3 for the result stream
\item[Toolchain] arm-none-eabi-gcc with CMake, CMSIS-DSP built from source, and a host Python environment with SciPy, NumPy and the library's own Python wrapper
\item[Operating system] Bare metal
\item[Difficulty] 4 of 5
\item[Effort] 4 evenings of about four hours
\item[Deliverable] A filter and spectrum module whose output is proved equal to a host reference within a tolerance stated before the first run, and a table of cycles per transform in which every row names its numeric format, its code placement, its data placement and the state of both caches
\end{keyfacts}

\subsection*{Why this project}

Every earlier chapter moved bytes. This one computes with them, and it is the
first chapter where being fast and being correct are separate claims that need
separate evidence. A transform that is quick and wrong is worse than no
transform at all, because the wrongness arrives as a plausible number rather
than as an error. So the order here is fixed: the acceptance test is written
first, the tolerance is written down before anything is measured, and only then
does anyone look at cycle counts.

The second reason this chapter exists is that the obvious question about this
part has no published answer. Ask whether a filter and a spectrum should be
computed in single-precision floating point or in 16-bit fixed point on a
Cortex-M7 at 280 MHz, and the answer that comes back is a ratio measured on a
Cortex-M4, on a part with a different memory system, quoted without its
conditions. That number does not transfer. A Cortex-M4 has no cache, no
tightly coupled memories worth the name in this sense, and a different
floating-point unit. On this part the interesting variable is probably not the
numeric format at all but where the code and the data sit, and that expectation
is itself the finding the chapter sets out to confirm or overturn.

The third reason is portfolio-shaped. A repository that says "runs an FFT"
proves nothing. A repository that says "the float32 path matches a
double-precision host reference to a relative root-mean-square error below 1e-5
on the shipped vector, the Q15 path reaches at least 55 dB signal to error
ratio on the same vector, and here are the cycle counts for eight placements
with the cache state named in every row" is an engineering artefact. The work
is the same. Only the reporting is different, and the reporting is the part
that is usually missing.

\begin{note}{The widely repeated ratio is not for this core}
The figure everyone quotes for fixed point against floating point on a
microcontroller transform comes from a Cortex-M4 measurement. It is not wrong
there; it simply does not describe this part. This chapter treats it as a
hypothesis rather than a result, and the acceptance criteria below require the
measurement to name its own conditions so that the next reader does not inherit
the same problem from this book.
\end{note}

\subsection*{Prior art and what to reuse}

\begin{table}[H]\centering\small
\begin{tabular}{p{34mm}p{46mm}p{38mm}p{20mm}}
\toprule
Source & What it gives & What it does not & Licence \\
\midrule
CMSIS-DSP & The whole computational layer: cascaded biquads in float32 and Q15, real FFTs in both formats, complex magnitude, statistics. Maintained, tested, and the reference implementation for this architecture & It does no filter design at all. Coefficients arrive from somewhere else, and choosing them is the engineering & Apache-2.0, entirely clean \\
The library's Python wrapper & The same functions callable from Python with an API that mirrors the C API, so one set of test vectors drives both sides and a disagreement is localised immediately & It is a wrapper, not a specification. Where the C API is awkward the Python is awkward in the same way & Apache-2.0 \\
The maintainer's learning path on building a pipeline against a reference implementation & The method this chapter follows: design in Python, check against a reference, then move the same computation to the device & It targets the development flow rather than any one part, so nothing in it is about this silicon & Vendor learning material, read and cite \\
The vendor's application note on system architecture and performance & Exactly the benchmark shape this chapter wants: one transform run with code and data in each memory, with the cache on and off, and a companion package that ships the linker scripts for those placements & It does not cover this part. Its numbers are the shape of the answer and not the answer, and its linker scripts are checked against RM0455 before use & Vendor documentation, read do not copy \\
SciPy and NumPy on the host & Filter design, the double-precision reference output, and the statistics that turn two arrays into a pass or fail & Nothing runs on the device. The host is the judge, not the subject & BSD-3-Clause \\
\bottomrule
\end{tabular}
\caption{Prior art for chapter 18. The computational library is clean enough to vendor outright, which is unusual in this book and worth saying plainly.}
\end{table}

What is left to write is the part that decides whether the result is
trustworthy: the coefficient design and its export, the test vector and its
reference output, the comparison and its stated tolerance, the cycle harness
that reports a placement and a cache state alongside every count, and the
linker script sections that put code and data where the experiment wants them.
The library supplies arithmetic. It supplies no evidence, and evidence is the
deliverable.

\subsection*{Parts from the inventory}

\begin{table}[H]\centering\small
\begin{tabular}{p{40mm}p{66mm}p{40mm}}
\toprule
Part & Role & Interface \\
\midrule
NUCLEO-H7A3ZI-Q & The whole chapter. The transform runs here and the cycle counter lives here & Micro USB to the host \\
A USB data cable & Power, programming and the result stream on one lead & Micro USB \\
Host PC & Filter design, reference computation, the Python wrapper, and the comparison that decides pass or fail & Python with NumPy, SciPy and the wrapper \\
Chapter 6's sampler, optionally & Replaces the stored vector with real samples at a proven rate when you want an end to end demonstration rather than a reproducible test & Already on the board \\
\bottomrule
\end{tabular}
\caption{Inventory items used in chapter 18. Nothing is wired and nothing is bought. The stored vector is deliberate: an acceptance test that depends on a live signal is not an acceptance test.}
\end{table}

\subsection*{System architecture}

\diagram{c18_arch}{The host owns design and judgement, the board owns computation and measurement. The same test vector crosses the link in one direction and a result block plus a cycle count comes back in the other.}

Three things run on the host and two on the board, and the split matters. On
the host: the filter design, which produces coefficients; the reference
computation in double precision, which produces the array the board must match;
and the same computation through the library's Python wrapper, which is the
second opinion. On the board: the computation under test, and the cycle counter
that times it. The reason for two host references rather than one is that they
answer different questions. A disagreement with the double-precision reference
means the arithmetic has drifted. A disagreement with the wrapper, which is the
same library, means the port is wrong: a coefficient order, a state buffer
size, a scaling shift.

\subsection*{Peripheral configuration}

\begin{table}[H]\centering\small
\begin{tabular}{p{22mm}p{26mm}p{24mm}p{30mm}p{30mm}}
\toprule
Peripheral & Mode & Clock source & Pins and function & Interrupt and transfers \\
\midrule
RCC and PWR & Bypass input then phase-locked loop, 280 MHz & 8 MHz from the probe & None & None \\
Floating-point unit & Enabled in the coprocessor access register before any float runs & Core clock & None & None \\
Cycle counter & Trace enable then counter enable, read directly & Core clock & None & None. Confirm it counts with no probe attached \\
L1 instruction cache & On or off, named in every row of the result table & Core clock & None & None \\
L1 data cache & On or off, named in every row of the result table & Core clock & None & Write-back by default \\
USART3 & Asynchronous, 115200 8N1 & Peripheral bus & Board manual pins & Polled here, the ring buffer of chapter 2 if the block stream grows \\
\bottomrule
\end{tabular}
\caption{Peripheral configuration. The last four rows are the experiment: everything else is scaffolding. Whether the cycle counter runs without a debug probe attached is item 25 of the confirm list and is checked before any number in this chapter is believed.}
\end{table}

\subsection*{Wiring}

\diagram{c18_wiring}{Nothing is wired. The signal is a vector compiled into flash, which is what makes the acceptance test reproducible on a different desk on a different day.}

\subsection*{Memory and timing budget}

\diagram{c18_mem}{The experiment as a grid. Code lives in flash or in instruction-side tightly coupled memory, data lives in the main AXI region or in data-side tightly coupled memory, and each cell is run with the caches on and off. Every cell is a measurement that has not been taken.}

\begin{table}[H]\centering\small
\begin{tabular}{p{44mm}p{28mm}p{28mm}p{30mm}}
\toprule
Quantity & Budget & Measured & Margin \\
\midrule
Flash for the library modules used & 24 kB & not measured & not measured \\
Flash for twiddle and coefficient tables & 16 kB & not measured & not measured \\
Static memory for buffers, float32 & 12 kB & not measured & not measured \\
Static memory for buffers, Q15 & 6 kB & not measured & not measured \\
Cycles per 1024-point real FFT, float32 & 40000 & not measured & not measured \\
Cycles per 1024-point real FFT, Q15 & 40000 & not measured & not measured \\
Cycles per 1024-sample biquad cascade & 20000 & not measured & not measured \\
Block period at the node's own rate & 256 ms & not measured & not measured \\
\bottomrule
\end{tabular}
\caption{The budget table. The two cycle budgets are deliberately identical: this chapter has no prior belief about which format wins on this part, and writing the same number in both rows is the honest statement of that. The instrument for the cycle rows is the core's own cycle counter; the flash and memory rows are read from the linker map file.}
\end{table}

The node's real workload has enormous margin and is not the interesting case. A
256-sample block at the 1 kHz rate proven in chapter 6 arrives every 256 ms,
and even a pessimistic transform finishes in under a millisecond. The
benchmark workload is different on purpose: the same transform repeated back to
back with nothing else running, which is the only way to get a cycles-per-
transform number that means anything. Both are reported, because a reader who
sees only the benchmark will over-engineer and a reader who sees only the node
workload will learn nothing about the part.

\subsection*{Firmware design (UML)}

\diagram{c18_uml}{One case of the acceptance run. The host sends a case identifier, the board runs the pipeline, times it and returns the output block with its cycle count, and the host decides.}

The board side is a table of cases and a loop. A case names a numeric format,
a placement for code, a placement for data and a cache state; the runner
configures what it can configure at run time, runs the pipeline over the
compiled-in vector, reads the cycle counter before and after, and streams the
result block back as hexadecimal words with the case identifier in front of it.
Placement is not a run-time choice, so the placement half of the grid is a
build-time matrix: the same source compiled several times with different
section attributes and linked with different scripts. Saying that plainly in
the design is better than pretending one image can do everything.

\subsection*{Data flow (ASCII)}

\begin{asciiart}
  host                                             board
  +----------------------------------+             +-------------------------------------+
  | scipy.signal.butter -> sos       |             | const float32_t vec[1024]  (.rodata) |
  |   |                              |             |   |                                 |
  |   v  export 5 coeffs per stage   |  build      |   v                                 |
  | coeffs.h --------------------------------->    | biquad cascade (df2T f32 / df1 q15) |
  |                                  |             |   |                                 |
  | numpy float64 reference          |             |   v                                 |
  |   |                              |             | rfft 1024 (f32 / q15)               |
  |   |  cmsisdsp wrapper, same API  |             |   |                                 |
  |   v                              |  USART3     |   v                                 |
  | compare: rel RMS, SNR, tolerance | <---------- | cmplx mag + DWT cycle count         |
  +----------------------------------+             +-------------------------------------+
        pass or fail, printed with the tolerance it was judged against
\end{asciiart}

\subsection*{Repository layout}

\begin{plaincode}
nucleo-h7a3-dsp-acceptance/
  CMakeLists.txt
  cmake/arm-none-eabi.cmake
  ld/placement_flash_axi.ld     # code in flash, data in the AXI region
  ld/placement_itcm_dtcm.ld     # code and data in tightly coupled memory
  third_party/CMSIS-DSP/        # vendored, Apache-2.0, LICENSE kept
  src/dsp_pipeline.c            # biquad + rfft + magnitude, one entry point
  src/dsp_pipeline_q15.c        # the fixed-point twin, same entry shape
  src/cycles.c                  # cycle counter enable and read
  src/cases.c                   # the case table: format, placement, cache
  src/main.c                    # runner and result stream
  gen/coeffs.h                  # generated by tools/design_filter.py
  gen/testvec.h                 # generated by tools/make_vector.py
  tools/design_filter.py        # SciPy design, export in library layout
  tools/make_vector.py          # the vector and the float64 reference
  tools/reference_wrapper.py    # the same pipeline through the Python wrapper
  tools/accept.py               # compares, prints the tolerance, exits nonzero
  results/cycles.csv            # one row per case, cache state in every row
  README.md
\end{plaincode}

\subsection*{Steps}

\begin{steps}
\item \textbf{Vendor the library and build only what you use.} The licence is
Apache-2.0, which means it can be copied into the repository with its licence
file intact. Build it from source rather than linking a prebuilt archive, so
the compiler flags are yours and visible.
\begin{shellcode}
git -C third_party clone --depth 1 https://github.com/ARM-software/CMSIS-DSP
ls third_party/CMSIS-DSP/LICENSE     # keep it, it is a condition of the licence
\end{shellcode}
Add only the source groups the pipeline touches. Pulling the whole library in
costs flash for nothing and makes the size rows of the budget table
meaningless.
\begin{makecode}
set(DSP_SRC
  Source/FilteringFunctions/arm_biquad_cascade_df2T_f32.c
  Source/FilteringFunctions/arm_biquad_cascade_df1_q15.c
  Source/TransformFunctions/arm_rfft_fast_f32.c
  Source/TransformFunctions/arm_rfft_fast_init_f32.c
  Source/ComplexMathFunctions/arm_cmplx_mag_f32.c)
\end{makecode}

\item \textbf{Enable the floating-point unit and the cycle counter, and say
which cache state you are in.} Without the first, every float is a fault.
Without the second there is no instrument. Without the third the numbers are
not comparable to anything, including to themselves a week later.
\begin{ccode}
#include "stm32h7xx.h"      /* CMSIS device header, for SCB and DWT */

void cycles_init(void)
{
    SCB->CPACR |= (3UL << 20) | (3UL << 22);    /* CP10 and CP11 full access */
    __DSB(); __ISB();
    CoreDebug->DEMCR |= CoreDebug_DEMCR_TRCENA_Msk;
    DWT->CYCCNT = 0;
    DWT->CTRL  |= DWT_CTRL_CYCCNTENA_Msk;
}

static inline uint32_t cycles_now(void) { return DWT->CYCCNT; }
\end{ccode}
Confirm on the bench that the counter increments with no debug probe attached.
If it does not, this chapter loses its cheapest instrument and the fallback is
a general-purpose timer at the core clock divided down, which is coarser and
must be said so in the results.

\item \textbf{Design the filter on the host.} The library does not do design,
and that is not a shortcoming: design is a modelling decision and belongs where
the model is. A fourth-order Butterworth low pass as second-order sections is
enough to make the point.
\begin{pycode}
import numpy as np
from scipy import signal

FS, FC, ORDER = 1000.0, 40.0, 4
sos = signal.butter(ORDER, FC / (FS / 2.0), btype="low", output="sos")

# CMSIS biquads want 5 coefficients per stage: b0, b1, b2, -a1, -a2.
# SciPy gives 6 per section: b0, b1, b2, 1, a1, a2. The sign of the
# denominator terms is the single most common porting mistake.
rows = [(s[0], s[1], s[2], -s[4], -s[5]) for s in sos]
with open("gen/coeffs.h", "w") as f:
    f.write("#define NSTAGE %d\n" % len(rows))
    f.write("static const float32_t coeffs[5 * NSTAGE] = {\n")
    for r in rows:
        f.write("    %.9ef, %.9ef, %.9ef, %.9ef, %.9ef,\n" % r)
    f.write("};\n")
\end{pycode}

\item \textbf{Generate the vector and the reference output.} The vector is
compiled into flash so the test is reproducible; the reference is computed in
double precision so the comparison has somewhere to stand.
\begin{pycode}
N = 1024
t = np.arange(N) / FS
x = (np.sin(2 * np.pi * 12.0 * t)
     + 0.35 * np.sin(2 * np.pi * 180.0 * t)
     + 0.02 * np.random.default_rng(20260920).standard_normal(N))

y_ref = signal.sosfilt(sos, x.astype(np.float64))
mag_ref = np.abs(np.fft.rfft(y_ref))
np.savez("gen/reference.npz", x=x, y=y_ref, mag=mag_ref)
\end{pycode}
The seed is fixed and written down. A test vector that changes between runs
turns a regression into an argument.

\item \textbf{Run the same pipeline through the library's Python wrapper.}
This is the step that saves an evening. The wrapper's API mirrors the C API, so
if the wrapper agrees with the double-precision reference and the board does
not, the problem is in the port and not in the arithmetic.
\begin{pycode}
import cmsisdsp as dsp

biquad = dsp.arm_biquad_casd_df1_inst_f32()
dsp.arm_biquad_cascade_df1_init_f32(biquad, len(rows),
                                    np.array(rows, dtype=np.float32).flatten(),
                                    np.zeros(4 * len(rows), dtype=np.float32))
y_wrap = dsp.arm_biquad_cascade_df1_f32(biquad, x.astype(np.float32))
print("wrapper against float64 reference:",
      np.sqrt(np.mean((y_wrap - y_ref) ** 2)) / np.sqrt(np.mean(y_ref ** 2)))
\end{pycode}
Read the exact instance and initialiser spelling from the wrapper's own
documentation and from the learning path rather than from memory: the names
track the C names closely but not blindly, and the transposed direct form has a
different state size from the direct form.

\item \textbf{Write the board pipeline once, twice.} One entry point, two
implementations, the same shape. The float32 side uses the transposed direct
form, which needs two state words per stage. The Q15 side uses the direct form,
which needs four, and carries a post shift because the coefficients live in
1.15 with headroom.
\begin{ccode}
#include "arm_math.h"
#include "coeffs.h"

static arm_biquad_cascade_df2T_instance_f32 bq;
static float32_t bq_state[2 * NSTAGE];
static arm_rfft_fast_instance_f32 fft;
static float32_t spec[1024];        /* packed real spectrum, in place */
static float32_t mag[513];

void dsp_init(void)
{
    arm_biquad_cascade_df2T_init_f32(&bq, NSTAGE, coeffs, bq_state);
    arm_rfft_fast_init_f32(&fft, 1024);
}

uint32_t dsp_run(const float32_t *in, float32_t *out_mag)
{
    uint32_t t0 = cycles_now();
    arm_biquad_cascade_df2T_f32(&bq, in, spec, 1024);
    arm_rfft_fast_f32(&fft, spec, spec, 0);   /* 0 = forward */
    /* The packed layout puts the two purely real bins first. */
    out_mag[0]   = fabsf(spec[0]);
    out_mag[512] = fabsf(spec[1]);
    arm_cmplx_mag_f32(&spec[2], &out_mag[1], 511);
    return cycles_now() - t0;
}
\end{ccode}
The packed output of the real transform catches everyone once. The first two
floats are not a complex pair: they are the zero-frequency bin and the Nyquist
bin, both purely real, packed together to keep the output the same length as
the input.

\item \textbf{Report a cycle count that carries its conditions.} A bare number
is not a measurement. Each row names the format, where the code is, where the
data is and what both caches were doing.
\begin{ccode}
printf("case=%s fmt=%s code=%s data=%s icache=%s dcache=%s cycles=%lu\n",
       c->name, c->fmt, c->code_where, c->data_where,
       c->icache ? "on" : "off", c->dcache ? "on" : "off",
       (unsigned long) cycles);
\end{ccode}
Run each case at least thirty-two times and report the median rather than the
minimum. The minimum is the best case with a warm cache and flatters exactly
the configurations this chapter is trying to tell apart.

\item \textbf{Move the code and the data.} This is the half of the experiment
that the run-time switches cannot reach. Give the pipeline its own sections and
let the linker script decide where they land.
\begin{ccode}
__attribute__((section(".itcm_text"))) uint32_t dsp_run(const float32_t *in,
                                                        float32_t *out_mag);
__attribute__((section(".dtcm_bss"))) static float32_t spec[1024];
\end{ccode}
\begin{ldcode}
SECTIONS
{
  .itcm_text : { *(.itcm_text*) } > ITCM AT> FLASH
  .dtcm_bss (NOLOAD) : { *(.dtcm_bss*) } > DTCM
}
\end{ldcode}
Every region name, origin and length in that script is checked against RM0455
for this part. The application note that inspired this experiment ships linker
scripts for the family, and they are read as a template and retyped against the
right manual, never inherited.

\item \textbf{Write the acceptance test and give it a tolerance before you run
it.} The tolerance is a decision, not an observation, and writing it after the
fact is how a test becomes a rubber stamp.
\begin{pycode}
TOL_F32_REL_RMS = 1e-5      # float32 path against the float64 reference
TOL_Q15_SNR_DB  = 55.0      # Q15 path, after undoing the transform's scaling

def rel_rms(a, b):
    return float(np.sqrt(np.mean((a - b) ** 2)) / np.sqrt(np.mean(b ** 2)))

def snr_db(a, b):
    err = a - b
    return float(10.0 * np.log10(np.sum(b ** 2) / np.sum(err ** 2)))

ok = rel_rms(board_f32, mag_ref) <= TOL_F32_REL_RMS
print("float32 rel RMS %.3e, tolerance %.3e -> %s"
      % (rel_rms(board_f32, mag_ref), TOL_F32_REL_RMS, "pass" if ok else "fail"))
\end{pycode}
The Q15 comparison needs one extra step that the float comparison does not: the
fixed-point real transform scales its output down as it goes, stage by stage,
so the board result is rescaled before it is compared. Doing that rescaling
inside the acceptance script, where it is visible, is better than hiding it in
the firmware where it becomes folklore.

\item \textbf{Collect the grid.} Eight cases at minimum: two formats, two code
placements, two data placements, with the cache state swept inside each. Put
them in a CSV file and let a script make the table, so the table cannot drift
from the data.
\begin{shellcode}
python tools/accept.py --port /dev/ttyACM0 --cases all --out results/cycles.csv
python tools/table.py results/cycles.csv > results/table.md
\end{shellcode}
\end{steps}

\subsection*{Build, flash and debug}

\diagram{c18_timing}{The two workloads on one time axis. The node workload leaves most of the block period idle; the benchmark workload removes the idle so that a cycles-per-transform figure means something. Every interval marked with a question mark is unmeasured.}

Four builds come out of one source tree, one per placement combination, and the
build system names them rather than leaving the reader to guess which image is
on the board.

\begin{shellcode}
cmake -B build-fw/flash_axi -G Ninja -DPLACEMENT=flash_axi \
  -DCMAKE_TOOLCHAIN_FILE=cmake/arm-none-eabi.cmake
cmake --build build-fw/flash_axi -j
arm-none-eabi-size -A build-fw/flash_axi/firmware.elf | grep -E 'itcm|dtcm|text'
cp build-fw/flash_axi/firmware.bin "$PROBE_DISK"/   # onto the probe disk
\end{shellcode}

When a placement build produces a board that does not start, the cause is
almost always the linker script rather than the code: a region that does not
exist on this part, a load address that was never set so the tightly coupled
code is never copied out of flash, or a section marked to be loaded when it
should not be. Recovery is the usual sequence, namely hold the reset, connect
under reset, and flash a known good image; the drag-and-drop disk of chapter 1
is the fastest way back to a working board when the debug connection itself has
become unreliable.

\begin{note}{When the fast placement is slower}
If moving the pipeline into tightly coupled memory makes it slower rather than
faster, check three things in order. The instruction cache may have been left
on, in which case flash was already fast and you have added a copy at startup
for nothing. The data may still be in the main region while only the code
moved, so every sample still crosses the bus matrix. Or the transform's tables,
which are constant data in flash, may now be the slowest thing in the loop
because everything else moved away from them. The third case is the interesting
one and it is why the table has a column for where the data is.
\end{note}

\subsection*{Verification and acceptance criteria}

\begin{itemize}
\item The float32 pipeline matches the double-precision host reference on the
shipped vector to a relative root-mean-square error at or below the tolerance
written in the script before the first run, and the script prints both the
value and the tolerance whether it passes or fails.
\item The Q15 pipeline reaches at least the stated signal to error ratio on the
same vector after the transform's own scaling has been undone, and the
rescaling is done in the visible part of the code.
\item The library's Python wrapper and the board agree to within the float32
tolerance on the same input, which separates a porting error from an arithmetic
one. A failure here names which of the two it is.
\item Every cycle figure carries its format, its code placement, its data
placement and the state of both caches. A row missing any of the four is not
reported.
\item Each case is run at least thirty-two times and the median is reported,
with the minimum and maximum alongside it, using the core's cycle counter as
the instrument. If the counter is found not to run without a probe attached,
the substitute instrument is named in the same table.
\item The size rows of the budget table are filled from the linker map file and
reconcile with the size output, in the same way as chapter 1.
\item The comparison script exits with a nonzero status on failure, so it can
become a regression test in chapter 12's rig without being rewritten.
\end{itemize}

\subsection*{Variants}

\begin{table}[H]\centering\small
\begin{tabular}{p{24mm}p{26mm}p{44mm}p{22mm}p{20mm}}
\toprule
Axis & Variant & What changes & Cost & Built in full in \\
\midrule
Execution model & Polled loop & The baseline here. The pipeline runs when called and nothing else is happening, which is what makes the cycle count clean & Nothing else runs & Chapter 3 \\
Execution model & DMA double buffered & Real samples arrive while the previous block is still being transformed, which is the shape any deployed version takes & A coherency problem, which is chapter 19 & Chapter 6 \\
Operating system & An RTOS task & The pipeline becomes one task among three and the cycle count acquires a distribution rather than a value & Scheduling jitter enters the measurement & Chapter 20 \\
Language & C++ & Templates over the block length and the numeric format remove the duplicated entry point, at the cost of reading generated assembly to confirm nothing was added & A decision to justify & Chapter 9 \\
Language & Rust & The computation moves to a pure-Rust transform crate or to bindings over this same library, and the acceptance script does not change at all & The library's test coverage is left behind & Chapter 17 \\
Language & MicroPython & The design and the reference can run on the board itself for exploration, which is pleasant and is not the measurement & Two orders of magnitude of speed & Chapter 1 and 9 \\
Peripheral substitution & The transform on the host & Send the block over the link and let the host compute, which is the honest comparison for a node with a host twenty centimetres away & Link energy replaces compute energy & Chapter 16 \\
Intelligence and reach & Transform plus a small statistical model & The spectrum feeds a decision tree or an anomaly score rather than a display, which for most sensor problems on this part beats a neural network per kilobyte & Accuracy on hard problems & Chapter 16 \\
\bottomrule
\end{tabular}
\caption{Variants for chapter 18. Only the first row is built here; the rest are named so that a reader can follow one axis across the book rather than reading front to back.}
\end{table}

\subsection*{Pitfalls}

\begin{itemize}
\item The denominator sign. The design tool gives the transfer function's
coefficients; the library wants the two feedback coefficients negated. A filter
built from unnegated coefficients does not fail loudly. It rings, or it grows
without bound after a few hundred samples, which looks like a data problem.
\item The state buffer size. The transposed direct form needs two words per
stage and the direct form needs four. Getting it wrong overwrites whatever is
next in memory, and the symptom appears in an unrelated variable.
\item The packed real spectrum. The first two outputs are the zero-frequency
and Nyquist bins, not a complex pair. Treating them as a pair puts a large
false value in the first magnitude bin and shifts everything after it.
\item The fixed-point transform's internal scaling. It divides as it goes. A
comparison against an unscaled reference reports a failure that is arithmetic
doing exactly what it documents.
\item Reporting the minimum cycle count. The minimum is the warm-cache case and
it hides the differences this chapter exists to expose. Report the median with
the range.
\item Measuring with the debugger halted between runs, or with a breakpoint
inside the timed region. Both perturb the cache and the pipeline.
\item Quoting a fixed-point against floating-point ratio measured on a
different core. It is the single most repeated error in this subject area.
\item Inheriting a linker script for the tightly coupled regions from the
better-known member of this family. The regions differ, and the failure looks
like a hardware fault.
\item Feeding the pipeline from a buffer that a transfer engine just filled
without the maintenance step. That is chapter 19, and it produces a spectrum of
stale data that looks entirely plausible.
\end{itemize}

\subsection*{Best practices applied}

\begin{itemize}
\item The tolerance is written before the first measurement and printed on
every run, pass or fail.
\item Two independent references, one in double precision and one through the
same library on the host, so a failure is localised rather than merely
detected.
\item Every measurement names its instrument and its conditions, and a
condition that is unknown is written as unknown rather than omitted.
\item A maintained, permissively licensed library does the arithmetic, and the
chapter says clearly what the library does not do, which is design.
\item The experiment's build matrix is in the build system, not in a comment,
so a result can be reproduced by name.
\item The test vector is fixed, seeded and committed, which is what makes this
a regression test rather than a demonstration.
\end{itemize}

\subsection*{Stretch goals}

\begin{itemize}
\item Add the optimised neural kernels from the same maintainer and measure a
small classifier on the spectrum. This core has the signal-processing
instruction extension but no vector extension, so the speed-ups reported for
later cores do not apply here; measuring that gap is a result worth having.
\item Compare the spectrum plus a small statistical model against a neural
network of the same flash budget on one real task, which is the argument
chapter 16 tests and the thesis this book returns to.
\item Take the window function and the magnitude calculation out of the loop by
folding the window into the biquad's final stage, and report whether the saving
survives the placement sweep.
\item Run the same acceptance script against the Rust implementation of chapter
17 without changing the script, which proves the acceptance test is about the
numbers and not about the language.
\item Extend the grid with the block length as a third axis, from 128 to 4096
samples, and plot cycles per sample rather than cycles per transform.
\end{itemize}

\subsection*{Roadmap and next steps}

Chapter 19 is the direct sequel and should be read next by anyone who intends
to feed this pipeline from a transfer engine rather than from a compiled-in
array, because the cache behaviour this chapter treats as a performance
variable becomes a correctness problem the moment a second bus master writes
the input buffer.

For going deeper, the maintainer's own learning path on building a pipeline in
Python against a reference implementation is the shortest route to fluency with
the wrapper, and it is written by the person who maintains both sides. The
vendor's application note on system architecture and performance is the model
for the placement grid; read it for its method and not for its numbers, which
belong to a different part. For the argument that a transform plus a
statistical model beats a small neural network on accuracy per kilobyte, there
is one peer-reviewed paper making the case directly, and the open textbook that
grew out of a university course is the closest thing to a canonical curriculum
for the machine-learning chapters that follow. Appendix H lists the courses and
their licences, including the two whose terms forbid reproduction in a
commercial work.

\subsection*{Portfolio evidence}

\begin{itemize}
\item A public repository whose README opens with the placement grid and states
the tolerance in the first paragraph, before any performance claim.
\item The acceptance output for both numeric formats, showing the measured
value next to the tolerance it was judged against.
\item The cycle table with one row per case, every row naming format,
placements and both cache states, generated from the CSV file rather than typed.
\item A short written note on what the placement sweep found, including the
case where the expected answer did not hold, because a sweep that confirms
everything was not a sweep.
\item The generated coefficient header next to the design script that produced
it, so a reader can regenerate rather than trust.
\end{itemize}

\subsection*{Sources}

Normative references:
\begin{itemize}
\item Reference manual RM0455, for the memory regions, their addresses and
which bus each sits behind. Not RM0433, which is its sibling.
\item The STM32H7A3xI datasheet, for the flash and memory sizes, and for
whether the floating-point unit on this part is double precision.
\item The Armv7-M architecture reference manual and the Cortex-M7 technical
reference manual, for the cycle counter, the caches and the signal-processing
instruction set extension.
\item The vendor's application note on system architecture and performance, for
the benchmark method and for the shape of the placement result. It does not
cover this part.
\end{itemize}

Reusable implementations:
\begin{itemize}
\item The computational library, Apache-2.0, vendored with its licence
file.\newline
\url{https://github.com/ARM-software/CMSIS-DSP}
\item The learning path that builds a pipeline in Python against a reference
implementation, written by the library's maintainer.\newline
\url{https://learn.arm.com/learning-paths/embedded-and-microcontrollers/cmsisdsp-dev-with-python/}
\item The optimised neural kernels for this architecture, Apache-2.0, which
reach this core through the signal-processing extension rather than through
vectors.\newline
\url{https://github.com/ARM-software/CMSIS-NN}
\item The small-model library for the statistical alternative, MIT, which does
decision trees, naive Bayes and anomaly detection from a very small flash
budget.\newline
\url{https://github.com/emlearn/emlearn}
\item The peer-reviewed argument for classical features over small neural
networks on parts of this class.\newline
\url{https://arxiv.org/pdf/2107.09448}
\end{itemize}
   % Transforms and filters with a numerical acceptance test
\project{19}{Caches, the MPU, and why DMA reads stale bytes}{NUCLEO-H7A3ZI-Q + DFRobot SEN0032 (ADXL345)}{D-cache maintenance, MPU regions, which engine reaches which memory}

\begin{keyfacts}
\item[Board] NUCLEO-H7A3ZI-Q with the ADXL345 breakout on the Arduino header, wired with jumpers
\item[Peripherals] SPI with transfers in both directions, the main transfer engine, the L1 data cache, the memory protection unit, GPIO for chip select and the sensor's interrupt line
\item[Toolchain] arm-none-eabi-gcc with CMake, a debug probe that can halt and dump memory, and a host script that counts failures
\item[Operating system] Bare metal
\item[Difficulty] 5 of 5
\item[Effort] 4 evenings of about four hours
\item[Deliverable] A repository that reproduces silent data corruption on purpose at a measured rate, removes it three separate ways, and ships a small declarative region table with cache attributes, which nothing permissive and maintained currently provides
\end{keyfacts}

\subsection*{Why this project}

This is the trap that costs this family of parts more engineering time than
anything else, and it is the one a reader is most likely to meet in a hurry, at
the end of a week, in somebody else's code. A transfer engine writes a buffer, the
processor reads it, and the values are wrong. Not obviously wrong: plausibly
wrong, sometimes, on some builds, and never under the debugger. The cause is
that the processor has a data cache and the transfer engine does not go through
it, so the two have different opinions about what is in memory. The fix is
well known and the failure is still shipped regularly, including in the silicon
vendor's own example projects, where it has been an open report for a long time,
and in at least two widely used open projects.

So the chapter does not begin with the fix. It begins by producing the failure
deliberately and measuring how often it happens, because a fault that cannot be
reproduced cannot be proved fixed. The harness fills the receive buffer with a
sentinel byte, starts the transfer, waits for the completion flag, and then
counts how many of the sentinel bytes survived into what the processor reads. A
surviving sentinel means the processor read a cache line that memory has since
moved on from. Run that a thousand times and the result is a rate, not an
anecdote.

Then it is removed three ways, and the three are different engineering
decisions rather than three spellings of one. Maintenance operations put the
burden at every call site and are exact. A protection-unit region marked
non-cacheable puts the burden in one table and costs performance on every
access to that memory. Placement decides the question once in the memory map.
There is also a fourth failure that is not a coherency problem at all and is
the most common version of this bug in this family: a buffer placed in a memory
the transfer engine cannot reach, which looks exactly like a configuration
error.

\begin{note}{The domains are not the ones in the popular manual}
The three-domain vocabulary that most writing about this family uses belongs to
RM0433 and the STM32H743. This part is documented by RM0455 and names its
domains differently. A blog post that tells you to put a buffer in "D3 SRAM"
is describing a memory this part does not call by that name, and following it
literally produces a linker script that either fails to link or places a buffer
somewhere useful only by accident. Every region name, base address and length
in this chapter is checked against RM0455 before it is written into a script.
\end{note}

\subsection*{Prior art and what to reuse}

\begin{table}[H]\centering\small
\begin{tabular}{p{34mm}p{46mm}p{38mm}p{20mm}}
\toprule
Source & What it gives & What it does not & Licence \\
\midrule
The vendor's application notes on the L1 cache, on the memory protection unit, and on using the cache for performance & The normative account: what the cache does, how regions are described, which attribute combinations are legal, and what each maintenance operation means & They are written for the family, so the region inventory and the domain names are checked against the manual for this part before use & Vendor documentation, read do not copy \\
A community article by one of the vendor's own engineers on using the protection unit to manage coherency & Exactly the three fixes, with complete configuration code for each, which makes it the single best starting page on the subject & It is one part's configuration, not a reusable component, and there is no build system or test around it & Community post, cite and re-derive \\
The same vendor's write-up of why transfers fail on this family & The fourth failure, which the coherency literature usually omits: a buffer in tightly coupled memory that the main engines cannot reach at all & It is a troubleshooting page rather than a design method & Community post, cite and re-derive \\
Three application briefs from another silicon vendor & The same three fixes, mapped one to one, and clearer than the originals because each brief covers exactly one fix & Written for that vendor's parts, so the register names differ & Vendor documentation, read do not copy \\
A three-part written series on the Cortex-M7 cache & The best teaching treatment in print: cache basics, then the coherency problem, then the maintenance operations, all at the architecture level where it is all true & No board, no build, no measurement & Blog series, cite \\
One video walking this through on an H7 with a non-cacheable region and a custom linker section & Proof that the linker-section approach works end to end, and a sanity check on the section attributes & Not this part, and a video is not a reference & Video, cite \\
\bottomrule
\end{tabular}
\caption{Prior art for chapter 19. This is the richest pool in the book, and the reason the chapter still has work to do is at the end of the next paragraph.}
\end{table}

What is left to write is the part every one of those sources leaves out. None
of them measures anything: they explain the failure and then assert the fix. So
the harness, the failure rate and the thousand-run comparison are new. And no
maintained permissively licensed library offers a declarative table of memory
regions with their cache attributes, checked at build time for the alignment
and size rules the protection unit imposes. That is a small, sharp, useful
piece of software that does not exist, it is about two hundred lines, and it is
this chapter's deliverable.

\subsection*{Parts from the inventory}

\begin{table}[H]\centering\small
\begin{tabular}{p{40mm}p{66mm}p{40mm}}
\toprule
Part & Role & Interface \\
\midrule
NUCLEO-H7A3ZI-Q & The processor, its caches, its protection unit and its transfer engines & Micro USB to the host \\
DFRobot SEN0032 (ADXL345) & Supplies a stream of real bytes at a known rate, with a 32-sample first in first out buffer that makes a burst transfer natural & SPI at up to 5 MHz, or I2C. 3.3 V only \\
Jumper wires & Six connections. Confirm that male to male wires exist in the drawer before planning the evening & Arduino header to breakout \\
A USB data cable & Power, programming and the failure count & Micro USB \\
Host PC & Runs the thousand iterations and tabulates the rate & Python over the virtual COM port \\
\bottomrule
\end{tabular}
\caption{Inventory items used in chapter 19. Nothing is bought. The sensor is chosen because it produces bursts, not because the acceleration matters; the experiment is about the bytes.}
\end{table}

\subsection*{System architecture}

\diagram{c19_arch}{Who reaches what. The processor sees memory through two caches for everything on the far side of the bus matrix and through neither for the tightly coupled memories. The transfer engines sit on the matrix and see memory directly, which is the whole problem in one picture.}

Read the figure as a reachability question rather than a performance one. The
processor's path to the tightly coupled memories does not pass the caches, which
is why a buffer there is never stale and why it is also unreachable by the main
transfer engines. Everything else the processor touches goes through the data
cache. The transfer engines are separate masters on the matrix: they neither
fill the cache nor look in it. Two masters, two views, no hardware arbitration
between them on this core. Once that is stated plainly the three fixes become
three places to put a decision.

\subsection*{Peripheral configuration}

\begin{table}[H]\centering\small
\begin{tabular}{p{22mm}p{26mm}p{24mm}p{30mm}p{30mm}}
\toprule
Peripheral & Mode & Clock source & Pins and function & Interrupt and transfers \\
\midrule
SPI on the Arduino header & Master, mode 3, 5 MHz maximum for this sensor & Peripheral bus & D11 to D13, instance confirmed in the board manual & Two transfer streams, one each way \\
GPIO chip select & Push-pull output, software controlled & Peripheral bus & D10, confirm & None \\
GPIO sensor interrupt & Input with edge detect & Peripheral bus & D2, confirm & External interrupt line \\
Main transfer engine & Peripheral to memory and memory to peripheral, single buffer & Bus clock & None & Complete and error interrupts \\
L1 data cache & On for the experiment. Off is one of the controls & Core clock & None & Write-back, write-allocate by default \\
Memory protection unit & One region per entry in the region table & Core clock & None & Memory management fault enabled \\
USART3 & Asynchronous, 115200 8N1 & Peripheral bus & Board manual pins & Carries the failure count \\
\bottomrule
\end{tabular}
\caption{Peripheral configuration. Enabling the memory management fault rather than leaving it to escalate is deliberate: a region table with a bad alignment should stop the board at the offending access with a named fault, not somewhere else later.}
\end{table}

\begin{note}{Three point three volts only}
The breakout is wired to the 3V3 pin and to ground, never to 5V, and no signal
from it reaches a processor pin at more than 3.3 V. Nothing on this bench can
be soldered, so if the breakout's bus selection turns out to need a solder
bridge moved, the chapter uses the other bus and says so rather than inventing
a workaround.
\end{note}

\subsection*{Wiring}

\diagram{c19_wiring}{Six wires. The interrupt line is not optional here: it is what makes the burst arrive on the sensor's schedule rather than the processor's, which is what gives the failure its intermittent character.}

\begin{table}[H]\centering\small
\begin{tabular}{p{34mm}p{34mm}p{44mm}p{40mm}}
\toprule
Breakout pin & Board pin & Function & Note \\
\midrule
VCC & 3V3 & Supply & Never 5V \\
GND & GND & Return & Common return, one wire \\
SCL or SCK & D13 & Serial clock & Instance confirmed in the board manual \\
SDA or SDI & D11 & Master out, slave in & Mode 3 for this sensor \\
SDO & D12 & Master in, slave out & Leave unconnected only in three-wire mode \\
CS & D10 & Chip select & Software controlled, not the peripheral's own \\
INT1 & D2 & Watermark interrupt & Rising edge, confirm the polarity bit \\
\bottomrule
\end{tabular}
\caption{Wiring table. Every row whose pin mapping comes from the Nucleo-144 convention rather than from the board manual is confirmed before the first power-on, which is confirm item 12 of the authoring list.}
\end{table}

\subsection*{Memory and timing budget}

\diagram{c19_mem}{The memory map, which is this chapter's silicon figure and the reference for the whole book. Sizes marked for confirmation are read from RM0455 and the datasheet for this part before any linker script uses them. The dashed outline is one protection-unit region.}

\begin{table}[H]\centering\small
\begin{tabular}{p{44mm}p{28mm}p{28mm}p{30mm}}
\toprule
Quantity & Budget & Measured & Margin \\
\midrule
Receive buffer, 32 samples of 6 bytes & 192 B, padded to 256 B & not measured & not measured \\
Alignment required for maintenance & 32 B, both ends & not measured & not measured \\
Protection-unit regions used & 4 of the available set & not measured & not measured \\
Failure rate, naive variant, 1000 runs & no budget, this is the result & not measured & not measured \\
Failure rate, each of the three fixes & 0 in 1000 runs & not measured & not measured \\
Added cycles per transfer, maintenance & 600 & not measured & not measured \\
Added cycles per transfer, non-cacheable region & 2000 & not measured & not measured \\
Region table code size & 2 kB & not measured & not measured \\
\bottomrule
\end{tabular}
\caption{The budget table. The two cycle rows are the honest cost of the two fixes and are the reason the chapter does not simply recommend one: maintenance is cheap and easy to forget, a non-cacheable region is impossible to forget and is not cheap. The instrument for both is the core's cycle counter, as in chapter 18.}
\end{table}

The alignment row is the second bug, the one that arrives after the first is
fixed. Maintenance acts on whole cache lines, and a line is 32 bytes on this
core. Invalidating a buffer that is not aligned to 32 bytes, or whose length is
not a multiple of 32, also invalidates the neighbouring bytes that share the
first and last lines, and any processor writes to those neighbours that have
not reached memory are lost. The symptom is corruption in an unrelated
variable.

\subsection*{Firmware design (UML)}

\diagram{c19_uml}{One run of the harness as a state machine. The four variants differ only in which of the three shaded states is entered, which is what makes the failure counts comparable.}

The harness is deliberately boring. Fill the buffer with the sentinel, arm the
sensor's watermark interrupt, start the transfer when it fires, wait for
completion, then read the buffer and count surviving sentinel bytes. The only
difference between the four variants is which optional state the run passes
through: none at all, an invalidate, nothing because the region is
non-cacheable, or nothing because the buffer lives in a memory that is not
cached. Keeping everything else identical is what lets the four failure counts
be compared without argument.

\subsection*{Data flow (ASCII)}

\begin{asciiart}
    processor                 L1 data cache            SRAM                 transfer engine
  +-------------+          +----------------+     +---------------+       +-----------------+
  | fill 0xA5   | -------> | line dirty A5  | --> | A5 A5 A5 A5   |       |                 |
  +-------------+          +----------------+     +---------------+       +-----------------+
                                                          ^                        |
  sensor watermark fires, the engine writes the FIFO burst |<-----------------------+
                                                          v
  +-------------+          +----------------+     +---------------+
  | read buffer | <------- | line STALE A5  |  x  | 1F 22 03 FE   |   the engine never
  +-------------+          +----------------+     +---------------+   touched the cache
        |
        +--> naive: sees A5, the run is counted as a failure
        +--> invalidate by address first: the line is dropped, the read reaches SRAM
        +--> region marked non-cacheable: there was never a line to go stale
        +--> buffer in DTCM: no cache on that path, and the main engine cannot reach it
\end{asciiart}

\subsection*{Repository layout}

\begin{plaincode}
nucleo-h7a3-cache-coherency/
  CMakeLists.txt
  cmake/arm-none-eabi.cmake
  ld/stm32h7a3zi.ld              # carries the .dma_buffers and .noncacheable sections
  include/mpu_table.h            # the declarative region table, the deliverable
  src/mpu_table.c                # one call configures every region in the table
  src/mpu_static_assert.h        # build-time checks: power of two, aligned to size
  src/adxl345.c                  # FIFO setup, watermark, burst read
  src/variant_naive.c            # no maintenance, the control
  src/variant_maintain.c         # clean before, invalidate after
  src/variant_region.c           # buffer pool inside a non-cacheable region
  src/variant_placement.c        # buffer in a memory the map already settles
  src/harness.c                  # 1000 runs, sentinel count, result line
  tools/run_matrix.py            # drives all four, tabulates, exits nonzero
  results/failure_rates.csv
  docs/region_map.md             # the table in prose, generated from the header
  README.md
\end{plaincode}

\subsection*{Steps}

\begin{steps}
\item \textbf{Bring the sensor up with no transfer engine at all.} Polled
register reads first, so that a later failure is known to be about coherency
and not about the bus. Confirm the device identification register reads the
documented value before anything else happens.
\begin{ccode}
uint8_t who = adxl_read_reg(0x00);      /* DEVID */
if (who != 0xE5) {                      /* documented value for this part */
    printf("sensor not answering: 0x%02X\n", who);
    for (;;) { }
}
\end{ccode}

\item \textbf{Configure the FIFO so a burst is natural.} The sensor holds 32
samples. Set the watermark, enable the interrupt on the first line, and let the
sensor decide when the burst is ready. A burst on the sensor's schedule is what
gives the failure its intermittent character, and an intermittent failure is
the one worth reproducing.

\item \textbf{Write the harness and make it fail.} This is the step that is
usually skipped and it is the one that makes the rest defensible.
\begin{ccode}
#define SENTINEL 0xA5u
static uint8_t rx_buf[256] __attribute__((aligned(32),
                                          section(".dma_buffers")));

static unsigned one_run(void)
{
    memset(rx_buf, SENTINEL, sizeof rx_buf);   /* dirties the lines */
    adxl_start_burst_dma(rx_buf, 192);
    while (!dma_complete_flag) { }
    dma_complete_flag = 0;

    unsigned survivors = 0;                    /* no maintenance here */
    for (unsigned i = 0; i < 192; i++)
        if (rx_buf[i] == SENTINEL) survivors++;
    return survivors;
}
\end{ccode}
A run with any survivor is a failure. Report the count of failing runs out of
a thousand and the distribution of survivors, because a run where all 192 bytes
survived and a run where 32 survived are different stories: the first is a
whole buffer of stale lines, the second is one line that was not evicted.

\item \textbf{Do not believe a run under the debugger.} Halting the processor
changes the cache. So does a breakpoint, and so does reading the buffer from
the debug view, which can fill lines that would not otherwise be filled. The
harness runs free and reports over the serial port. The debugger is used to
look at the state after a failure, never to produce the count.

\item \textbf{Fix it the first way: maintenance operations at the call site.}
Clean before the engine reads memory, invalidate after the engine has written
it. Both take an address and a length, and both round outward to whole lines,
which is why the alignment matters.
\begin{ccode}
/* memory to peripheral: the engine must see what the processor wrote */
SCB_CleanDCache_by_Addr((uint32_t *) tx_buf, (int32_t) tx_len);
dma_start_tx(tx_buf, tx_len);

/* peripheral to memory: the processor must not see what it cached before */
dma_start_rx(rx_buf, rx_len);
wait_complete();
SCB_InvalidateDCache_by_Addr((uint32_t *) rx_buf, (int32_t) rx_len);
\end{ccode}
Both buffers are declared with an alignment of 32 and a size rounded up to a
multiple of 32. Add a build-time check so the rule cannot be broken quietly.
\begin{ccode}
_Static_assert(sizeof rx_buf % 32u == 0u, "buffer must be whole cache lines");
_Static_assert(((uintptr_t) &rx_buf[0] % 32u) == 0u, "buffer must be aligned");
\end{ccode}

\item \textbf{Fix it the second way: a region marked non-cacheable.} The
protection unit describes a region by base address, size and attributes. Normal
memory that is not cached is one attribute combination; device memory is
another and is stricter than needed here. The rules that catch people are that
the size is a power of two, that the base is aligned to the size, and that a
later region wins where two overlap.
\begin{ccode}
#include "mpu_table.h"

static const mpu_region_t regions[] = {
  /* name             base          size        attributes            */
  { "flash",          0x08000000u,  MPU_2MB,    MPU_NORMAL_WT_RO      },
  { "axi_sram",       0x24000000u,  MPU_1MB,    MPU_NORMAL_WB_RW      },
  { "dma_buffers",    0x24080000u,  MPU_32KB,   MPU_NORMAL_NONCACHE   },
  { "peripherals",    0x40000000u,  MPU_512MB,  MPU_DEVICE_NX         },
};

void board_mpu_init(void) { mpu_apply(regions, MPU_COUNT(regions)); }
\end{ccode}
That header is the deliverable. It is small, it is declarative, it checks its
own alignment at build time, and nothing permissive and maintained currently
offers it. Every base address and length in the array above is a placeholder
until it has been checked against RM0455 for this part, and the header's
comment says so at the top so that a reader cannot copy it by accident.

\item \textbf{Fix it the third way: decide it once in the memory map.} Give the
buffer pool a section of its own, place that section in a known area, and let
the region table above cover exactly that area. Now no call site has to
remember anything, and a new buffer joins the scheme by being declared in the
right section.
\begin{ldcode}
MEMORY
{
  FLASH      (rx)  : ORIGIN = 0x08000000, LENGTH = 2048K
  AXI_SRAM   (rwx) : ORIGIN = 0x24000000, LENGTH = 512K   /* confirm RM0455 */
  DMA_POOL   (rw)  : ORIGIN = 0x24080000, LENGTH = 32K    /* confirm RM0455 */
}
SECTIONS
{
  .dma_buffers (NOLOAD) : ALIGN(32) { *(.dma_buffers*) } > DMA_POOL
}
\end{ldcode}
Say plainly what this fix does and does not do. Placement alone does not make
anything coherent: with the data cache on, the default memory map treats every
one of these memories as cacheable write-back. What placement buys is a
boundary the region table can describe in one line, so the decision lives in
the map instead of at four hundred call sites.

\item \textbf{Meet the fourth failure on purpose.} Move the buffer into the
data-side tightly coupled memory and watch the transfer fail to happen at all.
This is not a coherency problem: that memory is not cached from the processor's
side, so the bytes would have been correct. The main transfer engines simply
cannot reach it.
\begin{ccode}
/* This compiles, links, and does not work. It is the most common version
   of this bug in this family. The master engine can reach this memory;
   the engine driving the serial peripheral cannot. */
static uint8_t rx_buf_dtcm[256] __attribute__((aligned(32),
                                               section(".dtcm_bss")));
\end{ccode}
Record what the failure actually looks like, because that is the part nobody
writes down: whether the transfer never starts, whether it raises a transfer
error, or whether it completes with a count that never decremented.

\item \textbf{Write down the reachability rule and test it.} Three engines,
three different answers, and the difference is the reason the third one exists
at all.
\begin{table}[H]\centering\small
\begin{tabular}{p{30mm}p{46mm}p{34mm}p{34mm}}
\toprule
Engine & Reaches & Does not reach & Why it exists \\
\midrule
Main transfer engines & The main SRAM regions, the peripheral space, flash & The tightly coupled memories & The ordinary workhorse for peripheral traffic \\
Master transfer engine & Everything, including the tightly coupled memories & Nothing of consequence & Large block moves and the only path into tightly coupled memory \\
Low-power engine & One small memory in the low-power domain, and the peripherals in that domain & The main SRAM regions & It keeps running while the rest of the part sleeps, which is its entire purpose \\
\bottomrule
\end{tabular}
\caption{The reachability rule. Confirm each row against RM0455 for this part before relying on it; the equivalent table for the better-known sibling is not the same table and names domains this part does not have.}
\end{table}

\item \textbf{Run the matrix and publish the rates.} Four variants, a thousand
runs each, one table.
\begin{shellcode}
python tools/run_matrix.py --port /dev/ttyACM0 --runs 1000 \
  --variants naive,maintain,region,placement --out results/failure_rates.csv
\end{shellcode}
The script exits with a nonzero status if the naive variant did not fail at
all, because a control that passes means the experiment did not reproduce the
condition and the other three results prove nothing.
\end{steps}

\subsection*{Build, flash and debug}

\diagram{c19_timing}{The coherency figure. Three moments in one transfer: the processor's write leaves a dirty line, the engine's write leaves that line stale, and the processor's read then has two possible answers depending on what happened in between.}

\begin{shellcode}
cmake -B build -DVARIANT=naive -DCMAKE_TOOLCHAIN_FILE=cmake/arm-none-eabi.cmake
cmake --build build -j
probe-rs run --chip STM32H7A3ZITx build/firmware.elf
\end{shellcode}

When the board stops in a memory management fault, the protection unit is doing
its job and the region table is wrong. Read the fault status and the faulting
address, and check the two rules first: the size is a power of two and the base
address is a multiple of that size. When the board stops in a bus fault during
a transfer, suspect reachability rather than attributes. When the board does
not stop at all and simply produces wrong numbers, that is the original
condition and the harness should already be counting it.

\begin{note}{Why it always works under the debugger}
Halting the processor and inspecting memory from the debug view changes the
cache state, often in the direction that hides the problem. A team that debugs
only interactively can spend a week concluding the fault is in the sensor. The
harness exists so that the evidence is produced by a free-running board and
reported over a serial port, and so that a failure count can be compared
between two builds rather than argued about.
\end{note}

\subsection*{Verification and acceptance criteria}

\begin{itemize}
\item The naive variant fails at a rate greater than zero over a thousand runs,
reported as failing runs out of a thousand together with the distribution of
surviving sentinel bytes. A naive variant that does not fail is a failed
experiment and the harness says so.
\item Each of the three fixes produces zero failures in a thousand runs, with
the same sensor, the same rate and the same build options apart from the one
thing that differs.
\item The two cycle costs are measured with the core's cycle counter and
reported per transfer, so that the choice between maintenance and a
non-cacheable region can be made on numbers rather than on taste.
\item Every buffer in the repository is aligned to 32 bytes and sized to a
multiple of 32 bytes, and both facts are enforced by a build-time assertion
rather than by a comment.
\item The region table is applied at startup and a deliberately misaligned
entry stops the build, not the board.
\item The fourth failure is reproduced and described: what the transfer engine
actually does when it is pointed at a memory it cannot reach, in words, with
the register evidence.
\item The reachability table is checked row by row against RM0455 for this
part, and each row carries the manual section it came from.
\item The result table is generated from the CSV file, so it cannot drift from
the data it claims to summarise.
\end{itemize}

\subsection*{Variants}

\begin{table}[H]\centering\small
\begin{tabular}{p{24mm}p{26mm}p{44mm}p{22mm}p{20mm}}
\toprule
Axis & Variant & What changes & Cost & Built in full in \\
\midrule
Execution model & Interrupt with a flag & The control path here: the completion interrupt sets a flag and the main loop reads the buffer & Nothing else runs while waiting & Chapter 3 \\
Execution model & DMA & The subject of the chapter. Every other execution model in the book that moves a buffer this way depends on getting this right & The whole coherency problem & Chapter 4 \\
Execution model & DMA double buffered & Two buffers alternate, so the maintenance step has to name the right half, which is where this gets genuinely easy to get wrong & One more thing to get wrong & Chapter 6 \\
Operating system & An RTOS task & Maintenance moves into a driver layer with a lock around it, and the non-cacheable pool becomes a shared resource with an owner & Allocation policy & Chapter 20 \\
Peripheral substitution & The same experiment over I2C & Slower, so the window in which a line can be evicted is wider and the failure rate changes. A useful second data point & A second wiring & Chapter 17 \\
Language & Rust & The buffer pool becomes a type that cannot be constructed outside the non-cacheable section, which is the strongest version of this fix that exists & The library ecosystem here is thinner & Chapter 17 \\
Time and safety & Write-through instead of non-cacheable & Cures the direction where the processor writes and the engine reads, and leaves the other direction broken. A half fix that looks like a whole one & A false sense of correctness & Here, as a note \\
\bottomrule
\end{tabular}
\caption{Variants for chapter 19. The last row is included because it is the most common wrong answer given in forum threads on this subject and it deserves to be named.}
\end{table}

\subsection*{Pitfalls}

\begin{itemize}
\item Invalidating a buffer that shares its first or last cache line with
another variable. The neighbour's pending writes are dropped and the corruption
appears somewhere unrelated.
\item Cleaning when you meant to invalidate. Clean pushes the processor's view
out to memory; invalidate drops the processor's view. Using the wrong one
leaves the problem in place and adds cost.
\item Assuming a transfer that completed successfully means the data is
visible. The completion flag describes the engine's work, not the processor's
view of it.
\item Putting a buffer in tightly coupled memory because it is fast, and then
pointing the main transfer engine at it. The code compiles and links.
\item Copying a region table from material written for the better-known sibling
part. The domain names differ and the addresses differ.
\item Marking a region device memory when non-cacheable normal memory was
meant. Device memory forbids reordering and unaligned access, and the cost is
paid on every access forever.
\item Turning the data cache off to make the problem go away. It does, at a
price that is never measured and never revisited, and it removes the reason the
part was chosen.
\item Producing the failure count under the debugger. Halting changes the
cache, and a run that is watched behaves differently from a run that is not.
\item Believing that the transfer engine snoops the cache. On this core it does
not. There is no hardware arbitration between the two views.
\end{itemize}

\subsection*{Best practices applied}

\begin{itemize}
\item The failure is reproduced and measured before any fix is written, so the
fix has something to prove itself against.
\item The control variant is required to fail. An experiment whose control
passes reports nothing.
\item Three fixes are implemented and their costs are measured, rather than one
being recommended on the strength of a preference.
\item Rules the compiler can check are checked by the compiler: alignment,
size, and the power-of-two region rule.
\item The memory decisions live in one declarative table, in a header, next to
a generated prose description, instead of being distributed through the driver
layer.
\item Everything inherited from material about a sibling part is marked as
such and checked against the manual for this one before it is written as fact.
\end{itemize}

\subsection*{Stretch goals}

\begin{itemize}
\item Publish the region table as a small standalone library with its
build-time checks and a test suite, under a permissive licence. It is the
missing piece named in the gap list and it is a weekend of work.
\item Add a region-table linter that reads the linker map file and reports any
section that lands outside every declared region, which catches the whole class
of faults where a buffer moved and nobody noticed.
\item Reproduce the reported problem in the vendor's own example tree on this
board, confirm it is the same mechanism, and write it up with the register
evidence.
\item Add the master transfer engine as a fourth path and measure whether
moving into tightly coupled memory through it is faster than leaving the
buffer where the main engine can reach it.
\end{itemize}

\subsection*{Roadmap and next steps}

Chapter 20 is the next step for a reader who now has a correct transfer,
because the moment two tasks share that buffer the question changes from
coherency to ownership, and the answer is a different kind of design.
Chapters 13, 15 and 18 all move a buffer with a transfer engine and all depend
on this chapter being understood first.

For going deeper, read the three-part written series on the Cortex-M7 cache in
order: it is the clearest teaching treatment available and it sits at the
architecture level, so nothing in it goes stale when the silicon changes. Then
read the vendor's application notes on the cache and on the protection unit,
and the community article by their own engineer that puts the three fixes side
by side with complete configuration code. The three application briefs from
the other silicon vendor are worth reading afterwards precisely because they
cover one fix each. The community roadmap places this material in the firmware
track rather than the embedded Linux one, and its central advice, that projects
teach what reading does not, is why this chapter is an experiment rather than
an explanation. Appendix H lists the courses and their licences.

\subsection*{Portfolio evidence}

\begin{itemize}
\item A public repository containing the harness, the four variants and the
region table, with a README that opens with the coherency figure and states the
measured failure rate in the first paragraph.
\item The result table: four variants, a thousand runs each, failures and cycle
cost, generated from the CSV file.
\item The region table header with its build-time checks, published separately
under a permissive licence, since nothing maintained currently provides it.
\item A short written note on the fourth failure, describing exactly what the
transfer engine does when pointed at a memory it cannot reach, with register
evidence rather than a description.
\item A checked reachability table for this part, each row carrying its manual
section, which is a thing a reader currently has to assemble themselves.
\end{itemize}

\subsection*{Sources}

Normative references:
\begin{itemize}
\item Reference manual RM0455, for the memory map, the domain names, the region
inventory and which master reaches which memory. Not RM0433, whose domain
vocabulary this part does not use.
\item The STM32H7A3xI datasheet, for the memory sizes that RM0455 leaves to the
part-specific document.
\item The Armv7-M architecture reference manual, for the protection unit's
region rules and attribute encodings, and the Cortex-M7 technical reference
manual for the cache line size and the maintenance operations.
\item The vendor's application notes on the L1 cache, on the memory protection
unit, and on using the cache for performance.
\end{itemize}

Reusable implementations:
\begin{itemize}
\item The vendor engineer's community article putting the three fixes side by
side with complete configuration code.\newline
\url{https://community.st.com/t5/stm32-mcus/}\newline
\url{how-to-use-the-memory-protection-unit-to-manage-cache-coherency/ta-p/858387}
\item The same vendor's write-up of why transfers fail on this family, which
adds the fourth failure.\newline
\url{https://community.st.com/t5/stm32-mcus/dma-is-not-working-on-stm32h7-devices/ta-p/49498}
\item The three-part written series on the Cortex-M7 cache, which is the best
teaching treatment in print.\newline
\url{https://blog.feabhas.com/2020/10/}\newline
\url{introduction-to-the-arm-cortex-m7-cache-part-1-cache-basics/}
\item One video doing this on an H7 with a non-cacheable region and a custom
linker section.\newline
\url{https://www.youtube.com/watch?v=_K3GvQkyarg}
\item The open report against the vendor's own examples, which is the evidence
that this is a shipping problem and not a teaching example.\newline
\url{https://github.com/STMicroelectronics/STM32CubeH7/issues/153}
\end{itemize}
   % Caches, the MPU, and why DMA reads stale bytes
\project{20}{Three jobs at once: the node as an RTOS application}{NUCLEO-H7A3ZI-Q + X-NUCLEO-IKS4A1}{Static allocation, synchronisation, measured latency distribution}

\begin{keyfacts}
\item[Board] NUCLEO-H7A3ZI-Q with the X-NUCLEO-IKS4A1 on the Arduino header, one shield only
\item[Peripherals] I2C to the shield, one interrupt line for the watermark, a general timer as the synthetic release source, the system tick, USART3 for the result stream
\item[Toolchain] arm-none-eabi-gcc with CMake for the first kernel, and the other operating system's own build tool for the comparison
\item[Operating system] FreeRTOS, statically allocated, with a Zephyr variant built in full
\item[Difficulty] 5 of 5
\item[Effort] 6 evenings of about four hours
\item[Deliverable] The node of chapters 8 to 12 as a three-task application with no dynamic allocation anywhere, a written design that states its deadline before any measurement, and a latency distribution over a million events reporting the median, the 99.9th percentile and the maximum
\end{keyfacts}

\subsection*{Why this project}

This is the chapter the book exists to make possible. A job-fit note written
against real vacancies records one standing gap in the portfolio, and it is not
a language or a bus or a sensor: it is that there is no shipped application
built on a real-time kernel. Everything earlier in this book is a superloop
with interrupts, which is the right architecture for what it does and is not
what a job advertisement means when it asks for experience with a kernel. One
exercise on a board already owned closes the gap, and this is it.

The second reason is that the node of chapters 8 to 12 has outgrown its own
structure. Sensing, computing a feature and getting bytes onto a link are three
jobs with three different timing characters: one is driven by a sensor and must
not be late, one is compute-bound and can be interrupted, and one spends most
of its life waiting for a modem that answers when it feels like it. A superloop
can do all three, and the chapter that rebuilds them as three tasks is
worthwhile precisely because it must then prove that the rebuild did not make
the first one late.

The third reason is evidence. This subject area carries more confident numbers
with no method behind them than any other in the book, and the loudest is
retired below. The corrective is a deadline written down before the first run,
a distribution rather than an average, and a million events so that the 99.9th
percentile means something.

\begin{note}{A number this chapter retires}
The most repeated claim in this field is that waking a task by direct
notification is forty-five percent faster than by a semaphore. Its origin is a
vendor tutorial page that states no processor, no compiler, no optimisation
level and no measurement method. The kernel's own official book has since
dropped the figure and now says only that notifications are significantly
faster. The hard, current, falsifiable number is the memory cost, which is
five bytes per task for one notification entry. So this chapter states the
memory cost as a fact and turns the speed claim into an exercise measured here,
with its conditions printed next to it.
\end{note}

\subsection*{Prior art and what to reuse}

\begin{table}[H]\centering\small
\begin{tabular}{p{34mm}p{46mm}p{38mm}p{20mm}}
\toprule
Source & What it gives & What it does not & Licence \\
\midrule
The kernel and its free official book & The scheduler, the synchronisation objects, static allocation, and a book that is the normative description of all of it & No application architecture, no task watchdog abstraction, and no measurement method & MIT \\
The other operating system, with an upstream in-tree port for this exact board & A second implementation of the same application with a different primitive vocabulary, on a board definition somebody else already wrote and tested, which makes the comparison cheap rather than heroic & Its primitives do not map one to one, and neither project publishes the mapping & Apache-2.0 \\
The third commercial kernel & Silicon-level support for this family & No example for this board at all, so it is named here and not used & Vendor terms \\
The tick-based measurement library & A cycle-accurate stopwatch that deliberately uses the system tick rather than the cycle counter, because the tick is present on every part and in every build & It is an instrument, not an experiment. The histogram, the release source and the percentiles are the author's & Apache-2.0 \\
A two-part published study comparing two kernels & The method worth copying exactly: a stated release source, a large sample, and a report of median, high percentile and maximum rather than an average & Its platform is a Cortex-M4 at 80 MHz and not this core, so none of its numbers transfer. Only the method does & Article, cite \\
The active-object framework and its free course & The strongest available argument that tasks blocking on multiple objects is the wrong structure, and a complete framework that implements the alternative & The framework is GPL or commercial, its safety editions are commercial only, and the course code is AGPL-3.0, which is the most restrictive licence in this book's pool & GPL or commercial; course AGPL-3.0 \\
\bottomrule
\end{tabular}
\caption{Prior art for chapter 20. The last row is recommended and never quoted: its licence terms make reproduction in a commercial work impossible, and the chapter respects that by re-deriving the pattern rather than copying it.}
\end{table}

What is left to write is most of it. Neither operating system publishes a
mapping of its synchronisation primitives to the other's, and the third-party
attempts are partial, so the mapping tables below are checked against both
projects' own documentation and are original work. Nobody has implemented the
active-object pattern twice, by hand on ordinary queues and on the framework,
and compared them honestly on one board, so that comparison is new. And there
is no published latency distribution for either kernel on this core, which is
why the method is borrowed and the numbers are not.

\subsection*{Parts from the inventory}

\begin{table}[H]\centering\small
\begin{tabular}{p{40mm}p{66mm}p{40mm}}
\toprule
Part & Role & Interface \\
\midrule
NUCLEO-H7A3ZI-Q & The application, the scheduler and the instrument & Micro USB to the host \\
X-NUCLEO-IKS4A1 & The sensor source for the real workload. One shield at a time, never two & Arduino header, I2C, one interrupt line \\
A USB data cable & Power, programming and the result stream & Micro USB \\
Host PC & Collects the histogram, computes the percentiles and plots the tail & Python over the virtual COM port \\
\bottomrule
\end{tabular}
\caption{Inventory items used in chapter 20. Nothing is wired: the shield plugs onto the header. The one-shield-at-a-time rule is inherited from the kit labs and is about address collisions and the supply budget, not about convenience.}
\end{table}

\subsection*{System architecture}

\diagram{c20_arch}{The three tasks, the two interrupt sources, and every object that carries data between them. Nothing in this figure is allocated at run time: each object's storage is a named static array.}

Three tasks, and the priorities follow from the timing characters rather than
from importance. The sensor task is highest because it has a deadline. The
feature task is next because it has work but no deadline. The link task is
lowest because it spends its life waiting for a modem and because making it
higher would let a slow modem delay a sensor read. A fourth task, the
supervisor, runs at the top priority, does almost nothing, and exists to check
the liveness bitmask of chapter 12 and to service the independent watchdog. The
idle hook counts, which is the cheapest processor-load measurement there is.

\subsection*{Peripheral configuration}

\begin{table}[H]\centering\small
\begin{tabular}{p{22mm}p{26mm}p{24mm}p{30mm}p{30mm}}
\toprule
Peripheral & Mode & Clock source & Pins and function & Interrupt and transfers \\
\midrule
I2C to the shield & Fast mode, addresses from the shield manual & Peripheral bus & D14 and D15, confirm in the board manual & Interrupt driven, transfers optional \\
Watermark interrupt & Input, rising edge & Peripheral bus & Shield interrupt line, confirm & External interrupt, priority below the kernel's ceiling \\
General timer & Periodic at 1 kHz, the synthetic release source & Peripheral bus & None & Update interrupt, the release event \\
System tick & 1 kHz kernel tick, shared with the measurement library & Core clock & None & The kernel's own exception \\
USART3 & Asynchronous, 115200 8N1 & Peripheral bus & Board manual pins & Stream buffer from a task, never from an interrupt \\
Independent watchdog & Enabled, fed only by the supervisor & Low-speed internal & None & Reset on timeout \\
\bottomrule
\end{tabular}
\caption{Peripheral configuration. Every interrupt that calls a kernel function has a priority at or below the kernel's ceiling, and one that does not is the single most common way to produce a kernel that fails in a way nobody can reproduce.}
\end{table}

\begin{note}{Interrupt priority is the first thing to get right}
On this core a numerically lower priority value is more urgent, and the kernel
refuses to be called from any interrupt more urgent than its configured
ceiling. An interrupt above the ceiling that calls a kernel function does not
fail immediately: it corrupts a list and fails later, somewhere else. Build
with the kernel's assertion hook enabled during development so that the mistake
stops the board at the offending call.
\end{note}

\subsection*{Wiring}

\diagram{c20_wiring}{Nothing is wired. The shield plugs onto the header, and the only cable is the one to the host. The synthetic release source for the latency run is a timer on the board itself, which is what makes the release time known rather than assumed.}

\subsection*{Memory and timing budget}

\diagram{c20_mem}{Static allocation as a picture. Every stack, every control block and every object's storage is a named array in a known section, so the total is known at link time and the linker fails rather than the board.}

\begin{table}[H]\centering\small
\begin{tabular}{p{44mm}p{28mm}p{28mm}p{30mm}}
\toprule
Quantity & Budget & Measured & Margin \\
\midrule
Sensor task stack & 1 kB & not measured & not measured \\
Feature task stack & 4 kB & not measured & not measured \\
Link task stack & 2 kB & not measured & not measured \\
Supervisor task stack & 512 B & not measured & not measured \\
Kernel objects, all static & 2 kB & not measured & not measured \\
Notification storage per task & 5 B & fixed by the kernel & none needed \\
Release to sensor task running, median & 100 \textmu{}s & not measured & not measured \\
Release to sensor task running, 99.9th & 250 \textmu{}s & not measured & not measured \\
Release to sensor task running, maximum & 1 ms & not measured & not measured \\
Processor load at 104 Hz sensor rate & 20 percent & not measured & not measured \\
\bottomrule
\end{tabular}
\caption{The budget table, and the deadline is in it. The three latency rows are written before the first run and are the acceptance criteria for the whole chapter. The instrument is the tick-based measurement library, cross-checked against the core's cycle counter.}
\end{table}

Stack sizes are budgets until the high-water mark says otherwise. The kernel
reports, for each task, the smallest amount of free stack that task has ever
had, and that number is collected in the same result line as everything else.

\subsection*{Firmware design (UML)}

\diagram{c20_uml}{One event through the system, and the two points between which latency is defined. The definition matters more than the number: from the release edge to the first instruction of the task body, not to the end of the interrupt handler.}

The design is written before the code, which for a task set means four things
on one page: the tasks and their priorities, the objects between them and which
task owns each one, the blocking points of every task, and the deadline. The
last of these is what makes the rest checkable. A task set with no stated
deadline cannot be shown to work; it can only be shown to run.

Every task blocks on exactly one thing. That is a design rule and not an
accident of the kernel: a task that waits on several objects at once needs a
queue set or a polling object, both of which cost memory and make the worst
case harder to reason about. Where a task genuinely needs several sources, the
sources are merged into one queue of tagged events before the task sees them.
That rule is also the bridge to the active-object argument below, which says
the same thing more forcefully.

\subsection*{Data flow (ASCII)}

\begin{asciiart}
   timer 1 kHz            watermark IRQ
        |                      |
        v                      v
  +-----------------------------------+   notification (5 bytes per task)
  |  ISR: timestamp, notify, return   | ------------------+
  +-----------------------------------+                   |
                                                           v
  +------------------+  stream buffer  +------------------+   queue   +---------------+
  | sensor task      | --------------> | feature task     | --------> | link task     |
  | prio 3, I2C FIFO |   bytes         | prio 2, biquad   |  payload  | prio 1, AT    |
  +------------------+                 +------------------+           +---------------+
        |                                    |                              |
        +----------------+-------------------+------------------------------+
                         v  event group bit per task, one 32-bit word
                  +---------------------+       +-------------------------+
                  | supervisor, prio 4  | ----> | independent watchdog    |
                  +---------------------+       +-------------------------+
\end{asciiart}

\subsection*{Repository layout}

\begin{plaincode}
nucleo-h7a3-rtos-node/
  CMakeLists.txt
  cmake/arm-none-eabi.cmake
  third_party/FreeRTOS-Kernel/       # vendored, MIT, LICENSE kept
  third_party/perf_counter/          # vendored, Apache-2.0
  config/FreeRTOSConfig.h            # static only, assertions on
  src/tasks_sensor.c                 # highest application priority
  src/tasks_feature.c                # biquad and statistics, chapter 18
  src/tasks_link.c                   # AT engine of chapter 11, or the stub
  src/tasks_supervisor.c             # liveness bitmask, watchdog, chapter 12
  src/objects_static.c               # every queue, buffer and semaphore
  src/latency.c                      # histogram, percentiles, result line
  src/release_timer.c                # the synthetic 1 kHz release source
  docs/DESIGN.md                     # tasks, priorities, objects, deadline
  docs/PRIMITIVE_MAP.md              # the mapping between the two kernels
  zephyr/                            # the same application, second kernel
    prj.conf  CMakeLists.txt  src/
  active_object/                     # the pattern by hand, on plain queues
  tools/percentiles.py               # reads the histogram, plots the tail
  results/latency_million.csv
  README.md
\end{plaincode}

\subsection*{Steps}

\begin{steps}
\item \textbf{Write the design page before any code.} Four tasks, their
priorities, the objects between them, where each task blocks, and the deadline.
If that page cannot be written, the task set is not understood yet and the code
will encode the confusion.

\item \textbf{Turn dynamic allocation off and keep it off.} This is one
configuration line and it changes the character of the project. With it off,
every object's storage is visible, the total is known at link time, and there is
no allocator to fragment.
\begin{ccode}
#define configSUPPORT_STATIC_ALLOCATION     1
#define configSUPPORT_DYNAMIC_ALLOCATION    0
#define configUSE_MUTEXES                   1
#define configUSE_TASK_NOTIFICATIONS        1
#define configTASK_NOTIFICATION_ARRAY_ENTRIES 1   /* 5 bytes per task */
#define configCHECK_FOR_STACK_OVERFLOW      2
#define configASSERT(x) if ((x) == 0) { taskDISABLE_INTERRUPTS(); for(;;); }
\end{ccode}
With dynamic allocation off the kernel asks the application for the idle task's
and the timer task's memory. Providing those two hooks is not optional and the
link fails without them, which is the correct behaviour.

\item \textbf{Create the task set statically.} Each task owns a named stack
array and a named control block, so the memory figure above can be drawn from
the map file rather than from intention.
\begin{ccode}
static StackType_t  sensor_stack[256];          /* words, so 1 kB */
static StaticTask_t sensor_tcb;

TaskHandle_t sensor_task = xTaskCreateStatic(
    sensor_task_body, "sensor", 256, NULL,
    PRIO_SENSOR, sensor_stack, &sensor_tcb);
\end{ccode}

\item \textbf{Wake the sensor task from the interrupt, and define latency
before measuring it.} The interrupt timestamps the release and notifies; the
task timestamps again as its first act. Latency is the difference, which
includes the handler, the scheduler and the context switch, because all three
are between the event and the work.
\begin{ccode}
void TIM_Release_IRQHandler(void)             /* the synthetic release */
{
    release_ts = perf_counter_now();
    BaseType_t woken = pdFALSE;
    vTaskNotifyGiveFromISR(sensor_task, &woken);
    portYIELD_FROM_ISR(woken);
}

static void sensor_task_body(void *arg)
{
    for (;;) {
        ulTaskNotifyTake(pdTRUE, portMAX_DELAY);
        latency_record(perf_counter_now() - release_ts);   /* first act */
        sensor_read_fifo();
    }
}
\end{ccode}

\item \textbf{Record a distribution, not an average.} A million samples are not
stored. A fixed-bin histogram plus a running maximum is, and the percentiles
come from the histogram on the host.
\begin{ccode}
#define BINS 1024                       /* one bin per microsecond */
static uint32_t hist[BINS];
static uint32_t hist_over, hist_max, hist_n;

void latency_record(uint32_t cycles)
{
    uint32_t us = cycles / CYCLES_PER_US;
    if (us < BINS) hist[us]++; else hist_over++;
    if (us > hist_max) hist_max = us;
    hist_n++;
}
\end{ccode}
The overflow counter is not decoration. A distribution that silently discards
its tail reports a maximum that is a property of the histogram rather than of
the system.

\item \textbf{Run a million events with a known release source.} A timer at
1 kHz gives a million releases in about seventeen minutes, with a release time
that is known rather than inferred. The sensor-driven run follows and is the
real workload; the timer-driven run is what makes the distribution comparable
between builds.
\begin{shellcode}
python tools/percentiles.py --port /dev/ttyACM0 --events 1000000 \
  --out results/latency_million.csv
\end{shellcode}

\item \textbf{Measure the claim instead of repeating it.} Build the same
harness twice, once notifying and once giving a binary semaphore, change
nothing else, and report both distributions with the conditions printed beside
them.
\begin{ccode}
#if WAKE_BY_NOTIFICATION
    vTaskNotifyGiveFromISR(sensor_task, &woken);
#else
    xSemaphoreGiveFromISR(sensor_sem, &woken);
#endif
\end{ccode}
Print the processor, the clock, the compiler and its optimisation level, the
kernel version, the cache state and the release source in the result header.
Those are exactly the six things the widely repeated figure omits, which is why
it cannot be checked and should not be quoted.

\item \textbf{State what the mutex does and does not do.} This kernel
implements a basic form of priority inheritance. It assumes a task holds one
mutex at a time, it is not transitive, and there is no priority ceiling
protocol. That is a design constraint, not a defect, and a design that needs
more than it offers needs a different structure rather than a different
opinion.
\begin{ccode}
/* Correct use: one mutex, held briefly, released on every path. */
if (xSemaphoreTake(i2c_mutex, pdMS_TO_TICKS(10)) == pdTRUE) {
    i2c_transfer(...);
    xSemaphoreGive(i2c_mutex);
}
\end{ccode}

\item \textbf{Build the same application on the other kernel.} The upstream
in-tree board definition exists, so this costs an evening rather than a week.
The point is not that one is better: it is that writing the same design twice
exposes which parts were design and which were the first kernel's vocabulary.
\begin{shellcode}
west build -b nucleo_h7a3zi_q zephyr/
west flash
\end{shellcode}

\item \textbf{Write the mapping table, because nobody else has.} Neither
project documents the correspondence and the third-party attempts are partial.
Each row below is checked against both projects' own reference documentation,
and the rows with no equivalent are the interesting ones.

\begin{table}[H]\centering\small
\begin{tabular}{p{38mm}p{38mm}p{76mm}}
\toprule
This kernel & The other kernel & What actually differs \\
\midrule
Task, created statically & Thread with a statically defined stack & The other declares the stack with a macro that also handles guard regions and alignment \\
vTaskDelay & k\_sleep or k\_msleep & Equivalent. Both are relative delays \\
vTaskDelayUntil & No direct equivalent & The closest is a kernel timer or a work item rescheduled from its own handler \\
Binary semaphore & Semaphore with a limit of one & Equivalent \\
Counting semaphore & Semaphore with a limit of N & Equivalent \\
Mutex with priority inheritance & Mutex with priority inheritance & Both are basic forms. Neither offers a priority ceiling protocol \\
Queue, copies fixed-size items & Message queue, copies fixed-size items & Equivalent, and both have static variants \\
Queue of pointers & Queue or first-in-first-out list & The other's list is intrusive: the item carries the link word \\
\bottomrule
\end{tabular}
\caption{Primitive mapping, part one: threads, time and the basic objects. Checked against both projects' own reference documentation on Sunday 20 September 2026. Neither project publishes this table.}
\end{table}

\begin{table}[H]\centering\small
\begin{tabular}{p{38mm}p{38mm}p{76mm}}
\toprule
This kernel & The other kernel & What actually differs \\
\midrule
Stream buffer, bytes, one reader one writer & Pipe, bytes & The other's pipe does not assume a single reader and a single writer, so it takes a lock the stream buffer does not need \\
Message buffer, variable-length messages & No direct equivalent & Framing over a pipe, or a message queue sized for the largest message \\
Event group, 24 usable bits & Event object, 32 bits & The other has more bits and no reserved range \\
Direct task notification & No direct equivalent & The nearest is a per-thread semaphore or a poll signal. This is the sharpest difference between the two \\
Queue set & Poll & The other's poll is the general mechanism and is not limited to queues \\
Software timer & Kernel timer or delayable work & The other separates a timer callback in interrupt context from work deferred to a thread \\
Deferred interrupt handling to the timer task & Work submitted to the system work queue & The other's name for it is clearer and the mechanism is the same \\
Critical section macros & Interrupt lock or a spin lock & The other's spin lock is the multiprocessor-safe form and is the one to use \\
Tickless idle & Tickless kernel & Both are a configuration option. Chapter 7 is where the consequences were measured \\
\bottomrule
\end{tabular}
\caption{Primitive mapping, part two: data movement, events and deferral. The row with no equivalent on either side is the one worth reading twice, because it is where a design does not port by renaming.}
\end{table}

\item \textbf{Implement the active-object pattern twice.} By hand first: one
queue per task, one run-to-completion event handler per task, a static event
pool, and no task ever blocking anywhere except on its own queue. Then on the
framework. Compare code size, cycles per event and the same latency
distribution. Keep the two in separate repositories: the hand-written one can
carry a permissive licence, the framework one cannot.
\end{steps}

\subsection*{Build, flash and debug}

\diagram{c20_timing}{The schedule figure: two releases, four tasks, and the deadline drawn where it was written down. The marked interval at the left is what the million-event distribution measures.}

\begin{shellcode}
cmake -B build-fw -G Ninja -DWAKE=notification \
  -DCMAKE_TOOLCHAIN_FILE=cmake/arm-none-eabi.cmake
cmake --build build-fw -j
cp build-fw/firmware.bin "$PROBE_DISK"/   # onto the probe disk
\end{shellcode}

A debugger that understands the kernel shows the task list, each task's state
and each task's stack high-water mark, which turns a hung board into a reading.
When the board stops in the assertion hook the cause is usually one of three:
an interrupt above the kernel ceiling called a kernel function, a task returned
instead of looping forever, or a stack overflowed. All three are configuration
rather than logic.

\begin{note}{When the board runs and the distribution is wrong}
A tail that is far worse than the median usually has one of four causes, in
order of likelihood: an interrupt somewhere that runs for longer than anyone
thinks; a critical section held across something slow, such as a bus transfer;
a lower-priority task holding a mutex the sensor task needs, which is priority
inversion and is the one the kernel only partly protects against; and the
measurement itself, if the result line is printed from inside the timed path.
Print from a task, never from the interrupt, and never inside the interval
being measured.
\end{note}

\subsection*{Verification and acceptance criteria}

\begin{itemize}
\item The design page exists and states the deadline before any measurement
appears anywhere in the repository history.
\item No dynamic allocation anywhere: the configuration disables it, the
allocator is not linked, and the map file is the evidence.
\item A million release events are recorded with the timer source, and the
report gives the median, the 99.9th percentile and the maximum, with the
overflow count shown so the tail can be trusted.
\item All three latency figures meet the budget written in the budget table
above, or the table is updated with the reason and the design is revisited.
\item The same distribution is reported for the notification wake and the
semaphore wake, with the six conditions printed in the header, and the chapter
states the measured relationship rather than repeating the claim it retires.
\item Every task's stack high-water mark is reported in the same result line,
and no task has used more than eighty percent of its budget.
\item The application runs unmodified on both kernels, and the mapping table
names every place where the design had to change rather than be renamed.
\item The active-object implementation exists twice and is compared on code
size, cycles per event and the same latency distribution, with the counter
argument stated in the text.
\item The supervisor services the watchdog only when every task has set its
liveness bit since the last cycle, and a deliberately stalled task produces a
reset with a crash record, as in chapter 12.
\end{itemize}

\subsection*{Variants}

\begin{table}[H]\centering\small
\begin{tabular}{p{24mm}p{26mm}p{44mm}p{22mm}p{20mm}}
\toprule
Axis & Variant & What changes & Cost & Built in full in \\
\midrule
Execution model & An RTOS task set, statically allocated & The whole chapter. Four tasks, fixed priorities, no allocator & A scheduler in the timing path & Here \\
Synchronisation & Binary and counting semaphore & The signal from interrupt to task, and the resource count for the shield's bus & One object per signal & Here \\
Synchronisation & Mutex with priority inheritance & Exclusive access to the shared bus, with the kernel's basic inheritance and its stated limits & Inversion the kernel only partly prevents & Here \\
Synchronisation & Queue & Payloads from the feature task to the link task, copied rather than shared & A copy per message & Here \\
Synchronisation & Direct task notification & The cheapest wake, at five bytes per task, and the subject of this chapter's measured exercise & One signal per task at a time & Here \\
Synchronisation & Event group & The liveness bitmask the supervisor reads, one bit per task in one word & Bits are not counters & Here \\
Synchronisation & Stream buffer & Sensor bytes to the feature task without framing, one reader and one writer & That assumption is load bearing & Here \\
Operating system & The other kernel & The same design on an upstream in-tree port for this exact board, with a different primitive vocabulary & A second toolchain & Here, and referenced from chapter 13 \\
Operating system & Bare metal & What chapters 8 to 12 already are, and the thing this chapter is measured against & No scheduler, no latency to measure & Chapter 8 \\
Time and safety & Tickless idle & The kernel stops the tick when every task is blocked, which is what makes an RTOS node a low-power node & Wake latency & Chapter 7 \\
Time and safety & Independent watchdog & The supervisor is the only feeder, and it feeds only on a full liveness bitmask & A reset is a real outcome & Chapter 12 \\
Language & C++ & Tasks as objects with their stacks as members, which makes static allocation natural and the ownership explicit & Care with construction order & Chapter 9 \\
Intelligence and reach & Feature on the MCU & The feature task is chapter 18's pipeline, now sharing a processor with two other jobs & Jitter enters the compute path & Chapter 16 \\
\bottomrule
\end{tabular}
\caption{Variants for chapter 20. This chapter carries more of the book's variant matrix than any other, which is why it is last and why it is six evenings rather than four.}
\end{table}

\subsection*{Pitfalls}

\begin{itemize}
\item An interrupt above the kernel's ceiling calling a kernel function. It
does not fail at the call. It corrupts a list and fails later somewhere else.
\item Reporting an average latency. A real-time system is described by its
tail, and an average hides exactly the events that matter.
\item Sizing a histogram so that the interesting events land in the overflow
bin, and then reporting a maximum that is a property of the histogram.
\item Assuming the kernel's priority inheritance is a full protocol. It is a
basic form, it assumes one mutex at a time, it is not transitive, and there is
no priority ceiling.
\item Printing the result line from inside the interval being measured.
\item Letting a task return. A task function must never return; the kernel has
no sensible behaviour for it.
\item Quoting the forty-five percent figure. It has no stated processor,
compiler, optimisation level or method, and the kernel's own book has dropped
it.
\item Porting a design to the other kernel by renaming functions. Two rows of
the mapping tables have no equivalent, and those are where the design has to
change.
\item Moving a buffer into a task and forgetting that a transfer engine still
writes it. Chapter 19 does not stop applying because there is now a scheduler.
\end{itemize}

\subsection*{Best practices applied}

\begin{itemize}
\item The deadline is written down before the first measurement, and the
acceptance criteria are the budget table rather than a judgement afterwards.
\item Nothing is allocated at run time, so the memory total is a link-time
fact.
\item Every task blocks on exactly one object, which keeps the worst case
possible to reason about.
\item A claim with no method is measured rather than repeated, and the
conditions the original omitted are printed with the result.
\item A licence that forbids reproduction is respected by re-deriving the
pattern rather than copying the code, and the two implementations live in
separate repositories with separate terms.
\item The comparison between the two kernels is made by writing the same design
twice, which is the only way to find out which parts were design.
\end{itemize}

\subsection*{Stretch goals}

\begin{itemize}
\item Measure priority inversion directly: a low-priority task holding the bus
mutex while the sensor task waits, with and without inheritance enabled. No
published measurement of inversion blocking exists for this core.
\item Add the task watchdog abstraction that the other kernel ships and this
one does not, as a small permissive library, and compare it against the
supervisor of chapter 12.
\item Run the distribution again with the caches off, and again with the code
in tightly coupled memory, joining this chapter to chapters 18 and 19.
\item Publish the mapping tables as a standalone document. It does not exist
anywhere, and it is the kind of reference other people link to.
\end{itemize}

\subsection*{Roadmap and next steps}

This is the last chapter, so the roadmap is about what comes after the book.
The community roadmap that runs through this volume separates firmware from
embedded Linux from hardware and is explicit that projects teach what reading
does not, which is the thesis this chapter ends on.

For the kernel itself, the free official book is the normative description and
is worth reading cover to cover once, now that there is an application to read
it against. For the other kernel, the best free course comes from a different
silicon vendor and its concepts transfer directly. For the mathematics of
real-time scheduling, which almost nothing else teaches, there is one
university specialisation that actually covers it, and the free course whose
code is under a permissive licence and whose slides are freely licensed is the
safest material to build further exercises on, although its examples are not on
this board. The free course that walks a single problem from a superloop to an
event-driven architecture across fifty-six lessons is the best companion to
this chapter specifically, and its licence means it is recommended and never
quoted. Appendix H lists all of them with their terms, and appendix I sets out
what a reader does next.

\subsection*{Portfolio evidence}

\begin{itemize}
\item A public repository with the design page dated before the first result,
the static task set, and a README whose first figure is the schedule diagram.
\item The latency report: a million events, median, 99.9th percentile and
maximum, the overflow count, the tail plotted on a logarithmic axis, and the
six conditions in the header.
\item The two wake methods measured side by side, with a sentence saying what
the widely repeated figure omits and what this board actually did.
\item The stack high-water marks for every task next to their budgets.
\item The same application on the second kernel, in the same repository, with
the mapping tables as a document rather than as comments.
\item The active-object pattern implemented twice, with the counter argument
stated fairly and the comparison published, since nobody has published it.
\end{itemize}

\subsection*{Sources}

Normative references:
\begin{itemize}
\item The kernel's own official book, which is the normative description of the
scheduler, the objects and static allocation, and which no longer carries the
speed figure this chapter retires.
\item The other operating system's reference documentation for kernel services,
which is the authority for every row of the mapping tables on that side.
\item Reference manual RM0455 for the interrupt controller's priority grouping,
and the Armv7-M architecture reference manual for how priority values order.
\item The board user manual for MB1363 and the shield user manual UM3239, for
the header pins and the shield's bus topologies.
\end{itemize}

Reusable implementations:
\begin{itemize}
\item The kernel, MIT, vendored with its licence file.\newline
\url{https://github.com/FreeRTOS/FreeRTOS-Kernel}
\item The kernel's free official book, MIT.\newline
\url{https://github.com/FreeRTOS/FreeRTOS-Kernel-Book}
\item The other operating system's in-tree port for this exact board,
Apache-2.0.\newline
\url{https://docs.zephyrproject.org/latest/boards/st/nucleo_h7a3zi_q/doc/index.html}
\item The tick-based measurement library, Apache-2.0.\newline
\url{https://github.com/GorgonMeducer/perf_counter}
\item The two-part study whose method this chapter copies and whose numbers it
does not, since its platform is a different core.\newline
\url{https://www.ul.com/sis/blog/}\newline
\url{measuring-real-time-operating-system-performance-part-ii-comparing-freertos-vs-zephyr}
\item The active-object framework, GPL or commercial, recommended and never
quoted.\newline
\url{https://www.state-machine.com/products/qp}
\end{itemize}
   % Three jobs at once: the node as an RTOS application

\appendix
\clearpage
\section{The wider project catalogue, scored}

The twenty chapters were selected from a larger scored list written for the same
bench. That list is a ranking document; this volume is the build manual. Every
entry below keeps the identifier and the total score it carried in the original,
so a reader holding both documents can move between them without a translation
table. The score columns are impact, learnability and immediate usefulness, each
out of five, against an effort estimate in hours; the total is the original's own
weighting of those four and is reproduced rather than recomputed.

Three classes appear. Entries beginning with Z need no extra parts. Entries
beginning with N were added in the revision that produced the list. Entries
beginning with B need a purchase, and this volume buys nothing, so every B entry
stayed out for the same reason and is listed here only so that the numbering
survives and so that the reader can see what a budget would unlock.

\begin{table}[H]\centering\small
\begin{tabular}{p{9mm}p{70mm}p{11mm}p{55mm}}
\toprule
ID & Project as the original named it & Total & Where it went \\
\midrule
Z01 & H7A3 with the IKS5A1 industrial MEMS logger & 14.5 & Chapter 14 \\
Z02 & H7A3 with the IKS4A1 environment node & 11.5 & Chapter 13 \\
Z03 & H7A3 with the 8 by 8 time-of-flight sensor for occupancy and level & 13.3 & Chapter 15 \\
Z04 & IKS5A1 and the time-of-flight shield fused & 11.0 & Out: two shields at once, which the one-at-a-time rule forbids for address and supply reasons \\
Z05 & IKS4A1 and the time-of-flight shield fused & 9.0 & Out: same reason as Z04 \\
Z06 & STWIN.box wideband vibration logger & 18.7 & Out: the STWIN.box is the committed subject of a separate build and appears in no chapter here \\
Z07 & In-sensor classifier wake instead of a raw stream & 18.3 & Chapter 16, moved from the STWIN.box to the IKS4A1 so that it fits this bench \\
Z08 & Energy map of the STWIN.box modes & 18.7 & Out: STWIN.box. The method is chapters 7 and 10 \\
Z09 & Energy map of the H7A3 with the IKS5A1 & 15.2 & Chapter 10, with N04 \\
Z10 & Energy of time-of-flight ranging cadences & 13.3 & Referenced as a stretch goal in chapter 15 \\
Z11 & Transform library work on streamed inertial data & 15.0 & Chapter 18 \\
Z12 & On-box features then a small log, on the STWIN.box & 14.0 & Out: STWIN.box \\
Z13 & Host-side Python pipeline for the vendor's logging firmware & 18.5 & Out: it teaches the host, and the sibling volume owns that side of the cable \\
Z14 & Serial logger from the board to a Raspberry Pi, with bench notes & 13.7 & Absorbed into the rig of chapter 12 \\
Z15 & Narrowband command bring-up without a working subscription & 9.5 & Chapter 11 \\
Z16 & Application state machine: sense, feature, and a transmit that is a stub & 18.0 & Chapter 8 \\
Z17 & Two benches compared: laboratory inertial against field inertial & 11.0 & Out: STWIN.box on one side \\
Z18 & Occupancy policy engine with no camera & 12.3 & Absorbed into chapter 15, which decides on the microcontroller \\
Z19 & Vendor layer against register level, as a kata & 16.3 & Chapter 17 \\
Z20 & The one-page matrix of sensor, current and payload & 14.7 & Absorbed into chapter 10 and into appendix C \\
\bottomrule
\end{tabular}
\caption{The Z entries, which need no extra parts. Nine became chapters, three were absorbed into a chapter that covers the same ground more completely, and the rest stayed out for a stated reason.}
\end{table}

\begin{table}[H]\centering\small
\begin{tabular}{p{9mm}p{70mm}p{11mm}p{55mm}}
\toprule
ID & Project as the original named it & Total & Where it went \\
\midrule
N01 & Bearing envelope spectrum on the STWIN.box & 17.3 & Out: STWIN.box, and it is the committed subject of the separate condition-monitoring build \\
N02 & Energy as a regression test, in the host test framework & 17.7 & Chapter 12 \\
N03 & Payload codec with a Python twin, protocol discipline without a radio & 17.2 & Chapter 9 \\
N04 & Wake overhead split by labelling the current trace with marker pins & 18.2 & Chapter 10, with Z09 \\
N05 & The open narrowband stack running on a small board with a radio module and a software-defined receiver & 17.7 & Out: it needs a purchase and a different board. It is the first purchase to make when the budget exists \\
\bottomrule
\end{tabular}
\caption{The N entries, added in the revision that produced the list. Three of the five became chapters, which is the highest rate of any class and reflects that they were written later and with the bench in view.}
\end{table}

\begin{table}[H]\centering\small
\begin{tabular}{p{9mm}p{82mm}p{11mm}p{43mm}}
\toprule
ID & Project as the original named it & Total & Note \\
\midrule
B01 & Board with an 868 MHz radio module as an open narrowband endpoint & 13.3 & The radio subject this volume approaches through a modem instead \\
B02 & The official narrowband library on a small board & 16.3 & The highest-scoring purchase in the class \\
B03 & A subscription and a real narrowband attachment & 13.2 & Chapter 11 works without one, which is its point \\
B04 & A licensed command-set modem on the board's serial port & 15.7 & The command engine of chapter 11 is written to accept one \\
B05 & Industrial node with the IKS5A1 and an 868 MHz radio & 13.3 & Chapter 14 without the radio \\
B06 & Environment node with the IKS4A1, time of flight and 868 MHz & 10.3 & Two shields, so out twice over \\
B07 & STWIN.box with an external radio & 11.0 & STWIN.box \\
B08 & Transmit energy of 868 MHz bursts & 15.5 & The method is chapter 10 \\
B09 & A long-range contrast node, labelled as such & 12.2 & Useful as a comparison, not as a chapter \\
B10 & Software-defined base station on a Linux host & 9.7 & Host-side, and the sibling volume's territory \\
B11 & Mechanical mounting for the vibration node & 15.3 & Belongs with the condition-monitoring build \\
B12 & Primary-cell powered node & 12.3 & Chapter 7 produces the number that would size the cell \\
B13 & Additional sensors on the vibration board's flexible connector & 9.5 & STWIN.box \\
B14 & A second meter, only if the power instrument turns out to limit you & 6.0 & Appendix C states where the limit actually is \\
B15 & Radio front-end extras & 8.5 & Only after a basic radio works \\
B16 & An industrial fieldbus sibling project & 9.0 & The part has two of those interfaces and no transceiver on the board \\
B17 & A card socket on the expansion header, since the board has none & 13.2 & Confirms the inventory claim: there is no card socket \\
B18 & A coin cell and real time across a power cycle & 10.7 & The crystal is fitted; the coin cell is what is missing \\
B19 & A wireless debug bridge, labelled for debugging only & 9.5 & The coprocessor route named in the front matter \\
B20 & A complete ten-minute demonstration kit & 16.0 & What the twenty chapters add up to anyway \\
\bottomrule
\end{tabular}
\caption{The B entries, every one of which needs a purchase. This volume buys nothing, so all twenty stayed out. They are printed with their scores because the ranking is still useful the day a budget appears, and because the numbering has to survive.}
\end{table}

\section{Board reference}

Everything in this appendix is marked confirmed or to be confirmed. Confirmed
means it was read in a source that names this exact board, and two independent
machine-readable sources were used wherever possible. To be confirmed means it
is the convention of this board family and has not yet been read in the board
manual for MB1363. Nothing marked to be confirmed is written as fact anywhere
else in this book.

\begin{table}[H]\centering\small
\begin{tabular}{p{44mm}p{44mm}p{62mm}}
\toprule
Item & Value & Status \\
\midrule
Microcontroller & STM32H7A3ZIT6Q & Confirmed \\
Core and clock & Cortex-M7 at 280 MHz & Confirmed \\
Flash & 2 MB & Confirmed, dual bank to be confirmed \\
Static memory & Several regions, see chapter 19 & Sizes to be confirmed in the datasheet \\
Reference manual & RM0455, not RM0433 & Confirmed \\
Supply variant & Switched-mode, the Q suffix & Confirmed \\
Debug probe & ST-LINK V3E on board & Confirmed \\
High-speed clock & From the probe, bypass mode, 8 MHz & Confirmed from a working board definition \\
Low-speed crystal & 32.768 kHz, fitted & Confirmed from two independent sources \\
Ethernet & None. No controller and no media access controller & Confirmed from the upstream board description \\
Card socket & None & Confirmed \\
External flash & None & Confirmed \\
Bus transceiver for the automotive bus & None on the board & Confirmed \\
Cryptographic accelerator & None on this part & Confirmed from the peripheral list \\
Random number generator & Present & Confirmed from the peripheral list \\
\bottomrule
\end{tabular}
\caption{Board and silicon summary. The three entries that say none are the ones that most often surprise a reader arriving from material about the better-known member of this family.}
\end{table}

\begin{table}[H]\centering\small
\begin{tabular}{p{34mm}p{30mm}p{44mm}p{42mm}}
\toprule
Function & Pin & Note & Status \\
\midrule
LD1, green & PB0 & User LED & Confirmed, two sources agreeing \\
LD2, yellow & PE1 & User LED & Confirmed, two sources agreeing \\
LD3, red & PB14 & User LED & Confirmed, two sources agreeing \\
B1, user button & PC13 & Pull direction to be confirmed & Confirmed as the pin \\
Virtual COM port & USART3 & The peripheral is settled & Confirmed \\
Virtual COM port pins & PD8 and PD9 & The convention of this board family & To be confirmed in the board manual \\
Trace output & PB3 & Whether it reaches the probe at all & To be confirmed \\
\bottomrule
\end{tabular}
\caption{User interface and console. The peripheral is confirmed and its pins are not, which is exactly the distinction this book insists on.}
\end{table}

\begin{table}[H]\centering\small
\begin{tabular}{p{26mm}p{26mm}p{50mm}p{48mm}}
\toprule
Header pin & Expected port pin & Function used in this book & Status \\
\midrule
D14 and D15 & PB9 and PB8 & I2C1 to all three shields & To be confirmed \\
D13 & To be confirmed & Serial clock for the SPI chapters & To be confirmed \\
D12 & To be confirmed & Master in, slave out & To be confirmed \\
D11 & To be confirmed & Master out, slave in & To be confirmed \\
D10 & To be confirmed & Chip select, driven in software & To be confirmed \\
D0 and D1 & To be confirmed & A serial port for the modem chapters & To be confirmed \\
D2 to D7 & To be confirmed & Interrupt lines from the shields & To be confirmed \\
A0 to A5 & To be confirmed & Converter channels for chapter 6 & To be confirmed \\
Free Zio pins & To be confirmed & Marker pins for the power instrument & To be confirmed \\
IDD jumper & To be confirmed & What it isolates and whether the probe's own current is excluded & To be confirmed \\
\bottomrule
\end{tabular}
\caption{The Arduino-compatible header. Every row is unconfirmed, which is an honest statement of where this book currently stands: the chapters that depend on these rows say so, and the confirm list in the authoring contract carries the same items.}
\end{table}

\section{The two instruments}

This bench has no oscilloscope, no logic analyser and no Joulescope, confirmed on
Sunday 20 September 2026. Every timing and energy claim in this book therefore
comes from one of the two instruments below or from the board's own timers and
cycle counter.

\begin{table}[H]\centering\small
\begin{tabular}{p{44mm}p{58mm}p{48mm}}
\toprule
Property & nRF-PPK2 & MCC 118 on a Raspberry Pi \\
\midrule
What it is & Ammeter and source meter & Analogue acquisition hat \\
Channels & One current path & Eight single-ended inputs \\
Ceiling & 1 A, and that is a hard limit & Input range to be confirmed against the manual \\
Resolution & Dynamic across several ranges & 12 bit \\
Rate & About 100 thousand samples a second, to be confirmed & 100 thousand samples a second aggregate across the channels in use \\
Digital inputs & Eight, time-aligned with the current trace & None. The external trigger is a separate question \\
Supply mode & Can power the target, voltage to be confirmed against the board's needs & Powered from the host board \\
Host software & A Python interface under GPL-2.0: install it, do not fork it into a portfolio repository & The vendor's own library on the host board \\
Used in & Chapters 7, 10 and 12 & Chapters 6 and 12 \\
\bottomrule
\end{tabular}
\caption{The two instruments and their limits. The digital inputs of the power instrument are what make it this book's logic recorder for slow signals, which is the substitute wherever a chapter would otherwise want an oscilloscope.}
\end{table}

The traps are worth stating plainly, because each of them produces a plausible
number rather than an error.

\begin{itemize}
\item \textbf{Leakage through the debug probe.} With the probe connected, current
reaches the target through the debug connection as well as through the measured
path. A measurement taken with the probe attached is not the node's consumption.
Either isolate the target or state clearly that the figure includes the probe,
and never compare one of those numbers with the other.
\item \textbf{The one amp ceiling.} The power instrument measures up to 1 A. The
cellular modems on this bench ask for peaks of about 2 A. They are never fed
through it and never fed from the board; they get their own supply, and the
transmit energy question is answered by other means or not at all.
\item \textbf{Aggregate against per-channel rate.} The acquisition hat's hundred
thousand samples a second are shared across the channels in use. Eight channels
are eight times slower each, which is the arithmetic that turns a proof of a
1 kHz sampling rate into a measurement of something else.
\item \textbf{Time alignment.} The digital inputs are aligned with the current
trace by the instrument, not by the host clock. That alignment is what makes the
phase splitting of chapter 10 possible and is the reason a marker pin is worth
more here than a print statement.
\item \textbf{The export.} Whether the digital lines survive into the exported
file is to be confirmed before a chapter depends on post-processing them.
\end{itemize}

\section{Vocabulary, German and English}

\begin{table}[H]\centering\small
\begin{tabular}{p{40mm}p{40mm}p{72mm}}
\toprule
German & English & Note \\
\midrule
Taktbaum & clock tree & Chapter 1 draws this part's \\
Zwischenspeicher & cache & Chapter 19. Often left untranslated \\
Speicherabbild & memory map & Chapter 19 \\
Speicherschutzeinheit & memory protection unit & Chapter 19 \\
Direktspeicherzugriff & direct memory access & Chapters 4, 6 and 19 \\
Unterbrechung & interrupt & Chapter 3 \\
Unterbrechungsvektor & interrupt vector & Chapter 1 writes the table by hand \\
Ringpuffer & ring buffer & Chapter 2 \\
Warteschlange & queue & Chapter 20 \\
Stapelspeicher & stack & Chapter 20 budgets one per task \\
Halde & heap & Chapter 20 does not link one \\
Ablaufplaner & scheduler & Chapter 20 \\
Aufgabe & task & Usually left as the English word \\
Priorit\"atsumkehr & priority inversion & Chapter 20, measured as a stretch goal \\
Priorit\"atsvererbung & priority inheritance & Chapter 20, and only a basic form \\
Semaphor & semaphore & Chapter 20 \\
Ereignisgruppe & event group & Chapter 20 \\
Zeitgeber & timer & Chapter 6 \\
Abtastrate & sample rate & Chapter 6 proves this externally \\
Aufl\"osung & resolution & Chapter 6 \\
Wandler & converter & Chapter 6 \\
Quantisierung & quantisation & Chapter 18 \\
Festkommazahl & fixed-point number & Chapter 18 \\
Gleitkommazahl & floating-point number & Chapter 18 \\
Filterentwurf & filter design & Chapter 18, done on the host \\
Pr\"ufsumme & checksum & Chapter 5 \\
Rahmen & frame & Chapter 5 \\
Baudrate & baud rate & Chapter 3 \\
Latenz & latency & Chapter 20 measures a distribution of it \\
Echtzeit & real time & Chapter 20 \\
Betriebssystem & operating system & Chapter 20 builds two \\
Wachhund & watchdog & Chapter 12. Usually left as the English word \\
Stromaufnahme & current consumption & Chapters 7 and 10 \\
Ladung & charge & Chapter 7 reports charge per cycle \\
Ruhezustand & stop mode & Chapter 7 \\
Aufwachen & wake-up & Chapter 7 \\
Schwingquarz & crystal oscillator & Fitted on this board at 32.768 kHz \\
Steckverbinder & connector & Appendix B \\
Steckbrett & breadboard & Chapter 19 \\
L\"otbr\"ucke & solder bridge & Nothing here needs one moved \\
Datenblatt & datasheet & The part-specific document \\
Referenzhandbuch & reference manual & RM0455, not RM0433 \\
Anwendungshinweis & application note & Several, listed in appendix E \\
Abnahmekriterium & acceptance criterion & Every chapter has a list of them \\
Messger\"at & instrument & Appendix C \\
Sensorknoten & sensor node & Chapters 8 to 12 \\
Nutzlast & payload & Chapter 9 \\
Funkmodul & radio module & Not on this bench; chapter 11 uses a modem \\
\bottomrule
\end{tabular}
\caption{Working vocabulary, extended from the workbook's glossary. The terms that a German-speaking engineer normally leaves in English are marked, because using the German word where nobody does sounds stranger than the loan word.}
\end{table}

\section{The reference library}

Grouped by subject, and every entry was checked on Sunday 20 September 2026.
Three domains block automated checking, namely the silicon vendor's own site,
its wiki and the core designer's documentation site; the rows drawn from those
are marked as opened by hand on the same day. Licence categories follow the
provenance document: A is safe to depend on and to vendor, B is safe to depend
on but not to copy, C is read and do not copy, D is recommend and never quote at
length, and E means no licence file at all, which is all rights reserved.

\begin{table}[H]\centering\small
\begin{tabular}{p{28mm}p{80mm}p{22mm}p{18mm}}
\toprule
Subject & Source & Licence & Checked \\
\midrule
Normative & RM0455, the reference manual for this part & C & By hand \\
Normative & The STM32H7A3xI datasheet & C & By hand \\
Normative & The MB1363 board user manual & C & By hand \\
Normative & The Armv7-M architecture reference manual, and the Cortex-M7 technical reference manual & C & By hand \\
Toolchain & Lyubka's bare-metal programming guide & A, MIT & Yes \\
Toolchain & Oblaukhov's CMake package layer for this family & A, MIT & Yes \\
Toolchain & Majerle's configurator and editor template & A, MIT & Yes \\
Toolchain & The flashing and debugging tool whose manifest names this die & A & Yes \\
Toolchain & Two register-level projects for this exact board & E, none stated & Yes \\
Console & Memfault's article on printf for embedded systems & D & Yes \\
Buffering & Majerle's lightweight ring buffer & A, MIT & Yes \\
Buffering & Quantum Leaps' lock-free ring buffer with its test suite & A, MIT & Yes \\
Buffering & The bounded queue paper that gives the exact barrier placement & C & Yes \\
Transfers & Majerle's serial port and transfer-engine examples & A, MIT & Yes \\
Framing & Majerle's lightweight packet library & A, MIT & Yes \\
Framing & McQueen's byte stuffing library & A, MIT & Yes \\
Framing & The application note on the checksum unit's reversal settings & C & By hand \\
Timing & Ross's timer, converter and transfer progression on a Cortex-M7 board & A, MIT & Yes \\
\bottomrule
\end{tabular}
\caption{The reference library, part one: normative documents, toolchain, buffering, framing and timing. Checked Sunday 20 September 2026.}
\end{table}

\begin{table}[H]\centering\small
\begin{tabular}{p{28mm}p{80mm}p{22mm}p{18mm}}
\toprule
Subject & Source & Licence & Checked \\
\midrule
Power & The vendor's three power examples for this exact board & C & By hand \\
Power & The Python interface to the power instrument & B, GPL-2.0 & Yes \\
Power & A published method for measuring with that instrument & D & Yes \\
Codec & The smallest encoder, effectively public domain & A, CC0 & Yes \\
Codec & A schema-driven generator used in production & A, Apache-2.0 & Yes \\
Codec & A minimalist alternative & A, MIT & Yes \\
Modem & Majerle's cellular library, which requires an operating system & A, MIT & Yes \\
Test rig & The host unit test framework and its build tool & A & Yes \\
Test rig & The header-only fake framework, with its documented limits & A, MIT & Yes \\
Test rig & The published series on hardware-in-the-loop testing & D & Yes \\
Watchdogs & The two-layer watchdog article, and Ganssle on why naive ones fail & D & Yes \\
Sensors & The register-level driver collection for every sensor on both shields & A, BSD-3-Clause & Yes \\
Sensors & The other operating system's shield definitions, clearer than the manual & A, Apache-2.0 & Yes \\
Sensors & Ready-made in-sensor classifier and state-machine configurations & A, BSD-3-Clause & Yes \\
Sensors & The in-sensor processor examples, with no licence at the root & E & Yes \\
Sensors & The time-of-flight driver republished permissively & A, BSD-3-Clause & Yes \\
\bottomrule
\end{tabular}
\caption{The reference library, part two: power, the codec, the modem, the test rig and the shields. Checked Sunday 20 September 2026.}
\end{table}

\begin{table}[H]\centering\small
\begin{tabular}{p{28mm}p{80mm}p{22mm}p{18mm}}
\toprule
Subject & Source & Licence & Checked \\
\midrule
Transforms & The signal-processing library and its Python wrapper & A, Apache-2.0 & Yes \\
Transforms & The learning path that builds a pipeline against a reference & D & Yes \\
Transforms & The neural kernels that reach this core through the instruction extension & A, Apache-2.0 & Yes \\
Transforms & The small-model library for the statistical alternative & A, MIT & Yes \\
Cache & The application notes on the cache, on the protection unit and on cache performance & C & By hand \\
Cache & The vendor engineer's article putting the three fixes side by side & C & Yes \\
Cache & The vendor's own write-up of the fourth failure & C & Yes \\
Cache & The three-part written series on this core's cache & D & Yes \\
Cache & The open report against the vendor's own examples & A & Yes \\
Kernels & The kernel and its free official book & A, MIT & Yes \\
Kernels & The other operating system's in-tree port for this exact board & A, Apache-2.0 & Yes \\
Kernels & The tick-based measurement library & A, Apache-2.0 & Yes \\
Kernels & The two-part study whose method transfers and whose numbers do not & D & Yes \\
Kernels & The active-object framework and its video course & D, GPL and AGPL-3.0 & Yes \\
Reach & The coprocessor firmware that exposes messaging commands directly & A, Apache-2.0 & Yes \\
Reach & The transport-agnostic messaging client & A, MIT & Yes \\
Reach & The framed remote-call generator and the small payload encoder & A & Yes \\
Rust & The hardware crate, which routes this part through its twin & A, 0BSD & Yes \\
Rust & The asynchronous framework, which names this exact package & A & Yes \\
\bottomrule
\end{tabular}
\caption{The reference library, part three: inside the core, reach, and language support. Checked Sunday 20 September 2026. A row marked E has no licence file, which means all rights reserved: read it, do not vendor it.}
\end{table}

\section{The variant matrix}

Each variant is built in full exactly once, in the chapter where it teaches
most, and is referenced from everywhere else it applies. This table is the index
to that rule. A cell reading \emph{built} means that chapter builds that axis in
full; a number is a pointer to the chapter that does; a dash means the axis does
not apply to that chapter. The chapter's own variants table gives the detail
this one has no room for.

\begin{table}[H]\centering\small
\begin{tabular}{p{9mm}p{17mm}p{17mm}p{17mm}p{17mm}p{17mm}p{17mm}p{17mm}}
\toprule
No. & Execution & Sync & Time and safety & Peripheral & Operating system & Language & Reach \\
\midrule
1 & built & flag & SysTick & probe disk & bare metal & built & local \\
2 & built & built & 3 & 3 & bare metal & 9 & local \\
3 & built & built & 6 & 17 & bare metal & built & local \\
4 & built & 2 & 6 & 17 & bare metal & 9 & local \\
5 & 4 & 2 & 6 & built & bare metal & 9 & local \\
6 & built & 2 & built & 17 & bare metal & 9 & local \\
7 & 4 & 2 & built & 17 & bare metal & 9 & local \\
8 & 3 & 2 & 6 & 17 & bare metal & 9 & local \\
9 & 4 & 2 & 6 & 17 & bare metal & built & built \\
10 & 4 & 2 & 7 & 17 & bare metal & 9 & 9 \\
11 & 3 & 2 & 7 & 17 & bare metal & 9 & built \\
12 & 4 & 20 & built & 17 & bare metal & built & 9 \\
13 & 4 & 2 & 6 & built & 20 & 9 & 16 \\
14 & 4 & 2 & 6 & 17 & bare metal & 9 & 16 \\
15 & 4 & 2 & 6 & 17 & bare metal & 9 & 16 \\
16 & 4 & 2 & 6 & 17 & bare metal & 9 & built \\
17 & 3 & 2 & 6 & built & bare metal & built & local \\
18 & 3 & 2 & 6 & 16 & bare metal & 17 & 16 \\
19 & 4 & 2 & 6 & 17 & bare metal & 17 & local \\
20 & built & built & 7 & 17 & built & 9 & 16 \\
\bottomrule
\end{tabular}
\caption{The variant matrix. Read a column downward to follow one axis across the book; read a row across to see what one chapter touches. Chapter 20 carries more of this matrix than any other chapter, which is why it is last and why it is six evenings rather than four.}
\end{table}

What the \emph{built} cells mean, column by column. Execution: the polled loop
and the flag are chapter 3, the ring buffer is chapters 2 and 3, transfers are
chapter 4, double buffering is chapter 6, and the kernel task set is chapter 20.
Synchronisation: the volatile flag is chapters 2 and 3, and every kernel object,
from the semaphore to the stream buffer, is chapter 20. Time and safety: the
tick and the general timer are chapter 6, the low-power timer, the real-time
clock and tickless idle are chapter 7, and the two watchdogs are chapter 12.
Peripheral substitution is chapter 17, with the checksum unit in chapter 5 and
the shield bus topologies in chapter 13. Operating system: bare metal
throughout, and both kernels in chapter 20. Language: the interpreter path is
chapter 1, C++ and the host twin are chapter 9, Rust is chapter 17, and the
second host language is chapter 12. Reach: the framed link to an edge host is
chapters 9 and 11, and the three classifier placements are chapter 16.

\section{Languages on this part}

\begin{table}[H]\centering\small
\begin{tabular}{p{26mm}p{22mm}p{54mm}p{42mm}}
\toprule
Language & On this part & What is actually true & The board that does work \\
\midrule
C & Runs & The baseline of this book, verified on this exact device & This one \\
C++ & Runs & Same toolchain. Exceptions cost roughly 10 to 30 kB of tables, which on 2 MB of flash is a choice rather than a necessity & This one \\
Rust & Runs & The hardware crate routes this part through its twin's peripheral module and selects the right manual, and the asynchronous framework names this exact package and pin count & This one, although no worked example exists for it yet \\
MicroPython & Runs & The board definition is in the main branch with its own linker script, and prebuilt firmware is published. Its declared feature list is empty, so do not promise networking or host mode without building it & This one \\
Go & Sibling only & It compiles well, but the only board it offers in this family hardcodes a different member's memory size in its only linker script, and its own comments admit two memory domains are unmapped and therefore skipped by cache maintenance & The better-known member of this family, the H743 board, to be confirmed \\
Ada & Sibling only & The free toolchain and driver library support the previous generation of this core, on two development boards, and not this one & The two discovery boards of the previous generation, to be confirmed \\
C sharp & Sibling only & The small framework reaches this family only through a different member, so the family is supported and this part is not & The H743 board, to be confirmed \\
JavaScript & Does not & There is no board in this family at all, and the nearest supported parts are not even the same core & None in this family \\
Java & Does not & No free maintained virtual machine ports it here, and the one serious commercial option may not be used in a shipped product without terms this book will not assume & None that may be shipped \\
\bottomrule
\end{tabular}
\caption{What runs on this part, what runs only on a sibling, and what does not run at all. The middle group is the dangerous one: a search finds a project that names this family, and the linker script inside it is written for a different member.}
\end{table}

Two remarks on the middle group. First, a project that supports the family is
not the same as a project that supports this part, and the distinction is the
same one the front matter is about. Second, where a language reaches only a
sibling, the sibling named in the last column is the board to buy or borrow if
that language is the requirement, and the gap is worth filling: writing the
missing linker script and board definition for this part is a small, visible
contribution in every one of those three projects.

\section{Courses and learning paths}

\begin{table}[H]\centering\small
\begin{tabular}{p{62mm}p{44mm}p{46mm}}
\toprule
Course or path & Licence & How to use it \\
\midrule
The free course that walks one problem from a superloop to an event-driven architecture across fifty-six lessons & Code AGPL-3.0 & The best free companion to chapter 20. Recommend it, watch it, and never quote its code \\
The free course whose code is 0BSD and whose slides are CC BY 4.0 & 0BSD and CC BY 4.0 & The safest material in the pool to build an exercise on, although its examples are not on this board \\
The best free course on the other operating system, from a different silicon vendor & Free to register & Its concepts transfer directly. Chapter 20's second half is easier after it \\
The university specialisation that teaches the mathematics of real-time scheduling & Paid, audit available & Almost nothing else teaches this. Worth it before making a scheduling claim \\
The free-to-audit course on embedded machine learning & Free to audit & Background for chapter 16 \\
The free open textbook that grew out of a university course & Open & The closest thing to a canonical curriculum for the machine-learning chapters \\
The learning path that builds a signal-processing pipeline in Python against a reference & Vendor learning material & The shortest route to fluency with the wrapper chapter 18 uses \\
The most recommended paid course on this kernel and this silicon family & Paid & Thorough, and written against a different member of the family. Read it with the front-matter warning in mind \\
The silicon vendor's own free training modules & Vendor terms & They cover this family well, and there is no course specific to this part \\
\bottomrule
\end{tabular}
\caption{Courses, with their licences. Two of the best free courses in this list carry terms that forbid reproduction in a commercial work. They are recommended here and never quoted, and where this book needed what they teach it re-derived it.}
\end{table}

\section{Roadmaps}

Three roadmaps matter for a reader who finishes this volume, and they point in
different directions on purpose.

\begin{itemize}
\item \textbf{The community roadmap} separates firmware from embedded Linux from
hardware and says which to weight. It is the map to read before choosing what to
learn next, and its central claim, that projects teach what reading does not, is
this book's own thesis. The accompanying essay on bridging the gap between
reading and shipping is the shorter version of the same argument.
\item \textbf{The machine-learning curriculum.} The free open textbook that grew
out of a university course and its certificate is the closest thing to a
canonical path for what chapters 16 and 18 open up. With this part sitting at
the large end of the small-model envelope, memory is rarely the binding
constraint; throughput and energy are, and that reframing is what the textbook
supplies.
\item \textbf{The standards a reader eventually meets.} The consumer device
security baseline, the industrial component security series, the two evaluation
schemes, and the coding standard that matters once safety enters. None of them
is needed to build any chapter here. All of them appear in a job description
sooner or later, and knowing which one governs which question is worth an
afternoon.
\end{itemize}

Beyond the roadmaps there are the gaps this book would fill if its chapters were
built and published. There is no rigorous comparison of the two numeric formats
for transforms on this core; no peer-reviewed work measures all three placements
of a classifier on one task on one set of hardware; nobody has measured six
execution models against one problem on one board with a stated method; there is
no published measurement of priority-inversion blocking on this core; nobody has
implemented the active-object pattern twice and compared it honestly; no
maintained permissive library offers a declarative region table with cache
attributes; and no published method exists for the power instrument on this
bench. Each of those falls inside a chapter that is already planned, which is
the difference between a list of open questions and a list of excuses.

\end{document}
